%%%%%%%%%%%%%%%%%%%%%%%%%%%%%%%%%%%%%%%%%%%%%%%%%%%%%
%
\documentclass[11pt,english]{amsart}
%
\usepackage[english]{babel}
\usepackage{amsmath,amssymb,url,xspace,amsthm}
%
%%%%%%%%%%%%%%%%%%%% DIAGRAMS %%%%%%%%%%%%%%%%%%%%%%%\usepackage{xypic}
%\usepackage[pdftex]{graphicx}
%\usepackage{epstopdf}
%\usepackage{pstricks-add}
%\usepackage{pgf,tikz}
%\usetikzlibrary{arrows,automata}
%%%%%%%%%%%%%%%%%%%%%% FONTS %%%%%%%%%%%%%%%%%%%%%%%%%%%\usepackage[utf8]{inputenc}  
%\usepackage[T1]{fontenc}
%\mathversion{bold}
%\usepackage{pxfonts}
%\usepackage{kmath}
\usepackage{mathrsfs,eucal}
%\usepackage{fouriernc}
%\usepackage{mathpple}
%\usepackage{fourier}
\usepackage{times}
\usepackage{hyperref}
\usepackage{xcolor}
\pagecolor{white}
\hypersetup{
  colorlinks   = true, %Colours links instead of ugly boxes
  linkcolor={blue!50!black}, %Colour of internal links
  citecolor={blue!50!black}, %Colour of citations
  urlcolor={blue!80!black}, %Colour for external hyperlinks  
}
%\usepackage{etoolbox}
%\patchcmd{\thebibliography}{*}{}{}{}
%\pretocmd\thebibliography{\csname c@secnumdepth\endcsname=-2 }{}{}
%%%%%%%%%%%%%%%%%%%%%%%%%%%%%%%%%%%%%%%%%%%%%%%%%%%%%\renewcommand{\rm}{\mathrm}
\renewcommand{\sf}{\mathsf}
\renewcommand{\bf}{\mathbf}
\newcommand{\bb}{\mathbb}
\newcommand{\fr}{\mathfrak}
\newcommand{\cal}{\mathcal}
\newcommand{\scr}{\mathscr}
%%%%%%%%%%%%%%%%%%%%%%%%%%%%%%%%%%%%%%%%%%%%%%%%%%%%%
\renewcommand{\H}{\mathrm{H}}
\newcommand{\Bt}{\mathrm{B}}
\newcommand{\id}{\mathrm{id}}
\newcommand{\ab}{\mathrm{ab}}
\newcommand{\im}{\mathrm{im}}
\newcommand{\F}{\mathrm{F}}
\newcommand{\U}{\mathrm{U}}
\newcommand{\V}{\mathrm{V}}
\newcommand{\W}{\mathrm{W}}
\newcommand{\FT}{\mathfrak{F}}
\newcommand{\add}{\mathrm{s}}
\renewcommand{\sp}{\mathrm{sp}}
\newcommand{\tsp}{\widetilde{\mathrm{sp}}}
\newcommand{\Z}{\mathrm{Z}}
\newcommand{\diag}{\mathrm{diag}}
\newcommand{\vun}{\mathbf{1}}
\newcommand{\Fl}{\mathrm{Fl}}
%%%%%%%%%%%%%%%%%%%%%%%%%%%%%%%%%%%%%%%%%%%%%%%%%%%%%
\DeclareMathOperator*{\argmin}{argmin}
%%%%%%%%%%%%%%%%%%%%%%%%%%%%%%%%%%%%%%%%%%%%%%%%%%%%%\newcommand{\Rep}{\mathsf{Rep}}
\newcommand{\Mod}{\mathrm{M}}
\newcommand{\Dd}{\mathrm{D}}
\newcommand{\Rd}{\mathrm{R}}
\newcommand{\Ld}{\mathrm{L}}
\newcommand{\tX}{\tilde{X}}%
\newcommand{\RN}{\mathrm{RN}}
%%%%%%%%%%%%%%%%%%%%%%%%%%%%%%%%%%%%%%%%%%%%%%%%%%%%%
\newcommand{\llangle}{\langle\!\langle}
\newcommand{\rrangle}{\rangle\!\rangle}
%%%%%%%%%%%%%%%%%%%%%%%%%%%%%%%%%%%%%%%%%%%%%%%%%%%%%
%
\newcommand{\PnL}{\mathrm{P\&L}}
\newcommand{\ONPnL}{\mathrm{O\textrm{/}N~P\&L}}
\newcommand{\IDPnL}{\mathrm{Intraday~P\&L}}
%
%%%%%%%%%%%%%%%%%%%%%%%%%%%%%%%%%%%%%%%%%%%%%%%%%%%%%
%
\newcommand{\BS}{\mathrm{BS}}
\newcommand{\ODF}{\mathrm{1DF}}
\newcommand{\ATM}{\mathrm{ATM}}
\newcommand{\MTM}{\mathrm{MTM}}
\newcommand{\EOD}{\mathrm{EOD}}
\newcommand{\SOD}{\mathrm{SOD}}
\newcommand{\DCB}{\mathrm{DCB}}
%
%%%%%%%%%%%%%%%%%%%%%%%%%%%%%%%%%%%%%%%%%%%%%%%%%%%%%
%
\renewcommand{\AA}{\mathbb{A}}
\newcommand{\BB}{\mathbb{B}}
\newcommand{\CC}{\mathbb{C}}
\newcommand{\EE}{\mathbb{E}}
\newcommand{\FF}{\mathbb{F}}
\newcommand{\GG}{\mathbb{G}}
\newcommand{\II}{\mathbb{I}}
\newcommand{\LL}{\mathbb{L}}
\newcommand{\NN}{\mathbb{N}}
\newcommand{\PP}{\mathbb{P}}
\newcommand{\QQ}{\mathbb{Q}}
\newcommand{\RR}{\mathbb{R}}
\newcommand{\TT}{\mathbb{T}}
\newcommand{\VV}{\mathbb{V}}
\newcommand{\ZZ}{\mathbb{Z}}
%
\newcommand{\f}{\ell}
%
\newcommand{\Cov}{\mathrm{Cov}}
%
\newcommand{\disc}{\mathrm{disc}}
\newcommand{\proj}{\mathrm{proj}}
\newcommand{\spread}{\mathrm{spread}}

\newcommand{\Q}{\Xi}
\newcommand{\QZ}{{\Xi\mathrm{Z}}}
\newcommand{\Qto}{\mathrm{Q}}
%
\theoremstyle{definition}
\newtheorem{prop}{Proposition}
\newtheorem*{prop*}{Proposition}
\newtheorem{lemm}{Lemma}
\newtheorem{theo}{Theorem}
\newtheorem*{theo*}{Theorem}
\newcommand{\theoautorefname}{Theorem}
\newcommand{\propautorefname}{Proposition}
\newcommand{\lemmautorefname}{Lemma}
\newcommand{\remaautorefname}{Remark}
\newcommand{\algoautorefname}{Algorithm}
\newtheorem{coro}{Corrolary}
\newtheorem{defi}{Definition}
\newtheorem{rema}{Remark}
\newtheorem{exem}{Example}
\newtheorem*{hypo}{Hypothesis}
\newtheorem{algo}{Algorithm}
%
%\newtheorem{nota}{Notation}
%%%%%%%%%%%%%%%%%%%%%%%%%%%%%%%%%%%%%%%%%%%%%%%%%%%%%%%%%%%%%%%%%%%%%
\addtolength{\hoffset}{-2cm}
\addtolength{\voffset}{-2cm}
\setlength{\textwidth}{17cm}
\setlength{\textheight}{24cm}
%%%%%%%%%%%%%%%%%%%%%%%%%%%%%%%%%%%%%%%%%%%%%%%%%%%%%%%%%%%%%%%%%%%%%
\title{Notes on collateral discounting}%

%\title[Not so linear Linear rates]{Not so linear Linear rates \\ 
%(Notes on Discounting)}%
%\title{Rapport de stage}%\\ Sous la direction de Pascal Delanoë}
%\date {}
\author{Yassine Bounfour}

%\address{}
%\email{yassinebounfour@gmail.com}
%\urladdr{}
%\keywords{Discounting, collateral, cross-currency swap, FX forward, floored CSA, cheapest to deliver.}
%\altkeywords{D-modules, vanishing cycles, Verdier specialization, Sebastiani-Thom, convolution}

\begin{document}
%\def\smfbyname{}

%%%%%%%%%%%%%%%%%%%%%%%%%%%%%%%%%%%%%%%%%%%%%%%%%%%%%

\maketitle


\begin{abstract}
In these personal notes, we generalise the HJM framework\footnote{i.e. the arbitrage free pricing framework under stochastic discount rates} to multiple currencies and stochastic collateral. Then, we define the Linear Gaussian Collateral Model by specialising this framework to separable and deterministic volatility and constant funding spread thus generalising the Black-Scholes-Hull-White model to account for stochastic collateral. In this model, we derive the dynamic of all market observables (ZCBs, FX-forwards, FRAs, swap rates, currency basis spreads etc\ldots)  as well as their collateral convexity adjustments. Finally, we describe a procedure to calibrate this model accurately to market quotes and/or historical data.  
\end{abstract}

%\newpage

\subsection*{TODO LIST}%

All items below were addressed on 2026-10-09 and checked numerically (see the verification remarks in each section).
\begin{enumerate}
\item \emph{Done.} Gaussian approximation for swap rates and currency basis spreads: \autoref{prop:gauss_swap_rate}, \autoref{prop:gauss_cb_spread}, proofs in \autoref{sec:gaussian_approx_proofs}.
\item \emph{Done.} Reconstruction formulas in the Quasi-Gaussian model: \autoref{theo:qgcm} and \autoref{rema:qgcm_vfc}.
\item \emph{Done.} Algorithm for the computation of covariance integrals: \autoref{algo:covariance_integrals}.
\item \emph{Done.} Drifts and covariance matrices in the LGCM under risk neutral, $T$-forward and jumping num\'eraire measures: \autoref{prop:lgcm_transition_rn}, \autoref{prop:lgcm_drifts}.
\item \emph{Done.} Consistency of interpolation with model dynamic: \autoref{ssec:interpolation_consistency}.
\end{enumerate}
Remaining open points found on the way: the Gaussian approximations are for a single collateral basis block ($f\to d$); options on currency basis spreads with smile are not covered.


\tableofcontents




\newpage%
\setcounter{section}{-1}
\section{Introduction}%

\subsection{Motivation}
Discount and index curves are built from market quotes for the prices of liquid rates instruments. They serve as interpolation scheme for these prices, allowing a market participant to value non-standard instruments. 
These liquid instruments typically come with a standardised CSA which we call the \emph{native} collateral so the corresponding curves are only valid for the valuation of cashflows under the same CSA. 
In the OTC markets, non-native CSAs are extremely common but since hedging instruments, the standardised liquid instruments, have different CSAs, one has to rely on a model to value those instead. 

Since the 2008 financial crisis, market participants recognised that currency basis and funding\footnote{Libor-OIS spread or Euribor-EONIA spread for example.} basis spreads could not be neglected. This lead to a new paradigm where the USD Libor discounting practice was replaced by a multi-curve framework with a different discount curve for each payment currency and CSA. The current market practice for building these curves relies on the assumption that funding spreads are deterministic. The consensus around this practice comes from the following very appealing property: it allows to account for funding cost with very minimal changes to exotic pricing models: namely replacing the "risk-free" discount curve by the CSA-specific ones.

However, at the same time, currency basis and funding spreads also became more volatile which generated demand for models with stochastic spreads to quantify and hedge the associated risks. But, being based on deterministic spread assumptions, the CSA discount curves are fundamentally inconsistent with such models and can no longer be used as inputs. So stochastic spreads models require a new change of paradigm. Until now, derivatives modelling was always assumed to be linear and incremental when it turns out to be a circle! The recipe used to be: first build the curves for linear products, then the vol surfaces for vanilla options and finally dynamic models for exotics. And the CSA-discounting market practice allowed to maintain this hierarchy. But, under non-native collateral, even the simplest linear products are actually exotic streams of cash-flows whose valuation depend on the covariance between the underlying coupons and the collateral. So non-native discount and index curves need to be built at the same time as we calibrate the stochastic spread model. 


%In order to calibrate stochastic spread models, a more refined approach has to be taken where the curve building is part of the model calibration. 

%While derivatives modelling always assumed the problem to be linear: first build the curves, then the vol surfaces then the dynamic models, it turns out to be a circle! We need to build the curves at the same time as we calibrate the model. 

\subsection{Examples}
Let us give a few examples where stochastic currency basis risk is material.


\begin{exem}[Currency basis options] 
Just like the time value of a swaption comes from the volatility of the swap rate:
\begin{eqnarray}
  \MTM^{\rm{Swaption}}_t &=& A^{d|c}_t \EE_t^{A^{d|c},T;M}\left[ \left(S^{d|c}_{T;M} - K\right)_+ \right] 
\end{eqnarray}
 the time value of an option to enter a currency basis swap depends on the volatility of the currency basis spread
 \begin{eqnarray}
  \MTM^{\rm{CBSwaption}}_t &=& A^{f|c}_t \EE_t^{A^{f|c},T;M}\left[ \left(B^{f/d|c}_{T;M} - K\right)_+ \right] .
\end{eqnarray}
\end{exem}


\begin{exem}[CVA]
The CVA can be seen as an option on a netting set of trades. As a result, any netting set involving several currencies has currency basis volatility risk. 
\end{exem}


\begin{exem}[FVA] 
FVA is defined as the cost of hedging a non-collateralised trade with collateralised instruments. By hedging with collateralised trades, the trader is subject to margin calls that he needs to fund by borrowing from the treasury. The FVA is thus defined as the difference between the price under CSA discounting and under treasury rate discounting. If either of those is correlated to the underlying, the FVA is exposed to collateral convexity.  
\end{exem}




\subsection{Structure of these notes}

\begin{enumerate}
\item \emph{Valuation of collateralised trades}: reminder on risk-neutral pricing to set notations.
\item \emph{Market Models}: tentative formal definition of the class of models we are considering.
\item \emph{Deterministic Collateral Basis Model}: short description of the current standard market practice for the valuation of fixed and floating rates cash-flows (swaps, fx-forwards, cross-currency swaps, etc...). 
\item \emph{Arbitrage free models with stochastic collateral}: extension of the HJM framework to fixed cash-flows under stochastic collateral.
\item Specialisation to deterministic volatility to get explicit formulas\footnote{This could easily be extended to encompass displaced lognormal diffusions thus capturing the value of the forward, the atm vol and the atm skew of the calibration targets.} for convexity adjustments.
\item Specialisation to separable volatility to get finite dimensional Markovian dynamic.
\item \emph{Specification of the LGCM}. By combining the two previous sections we generalise the IR Linear Gaussian model and FX Black-Scholes model to account for stochastic collateral. The main twist is the closed form approximation for the volatility of the currency basis spread. 
\item \emph{Calibration of the LGCM}. The main point is that the EUR$|$USD discount curve and the EUR collateral basis spot volatility should be \emph{jointly} bootstrapped. 
%\item 
\end{enumerate}

Appendices contain details of mathematical proofs. 

%%%%%%%%%%%%%%%%%%%%%%%%%%%%%%%%%%%%%%%%%%%%
%% Risk neutral pricing under perfect collateral
%%%%%%%%%%%%%%%%%%%%%%%%%%%%%%%%%%%%%%%%%%%%
\newpage
\section{Valuation of collateralized trades}%


\subsubsection*{Notations}%
In these notes, we will use the following notations:
\begin{itemize}
\item $u$, reference currency (USD)
\item $e$, generic currency
\item $d,f$, generic currency pair
\item $c$, generic collateral
\item $i,j$, generic collateral currencies
\item $t$, valuation date
\item $T$, expiry date  
\item $\theta$, tenor of an index rate (typically 3M for a Libor)
\item $M$, tenor of a swap
\item $\ell$, frequency of a swap (typically $\ell = 4$ if $\theta$ is 3M). 
\end{itemize}


\subsection{Risk-neutral pricing of collateralized trades}

We summarise the main results about the valuation of derivatives under perfect collateral in a perfect Brownian market (see Piterbarg \cite{Piterbarg:Cooking} and McCloud \cite{McCloud:CollateralVolatility}).

We will always assume that the following hypothesis holds. 
\begin{hypo}[Frictionless market]
We will say that the market is \emph{frictionless} if the following properties are satisfied:
\begin{enumerate} 
\item Absence of arbitrage.
\item Assets are infinitely divisible
\item There are no restriction on transactions 
\item Transactions are instantaneous
\item There are no jumps 
\end{enumerate}
\end{hypo}

[TODO: Compare to hypothesis in Kjaer's Consistent XVA metrics]

We will always assume that the following hypothesis holds. 
\begin{hypo}[Brownian market]
A Brownian market is a 
\begin{enumerate} 
\item Brownian assumption: under the historical measure $\PP$, asset prices are Ito processes driven by a finite dimensional Brownian motion $W^\PP$. 
\item Absence of arbitrage.
\end{enumerate}
\end{hypo}

Under these hypothesis, one can derive the whole risk-neutral pricing framework for perfectly collateralised trades as follows: 
\begin{enumerate}%{alpha}
\item For each currency $e$, there exist a unique measure $\QQ^{e,\RN}$ equivalent to the historical measure $\PP$ such that, the value $V^{e|c}$ of a self-financed portfolio in currency $e$ under domestic collateral rate $r^c$ satisfies the martingale property:
\begin{eqnarray}
	\EE^{e,\RN}_t\left[ e^{-\int_t^T r^c_u du} V^{e|c}_T \right] &=& V^{e|c}_t
\end{eqnarray} 
i.e. $V^{e|c}/M^{c}$ is a $\QQ^{e,\RN}$-martingale where $M^{c}_t = e^{\int_0^t r^c_u du}$. This measure\footnote{This is the measure associated with the numeraire $M^{c}$. And more generally, for each numeraire $N$ satisfying an integrability condition, there is an associated measure $\QQ^N$.} is called the \emph{risk-neutral measure} for currency $e$. 
\item For non-native collateral $i$, the \emph{effective discount rate} $r^{e|i}_t$ is \emph{defined} as the rate such that the formula
\begin{eqnarray}
	\EE^{e,\RN}_t\left[ e^{-\int_t^T r^{e|i}_u du} V^{e|i}_T \right] &=& V^{e|i}_t
\end{eqnarray} 
is still valid\footnote{We could extend this to non-cash collateral. If a security $S_t$ is posted as collateral, and  it is fully rehypothecable, then the party holding the collateral can repo it, earning the corresponding repo rate and the exact same reasoning applies. As suggested by the usual FX-EQ analogy, posting foreign cash as collateral is essentially equivalent to posting a domestic security with repo rate $r^{f|d}_t = r^{d|d}_t - \mu^{df}_t$.}. % $\frac{dS_t}{S_t} = (r^{d|d}_t - r^{S,\mathrm{repo}}_t) dt + \sigma^S_t \,dW^{d,\RN}_t$.}.
\item The drift of the FX spot rate under the domestic risk-neutral measure is defined as
\begin{eqnarray}
	\mu^{df}_t &:=& \EE^{d,\RN}_t\left[\tfrac{dx^{df}_t}{x^{df}_t}\right]
\end{eqnarray} 
It satisfies the relations\footnote{Very wonky remark!!! Mathematicians will recognise these relations as \emph{cocycle} conditions. Let us write $K^n$ for the vector space of  $\fr{X}^n$ valued processes $c = (c^x_t)_{x\in \fr{X}^n}$. Define $d: K^n \to K^{n+1}$ by $dc^{x}_t = \sum_i (-1)^i c^{x^{(i)}}_t$ where $x^{(i)} = (x^1,\ldots,\widehat{x}^i,\ldots,x^n)$ is obtained by dropping the $i$-th component. This defines a co-homological complex. Moreover this complex has trivial cohomology. In other words, the cocycle relation implies that there exists $(r^x_t)_{x\in \fr{X}}$ such that $\mu^{df} = r^d_t - r^f_t$ for all $d,f$. However, the choice of the processes $(r^x_t)_{x\in \fr{X}}$ is not unique since adding a process (a collateral basis) $c_t$ to all rates cancels out when taking the difference. In summary, one can interpret the FX-drifts $(\mu^{df})_{df\in \fr{X}^2}$ as a $2$-cocycle while the collateral basis $q^{i\to j}$ correspond to a coboundary.} 
\begin{eqnarray}
	\mu^{df}_t + \mu^{fd}_t  &=& 0 \\
	\mu^{de}_t + \mu^{ef}_t + \mu^{fd}_t  &=& 0 .
\end{eqnarray} 
\item The change of payment currency formula for effective discount rates is: 
\begin{eqnarray}
	r^{d|i}_t &=& \mu^{df}_t + r^{f|i}_t
\end{eqnarray} 
\item The short collateral basis is defined as 
\begin{eqnarray}
	q^{i\to j}_t &=& r^{e|j}_t  - r^{e|i}_t
\end{eqnarray} 
It does not depend on the payment currency $e$. So, for any currency $e$, we have the change of collateral formula:
\begin{eqnarray}
	r^{e|j}_t &=&  r^{e|i}_t + q^{i\to j}_t
\end{eqnarray} 
Moreover, we have the cocycle relations:
\begin{eqnarray}
	q^{i\to j}_t + q^{j\to i}_t &=& 0 \\ 
	q^{i\to j}_t + q^{j\to k}_t + q^{k\to i}_t  &=& 0 
\end{eqnarray} 

\item As a result, the dynamic of the spot FX rate is 
\begin{eqnarray}
	\frac{dx^{df}_t}{x^{df}_t}&=& \left( r^{d|d}_t - r^{f|f}_t - q^{f\to d}_t\right) dt + \sigma^{df}_t\cdot dW^{d,\RN}_t 
\end{eqnarray} 
\item The change of risk neutral measure is the usual one\footnote{Given a semi-martingale $L_t$, we write $\mathcal{E}(L)_t$ for the corresponding Doleans-Date exponential martingale: $d\mathcal{E}(L)_t = \mathcal{E}(L)_{t_-}dL_t$. Since we are only considering continuous semi-martingales, $\mathcal{E}(L)_t = e^{L_t - L_0 - \frac{1}{2}\langle L\rangle_t}$.}:
\begin{eqnarray}
	\EE^{d,\RN}_t\left[\frac{d\QQ^{f,\RN}}{d\QQ^{d,\RN}}\right] 
	~=~ e^{-\int_0^t \mu^{df}_u\, du} \frac{x^{df}_t}{x^{df}_0} 
	~=~ \mathcal{E}\left(\int_0^\cdot \sigma^{df}_v \cdot dW^{d,\RN}_v \right)_t 
\end{eqnarray} 
\end{enumerate}




%%%%%%%%%%%%%%%%%%%%%%%%%%%%%%%%%%%%%%%%%%%%
%% Instruments
%%%%%%%%%%%%%%%%%%%%%%%%%%%%%%%%%%%%%%%%%%%%
%\newpage%
\subsection{Standard instruments and market observables}\label{sec:market_observables_list}

We now list standard instruments of interest to us as well as the corresponding market quotes and set our notations:
\begin{itemize}
\item Zero coupon bonds: the value at $t$ in currency $e$ of the bond paying one unit of currency $e$ at $T$ under collateral $i$ is denoted by: 
\begin{eqnarray}
  Z^{e|i}_{t,T} &=& \EE^{e,\RN}_t \left[ e^{-\int_t^T r^{e|i}_u \,du} \right]
\end{eqnarray} 
The measure associated with this num\'{e}raire is denoted by $\QQ^{e|i,T}$. The instantaneous forward rate is defined by 
\begin{eqnarray}
  r^{e|i}_{t,T} &=& - \partial_T \log Z^{e|i}_{t,T}  
\end{eqnarray} 
\item Forward zero coupon bonds: 
\begin{eqnarray}
	Z^{e|i}_{t,S,T} ~=~ \frac{Z^{e|i}_{t,T}}{Z^{e|i}_{t,S}} ~=~ e^{-\int_S^T r^{e|i}_{t,u}\,du}
\end{eqnarray} 
\item Linear discount rates: $R^{e|i}_{t,T;\theta}$ is defined by 
\begin{eqnarray}
  1 + \theta R^{e|i}_{t,T;\theta} &=& \frac{Z^{e|i}_{t,T}}{Z^{e|i}_{t,T+\theta}} 
  ~=~ e^{\int_T^{T+\theta} r^{e|i}_{t,u}\,du}
\end{eqnarray} 
%\begin{eqnarray}
%	1 + \theta R^{e|i}_{t,T;\theta} &=& \frac{1}{\theta}\left( \frac{Z^{e|i}_{t,S}}{Z^{e|i}_{t,S+\theta}} - 1 \right) 
%\end{eqnarray} 
\item Forward Libors. The par strike of an FRA in currency $e$ under collateral $i$ is
\begin{eqnarray}
  	L^{e|i}_{t,T;\theta}	&=& \EE^{e|i,T+\theta}_t \left[L^e_{T;\theta}\right] 
\label{eq:forward_libor}
\end{eqnarray} 
where $L^e_{T;\theta}$ is the value of the Libor index with tenor $\theta$, fixing as $T$ and paid at $T+\theta$. 
\item Futures. A future contract is economically equivalent to a FRA with domestic collateral earning a collateral rate of zero:
\begin{eqnarray}
  \mathrm{Fut}^{e}_{t,T;\theta} &=& \EE^{e,\RN}_t \left[L^e_{T;\theta}\right] .
\label{eq:future}
\end{eqnarray} 
\item Funding spreads. The additive funding spread\footnote{The Libor index corresponds to unsecured lending which means that its value encompasses credit risk. On the other hand, the discount rates correspond to perfectly collateralised lending so the funding spread quantifies the credit risk.} in currency $e$ with collateral $i$ is\footnote{This funding spread is negative to be consistent with conventions used at JPMorgan.}:
\begin{eqnarray}
	S^{e|i,-}_{t,T;\theta} &=& R^{e|i}_{t,T;\theta}  - L^{e|i}_{t,T;\theta} 
	\label{eq:additive_funding_spread}
\end{eqnarray} 
For mathematical reasons, we will mostly deal with the multiplicative funding spread\footnote{This definition and notation can seem a bit strange but they are meant to be consistent with the negative funding spread convention. One could flip the signs: $S^{e|i,+}_{t,T;\theta} = -S^{e|i,-}_{t,T;\theta}$, $S^{e|i,\times}_{t,T;\theta} =  - S^{e|i,\div}_{t,T;\theta}$ and all the results would remain valid. However, they would break if we flipped the ratio in this definition. This is because, while both Discount and Libor forwards rates are martingales under the payment forward measure, there do not have symmetric economic roles. Discount forwards are determined by ratios of ZCBs which are prices of tradeable assets but there is no analog for Libor forwards. While we can always define pseudo-discount factors corresponding Libor forwards, these are not observable and not constrained by AOA. Specifying their dynamic requires additional model assumptions. For example, one could assume that the ratios $Z^{e|e,\mathrm{spread}}_{t,T} := \frac{Z^{e|e}_{t,T}}{Z^{e|e,\mathrm{proj}}_{t,T}}$  of pseudo-discount factors or the instantaneous forward spreads $s^{e|e}_{t,T} = r^{e|e,\mathrm{proj}}_{t,T} - r^{e|e}_{t,T}$ are constant but it is not clear that this is economically or mathematically justified or that this leads to significant mathematical simplifications.} instead as defined by:
\begin{eqnarray}
	1 - \theta S^{e|i,\div}_{t,T;\theta} &=& \frac{1+\theta L^{e|i}_{t,T;\theta}}{1+\theta R^{e|i}_{t,T;\theta}}  
	\label{eq:multiplicative_funding_spread}
\end{eqnarray} 
\item Swap rates. The forward swap rate in currency $e$ with schedule $T_0,\ldots,T_{\f M}$ under collateral $i$ is 
\begin{eqnarray}
  Y^{e|i}_{t,T;M}	&=&  \frac{\sum_{j=1}^{\f M} Z^{e|i}_{t,T_0,T_j} \theta_j L^{e|i}_{t,T_{j-1},T_j}}{\sum_{k=1}^{m} \theta_k Z^{e|i}_{t,T_0,T_k} }
    \label{eq:forward_swap_rate}
\end{eqnarray} 
\item FX forwards. The par strike of an FX forward under collateral $i$ is 
\begin{eqnarray}
  x^{df|i}_{t,T } &=& \EE^{d|i,T}_t \left[x^{df}_{T}\right]  = x^{df}_{t}\frac{Z^{f|i}_{t,T}}{Z^{d|i}_{t,T}}. 
  \label{eq:fx_forward}
\end{eqnarray} 
\item Resetting forward Libor:  $L^{f/d|i}_{t,T;\theta}$ is the par strike $K$ of a contract with collateral $i$ exchanging 
\begin{itemize}
\item at $T$, receive $N_d = x^{df}_T N_f$ in $d$, pay $N_f$ in $f$
\item at $T+\theta$, pay $N_d(1+\theta L^d_{T;\theta})$ in $d$, receive $N_f(1+\theta K)$ in $f$.   
\end{itemize}
So, $L^{f/d|i}_{t,T;\theta}$ is characterised by 
\begin{eqnarray*}
   1 + \theta L^{f/d|i}_{t,T;\theta} 
   &=& \EE^{f|i,T+\theta}_t\left[ \frac{x^{fd}_{T+\theta}}{x^{fd}_{T}} (1+\theta L^{d}_{T;\theta}) \right] 
\end{eqnarray*} 
%As a result, the par spread of a resetting cross-currency swaplet can be reformulated as 
%\begin{eqnarray*}
%   B^{f/d|i,r}_{t,T;\theta}  &=& L^{f/d|i}_{t,T;\theta} - L^{f|i}_{t,T;\theta}
%\end{eqnarray*} 
%while the general swap 
%\begin{eqnarray*}
%   B^{f/d|i,r}_{t,T;M}  &=& \frac{ \sum_i Z^{f|i}_{t,T_i} \theta (L^{f/d|i}_{t,T_{i-1};\theta} - L^{f|i}_{t,T_{i-1};\theta}) } { \sum_i Z^{f|i}_{t,T_i} \theta  }
%\end{eqnarray*} 
\item Currency basis spread. A resetting currency basis swap has
\begin{itemize}
\item foreign leg: pays in $f$, Libor plus spread $L^{f}_k + K$ and exchange of notionals $N^f_k$.
\item domestic leg: pays in $d$, Libor $L^{d}_k$ and exchange of notionals resetting at the beginning of each period to match the corresponding foreign notional $N_k^d = x^{df}_{T_{k-1}}N_{k}^f$. 
\end{itemize}
It follows that, under collateral $i$, the par-spread of each swaplet is
\begin{eqnarray}
B^{f/d|i,r}_{t,T_{k};\theta} &=&  L^{f/d|i,r}_{t,T_{k};\theta}  -  L^{f|i}_{t,T_{k};\theta}  \label{eq:forward_cb_swaplet_resetting}
\end{eqnarray}
while for the whole swap\footnote{For a generic non-resetting cross-currency swap, we can introduce $x^{df}_{t_{\mathrm{ref},i}} = N^d_i/N^f_i$. The previous formula generalises to 
\begin{eqnarray*}
B^{f/d|i,*}_{t,T;M} &=& 
  \frac{
  \sum_{k=1}^{\f M} \left( Z^{d|i}_{t,T_k} \EE^{d|i,T_k}_t\left[ x^{df}_{\mathrm{ref},k} (1+\theta L^{d|i}_{T_{k-1};\theta})\right] - Z^{d|i}_{t,T_{k-1}} \EE^{d|i,T_{k-1}}_t\left[ x^{df}_{\mathrm{ref},k}\right] \right) - 
    \sum_{k=1}^{\f M} \left(  Z^{f|i}_{t,T_k} \left(1+\theta L^{f|i}_{t,T_{k-1},T_k} \right) - Z^{f|i}_{t,T_{k-1}} \right) }
   {\sum_{k=1}^{\f M}  \theta Z^{f|i}_{t,T_k}}
  \label{eq:forward_cb_non_resetting}
\end{eqnarray*}
In the case of constant notionals, i.e. deterministic $x^{df}_{t_{\mathrm{ref},i}} $, the par-spread is: 
\begin{eqnarray*}
B^{f/d|i,nr}_{t,T;M} &=& 
  \frac{
    \sum_{k=1}^{\f M} \frac{x^{df}_{\mathrm{ref},k}}{x^{df}_t} \left( Z^{d|i}_{t,T_k} (1+\theta L^{d|i}_{T_{k-1};\theta}) - Z^{d|i}_{t,T_{k-1}} \right) - 
    \sum_{k=1}^{\f M} \left(  Z^{f|i}_{t,T_k} \left(1+\theta L^{f|i}_{t,T_{k-1},T_k} \right) - Z^{f|i}_{t,T_{k-1}} \right) }
   {\sum_{k=1}^{\f M}  \theta Z^{f|i}_{t,T_k}}
  \label{eq:forward_cb_non_resetting}
\end{eqnarray*}} 
it is 
\begin{eqnarray}
B^{f/d|i,r}_{t,T;M} &=& 
  \frac{\sum_{k=1}^{\f M} Z^{d|i}_{t,T_k} \theta B^{f/d|i,r}_{t,T_{k-1};\theta}}
   {\sum_{k=1}^{\f M}  Z^{f|i}_{t,T_k}  \theta}
  \label{eq:forward_cb_resetting}
\end{eqnarray}
%The currency basis spread depends on the FX-forward curve. 
%\item Non-resetting currency basis spread. A non-resetting currency basis swap is a swap consisting of two floating legs with notional exchanges in different currencies. Assuming the    \\
%Resetting cross-currency swap: the notional of the $d$ leg fixes at the beginning of each period to match the notional of the non-resetting $f$ leg. The spread is on the non-resetting leg. It follows that, under collateral $i$, the par-spread is 
%\begin{eqnarray}
%B^{f/d|i,r}_{t,T;M} &=& 
%  \frac{
%   \sum_{k=1}^{\f M} Z^{d|i}_{t,T_k} \theta \left(  L^{f/d|i,r}_{t,T_{k-1};\theta}  -  L^{f|i}_{t,T_{k-1};\theta} \right) }
%   {\sum_{k=1}^{\f M}  \theta Z^{f|i}_{t,T_k}}
%  \label{eq:forward_cb_resetting}
%\end{eqnarray}
%where the resetting forward Libor is defined by
%\begin{eqnarray*}
%   1 + \theta L^{f/d|i,r}_{t,T;\theta} 
%   &=& \EE^{f|i,T+\theta}_t\left[ \frac{x^{fd}_{T+\theta}}{x^{fd}_{T}} (1+\theta L^{d}_{T;\theta}) \right] 
%   \label{eq:resetting_forward_libor}
%\end{eqnarray*} 
\end{itemize}

[TODO: Add options, i.e. Swaptions, Caps, FX options and currency basis swaptions].

%
%\begin{rema}[Notations]
%A few rules to understand the notations:
%\begin{itemize}
%\item The first letter in the exponent is the native payment currency: 
%\item A letter after a vertical bar denotes the collateral: $Z^{d|c}_{t,T}$ has collateral $c$ 
%\item the first subscript is always the valuation time of the variable (typically $t$ but sometimes $u$ or $v$ when the variable is integrated over time)
%\item the second subscript denotes the expiry: $Z^{d|c}_{t,T}$, $R^{d|c}_{t,T,U}$,  $Y^{d|c}_{t,T;M}$ correspond to expiry $T$.
%\item the third subscript correspond to the tenor if preceded by a semi-colon and to the maturity if preceded by a colon. 
%\item $Z$ denotes pseudo-discount factors
%\item $R$ denotes pseudo-discount rates
%\item $L$ denotes linear (Ibor) rates
%\item $Y$ denotes swap rates
%\item $S$ denotes funding spreads
%\item $B$ denotes currency basis spreads
%\item $x$ denotes FX rates
%\end{itemize}
%\end{rema}

%
%\subsection{Liquid instruments}%
%
%\subsubsection{Forward Rate Agreements}%
%
%
%\subsubsection{Interest Rate Swaps}%
%
%
%\subsubsection{FX Forward Rate Agreements}%
%
%
%\subsubsection{Resetting cross-currency swap}%
%
%
%\subsubsection{Interest rate swaptions}%
%
%
%\subsubsection{FX options}%









%%%%%%%%%%%%%%%%%%%%%%%%%%%%%%%%%%%%%%%%%%%%
%% Market models
%%%%%%%%%%%%%%%%%%%%%%%%%%%%%%%%%%%%%%%%%%%%
\newpage%
\section{Market models}%

%\subsection{Market models}%
%\begin{defi}
%A \emph{market model} for collateralized fixed cashflows is the following data:
%\begin{enumerate}
%\item a set of currencies $\fr{X}$ 
%\item the following set of stochastic processes\footnote{Defined on a filtered probability space $(\Omega,\mathcal{A},\mathcal{F})$}:
%\begin{itemize}
%\item for each currency $c\in \fr{X}$, a money market collateral account $N^{c|c}_t$ earning collateral rate $r^{c|c}_t$.
%\item for each pair of currencies $(d,f)$, an FX spot rate $x^{df}_t$.
%\item for each pair of currencies $x,c$ and each maturity $T$, a price process $Z^{x|c}_{t,T}$ for the value in currency $x$ of the contract paying one unit of $x$ at $T$ under collateral $c$.
%\end{itemize}
%\end{enumerate}
%\end{defi}

\subsection{Definition}

\begin{defi}
A \emph{market model} consists of the following data:
\begin{enumerate}
\item A set of assets $S^i$ 
\item A set of parameters $\theta \in \Theta$
\item A stochastic process $(X^\theta_t)_{t\in [0,T],\theta \in \Theta}$ 
\item Reconstruction formulas for asset prices, i.e. deterministic functions s.t. $S^i_t = S^i(\theta,t,X^\theta_t)$. 
\end{enumerate}
\end{defi}

The model provides a dynamic for the value of any portfolios $\Pi_t = \sum_i w_{i,t} S^i_t$. 


\subsubsection*{Calibration}
The process of calibrating of the model corresponds to finding $\theta^* \in \Theta$ such that the model delivered prices of benchmark instruments match the market observed ones. 


\subsection{Fixed cash-flows}

Our focus is on fully collateralised assets. The most basic ones are fixed cash-flows with cash collateral in a single currency. 

A market model for collateralised fixed cash-flows consists of the following data:
\begin{enumerate}
\item a set of currencies $\fr{X}$ 
\item the following set of stochastic processes\footnote{Defined on a filtered probability space $(\Omega,\mathcal{A},\mathcal{F})$}:
\begin{itemize}
\item for each currency $c\in \fr{X}$, a reference native collateral rate $r^{c|c}_t$.
\item for each pair of currencies $(d,f)$, an FX spot rate $x^{df}_t$.
\item for each pair of currencies $x,c$ and each maturity $T$, a price process $Z^{x|c}_{t,T}$ for the value in currency $x$ of the contract paying one unit of $x$ at $T$ under collateral $c$.
\end{itemize}
\end{enumerate}


We will require that the model satisfies the following properties:
\begin{enumerate}
\item Absence of arbitrage
\item Calibration, i.e. the model reproduces prices of liquid market instruments at time $t = 0$. 
\item Regularity: discount curves $T\mapsto Z^{x|c}_{t,T}$ are derivable\footnote{They need to be at least absolutely continuous for instantaneous forward rates to exist.} to allow for risk management through computation of sensitivities.
\end{enumerate}


\subsection{Floating rates instruments}
Liquid long dated instruments are vanilla swaps, cross-currency swaps, tenor basis swaps etc\ldots They all involve floating cash-flows associated with IR indices. So, in order to calibrate our market model, we will have to extend its scope to include these indices. There are two options for that:
\begin{enumerate}
\item Assume that funding and tenor basis spreads are constant\footnote{Or deterministic or at least independent of all other processes.}.
\item Develop a model with stochastic funding and tenor basis spreads.
\end{enumerate} 

The second option adds a lot of complexity. Obviously the increased number of stochastic factors is one but more fundamentally, one can prove that the drift of instantaneous forward projection rates are not uniquely determined by AOA in a model with stochastic funding basis and there is no natural choice. Models of LMM type, on the other hand, do not have this problem because the dynamic of the market observable projection rates is modelled directly. This suggests that the right models would be of LMM type for the projection rate and funding spread but with a short discount rate\footnote{It is not surprising that discount rates and index forwards require different modelling since they are fundamentally different.} which makes them much more complex and high-dimensional. 


We will focus on the first option instead. 
\begin{hypo}[Constant funding spreads]
For every currency $e$ and tenor $\theta$, the native multiplicative funding spread
\begin{eqnarray}
  1 - \theta S^{e|e,\div}_{t,T;\theta} &=& \frac{1+\theta L^{e|e}_{t,T;\theta}}{1+\theta R^{e|e}_{t,T;\theta}}  
  \label{eq:constant_funding_spread}
\end{eqnarray} 
is constant. 
\end{hypo}
From now on, we will assume the constant funding spreads hypothesis holds.

%%%%%%%%%%%%%%%%%%%%%%%%%%%%%%%%%%%%%%%%%%%%
%% Deterministic model
%%%%%%%%%%%%%%%%%%%%%%%%%%%%%%%%%%%%%%%%%%%%
\newpage%
\section{Market practice: The Deterministic Collateral Basis Model}%

\setcounter{subsection}{-1}
\subsection{Introduction}%
In this section, we detail the assumptions currently used by most market participants to value linear products. 
We then explain how they allow them to build discount curves and projection curves to value any linear product under any single currency cash collateral. 



\subsection{FX forward invariance}%

\begin{hypo}[Constant collateral basis]
All collateral basis spreads $q^{e|i\to j}_{t,T}$ are constant:
\begin{eqnarray}
  \frac{Z^{e|i}_{t,T}}{Z^{e|j}_{t,T}} &=& \frac{Z^{e|i}_{0,T}}{Z^{e|j}_{0,T}} 
  \label{hypo:constant_collateral_basis}
\end{eqnarray}
\end{hypo}


This is equivalent to the following.
\begin{hypo}[FX forward invariance]
FX forward rates do not depend on the collateral: 
\begin{eqnarray}
  x^{df|i}_{t,T} &=& x^{df|j}_{t,T} 
  \label{eq:fx_forward_invariance}
\end{eqnarray}
\end{hypo}

\begin{rema}
In conjunction with the constant funding spread assumption, this implies that all spreads constant. Note however that these assumption only involve forward rates. We do not assume anything about the initial discount curves $Z^{f|d}_{0,T}$. In fact, we assume that the impact of the change of collateral on FX forwards an FRAs is entirely accounted for by the change of discount factor.  

One could weakened the assumption a bit: our results would still hold if we assumed that all the basis spreads are independent of all the Libor and FX rates instead. 
\end{rema}

\subsection{Valuation of resetting cross-currency swaps}%

\subsubsection{Discount curves under reference collateral}
The first problem to define a market model is that for each new cross-currency swap expiry, we have:
\begin{enumerate}
\item One new constraint: the master cross-currency equation (\ref{eq:forward_cb_resetting}).
\item Three new unknowns\footnote{This is assuming that we can observe the currency basis spread $B^{f/d|d,r}_{t,T_{k};\theta}$ for each swaplet which is not case. In practice, we only observe spot starting currency basis spreads for a given set of expiries. So some interpolation procedure is always required at the time of the bootstrapping.}: $Z^{f|d}_{t,T_{k}}$ and $L^{f|d}_{t,T_{k};\theta}$ and $L^{f/d|d,r}_{t,T_{k};\theta}$. 
\end{enumerate}
This system of equations is clearly under-determined so we need to impose some constraints and/or include additional parameters.


\subsubsection{Deterministic collateral basis model}
The standard market practice is to assume that all collateral basis spreads are constant. This ensures that all collateral convexity adjustments in the valuation of cross-currency swaps can be neglected:
\begin{enumerate}
\item The convexity adjustment due to non-standard collateral for foreign Libors: $L^{f|d,\DCB}_{t,T_{i-1},T_i} = L^{f|f}_{t,T_{i-1},T_i}$.
\item The convexity adjustment due to the resetting notional $N^d_k = x^{df|d}_{t,T_{k-1}}N^f_k$ captured by $L^{f/d|d}_{t,T_{k-1};\theta}$. 
\end{enumerate}


\subsubsection{Main theorem}%

\begin{theo}[DCB]
Under the above assumptions, 
\begin{eqnarray}
  R^{f|d}_{t,T;\theta} &=& L^{f|f}_{t,T;\theta} + B^{f/d|d,r}_{t,T;\theta}  + S^{d|d,-}_{t,T;\theta}  .
  \label{eq:differential_discounting}
\end{eqnarray} 
\end{theo}

\begin{rema}
This formula is quite pleasing as it has an intuitive hedging interpretation. To hedge fixed cash-flows in foreign currency with domestic collateral, one should enter a foreign vanilla swap to swap the fixed foreign cash-flows into floating ones, a currency basis swap and a domestic funding basis swap\footnote{Can we make the replication argument more precise? If I enter, 
\begin{itemize}
\item a foreign vanilla swap to swap the fixed foreign cash-flows into floating ones
\item a currency basis swap to swap the foreign Ibor cash-flows into domestic ones 
\item a domestic funding basis swap to swap the domestic Ibor cash-flows into domestic O/N ones. 
\end{itemize}
How should I set the notionals? What is the residual risk? Am I really left with a foreign fixed cash flow vs domestic OIS? Is this true at first order? 
}
\end{rema}



\subsubsection{Cross-currency curve stripping\label{ssec:dd_curve_stripping}}%
Under these simplifying assumptions, the master equation for the cross-currency discount curve is 
\begin{eqnarray}
  & & \sum_{k=1}^{\f M} \left(  Z^{f|d}_{t,T_k} \left(1+\theta (L^{f|d}_{t,T_{k-1},T_k} + B^{f/d|d,r}_{t,T;M}) \right) - Z^{f|d}_{t,T_{k-1}} \right) \nonumber \\
  &=& 
  \sum_{k=1}^{\f M} x^{df|d}_{t,T_{k-1}} \left( Z^{d|d}_{t,T_k} (1+\theta L^{d|d}_{T_{k-1};\theta}) - Z^{d|d}_{t,T_{k-1}} \right) 
  \label{eq:dcb_xccy_master}
\end{eqnarray}
%\begin{eqnarray*}
%B^{f/d|i,nr}_{t,T;M} &=& 
%  \frac{
%    \sum_{k=1}^{\f M} \frac{x^{df}_{\mathrm{ref},k}}{x^{df}_t} \left( Z^{d|i}_{t,T_k} (1+\theta L^{d|i}_{T_{k-1};\theta}) - Z^{d|i}_{t,T_{k-1}} \right) - 
%    \sum_{k=1}^{\f M} \left(  Z^{f|i}_{t,T_k} \left(1+\theta L^{f|i}_{t,T_{k-1},T_k} \right) - Z^{f|i}_{t,T_{k-1}} \right) }
%   {\sum_{k=1}^{\f M}  \theta Z^{f|i}_{t,T_k}}
%  \label{eq:forward_cb_non_resetting}
%\end{eqnarray*}
In order to bootstrap the cross-currency discount curve, one writes this equation for two consecutive tenors $M$ and $N$, and take the difference. 

The problem is still under-determined since we only have one equation but each of the $Z^{f|d}_{t,T_k}$ for $T_k \in [T_{\f M},T_{\f N}[$ is unknown. 
 
A solution is to specify an interpolation scheme for the discount curve on the relevant segment. Typical choices are:
\begin{itemize}
\item Flat instantaneous forward rates: $r^{f|d}_{0,T}$ is constant for $T \in [T_{\f M},T_{\f N}[$. 
\item Parabolic instantaneous forward rates: $r^{f|d}_{0,T}$ is a second order polynomial on $T \in [T_{\f M},T_{\f N}]$ and it is derivable at $T = \f M$.  
\end{itemize}
The main point is that such a parametrisation reduces the knowledge of the discount factors in the relevant period to a single unknown:
\begin{eqnarray}
  \overline{r}^{f|d}_{0,[T_{\f M},T_{\f N}]} &=& \frac{1}{T_{\f M} - T_{\f N}} \int_{T_{\f M}}^{T_{\f N}} r^{f|d}_{0,u} \, du
\end{eqnarray}
which is the average level of the instantaneous forward rate on this period. 

\begin{rema}
The curve bootstrapping has some limitations. For example, it assumes that the schedules of the instruments are aligned. An alternative solution is to consider a global calibration. This allows the user to include more diversity in the calibration targets, and more granularity in the pillars. It is however more complex and it does not guarantee that risk is local: an instrument with maturity $T$ can have exposure to benchmark instruments with arbitrary maturity.  
\end{rema}

\subsection{Off-market discount curves}%


\subsubsection{Discount curves under general collateral}
The second problem is that cross-currency swaps are only liquid against USD. So, there are no quotes to build a discount curve with non-native non-USD collateral such as $Z^{\mathrm{USD}|\mathrm{EUR}}$ or $Z^{\mathrm{EUR}|\mathrm{JPY}}$. 


\subsubsection{Off-market USD discount curve}

\begin{prop}[Off-market discount curve]
Under the FX-forward invariance assumption, the off-market discount curves $Z^{u|i}_{t,T}$ is characterised by:
\begin{eqnarray}
  Z^{u|i}_{t,T} &=& Z^{u|u}_{t,T} \frac{Z^{i|i}_{t,T}}{Z^{i|u}_{t,T}} 
  \label{eq:usd_off_market_discount_curve}
\end{eqnarray}
\end{prop}
This provides a way to build the curve.



\subsubsection{Off-market non-USD discount curves}


\begin{prop}[Off-market discount curve]
Under the FX-forward invariance assumption, the off-market discount curves $Z^{e|i}_{t,T}$ is characterised by:
\begin{eqnarray}
  Z^{e|i}_{t,T} &=& Z^{e|u}_{t,T} \frac{Z^{i|i}_{t,T}}{Z^{i|u}_{t,T}} 
  \label{eq:off_off_market_discount_curve}
\end{eqnarray}
\end{prop}




%%%%%%%%%%%%%%%%%%%%%%%%%%%%%%%%%%%%%%%%%%%%
%% Brownian HJM model
%%%%%%%%%%%%%%%%%%%%%%%%%%%%%%%%%%%%%%%%%%%%
\newpage%
\section{Arbitrage free models with stochastic collateral}


\setcounter{subsection}{-1}
\subsection{Introduction}%
The current market practice described in the previous section shows how to build a market model with \emph{deterministic basis}. But the construction relies on assumptions that are inconsistent with any stochastic spreads between discount rates. As a result, the corresponding curves cannot be used as inputs to models with stochastic spreads without introducing arbitrage. 

Our first objective in this section is to identify the arbitrage constraints on a stochastic multi-currency market model for collateralized fixed cashflows. 

%Under a deterministic tenor basis spread assumptions, we then proceed to derive the dynamic of currency %basis spreads in this class of models.  

%\begin{itemize}
%\item Define a market model for collateralized fixed cashflows with stochastic collateral. 
%\item Compute the mean and covariance of market observables in the model
%\item Provide a calibration procedure for the model.
%\item Provide analytic solution for the discount curves.  
%\end{itemize}

We assume that tenor basis (aka funding) spreads are deterministic (cf.  \autoref{hypo:constant_collateral_basis}) so the stochasticity of collateral is essentially equivalent to the stochasticity of the currency basis. 



\subsection{Model specification}%


\begin{theo}
Let $\{ Z^{x|c}_{t,T} \}_{t\le T, x,c\in\fr{X}}$ be an arbitrage-free market model for fixed collateralized cash-flows such that $T \mapsto Z^{x|c}_{t,T}$ are absolutely continuous. 
Then, given a reference currency $d \in \fr{X}$, the model $\{ Z^{x|c}_{t,T} \}_{t\le T, x,c\in\fr{X}}$ is uniquely specified by the following data:  for each currency $f \in \fr{X}$,
\begin{enumerate}
\item Initial data for
\begin{itemize}
\item native discount curves $\{ Z^{f|f}_{t_0,T} \}_{f\in\fr{X}}$.
\item collateral basis curves $\{ \QZ^{f|f\to d}_{t_0,T} \}_{f\in\fr{X}}$ or equivalently cross-currency discount curves $\{ Z^{f|d}_{t_0,T} \}_{f\in\fr{X}}$.
%\item collateral basis curves $\{ \QZ^{d|d\to f}_{t_0,T} \}_{c\in\fr{X}}$ or equivalently $\{ Z^{d|f}_{t_0,T} \}_{c\in\fr{X}}$.
\item FX spot rates $x^{df}_{t_0}$. 
\end{itemize}
\item Spot volatility processes for 
\begin{itemize}
\item instantaneous forward discount rates denoted by $\{ \sigma^{f|f}_{\cdot,T} \}$ 
\item instantaneous forward collateral basis denoted by$\{ \sigma^{f\to d}_{\cdot,T} \}$ (and equal to 0 if $f=d$)
\item FX spot rates denoted by $\sigma^{df}_{\cdot}$ (and equal to 0 if $f=d$)
\end{itemize}
\end{enumerate}
\end{theo}

\begin{rema}
The theorem essentially says that arbitrage conditions imply two things: 
\begin{enumerate}
\item (Triangle constraints) the dynamic can be decomposed\footnote{In mathematical terms, the market is a graph whose nodes correspond to currencies. Moreover this graph is complete. Once we choose a reference currency, there is an obvious spanning tree whose root is that reference currency. This is a very pedantic way to state things but it is interesting to remember when we consider the implementation.} with respect to a reference currency $d$. 
\item (Drift constraints) the risk-neutral drifts of the assets are entirely specified by their spot volatilities\footnote{This is a generalisation of the standard HJM result.}. 
\end{enumerate}

More explicitly, let us write, 
\begin{eqnarray}
  \Sigma^{\alpha}_{t,T} &=& \int_t^{T} \sigma^{\alpha}_{t,u}\,du \\
  \Sigma^{\alpha}_{t,T;\theta} &=& \int_T^{T+\theta} \sigma^{\alpha}_{t,u}\,du 
\end{eqnarray} 
Then, we have
\begin{eqnarray}
  dr^{e|e}_{t,T} &=& \left(\sigma^{e|e}_{t,T}\cdot \Sigma^{e|e}_{t,T} \right) dt +  \sigma^{e|e}_{t,T}\cdot \left(dW^{d,\RN}_t -\sigma^{de}_t \,dt\right)  \\
  dq^{f|f\to d}_{t,T} &=& \left(\sigma^{f\to d}_{t,T}\cdot \Sigma^{f\to d}_{t,T} + \sigma^{f\to d}_{t,T}\cdot \Sigma^{f|f}_{t,T} + \sigma^{f|f}_{t,T}\cdot \Sigma^{f\to d}_{t,T}\right) dt +  \sigma^{f\to d}_{t,T} \cdot \left( dW^{d,\RN}_t  - \sigma^{df}_t \, dt\right)
%  dq^{d|d\to i}_{t,T} &=& \left(\sigma^{d\to i}_{t,T}\cdot \Sigma^{d\to i}_{t,T} + \sigma^{d\to i}_{t,T}\cdot \Sigma^{d|d}_{t,T} + \sigma^{d|d}_{t,T}\cdot \Sigma^{d\to i}_{t,T}\right) dt +  \sigma^{d\to i}_{t,T} \cdot dW^{d,\RN}_t  
\end{eqnarray}
We define the short rate and short collateral basis in the obvious way\footnote{Note that instantaneous forward collateral basis does depend on the payment currency even if the short collateral basis does not! This is because the forward accounts for covariance between payment currency and collateral basis across time. This covariance tends to zero as the time interval goes to zero.}:
\begin{eqnarray}
  r^{e|e}_{t} &=& r^{e|e}_{t,t}  \\
  q^{e\to d}_{t} &=& q^{e|e\to d}_{t,t}
\end{eqnarray}
Then, the full dynamic is specified by
\begin{eqnarray}
  \frac{dZ^{d|i}_{t,T}}{Z^{d|i}_{t,T}} &=& \left(r^{d|d}_t + q^{d\to i}_t \right) dt - \left(\Sigma^{d|d}_{t,T} + \Sigma^{d\to i}_{t,T}  \right) dW^{d,\RN}_t  \\
   \frac{dx^{df}_{t}}{x^{df}_{t}} &=& \left(r^{d|d}_t - r^{f|f}_t + q^{d\to f}_t \right) dt + \sigma^{df}_{t} \cdot dW^{d,\RN}_t  \end{eqnarray}
for domestic observables and 
\begin{eqnarray}
  \frac{dZ^{e|i}_{t,T}}{Z^{e|i}_{t,T}} 
  &=&  \left(r^{e|e}_t + q^{d\to i}_t-q^{d\to e}_t \right)  dt 
  - \left(\Sigma^{e|e}_{t,T} + \Sigma^{d\to i}_{t,T} -  \Sigma^{d\to e}_{t,T}  \right) \cdot \left( dW^{d,\RN}_t -\sigma^{de}_tdt\right) \\
   \frac{dx^{ef}_{t}}{x^{ef}_{t}} 
   &=& \left(r^{e|e}_t - r^{f|f}_t + q^{d\to f}_t-q^{d\to e}_t \right) dt 
   + \left(\sigma^{df}_{t} - \sigma^{de}_{t}\right)  \cdot \left(dW^{d,\RN}_t - \sigma^{de}_{t}\,dt \right)  
\end{eqnarray}
for foreign observables. 
\end{rema}

\begin{rema}
For short, we will  write 
\begin{eqnarray}
  \sigma^{d|i}_{t,T} &=& \sigma^{d|d}_{t,T}  + \sigma^{d\to i}_{t,T} \\
  \Sigma^{d|i}_{t,T} &=& \Sigma^{d|d}_{t,T}  + \Sigma^{d\to i}_{t,T} 
\end{eqnarray} 
and use the same convention for the integrated vols. We will also write 
\begin{eqnarray}
  \Sigma^{j\to i}_{t,T} &=& - \Sigma^{i\to j}_{t,T} \\
  \Sigma^{j\to i}_{t,T;\theta} &=& - \Sigma^{i\to j}_{t,T;\theta}  
\end{eqnarray} 
so that 
\begin{eqnarray}
	\Sigma^{d|j}_{t,T;\theta} &=& \Sigma^{d|i}_{t,T;\theta}  + \Sigma^{i\to j}_{t,T;\theta}  
\end{eqnarray} 
\end{rema}


%%%%%%%%%%%%%%%%%%%%%%%%%%%%%%%%%%%%%%%%%%%%
%% Quanto convexity adjustments
%%%%%%%%%%%%%%%%%%%%%%%%%%%%%%%%%%%%%%%%%%%%

\subsection{Deterministic collateral basis model}

We want to show that if the collateral basis volatilities $\sigma^{d\to i}_{t,T}$ are zero, we recover the DCB model.   

This is straightforward since discount curve become deterministic. 

%%%%%%%%%%%%%%%%%%%%%%%%%%%%%%%%%%%%%%%%%%%%
%% Quanto convexity adjustments
%%%%%%%%%%%%%%%%%%%%%%%%%%%%%%%%%%%%%%%%%%%%

\subsection{Quanto adjustments}

Quanto adjustments are convexity adjustments associated with the payment of an index in a non-native currency.  


\subsubsection{Change of currency formula}%
 
%Now we consider explicit computations of these convexity adjustments in the HJM framework.
 
\begin{prop}
The Radon-Nikodym derivative of the change of $T$-forward measure is characterised by the exponential martingale 
\begin{eqnarray}
  \Q^{df|i}_{t,T}  := \EE^{d|i,T}_t\left[ \frac{d\QQ^{f|i,T}}{d\QQ^{d|i,T}} \right] 
  = \mathcal{E}\left( - \int_0^\cdot \Sigma^{df}_{u,T} \cdot dW^{d|i,T}_u \right)_t
  \label{eq:quanto_radon_nikodym}
\end{eqnarray}
where 
\begin{eqnarray}
  \Sigma^{df}_{u,T} 
  = \Sigma^{d|d}_{u,T} - \Sigma^{f|f}_{u,T} - \Sigma^{f\to d}_{u,T}  + \sigma^{df}_{u} 
\end{eqnarray}
is the volatility of the FX-forward.
\end{prop}

\subsubsection{Quanto adjustments}%
 

\begin{prop}

\end{prop}




%\begin{defi}
%For any payment currency $e$, collateral $i,j$ and maturity $T$, we define the corresponding \emph{discount factor  collateral convexity adjustment} as 
%\begin{eqnarray}
%  \Qto^{d\to f|i}_{t,T}[v^e_T]   &=& \frac{ \EE^{f|i}_{t,T}[v^d_T]  } { \EE^{e|i}_{t,T}[v^d_T]  } 
%\end{eqnarray}
%\end{defi}







\newpage%
\subsection{Dynamic of market observables}
We now derive the dynamic of market observables in order to calibrate the model.



\subsubsection{Discount factors}

\begin{prop}
The ZCBs have dynamic:
\begin{eqnarray}
  \frac{dZ^{e|i}_{t,T}}{Z^{e|i}_{t,T}} 
  &=&  \left(r^{e|e}_t + q^{e\to i}_t\right)  dt 
  - \Sigma^{e|i}_{t,T} \cdot \left( dW^{d,\RN}_t -\sigma^{de}_t\, dt\right) 
\end{eqnarray}
\end{prop}

\subsubsection{Forward discount rates}

\begin{prop}
The forward discount rate $R^{e|i}_{t,T;\theta}$ is a shifted exponential martingale under is a $\QQ^{e|i,T+\theta}$:
\begin{eqnarray}
	\frac{d (1+\theta R^{e|i}_{t,T;\theta})}{1+\theta R^{e|i}_{t,T;\theta}}  
	&=&   \Sigma^{e|i}_{t,T;\theta} \cdot dW^{e|i,T+\theta}_t .
\end{eqnarray} 
\end{prop}


\subsubsection{FX forwards}

\begin{prop}
The FX forward $x^{df|i}_{t,T}$ is an exponential martingale under $\QQ^{d|i,T}$:  
\begin{eqnarray}
  \frac{dx^{df|i}_{t,T}}{x^{df|i}_{t,T}} &=& \Sigma^{df}_{t,T} \cdot dW^{d|i,T}_{t} 
\end{eqnarray}
where 
\begin{eqnarray}
  \Sigma^{df}_{t,T} &=& \Sigma^{d|d}_{t,T} - \Sigma^{f|f}_{t,T} +\Sigma^{d\to f}_{t,T}  + \sigma^{df}_t  .
\end{eqnarray}
\end{prop}

\begin{rema}
Note that, a priori, the collateral does not explicitly appear in the volatility process of the FX-forward. The collateral only appears in the drift or equivalently in the initial value of the FX-forward when the model is calibrated to market quotes.  

This is a pleasant surprise. Indeed, a priori, we should have expected one implied volatility surface for each FX pair and each collateral. And it is indeed the case if we consider the volatility smile as a function of absolute strikes. However, if we consider our volatility surface as a function $(\tau,y)\mapsto \sigma^{df}_{\BS,t}(\tau,y)$ of time to expiry $\tau = T-t$ and log-moneyness forward $y = \ln (x^{df}_t / x^{df|i}_{t,t+\tau})$, then the model suggests that $\sigma^{df}_{\BS,t}$ should be independent of the collateral $i$ unless the volatility of the FX-forward is stochastic and correlated with the collateral. 

One can reformulate this as saying that collateral convexity impacts the intrinsic value of an FX option but not its time value, at zeroth order in vol of vol\footnote{Q: Can we get some estimate in the LV model where instantaneous correlation of the FX vol with the collateral basis is equal to that of the FX spot? }.  
\end{rema}


\subsubsection{Forward Libors}

\begin{prop}
The forward Libor $L^{e|i}_{t,T;\theta}$ is a shifted exponential martingale under $\QQ^{f|i,T+\theta}$: 
\begin{eqnarray}
  \frac{d (1+\theta L^{e|i}_{t,T;\theta})}{1+\theta L^{e|i}_{t,T;\theta}}  
  &=&   \Sigma^{e|e}_{t,T;\theta}  \cdot  dW^{e|i,T+\theta}_t  
\end{eqnarray} 
\end{prop}

\begin{rema}
Similarly to FX-forwards, this suggests that the Libor volatility is independent of the collateral, we can ignore collateral convexity on implied volatilities.
\end{rema}




\subsubsection{Funding basis spreads}

\begin{prop}
The multiplicative funding spread $S^{e|i,\div}_{t,T;\theta}$ is a shifted exponential martingale\footnote{Note that this is the $T$-forward measure and not the  he $(T+\theta)$-forward measure like $1+\theta L^{e|i,\div}_{t,T;\theta}$ are $1+\theta R^{e|i,\div}_{t,T;\theta}$. Also note that no such result holds for the additive funding spread (although they should be very close in practice).} under $\QQ^{e|i,T}$: 
%\begin{eqnarray}
%  1 - \theta S^{e|i,\div}_{t,T;\theta}
%  &=&  \frac{1+\theta L^{e|i,\div}_{t,T;\theta}}{1+\theta R^{e|i,\div}_{t,T;\theta}} 
%\end{eqnarray} 
\begin{eqnarray}
  \frac{d(1 - \theta S^{e|i,\div}_{t,T;\theta})}{1 - \theta S^{e|i,\div}_{t,T;\theta}} 	
  &=& -\Sigma^{e\to i}_{t,T;\theta} \cdot dW^{e|i,T}_t .
\end{eqnarray} 
\end{prop}


\begin{proof}
By definition, 
\begin{eqnarray}
  1 - \theta S^{e|i,\div}_{t,T;\theta}
  ~=~  \frac{1+\theta L^{e|i,\div}_{t,T;\theta}}{1+\theta R^{e|i,\div}_{t,T;\theta}}  
  ~=~  \frac{Z^{e|i}_{t,T+\theta}(1+\theta L^{e|i,\div}_{t,T;\theta})}{Z^{e|i}_{t,T}} 
  ~=~  \frac{\mathrm{FRA}^{e|i}_{t,T;\theta}}{Z^{e|i}_{t,T}}  .
\end{eqnarray} 
This is the ratio of the PV of a tradeable asset by the num\'eraire $Z^{e|i}_{t,T}$ so it is a martingale under the corresponding measure $\QQ^{e|i,T}$. It remains to compute its volatility. Ignoring finite variation terms, we have
\begin{eqnarray*}
  \frac{d(1 - \theta S^{e|i,\div}_{t,T;\theta})}{1 - \theta S^{e|i,\div}_{t,T;\theta}} 	
  &=& \frac{d(1 + \theta L^{e|i}_{t,T;\theta})}{1 + \theta L^{e|i}_{t,T;\theta}}  - \frac{d(1 + \theta R^{e|i}_{t,T;\theta})}{1 + \theta R^{e|i}_{t,T;\theta}}  + O(dt) \\
  &=& \Sigma^{e|e}_{t;T;\theta} \cdot dW^{e|i,T+\theta}_t  - \Sigma^{e|i}_{t;T;\theta} \cdot dW^{e|i,T+\theta}_t + O(dt) \\
  &=& - \Sigma^{e \to i}_{t;T;\theta} \cdot dW^{e|i,T+\theta}_t  + O(dt) \\
  &=& - \Sigma^{e \to i}_{t;T;\theta} \cdot dW^{e|i,T}_t 
\end{eqnarray*} 
\end{proof}

\begin{rema}
Note that even if $1 - \theta S^{d|d,\div}_{t,T;\theta}$ and $1 - \theta S^{f|f,\div}_{t,T;\theta}$ are constant,  $1 - \theta S^{f|d,\div}_{t,T;\theta}$ is not, unless the collateral basis vol is zero.
\end{rema}




\subsubsection{Resetting forward Libors}

%\begin{prop}
%The resetting forward Libor $L^{f/d|i}_{t,T;\theta} $ is a $\QQ^{f|i,T+\theta}$-martingale:
%\begin{eqnarray}
%   1 + \theta L^{f/d|i}_{t,T;\theta} 
%   &=& \EE^{f|i,T+\theta}_t\left[ \frac{x^{fd}_{T+\theta}}{x^{fd}_{T}} (1+\theta L^{d}_{T;\theta}) \right] .
%\end{eqnarray} 
%\end{prop}


\begin{prop}
The resetting forward Libor $L^{f/d|i}_{t,T;\theta} $ is a $\QQ^{f|i,T+\theta}$-martingale:
\begin{eqnarray}
   1 + \theta L^{f/d|i}_{t,T;\theta} 
   ~~=~~ \EE^{f|i,T+\theta}_t\left[ \frac{x^{fd}_{T+\theta}}{x^{fd}_{T}} (1+\theta L^{d}_{T;\theta}) \right]  %\\
   ~~=~~  (1+\theta R^{f|i}_{t,T;\theta})  \EE^{f|i,T}_t\left[ 1-\theta S^{d|i,\div}_{T;\theta} \right]  
\end{eqnarray} 
In particular, if native funding spreads are constant, 
\begin{eqnarray}
   1 + \theta L^{f/d|d}_{t,T;\theta} 
   &=&  (1+\theta R^{f|d}_{t,T;\theta}) ( 1-\theta S^{d|d,\div}_{t,T;\theta} )
\end{eqnarray} 
and 
\begin{eqnarray}
   \frac{d(1 + \theta L^{f/d|d}_{t,T;\theta})}{1 + \theta L^{f/d|d}_{t,T;\theta} }
   &=&  \Sigma^{f|d}_{t,T;\theta}\cdot  dW^{f|d,T+\theta}
\end{eqnarray} 
\end{prop}


\begin{proof}
The martingale property is clear from the definition. The more interesting part is the relation to the discount rate. We can rewrite the definition in terms in terms of domestic $T$-martingales: 
\begin{eqnarray*}
   1 + \theta L^{f/d|i}_{t,T;\theta} 
   &=& \EE^{f|i,T+\theta}_t\left[ \frac{x^{fd}_{T+\theta}}{x^{fd}_{T}} (1+\theta L^{d}_{T;\theta}) \right] \\
   &=& x^{fd|i}_{t,T+\theta} \EE^{d|i,T+\theta}_t\left[ x^{df}_{T} (1+\theta L^{d}_{T;\theta}) \right]  \\
   &=& x^{fd|i}_{t,T+\theta} \EE^{d|i,T}_t\left[\frac{Z^{d|i}_{T,T+\theta}}{Z^{d|i}_{t,T,T+\theta}} x^{df}_{T} (1+\theta L^{d}_{T;\theta}) \right]  \\
   &=&  \frac{x^{fd|i}_{t,T+\theta}}{Z^{d|i}_{t,T,T+\theta}} \EE^{d|i,T}_t\left[ x^{df}_{T} (1-\theta S^{d|i,\div}_{T;\theta}) \right]  \\
   &=&   \frac{x^{fd|i}_{t,T+\theta}}{Z^{d|i}_{t,T,T+\theta}}  \EE^{d|i,T}_t\left[ x^{df}_{T} (1-\theta S^{d|i,\div}_{T;\theta}) \right]  \\
   &=&   \frac{x^{fd|i}_{t,T+\theta}}{Z^{d|i}_{t,T,T+\theta}}   (1-\theta S^{d|i,\div}_{t,T;\theta}) \frac{ \EE^{d|i,T}_t\left[ x^{df}_{T} (1-\theta S^{d|i,\div}_{T;\theta}) \right]  }  { x^{df|i}_{t,T} (1-\theta S^{d|i,\div}_{t,T;\theta})  }  \\
   &=&  (1+\theta R^{f|i}_{t,T;\theta})  \EE^{f|i,T}_t\left[ 1-\theta S^{d|i,\div}_{T;\theta} \right]  
\end{eqnarray*} 
If $i=d$, $\EE^{f|d,T}_t\left[ 1-\theta S^{d|d,\div}_{T;\theta} \right]  = 1-\theta S^{d|d,\div}_{0,T;\theta}$ since funding spreads are constant and it follows that 
\begin{eqnarray*}
   \frac{d(1 + \theta L^{f/d|d}_{t,T;\theta})}{1 + \theta L^{f/d|d}_{t,T;\theta} }
   ~~=~~  \frac{d(1 + \theta R^{f|d}_{t,T;\theta})}{1 + \theta R^{f|d}_{t,T;\theta} }
   ~~=~~  \Sigma^{f|d}_{t,T;\theta}\cdot dW^{f|d,T+\theta}. 
\end{eqnarray*} 
\end{proof}


\begin{rema}
Under domestic collateral, the result can be easily extended to stochastic funding spreads by introducing the convexity adjustment due to the covariance of the funding basis and FX forward:
\begin{eqnarray}
   1 + \theta L^{f/d|d}_{t,T;\theta} 
   &=&  (1+\theta R^{f|d}_{t,T;\theta}) ( 1-\theta S^{d|d,\div}_{t,T;\theta} ) e^{\int_t^T \Sigma^{d,\mathrm{spread}}_{u,T;\theta} \cdot \Sigma^{df}_{u,T} \, du}
\end{eqnarray}
\end{rema}


%\newpage%
\subsubsection{Currency basis spread}


\begin{prop}\label{prop:aoa_currency_basis_spread}
The currency basis spread $B^{f/d|i,*}_{t,T;M}$ (resetting or non-resetting) is a martingale under the measure $\QQ^{f|i,T;M}$ associated with the annuity num\'{e}raire $A^{f|i,*}_{t,T;M}$. 
\end{prop}

\begin{proof}
Let us denote by $\mathrm{CCS}^{f/d|i,*}_{t,T;M}$ the MTM at $t$ in currency $f$ of a cross-currency swap with zero spread. By definition, $B^{f/d|i,*}_{t,T;M} = - \frac{ \mathrm{CCS}^{f/d|i,*}_{t,T;M} }{ A^{f|i,*}_{t,T;M} }$ so it is a martingale under the measure $\QQ^{f|i,T;M}$ associated with the annuity num\'{e}raire $A^{f|i}_{t,T;M}$. 
\end{proof}

In order to relate the dynamic of resetting currency basis spreads to that of discount factors and market observables (FX and Libor rates), we first consider the case of a swaplet. We restrict ourselves to the case of domestic collateral for now. 

\begin{theo}\label{theo:currency_basis_spread_constant_funding_basis}
The currency basis spread of a resetting cross-currency swaplet is:
\begin{eqnarray}
	B^{f/d|i,r}_{t,T;\theta} 
	&=& \frac{1}{\theta}\left(1+\theta R^{f|d}_{t,T;\theta} \right)  \left( \EE^{f|i,T}_t\left[ \frac{1+\theta L^{d|i}_{T;\theta}}{ 1+ \theta R^{d|i}_{T;\theta}} \right] - \frac{1+\theta L^{f|i}_{t,T;\theta}}{ 1+ \theta R^{f|i}_{t;T;\theta}}  \right)  
	\label{eq:cb_r_swaplet}
\end{eqnarray} 
In particular, 
\begin{eqnarray}
  B^{f/d|d,r}_{t,T;\theta} 
  &=& \left(1+\theta R^{f|d}_{t,T;\theta} \right)  \left(S^{f|d,\div}_{t,T;\theta} - S^{d|d,\div}_{t_0,T;\theta} \right)  
	\label{eq:cb_r_swaplet}
\end{eqnarray} 
%Equivalently,
%\begin{eqnarray}
%	Z^{f|d}_{t,T,T+\theta}B^{f/d|d,r}_{t,T;\theta} 
%	&=&  S^{f|d,\div}_{t,T;\theta} - S^{d|d,\div}_{t_0,T;\theta} 
%	\label{eq:cb_r_swaplet2}
%\end{eqnarray} 
\end{theo}


%\begin{rema}
%The fact that non-native funding spreads are stochastic explains why it is not trivial to extend the theorem to the case of non-native collateral. 
%
%Indeed, in order to compute the value of the resetting Libor forward, we have two convexity adjustments applied to the FX spot: the multiplication by the Libor and the delay between its fixing and payment. In the case of domestic collateral, our constant funding spread assumption ensures that the payment delay of the FX is exactly compensated by the accrual by the Libor (up to a multiplicative constant). But in the non-native collateral case, this is no longer the case. 
%
%The collateral convexity of discount factors and Libors are fundamentally different because the former are asset prices while the latter are forward rates.
%
%In the next section, we will extend the theorem to any collateral currency under the assumptions that volatilities are deterministic.
%\end{rema}


%
%\begin{coro}[Dynamic of currency basis]
%The currency basis spread of a swaplet is a martingale under $\QQ^{f|d,T+\theta}$ with dynamic\footnote{FIXME!!!}:
%\begin{eqnarray}
%  d B^{f/d|d,r}_{t,T;\theta}
%  &=&  \left( ???  \right)    \cdot dW^{f|d,T+\theta}_t  
%  %&=&  \left( \Sigma^{f|f}_{t,T;\theta} - \Sigma^{d\to f}_{t,T;\theta}   \right)    \cdot dW^{f|d,T+\theta}_t 
%\end{eqnarray} 
%\end{coro}
%
%\begin{proof}
%Since multiplicative funding spreads are deterministic, we have
%\begin{eqnarray*}
%	\frac{d B^{f/d|d,r}_{t,T;\theta}}{B^{f/d|d,r}_{t,T;\theta}}   =  \frac{d (1+\theta R^{f|d}_{t,T;\theta})}{1+\theta R^{f|d,r}_{t,T;\theta}} = \Sigma^{f|d}_{t,T;\theta}  \cdot dW^{f|d,T+\theta}_t 
%\end{eqnarray*} 
%\end{proof}



\begin{coro}
The currency basis spread of a resetting cross-currency swap is:
\begin{eqnarray}
  B^{f/d|d,r}_{t,T;M}  
  &=& \frac{\sum_{j=1}^{M}   Z^{f|d}_{t,T,T_j} B^{f/d|d,r}_{t,T_{j-1},\theta} }{\sum_{j=1}^{M}  Z^{f|d}_{t,T_0,T_j}} \nonumber \\
  &=& \frac{ \sum_{j=1}^{M}   Z^{f|d}_{t,T_0,T_{k}} (L^{f/d|d,r}_{t,T_{j-1},\theta}  - L^{f|d,r}_{t,T_{j-1},\theta}  ) }{\sum_{j=1}^{M}  Z^{f|d}_{t,T_0,T_j}} \nonumber \\
  &=& \frac{ \sum_{j=1}^{M}   Z^{f|d}_{t,T_0,T_{j-1}} (S^{f|d,\div}_{t,T_{j-1},\theta}  - S^{f/d|d,\div}_{t,T_{j-1},\theta} ) }{\sum_{j=1}^{M}  Z^{f|d}_{t,T_0,T_j}} 
\end{eqnarray}
\end{coro}


\begin{rema}
The above expressions correspond to different representations $B^{f/d|d,r}_{t,T;M} = f(t,X_t)$ of the currency basis spread as function of a stochastic process. We use those to get approximations for the dynamic of $B^{f/d|d,r}_{t,T;M}$ by using Ito-Taylor expansions. The quality of the approximation is controlled by how close the function $f$ is to a polynomial of low degree. For degree $\le 1$, the dynamic ends up being Gaussian if $X$ is. This technique is often known as weight freezing because the approximation amounts to replacing the stochastic weights in the exact expansion by their initial value.

The two representations as weighted sums above provide two opportunities to apply this technique to obtain Gaussian approximations. It would be interesting to compare those. 
\end{rema}






\begin{rema}[Triangle relation for currency basis spreads]
Let $d,e,f$ be three currencies. Then, non-resetting currency basis spreads satisfy:
\begin{eqnarray}
   B^{d/f|i,nr}_{t,T;M} &=& B^{d/e|i,nr}_{t,T;M} - \frac{A^{f|i}_{t,T;M}}{A^{d|i}_{t,T;M}} B^{f/e|i,nr}_{t,T;M}    .
\end{eqnarray}
This follows from the fact that an $f/e$ non-resetting cross-currency swap can be synthesised from $d/f$ and $d/e$ swaps. There is no analog for the resetting currency basis spreads. 
\end{rema}





%%%%%%%%%%%%%%%%%%%%%%%%%%%%%%%%%%%%%%%%%%%%
%% Collateral convexity adjustments
%%%%%%%%%%%%%%%%%%%%%%%%%%%%%%%%%%%%%%%%%%%%

\newpage%
\section{Collateral convexity adjustments and P\&L}




%When volatilities are deterministic, this provides closed formulas for collateral convexity adjustments:
%\begin{coro}
%Let
%\begin{itemize}
%\item $V^d_S$ be a random positive cash-flow fixing at $S$ and paid at $T$ in currency $d$.
%\item $V^{d|i}_t = \EE^{d|i,T}_t[ V^d_S]$ be the $T$-forward value of this cashflow under collateral $i$
%\item $\sigma^{V,d|i}_t$ be its volatility. 
%\end{itemize}
%If both $\sigma^{V,d|i}_t$ and $\Sigma^{i\to j}_{u,T} $ are deterministic, then 
%\begin{eqnarray}
%  \frac{V^{d|j}_{t}}{V^{d|i}_t}
%  = \exp \left( \Cov_t^{d|i,T}\left(  \ln \Q^{d|ij}_{S,T} , \ln V^{d}_S   \right) \right)  
%  = \exp \left( - \int_t^S \sigma^{V,d|i}_u \cdot \Sigma^{i\to j}_{u,T} \, du \right)
%\end{eqnarray}
%Moreover, the volatility of the forward value $V^{d|j}_t$ does not depend on the collateral: $\sigma^{V,d|j}_t = \sigma^{V,d|i}_t$. 
%\end{coro}

% 
%\subsubsection{Collateral convexity of FX forwards}%
%
%
%%The following proposition is a generalisation of FX-forward invariance under stochastic collateral rates. 
%\begin{prop}[Collateral convexity of FX forwards]
%Under deterministic volatility assumptions:
%\begin{eqnarray}
%	 \frac{x^{df|j}_{t,T}}{x^{df|i}_{t,T}}
%	&=& \exp\left\{ -\int_t^T \left( \Sigma^{d|d}_{u,T} - \Sigma^{f|f}_{u,T} + \Sigma^{d\to f}_{u,T} + \sigma^{df}_u  \right) \cdot \Sigma^{i\to j}_{u,T} \, du \right\}
%\end{eqnarray} 
%\end{prop}
%
%
%\begin{rema}
%This is the generalisation of FX forward invariance. Indeed, it implies that off-market discount curves are entirely specified by:
%\begin{eqnarray}
%	Z^{d|f}_{t,T} &=& Z^{d|d}_{t,T} \frac{Z^{f|f}_{t,T}}{Z^{f|d}_{t,T}} 
%	\exp\left\{ \int_t^T \left( \Sigma^{d|d}_{u,T} - \Sigma^{f|f}_{u,T} + \Sigma^{d\to f}_{u,T} + \sigma^{df}_u \right) \cdot \Sigma^{d\to f}_{u,T}\, du  \right\}
%\end{eqnarray} 
%\end{rema}
%
%
% 
%\subsubsection{Collateral convexity of Libors}%
%
%[FIXME: remove deterministic vol assumption]
%
%\begin{prop}
%Under the deterministic multiplicative funding spread hypothesis, 
%\begin{eqnarray}
%	1+\theta L^{e|c}_{t,T;\theta}  
%	&=& (1+\theta L^{e|e}_{t,T;\theta})  \exp\left(  - \int_t^T   \Sigma^{e|e}_{u,T;\theta}  \cdot \Sigma^{e\to c}_{u,T+\theta} \, du \right)
%\end{eqnarray} 
%\end{prop}

In this section, we describe a basic framework to decompose the convexity adjustments associated with a change of collateral. 

While our presentation is different, the results are equivalent to that of McCloud in \cite{McCloud:CollateralVolatility}.


\subsection{Change of collateral formula}%
 
%Now we consider explicit computations of these convexity adjustments in the HJM framework.
 
\begin{prop}
The Radon-Nikodym derivative of the change of $T$-forward measure is characterised by the exponential martingale 
\begin{eqnarray}
  \xi^{d|i\to j}_{t,T}  := \EE^{d|i,T}_t\left[ \frac{d\QQ^{d|j,T}}{d\QQ^{d|i,T}} \right] 
  = \mathcal{E}\left( - \int_0^\cdot \Sigma^{i\to j}_{u,T} \cdot dW^{d|i,T}_u \right)_t
  \label{eq:collateral_radon_nikodym}
\end{eqnarray}
\end{prop}


%Remember that, for any positive cashflow $v^e_T$ paid at $T$ in currency $e$, we defined by
%\begin{eqnarray}
%  \Q^{e|i\to j}_{t,T}[v^e_T] &=& \frac{\EE^{e|j,T}_t\left[ v^e_T \right] } {\EE^{e|i,T}_t\left[ v^e_T \right] }  
%\end{eqnarray}

%\begin{prop}
%For any positive cashflow $v^e_T$ paid at $T$ in currency $e$, 
%\begin{eqnarray}
%  \Q^{e|i\to j}_{t,T}[v^e_T] &=&  \exp\left(  - \int_t^T  \EE^{e|i,T} \left[ \sigma^{v|e}_{u,T}  \cdot \Sigma^{i\to j}_{u,T} \right]  \, du \right)
%\end{eqnarray}
%\end{prop}
%
%[TODO: Check this for stochastic vols]

\subsection{Collateral convexity adjustments}
Discount factors play a specific role in quantitative finance. They correspond to the price of tradeable assets (fixed cashflow of one unit of currency) but they are not directly observables: zero coupon bonds are not liquid assets. Nevertheless, they are ubiquitous because they provide a way to "factor out" inter-temporality: we can decompose the \emph{present value} of a cash-flow into the product of the \emph{discount factor} and the \emph{forward value} of the cash-flow. 

One can apply the same logic to collateral convexity adjustments, decomposing them into:
\begin{enumerate}
\item Convexity adjustment to discount factors
\item Convexity adjustment to forward rates.
\end{enumerate} 

\begin{rema}
The computation of convexity adjustment typically goes as follows:  
\begin{enumerate}
\item Consider a market quote
\item Express it in terms of PV of tradeable assets (ideally a ratio).
\item Derive its drift under a well-chosen measure (in case of a ratio, drift is zero under the measure associated with the denominator).
\item Compute the collateral convexity adjustment as a covariance with the change of measure (if the drift is deterministic then the convexity adjustment only depends on the volatility).  
\end{enumerate}
\end{rema}


\subsection{Convexity adjustment to domestic discount factors}
We first consider fixed cash-flows/discount factors.


\begin{defi}
For any payment currency $e$, collateral $i,j$ and maturity $T$, we define the corresponding \emph{discount factor  collateral convexity adjustment} as 
\begin{eqnarray}
  \QZ^{e|i\to j}_{t,T}   &:=& \frac{ Z^{e|j}_{t,T} } { Z^{e|i}_{t,T} } 
\end{eqnarray}
\end{defi}

\begin{rema}
By definition of the instantaneous collateral basis, it can be written as
\begin{eqnarray}
  \QZ^{e|i\to j}_{t,T} &=& \EE^{e|i,T}_t\left[ e^{-\int_t^T q^{i\to j}_u \,du}  \right]  .
\end{eqnarray}
\end{rema}






\subsection{Convexity adjustment to forward rates}

%We can generalise the previous definition.

\begin{defi}
For any positive cashflow $v^e_T$ paid at $T$ in currency $e$, we define the corresponding \emph{forward collateral convexity adjustment} as:
\begin{eqnarray}
  \Q^{e|i\to j}_{t,T}[v^e_T] &=& \frac{\EE^{e|j,T}_t\left[ v^e_T \right] } {\EE^{e|i,T}_t\left[ v^e_T \right] }  
\end{eqnarray}
\end{defi}


\begin{rema}
By definition of the instantaneous collateral basis, it can be computed as
\begin{eqnarray}
  \Q^{e|i\to j}_{t,T}[v^e_T] &:=& \frac{\EE^{e|i,T}_t\left[  \xi^{e|i\to j}_{T,T} v^e_T  \right]}{\EE^{e|i,T}_t\left[  \xi^{e|i\to j}_{T,T}   \right]}  .
\end{eqnarray}
\end{rema}

\begin{rema}
Note that the symbol $\Q^{e|i\to j}_{t,T}$ is neither linear nor a multiplicative operator\footnote{For mathematicians, it can be seen as the differential/coboundary operator in cohomology.
%We have $d: \bigoplus E^i \to \bigoplus_{i,j} E^{i,j}$, $d (i, v^e_T) = \log  \EE^{e|j,T}_t\left[ v^e_T \right] - \log  \EE^{e|i,T}_t\left[ v^e_T \right] $  
}. However, it provides a nice way to express convexity adjustments in the case of lognormal variables. This property makes it especially relevant to the extension of the BS and LGM models where the standard changes of measures are given by lognormal processes.   
\end{rema}

\subsection{Change of collateral}


\subsubsection{Multiplicative property}
The first (and trivial) property of collateral convexity adjustments is their multiplicative property with respect to collateral.

\begin{prop}[Change of collateral]
We have
\begin{eqnarray}
  \QZ^{e|i\to j}_{t,T} \QZ^{e|j\to k}_{t,T} &=& \QZ^{e|i\to k}_{t,T} \\
  \Q^{e|i\to j}_{t,T}[v^e_T] \Q^{e|j\to k}_{t,T}[v^e_T]  &=& \Q^{e|i\to k}_{t,T}[v^e_T]
\end{eqnarray}
\end{prop}

\begin{rema}
This implies that it is enough to know the change of collateral from the payment currency. The others can be deduced according to: 
\begin{eqnarray}
  \QZ^{e|i\to j}_{t,T} &=& \frac{\QZ^{e|e\to j}_{t,T}}{\QZ^{e|e\to i}_{t,T}} \\
  \Q^{e|i\to j}_{t,T}[v^e_T] &=& \frac{\Q^{e|e\to j}_{t,T}[v^e_T]}{\Q^{e|e\to i}_{t,T}[v^e_T]}
\end{eqnarray}
\end{rema}

\subsubsection{FX forwards}

\begin{prop}
Absence of arbitrage implies 
\begin{eqnarray}
   x^{df|c}_{t,T} &=& x^{df}_{t} \frac{Z^{f|c}_{t,T}}{Z^{d|c}_{t,T}}
\end{eqnarray}
And, more generally, 
\begin{eqnarray}
   \EE^{d|c,T}_t\left[ x^{df}_{T}v^{f}_T\right]  &=&  x^{df|c}_{t,T} \EE^{f|c,T}_t\left[ v^{f}_T\right] .
\end{eqnarray}
\end{prop}


%\begin{proof}
%An FX forward contract corresponds to a single exchange of cashflows.  At $T$, $A$ receive $N^f$ units of $f$ and pays $N^d = K N^f$ units of $d$. The value of this contract at $T$ in currency $d$ is 
%\begin{eqnarray}
%   V^d_T &=& x^{df}_{T}N^f - N^d  = (x^{df}_T - K)N^f 
%\end{eqnarray}
%We can replicate these cash-flows by buying $N^f$ foreign ZCB and selling $N^d$ domestic ones:
%\begin{eqnarray}
%   V^d_t &=& x^{df}_{t}N^f Z^{f|c}_{t,T} - N^dZ^{d|c}_{t,T}   = Z^{d|c}_{t,T} N^f\left(x^{df}_{t}\frac{Z^{f|c}_{t,T}}{Z^{d|c}_{t,T}} - K \right)
%\end{eqnarray}
%So, $x^{df|c}_{t,T} = x^{df}_{t} \frac{Z^{f|c}_{t,T}}{Z^{d|c}_{t,T}}$. 
%\end{proof}

\begin{prop}[Triangular relation for FX forwards]
For any triple of currencies $d,e,f$ and any collateral $i$, we have 
\begin{eqnarray}
  x^{df|i}_{t,T} &=& x^{de|i}_{t,T}x^{ef|i}_{t,T}
\end{eqnarray}
\end{prop}

\begin{proof}
Indeed,
\begin{eqnarray}
  x^{de|i}_{t,T}x^{ef|i}_{t,T} 
  =  x^{de}_{t} \frac{Z^{e|i}_{t,T}}{Z^{d|i}_{t,T}} x^{ef}_{t} \frac{Z^{f|i}_{t,T}}{Z^{e|i}_{t,T}} 
  =  x^{de}_{t}x^{ef}_{t} \frac{Z^{f|i}_{t,T}}{Z^{d|i}_{t,T}}   
  =  x^{df|i}_{t,T}.
\end{eqnarray}
\end{proof}


The FX triangle relation implies that the collateral convexity of the collateral convexity of FX forwards\footnote{To continue with the cohomological point of view, FX forwards are trivial cocycles so they decompose into images of discount factors under the differential operators.}\footnote{Note: This does not extend to the case of a quanto FX forwards as the triangle relationship on the FX vols becomes relevant rather than just the forwards.} is entirely determined by that of discount factors. % $\QZ^{d|i\to j}_{t,T}$.
\begin{prop}
For any currencies $d,f$ and any collaterals $i,j$, we have 
\begin{eqnarray}
  \Q^{d|i\to j}_{t,T} [ x^{df}_T ] ~=~ \frac{x^{df|j}_{t,T}}{x^{df|i}_{t,T}} ~=~ \frac{\QZ^{f|i\to j}_{t,T}}{\QZ^{d|i\to j}_{t,T}} 
\end{eqnarray}
\end{prop}

\begin{proof}
We have 
\begin{eqnarray}
 \Q^{d|i\to j}_{t,T} [ x^{df}_T ]  
  = \frac{\EE^{d|j,T}_t\left[ x^{df}_T \right] } {\EE^{d|i,T}_t\left[ x^{df}_T \right] }  
  = \frac{x^{df|j}_{t,T}}{x^{df|i}_{t,T}} 
  = \frac{x^{df}_{t}Z^{f|j}_{t,T}/Z^{d|j}_{t,T}}{x^{df}_{t}Z^{f|i}_{t,T}/Z^{d|i}_{t,T}} 
  = \frac{Z^{f|j}_{t,T}}{Z^{f|i}_{t,T}} \Big/ \frac{Z^{d|j}_{t,T}}{Z^{d|i}_{t,T}} 
  =  \frac{\QZ^{f|i\to j}_{t,T}}{\QZ^{d|i\to j}_{t,T}} 
\end{eqnarray}
\end{proof}

\subsection{Change of  payment currency}

The property of FX forwards generalises as follows.

\begin{prop}[Change of currency]
For any positive cashflow $v^f_T$ paid at $T$ in currency $f$, 
\begin{eqnarray}
  \Q^{d|i\to j}_{t,T}[x^{df}_Tv^f_T]  ~=~  \Q^{d|i\to j}_{t,T}[x^{df}_T] \Q^{f|i\to j}_{t,T}[x^{df}_Tv^f_T]   
  ~=~ \frac{\QZ^{f|i\to j}_{t,T}}{\QZ^{d|i\to j}_{t,T}} \Q^{f|i\to j}_{t,T}[v^f_T] .
\end{eqnarray}
Equivalently 
\begin{eqnarray}
  \QZ^{d|i\to j}_{t,T}\Q^{d|i\to j}_{t,T}[x^{df}_Tv^f_T]  &=& \QZ^{f|i\to j}_{t,T} \Q^{f|i\to j}_{t,T}[v^f_T] .
\end{eqnarray}
\end{prop}

\begin{proof}
\begin{eqnarray*}
  \Q^{d|i\to j}_{t,T}[x^{df}_Tv^f_T] 
  = \frac{\EE^{d|j,T}_t\left[ x^{df}_Tv^f_T\right] } {\EE^{d|i,T}_t\left[ x^{df}_Tv^f_T \right] }   
  = \frac{x^{df|j}_{t,T} \EE^{f|j,T}_t\left[ v^f_T\right] } {x^{df|i}_{t,T} \EE^{f|i,T}_t\left[ v^f_T \right] }   
  = \Q^{d|i\to j}_{t,T}[x^{df}_T]  \Q^{f|i\to j}_{t,T}[v^f_T]  
  = \frac{\QZ^{f|i\to j}_{t,T}}{\QZ^{d|i\to j}_{t,T}} \Q^{f|i\to j}_{t,T}[v^f_T] 
\end{eqnarray*}
\end{proof}

%\begin{rema}
%This implies that is enough to know the change of collateral from the payment currency. The others can be deduced from: 
%\begin{eqnarray}
%  \QZ^{e|i\to j}_{t,T} &=& \frac{\QZ^{e|e\to j}_{t,T}}{\QZ^{e|e\to i}_{t,T}} \\
%  \Q^{e|i\to j}_{t,T}[v^e_T] &=& \frac{\Q^{e|e\to j}_{t,T}[v^e_T]}{\Q^{e|e\to i}_{t,T}[v^e_T]}
%\end{eqnarray}
%\end{rema}

\subsection{Relation to McCloud's definition}


\begin{rema}
In \cite{McCloud:CollateralVolatility}, McCloud defines convexity adjustment curves as\footnote{The original notation is $Q^{ji|e}_{t,T} [ v^e_T ]$.} 
\begin{eqnarray}
  \widehat{Q}^{e|i\to j}_{t,T} [ v^e_T ]
  %&=&  \frac{\EE^{e|j,T}_t\left[ v^e_T \right] } {\EE^{e|i,T}_t\left[ v^e_T \right] }  
   &:=& \frac{\EE^{i|i,T}_t\left[ e^{-\int_t^T q^{i\to j}_u \,du} x^{ie}_Tv^e_T \right] } {\EE^{i|i,T}_t\left[ e^{-\int_t^T q^{i\to j}_u \,du}  \right]  \EE^{i|i,T}_t\left[ x^{ie}_Tv^e_T \right]  }  
\end{eqnarray}
How does this relate to our definition?
\begin{eqnarray*}
  \widehat{Q}^{e|i\to j}_{t,T} [ v^e_T ]
  %&=&  \frac{\EE^{e|j,T}_t\left[ v^e_T \right] } {\EE^{e|i,T}_t\left[ v^e_T \right] }  
   &=& \frac{\EE^{i|i,T}_t\left[ e^{-\int_t^T q^{i\to j}_u \,du} x^{ie}_Tv^e_T \right] } {\EE^{i|i,T}_t\left[ e^{-\int_t^T q^{i\to j}_u \,du}  \right]  \EE^{i|i,T}_t\left[ x^{ie}_Tv^e_T \right]  }  \\
   &=& \frac{Z^{i|i}_{t,T}\EE^{i|i,T}_t\left[ e^{-\int_t^T q^{i\to j}_u \,du} x^{ie}_Tv^e_T \right] } {Z^{i|i}_{t,T}\EE^{i|i,T}_t\left[ e^{-\int_t^T q^{i\to j}_u \,du}  \right]  \EE^{i|i,T}_t\left[ x^{ie}_Tv^e_T \right]  }  \\
   &=& \frac{\EE^{i,\RN}_t\left[ e^{-\int_t^T r^{i|i}_u + q^{i\to j}_u \,du} x^{ie}_Tv^e_T \right] } {\EE^{i,\RN}_t\left[ e^{-\int_t^T r^{i|i}_u + q^{i\to j}_u \,du}  \right]  \EE^{i|i,T}_t\left[ x^{ie}_Tv^e_T \right]  }  \\
   &=& \frac{Z^{i|j}_{t,T}\EE^{i|j,T}_t\left[  x^{ie}_Tv^e_T \right] } {Z^{i|j}_{t,T} \EE^{i|i,T}_t\left[ x^{ie}_Tv^e_T \right]  }  \\
   &=& \Q^{i|i\to j}_{t,T} [ x^{ie}_Tv^e_T ] \\
   &=& \frac{\QZ^{e|i\to j}}{\QZ^{i|i\to j}} \Q^{e|i\to j}_{t,T} [ v^e_T ]
\end{eqnarray*}
So McCloud's convexity adjustments can be expressed in terms of ours according to:
\begin{eqnarray}
   \widehat{Q}^{e|i\to j}_{t,T} [ v^e_T ] 
   %&=& \Q^{i|i\to j}_{t,T} [ x^{ie}_T v^e_T ] 
   &=& \frac{\QZ^{e|i\to j}}{\QZ^{i|i\to j}}   \Q^{e|i\to j}_{t,T} [ v^e_T ] 
\end{eqnarray}
and, conversely, our convexity adjustments can be expressed in terms of McCloud's according to:
\begin{eqnarray}
  \frac{\widehat{Q}^{e|i\to j}_{t,T} [ v^e_T ]}{\widehat{Q}^{e|i\to j}_{t,T} [ 1^e_T ]}
   &=& \Q^{e|i\to j}_{t,T} [ v^e_T ]  
\end{eqnarray}

Our definition splits McCloud's definition into two parts to emphasise the different roles played by discount factors and forward rates.
 
His formulation is interesting when one choses $i$ to be the domestic currency of the model as it decomposes the convexity adjustment into expectations under the native domestic $T$-forward measure. 
\end{rema}

%\begin{rema}
%In terms of MTM, the formula can be rewritten:
%\begin{eqnarray}
%  v^{e|i}_t &=& v^{e|e}_t s^{e\to i}_{t,T} \Q^{e|e\to i}_{t,T}[v^e_T] .
%\end{eqnarray}
%\end{rema}


%\subsubsection{Forward Libors}
%
%\begin{eqnarray}
%  \Q^{e|i\to j}_{t,T+\theta}[1+\theta L^{e}_{T;\theta}] 
%  &=& \frac{1+\theta L^{e|j}_{t,T;\theta}} {1+\theta L^{e|i}_{t,T;\theta} }  
%\end{eqnarray}
%
%
%\subsubsection{Resetting Forward Libors}
%
%\begin{eqnarray}
%  \Q^{f|i\to j}_{t,T+\theta}\left[ \frac{x^{fe}_{T+\theta}}{x^{fe}_T}(1+\theta L^{e}_{T;\theta}) \right] 
%  &=&  \frac{ \EE^{f|j,T+\theta}_t\left[ \frac{x^{fe}_{T+\theta}}{x^{fe}_T} (1+\theta L^{e}_{T;\theta}) \right] 
%  } {
%  \EE^{f|i,T+\theta}_t\left[ \frac{x^{fe}_{T+\theta}}{x^{fe}_T} (1+\theta L^{e}_{T;\theta}) \right]  }  
%\end{eqnarray}

 
\subsection{P\&L impact}%

\subsubsection{The P\&L equation}% 

\begin{eqnarray}
  V^{e|i}_t &=&  \EE^{e,\RN}_t\left[ e^{-\int_t^T r^{e|i}_u\,du} V_T \right]  \\
  &=&  Z^{e|i}_{t,T} \EE^{e|i,T}_t[V^{e|i}_T]  \\
\end{eqnarray}


Let's say the value of the portfolio under standard collateral is $X$, a function of the market observables $S$. Let $Q$ be the collateral basis so that 
\begin{eqnarray}
  \MTM_t &=&  \EE_t\left[ Q X \right]  \\
  &=& \EE_t\left[ Q \right] \EE_t\left[ X \right]  + \Cov_t\left[ Q,X \right]  
\end{eqnarray}

%\begin{eqnarray}
%  \MTM_t &=&  \EE\left[ Q X \right]  \\
%  &=& \EE_t\left[ Q \right] \EE_t\left[ X \right]  + \Cov_t\left[ Q,X \right]  
%\end{eqnarray}

\begin{eqnarray}
  \EE^\PP_t[d\PnL_t]&=&  \frac{1}{2}(\partial^2_SV) \left( \sigma^{S,2}_{t,\rm{real}} -   \sigma^{S,2}_{t,\rm{mdl}} \right)  \\
  & & + \frac{1}{2}(\partial^2_QV) \left(  \sigma^{Q,2}_{t,\rm{real}} -   \sigma^{Q,2}_{t,\rm{mdl}} \right)   \\
  & & + (\partial^2_{S,Q}V) \left( \rho^{S,Q}_{t,\rm{real}} \sigma^S_{t,\rm{real}} \sigma^Q_{t,\rm{real}} -  \rho^{S,Q}_{t,\rm{mdl}} \sigma^S_{t,\rm{mdl}} \sigma^Q_{t,\rm{mdl}} \right)  .
\end{eqnarray}
In a model with deterministic basis $Q$, we have $\sigma^Q_{t,\rm{mdl}} = 0$ so 
\begin{eqnarray}
  \EE^\PP_t[d\PnL_t] &=&  \frac{1}{2}(\partial^2_SV) \left( \sigma^{S,2}_{t,\rm{real}} -   \sigma^{S,2}_{t,\rm{mdl}} \right)  \\
  & & + \frac{1}{2}(\partial^2_QV)  \sigma^{Q,2}_{t,\rm{real}}  + (\partial^2_{S,Q}V)  \rho^{S,Q}_{t,\rm{real}} \sigma^S_{t,\rm{real}} \sigma^Q_{t,\rm{real}} .
\end{eqnarray}
The basis Gamma and cross-Gamma ($\Gamma\times\Gamma$) terms are unhedged which leads to a drift in the P\&L. From the quant point of vue they are more often called "P\&L leaks" or "P\&L bleed". Traders call these terms "deliveries": I hedge both my deltas, but when market moves, the difference between my predict and the realised P\&L is not zero on average. Instead, I get "delivered" some additional P\&L. 

If the Gamma and cross-Gamma are negative, these deliveries are systematically negative and the desk hedging with this model will bleed P\&L. 

So, one should know either
\begin{enumerate}
\item have a full model with stochastic $Q$ factor
\item know the sign of the $\Gamma\times\Gamma$ term to ensure the position is conservative. 
\end{enumerate}







\newpage%
\section{Deterministic volatility models}


\subsection{Model specification}%

We now consider the case where the covariance of instantaneous forward rates is deterministic. This implies that the main market observables are deterministic functions of (displaced) lognormal processes so that we have explicit formulas for their collateral convexity adjustments.

Unless specified otherwise. in this section we assume that:
\begin{enumerate}
\item Funding spreads are constant.
\item Volatility matrices are deterministic.
\end{enumerate}


\subsection{Dynamic of market observables}
We now derive the dynamic of market observables. % in order to calibrate the model.


\subsubsection{FX forwards}

\begin{prop}
The FX forward $x^{df|i}_{t,T}$ is a lognormal martingale under $\QQ^{d|i,T}$:  
\begin{eqnarray}
  \frac{dx^{df|i}_{t,T}}{x^{df|i}_{t,T}} &=& \Sigma^{df}_{t,T} \cdot dW^{d|i,T}_{t} 
\end{eqnarray}
where 
\begin{eqnarray}
  \Sigma^{df}_{t,T} &=& \Sigma^{d|d}_{t,T} - \Sigma^{f|f}_{t,T} +\Sigma^{d\to f}_{t,T}  + \sigma^{df}_t  .
\end{eqnarray}
\end{prop}

\begin{rema}
In the case of deterministic volatilities, FX forwards are lognormal processes with deterministic variance independent of the collateral. However, there is a convexity adjustment associated with the change of collateral. In other words, FX implied volatilities remains the same but there is a shift in the moneyness of the options corresponding to the covariance between the collateral basis and the log of the FX forward\footnote{Looks like a repeat of a previous remark. Avoid duplication.}.
\end{rema}


\begin{prop}
FX options prices are given by the BS formula\footnote{Here the BS function is the function of the log-moneyness forward $y$ and total variance $w$ defined by 
\begin{eqnarray}
   \BS\left( y , w \right) &=& \mathcal{N}(d_+(y,w)) - e^y \mathcal{N}(d_-(y,w)) \\
   d_{\pm}(y,w) &=& -\frac{\ln y}{\sqrt{w}} \pm \frac{1}{2} \sqrt{w}
\end{eqnarray}}:
\begin{eqnarray}
  \mathrm{Call}^{df|i}_t(T,K)
  %Z^{d|i}_{t,T}\EE^{d|i,T}_t\left[\max\left(x^{df}_T-K,0\right)\right] 
  &=& Z^{d|i}_{t,T} x^{df|i}_{t,T} \BS\left( y = \ln(K/x^{df|i}_{t,T}), w = w^{df}_{t}(T) \right) 
\end{eqnarray} 
where total variance is 
\begin{eqnarray}
  w^{df}_{t}(T) &=& \int_t^T \left| \Sigma^{d|d}_{u,T} - \Sigma^{f|f}_{u,T}  + \Sigma^{d\to f}_{u,T}  + \sigma^{df}_u \right|^2 \, du 
\end{eqnarray} 
regardless of the collateral $c$.
\end{prop}

%\begin{rema}
%Note that the total variance is unchanged by the collateral but the moneyness is not. 
%\end{rema}

%The convexity adjustment on option prices is given by the product of the FX forward convexity adjustment and a BS moneyness adjustment:
%\begin{eqnarray}
%  \frac{  \BS\left( y = \ln(K/x^{df|j}_{t,T}) , w = w^{df}_{t,T}\right) } 
%  {  \BS\left( y = \ln(K/x^{df|i}_{t,T}) , w = w^{df}_{t,T}\right) } 
%\end{eqnarray} 
%with shift $\Delta y = \ln \Q^{d|i\to j}[x^{df}_T] = $.

\subsubsection{Caplets}

\begin{coro}
Caplets prices are given by the BS formula
\begin{eqnarray}
  %& & \mathrm{Caplet}^{e|i,\theta}_{t}(T,K)  \nonumber \\
  %& & Z^{e|i}_{t,T+\theta}\EE^{e|i,T+\theta}_t\left[\max\left(L^{e}_{T,T;\theta}-K,0\right)\right]  \nonumber \\
  \mathrm{Caplet}^{e|i,\theta}_{t}(T,K) 
  &=&  Z^{e|i}_{t,T+\theta} \left( 1+\theta L^{e|i}_{t,T;\theta} \right)\BS\left(y = \ln \left( \tfrac{1+\theta K}{ 1+\theta L^{e|i}_{t,T;\theta} } \right), w = w^{e,\theta}_t(T,K) \right) .
\end{eqnarray} 
where total variance is [FIXME: do not use the LGM parametrisation yet]
\begin{eqnarray}
  w^{e,\theta}_t(T,K)
  &=&  \sum_{a=1}^{N_e}  \sum_{b=1}^{N_{e}} \int_t^T  \rho^{(e,a),(e,b)}_u \sigma^{e|e,a}_u \sigma^{e|e,b}_u  G^{e|e,a}_{u,T;\theta} G^{e|e,b}_{u,T;\theta} \, du  
\end{eqnarray} 
%The dynamic of the caplet:  
%\begin{eqnarray}
%	\frac{dC^{f|d}_t(T,K)}{C^{f|d}_t(T,K)}
%	&=&  r^{f|d}_t dt + \frac{F\partial_F\BS}{\BS}  \Sigma^{f|d}_{t,T;\theta}\cdot dW^{f,\RN}_t + \Sigma^{f|d}_{t,T +\theta}\cdot dW^{f,\RN}_t
%\end{eqnarray} 
\end{coro}

\begin{proof}
\begin{eqnarray*}
  & & Z^{e|i}_{t,T+\theta}\EE^{e|i,T+\theta}_t\left[\max\left(L^{e}_{T,T;\theta}-K,0\right)\right]  \nonumber \\
  &=& Z^{e|i}_{t,T+\theta}\EE^{e|i,T+\theta}_t\left[\max\left((1+\theta L^{e}_{T,T;\theta}) - (1+\theta K),0\right)\right]  \\
  &=& Z^{e|i}_{t,T+\theta} \left( 1+\theta L^{e|i}_{t,T;\theta} \right)\BS\left(y = \ln \left( \tfrac{1+\theta K}{ 1+\theta L^{e|i}_{t,T;\theta} } \right), w = w^{e,\theta}_t(T,K) \right) .
\end{eqnarray*}  
where total variance is [FIXME: do not use the LGM parametrisation yet]
\begin{eqnarray*}
  & & w^{e,\theta}_t(T,K) \\
  &=&   \int_t^T  \Sigma^{e|e,\top}_{u,T;\theta}  \Sigma^{e|e}_{u,T;\theta} \, du \\
  &=&   \int_t^T  \left(\Sigma^{e|e}_{u,T+\theta} - \Sigma^{e|e}_{u,T}\right)^\top  \left(\Sigma^{e|e}_{u,T+\theta}  - \Sigma^{e|e}_{u,T}\right) \, du \\
  &=&   \int_t^T  \left(\Gamma_u \Pi^{e|e}\sigma^{e|e}_{u}H^{e|e,-1}_u(H^{e|e}_{T+\theta} - H^{e|e}_{T})\right)^\top \left(\Gamma_u \Pi^{e|e}\sigma^{e|e}_{u}H^{e|e,-1}_u(H^{e|e}_{T+\theta} - H^{e|e}_{T})\right) \, du \\
  &=&   \int_t^T  \left(\sigma^{e|e}_{u}H^{e|e,-1}_u(H^{e|e}_{T+\theta} - H^{e|e}_{T})\right)^\top \Pi^{e|e,\top} \Gamma_u \Gamma_u \Pi^{e|e} \left( \sigma^{e|e}_{u}H^{e|e,-1}_u(H^{e|e}_{T+\theta} - H^{e|e}_{T})\right) \, du \\
  &=&  \sum_{a=1}^{N_e}  \sum_{b=1}^{N_{e}} \int_t^T  \rho^{(e,a),(e,b)}_u \sigma^{e|e,a}_u \sigma^{e|e,b}_u  G^{e|e,a}_{u,T;\theta} G^{e|e,b}_{u,T;\theta} \, du  
\end{eqnarray*} 
%\begin{eqnarray*}
%	d(C^{f|d}_t(T,K)/Z^{f|d}_{t,T+\theta})
%	&=&  (F\partial_F\BS)  \Sigma^{f|d}_{t,T;\theta}\cdot dW^{f|d,T+\theta}_t 
%\end{eqnarray*} 
\end{proof}







\subsubsection{Dynamic of resetting forward Libors}

\begin{prop}
The resetting forward Libors $L^{f/d|i}_{t,T;\theta}$ is a shifted lognormal martingale under $\QQ^{f|i,T+\theta}$:
\begin{eqnarray}
   \frac{d(1 + \theta L^{f/d|i}_{t,T;\theta})}{1 + \theta L^{f/d|i}_{t,T;\theta}}
   &=& \Sigma^{f|d}_{t,T;\theta} \cdot dW^{f|i,T+\theta}_t 
\end{eqnarray} 
so that
\begin{eqnarray}
   1 + \theta L^{f/d|i}_{t,T;\theta} 
  &=& (1+\theta R^{f|i}_{t,T;\theta}) (1-\theta S^{d|i,\div}_{t,T;\theta}) e^{ - \int_t^T \Sigma^{d\to i}_{u,T;\theta} \cdot \Sigma^{df|i}_{u,T} \,du }  
\end{eqnarray} 
\end{prop}


\begin{proof}
We use the dynamic of the funding spread to compute:
\begin{eqnarray*}
   1 + \theta L^{f/d|i}_{t,T;\theta} 
   &=&  (1+\theta R^{f|i}_{t,T;\theta})  \EE^{f|i,T}_t\left[ 1-\theta S^{d|i,\div}_{T;\theta} \right]   \\
   &=& (1+\theta R^{f|i}_{t,T;\theta})  \EE^{f|i,T}_t\left[ (1-\theta S^{d|i,\div}_{t,T;\theta}) \mathcal{E}\left(- \int_t^T \Sigma^{d\to i}_{u,T;\theta} \cdot dW^{d|i,T}_u \right) \right]   \\
  &=& (1+\theta R^{f|i}_{t,T;\theta}) (1-\theta S^{d|i,\div}_{t,T;\theta}) e^{ - \int_t^T \Sigma^{d\to i}_{u,T;\theta} \cdot \Sigma^{df}_{u,T} \,du }  
\end{eqnarray*} 
We can now compute the volatility, knowing this is a $\QQ^{f|i,T+\theta}$-martingale:
\begin{eqnarray*}
    \frac{d(1 + \theta L^{f/d|i}_{t,T;\theta})}{1 + \theta L^{f/d|i}_{t,T;\theta} } 
   &=& \frac{d(1 + \theta R^{f|i}_{t,T;\theta})}{1 + \theta R^{f|i}_{t,T;\theta}} 
   +  \frac{d(1-\theta S^{d|i,\div}_{t,T;\theta})}{1-\theta S^{d|i,\div}_{t,T;\theta}} + O(dt) \\
   &=& \Sigma^{f|i}_{t,T;\theta} \cdot dW^{f|i,T+\theta}_t  -  \Sigma^{d\to i}_{t,T;\theta} \cdot dW^{d|i,T}_t  + O(dt) \\
   &=& \left( \Sigma^{f|f}_{t,T;\theta} + \Sigma^{f\to i}_{t,T;\theta} -   \Sigma^{d\to i}_{t,T;\theta}\right) \cdot dW^{f|i,T+\theta}_t  \\
   &=&  \Sigma^{f|d}_{t,T;\theta} \cdot dW^{f|i,T+\theta}_t  
\end{eqnarray*} 

\end{proof}


%\begin{rema}
%Note that, once again, the volatility of the resetting forward Libor does not depend on the collateral. 
%\end{rema}


\begin{rema}
Note that the volatility of the resetting forward Libor does not depend on the native rates volatility. This is quite surprising since the domestic Libor $L^d$ appears explicitly in the definition. 

The explanation is the following. Due to the resetting notional, the domestic discount rate volatility risk is compensated by the domestic component in the volatility of the FX-forward, or equivalently in the drift of the FX spot. The FX spot volatility also cancels out due to the fact that the resetting forward is paid in the foreign currency. The remaining volatility risks are the ones from foreign rates and collateral basis volatility\footnote{and the domestic funding basis which we assumed to be zero.}. 
\end{rema}

\subsubsection{Dynamic of currency basis spreads}

%By definition, 
%\begin{eqnarray*}
%  B^{f/d|c,r}_{t,T;\theta}  &=&  L^{f/d|c,r}_{t,T;\theta}  - L^{f|c}_{t,T;\theta} 
%\end{eqnarray*} 
%What is its dynamic under our assumptions? 

\begin{prop}
The currency basis spread is the difference of two exponential $\QQ^{f|c,T+\theta}$-martingales:
\begin{eqnarray}
  B^{f/d|c,r}_{t,T;\theta} &=& L^{f/d|c,r}_{t,T;\theta} -  L^{f|c}_{t,T;\theta} \\
  &=& \frac{1+\theta L^{f/d|c,r}_{t,T;\theta}}{\theta}  - \frac{1+\theta L^{f|c}_{t,T;\theta}}{\theta} \nonumber
\end{eqnarray} 
with dynamic:
\begin{eqnarray*}
  d B^{f/d|c,r}_{t,T;\theta} 
  &=& ??
\end{eqnarray*} 
\end{prop}

\begin{rema}
Do we gain anything by deriving the full dynamic?

Even with deterministic spot covariance, the currency basis spread is the difference of two shifted lognormal processes which does not lead to a nice dynamic. In particular, computing the delivered volatility of a currency basis spread is equivalent to computing the implied volatility of a spread option in a two asset BS model. There is no exact formula and this will typically require 2D numerical integration. Instead, we will always approximate the dynamic of each Libor by a Gaussian one. These approximations are very accurate for low volatility regimes but they may fail for EM markets or stress scenarios.
\end{rema}





\begin{rema}
Instead of the additive currency basis spread, it can be mathematically easier to consider the multiplicative basis spread defined by:
\begin{eqnarray}
  1 - \theta B^{f/d|i,r,\div}_{t,T;\theta} 
  &=& \frac{ 1+\theta L^{f|i}_{t,T;\theta}}{1+\theta (L^{f|i}_{t,T;\theta} + B^{f/d|i,r}_{t,T;\theta})}  \\
  &=& \frac{ 1+\theta L^{f|i}_{t,T;\theta}}{1+\theta L^{f/d|i,r}_{t,T;\theta}}  \\
  &=& \frac{1-\theta S^{f|i,\div}_{t,T;\theta} } { \EE^{f|i,T}_t \left[  1-\theta S^{d|i,\div}_{T,T;\theta} \right] }  \end{eqnarray} 
If $i = d$, then $\EE^{f|i,T}_t \left[  1-\theta S^{d|i,\div}_{T,T;\theta} \right] = 1-\theta S^{d|d,\div}_{t,T;\theta}$ is constant and $1 - \theta B^{f/d|d,r,\div}_{t,T;\theta}$ is an exponential martingale under $\QQ^{f|d,T}$. 

The additive currency basis spread can be recovered from the multiplicative one and $1+\theta L^{f|i}_{t,T;\theta}$. 
\end{rema}


\begin{rema}
In the DCB model, the currency basis spread was used to bootstrap the cross-currency discount curve $Z^{f|d}_{t,T}$. \\
How can we build this curve now? \\
Can we still bootstrap it? \\
The answer is yes, we can still bootstrap the cross-currency discount curve, but we have to bootstrap the collateral basis spot volatility at the same time\footnote{The collateral basis volatility is integrated up to the payment date $T+\theta$, as in the collateral convexity of Libors. An earlier version had $\Sigma^{f\to d}_{u,T;\theta}$, which misses most of the adjustment: about $94\%$ in the numerical example of the verification of \autoref{theo:qgcm}.}:
\begin{eqnarray}\label{eq:xccy_swaplet_master_equation}
  \theta B^{f/d|d,r}_{t,T;\theta} 
  &=& (1+\theta R^{f|d}_{t,T;\theta})(1 - \theta S^{d|d,\div}_{t,T;\theta})   
  - (1+\theta L^{f|f}_{t,T;\theta}) e^{-\int_t^T \Sigma^{f|f}_{u,T;\theta} \cdot \Sigma^{f\to d}_{u,T+\theta}\, du}. 
\end{eqnarray} 
so 
\begin{eqnarray}
  1+\theta R^{f|d}_{t,T;\theta} 
  &=& \frac{1+\theta R^{d|d}_{t,T;\theta}}{1+\theta L^{d|d}_{t,T;\theta}}  \left\{    
  (1+\theta L^{f|f}_{t,T;\theta}) e^{-\int_t^T \Sigma^{f|f}_{u,T;\theta} \cdot \Sigma^{f\to d}_{u,T+\theta} \, du} + \theta B^{f/d|d,r}_{t,T;\theta} 
  \right\}
\end{eqnarray} 

If we assuming that we can observe the currency basis spread for all relevant expiries, either directly from market quotes or through an interpolation scheme, then the only unknown is the collateral basis spot vol which needs to be calibrated.
\end{rema}



\newpage%
\subsection{Collateral convexity adjustments}%
We derive closed form formulas for collateral convexity adjustments. 
 
\subsubsection{Change of collateral formula}%
 
The following is a key formula for our computations.

\begin{lemm}
Let $(X,Y)$ be a Gaussian vector then 
\begin{eqnarray}
  \frac{\EE[e^X e^Y] } {\EE[e^X] \EE[e^Y] } &=& \exp\left( \Cov[X,Y] \right)  
\end{eqnarray} 
\end{lemm}
 
%The following proposition is a generalisation of FX-forward invariance under stochastic collateral rates. 
\begin{prop}[Collateral convexity]
Let $v^{e|i}_t$ be a lognormal process under $\QQ^{e|i,T}$ 
\begin{eqnarray}
  \frac{dv^{e|i}_{t}}{v^{e|i}_{t}} &=& \sigma^{v}_{t,T} \cdot dW^{e|i,T}_t 
\end{eqnarray} 
then the collateral convexity adjustment is given by 
\begin{eqnarray}
  \Q^{e|i\to j}_{t,T}[v^e_T] = \frac{\EE^{e|j,T}_t\left[ v^e_T \right] } {\EE^{e|i,T}_t\left[ v^e_T \right] } 
  &=& \exp\left\{ -\int_t^T \sigma^{v}_{u,T} \cdot \Sigma^{i\to j}_{u,T} \, du \right\}
\end{eqnarray} 
\end{prop}

\begin{proof}
We apply the previous lemma with $X = \ln v^{e|i}_{T,T}$ and $Y = \ln \Q^{e|i\to j}_{T,T}$. 
\end{proof}

\subsubsection{Collateral convexity of FX forwards}%
Applying the collateral convexity formula to FX forwards, we find.

%The following proposition is a generalisation of FX-forward invariance under stochastic collateral rates. 
\begin{prop}[Collateral convexity of FX forwards]
\begin{eqnarray}
        \Q^{d|i\to j}_{t,T}[x^{df}_T] =
	 \frac{x^{df|j}_{t,T}}{x^{df|i}_{t,T}} 
	&=& \exp\left\{ -\int_t^T \Sigma^{df}_{u,T}  \cdot \Sigma^{i\to j}_{u,T} \, du \right\}
\end{eqnarray} 
\end{prop}


\begin{rema}[Generalised FX forward invariance]
The previous result provides a generalisation of FX forward invariance. Indeed, under zero collateral basis volatility, the formula reduces to $x^{df|j}_{t,T}= x^{df|i}_{t,T}$. Under non zero collateral volatility, it still provides a formula for off-market discount curves by setting:
\begin{eqnarray}
	Z^{d|f}_{t,T} &=& Z^{d|d}_{t,T} \frac{Z^{f|f}_{t,T}}{Z^{f|d}_{t,T}} 
	\exp\left\{ \int_t^T \left( \Sigma^{df}_{u,T} - \Sigma^{f|f}_{u,T} + \Sigma^{d\to f}_{u,T} + \sigma^{df}_u \right) \cdot \Sigma^{d\to f}_{u,T}\, du  \right\}
\end{eqnarray} 
The first part is the usual formula from deterministic basis model. The exponential term quantifies the convexity adjustment ignored by the FX forward invariant construction. 

But one has to be careful though as the discount curve $Z^{f|d}_{t,T}$ that appears in this formula is not exactly the same as the curve $Z^{f|d,\DCB}_{t,T}$ built with the DCB model, because the former takes into account collateral convexity adjustment of Libors $L^{f|d}_{t,T;\theta}$ while the DCB model ignores it. 
\footnote{TODO: Q: Do the two convexity adjustments on $Z^{f|d}_{t,T}$ and $Z^{d|f}_{t,T} $ compensate each other or do they compound to make $Z^{f|d}_{t,T}$ diverge even more from $Z^{f|d,\DCB}_{t,T}$?}
\end{rema}


 
\subsubsection{Collateral convexity of Libors}%

Applying the collateral convexity formula to forward Libors, we find.

\begin{prop}
Under the deterministic multiplicative funding spread assumption, 
\begin{eqnarray}
 \Q^{e|i\to j}_{t,T+\theta} \left[1+\theta L^{e}_{t,T;\theta} \right] 
 = \frac{1+\theta L^{e|j}_{S,T;\theta} }{ 1+\theta L^{e|i}_{t,T;\theta} }
 &=&  \exp\left\{  - \int_t^T \Sigma^{e|e}_{u,T;\theta}  \cdot  \Sigma^{i\to j}_{u,T+\theta}\, du \right\}
\end{eqnarray} 
\end{prop}


%\begin{rema}
%In terms of pseudo discount factors, we find 
%\begin{eqnarray}
%	Z^{e|c}_{t,T;\theta}  
%	&=& Z^{e|e}_{t,T;\theta}   \exp\left(   \int_t^S (\Sigma^{e|e}_{u,S+\theta} - \Sigma^{e|e}_{u,S} )  \cdot  \Sigma^{e\to c}_{u,S+\theta}\, du \right) \\
%	&=& Z^{e|e}_{t,T;\theta}   \exp\left(   \int_t^S g^{e|e}(u) \left(\int_u^{S+\theta} h^{e|e}_v \, dv -\int_u^{S} h^{e|e}_v \, dv \right)  \cdot g^{e|e}(u) \int_u^{S+\theta} h^{e\to c}_{v}\,dv \, du \right) 
%\end{eqnarray} 
%\end{rema}

\begin{rema}[Collateral convexity of pseudo-discount factors]
Let us assume the existence of projection pseudo-discount factors:
\begin{eqnarray}
  1+\theta L^{e|c}_{t,T;\theta}   &=& \frac{ Z^{e,\theta|c}_{t,T}  }{ Z^{e,\theta|c}_{t,T+\theta} } .
\end{eqnarray} 
Can we decompose the convexity adjustment of forward Libors as a convexity adjustment of the pseudo-discount factors? Namely, can we find $\QZ^{e,\theta|e\to c}_{t,T}$ such that  
\begin{eqnarray}
  \QZ^{e,\theta|e\to c}_{t,T} = \frac{ Z^{e,\theta|c}_{t,T}  }{ Z^{e,\theta|e}_{t,T} } 
\end{eqnarray} 
and the above relation remains valid for all $T$? 

In other words, can we find $\QZ^{e,\theta|e\to c}_{t,T}$ such that 
\begin{eqnarray}
    \frac{ \QZ^{e,\theta|e\to c}_{t,T}  }{ \QZ^{e,\theta|e\to c}_{t,T+\theta} }  
    &=&  \exp\left(  - \int_t^T \Sigma^{e|e}_{u,T;\theta}  \cdot  \Sigma^{e\to c}_{u,T+\theta} \, du \right)
\end{eqnarray} 

Note that AOA does not determine the drift of projection pseudo-discount factors uniquely as they are not uniquely determined by asset prices. So we need to add some constraint on them. 

A natural extension of our funding spread hypothesis is to assume that the following ratio of \emph{native} (pseudo)-discount factors is constant:
\begin{eqnarray}
  \frac{ Z^{e,\theta|e}_{t,T}   }{ Z^{e|e}_{t,S}  }   &=& \frac{ Z^{e,\theta|e}_{t_0,T}   }{ Z^{e|e}_{t_0,S}  }  .
\end{eqnarray} 
%Note that this is only assumed to be valid under native collateral.
This assumption is equivalent to imposing the following dynamic\footnote{
We can check that this is compatible with the AOA condition on forward Libors, Namely, 
\begin{eqnarray}
  \frac{ d (1+\theta L^{e|e}_{t,T;\theta})   }{ (1+\theta L^{e|e}_{t,T;\theta}) }   &=&  \frac{ dZ^{e,\theta|e}_{t,T}   }{ Z^{e,\theta|e}_{t,T} }  - \frac{ dZ^{e,\theta|e}_{t,T+\theta}   }{ Z^{e,\theta|e}_{t,T+\theta} } - \left\langle \frac{ dZ^{e,\theta|e}_{t,T+\theta}   }{ Z^{e,\theta|e}_{t,T+\theta} } ,   \frac{ dZ^{e,\theta|e}_{t,T}   }{ Z^{e,\theta|e}_{t,T} }  - \frac{ dZ^{e,\theta|e}_{t,T+\theta}   }{ Z^{e,\theta|e}_{t,T+\theta} }  \right\rangle \\
  &=& \Sigma^{e|e}_{t,T;\theta} \cdot dW^{e,\RN}_t +  \Sigma^{e|e}_{t,T+\theta} \cdot  \Sigma^{e|e}_{t,T;\theta} \cdot dt \\
  &=& \Sigma^{e|e}_{t,T;\theta} \cdot \left( dW^{e,\RN}_t +  \Sigma^{e|e}_{t,T+\theta} dt  \right) \\
  &=& \Sigma^{e|e}_{t,T;\theta} \cdot dW^{e|e,T+\theta}_t  .
\end{eqnarray}
so forward Libors are indeed $\QQ^{e|e,T+\theta}$ as required by AOA.  
}:
\begin{eqnarray}
  \frac{ dZ^{e,\theta|e}_{t,T}   }{ Z^{e,\theta|e}_{t,T} }   ~=~  \frac{ dZ^{e|e}_{t,T}   }{ Z^{e|e}_{t,T} } ~=~ r^{e|e}_t \, dt - \Sigma^{e|e}_{t,T} \cdot dW^{e,\RN}_t 
\end{eqnarray}
However, this is not necessarily the most relevant or useful assumption.

TODO: Complete this investigation.
\end{rema}





 
\subsubsection{Collateral convexity of resetting forward Libors}%


\begin{prop}
The convexity adjustment on resetting forward Libors $L^{f/d|i}_{t,T;\theta}$ is
\begin{eqnarray}
   \Q_{t,T+\theta}^{f | i\to j}\left[ 1 + \theta L^{f/d|i}_{t,T;\theta}) \right] 
   &=&  \exp\left\{  - \int_t^S \Sigma^{f|d}_{u,T;\theta}  \cdot  \Sigma^{i\to j}_{u,T+\theta}\, du \right\}  \quad ???\end{eqnarray} 

\end{prop}

%
%\begin{proof}
%We use the dynamic of the funding spread to compute:
%\begin{eqnarray*}
%   1 + \theta L^{f/d|i}_{t,T;\theta} 
%   &=&  (1+\theta R^{f|i}_{t,T;\theta})  \EE^{f|i,T}_t\left[ 1-\theta S^{d|i,\div}_{T;\theta} \right]   \\
%   &=& (1+\theta R^{f|i}_{t,T;\theta})  \EE^{f|i,T}_t\left[ (1-\theta S^{d|i,\div}_{t,T;\theta}) \mathcal{E}\left(- \int_t^T \Sigma^{d\to i}_{u,T;\theta} \cdot dW^{d|i,T}_u \right) \right]   \\
%  &=& (1+\theta R^{f|i}_{t,T;\theta}) (1-\theta S^{d|i,\div}_{t,T;\theta}) e^{ - \int_t^T \Sigma^{d\to i}_{u,T;\theta} \cdot \Sigma^{df}_{u,T} \,du }  
%\end{eqnarray*} 
%We can now compute the volatility, knowing this is a $\QQ^{f|i,T+\theta}$-martingale:
%\begin{eqnarray*}
%    \frac{d(1 + \theta L^{f/d|i}_{t,T;\theta})}{1 + \theta L^{f/d|i}_{t,T;\theta} } 
%   &=& \frac{d(1 + \theta R^{f|i}_{t,T;\theta})}{1 + \theta R^{f|i}_{t,T;\theta}} 
%   +  \frac{d(1-\theta S^{d|i,\div}_{t,T;\theta})}{1-\theta S^{d|i,\div}_{t,T;\theta}} + O(dt) \\
%   &=& \Sigma^{f|i}_{t,T;\theta} \cdot dW^{f|i,T+\theta}_t  -  \Sigma^{d\to i}_{t,T;\theta} \cdot dW^{d|i,T}_t  + O(dt) \\
%   &=& \left( \Sigma^{f|f}_{t,T;\theta} + \Sigma^{f\to i}_{t,T;\theta} -   \Sigma^{d\to i}_{t,T;\theta}\right) \cdot dW^{f|i,T+\theta}_t  \\
%   &=&  \Sigma^{f|d}_{t,T;\theta} \cdot dW^{f|i,T+\theta}_t  
%\end{eqnarray*} 
%
%\end{proof}





 
\subsubsection{Collateral convexity of currency basis spreads}%


%\subsubsection{Collateral convexity for swap rates}
%
%Obviously we can compute the convexity adjustment for annuities and swap rates by decomposing them into bonds and Libors. Can we get something a bit more intuitive?
%
%[TODO]






\newpage%
\section{The Quasi-Gaussian Collateral Model}


\setcounter{subsection}{-1}
\subsection{Introduction}
The state variable in the generic HJM model is the whole forward curve which means that the dynamic is not a finite dimensional process in general\footnote{Here is a simple situation showing that the process $(r(t))$ is usually not Markov.

Assume that $\sigma(s,t)=t-s$ for every $s\leqslant t$ and every $\nu$, then elementary computations yield $$r(t)=F(0,t)+\tfrac18t^4+X_t$$ with $$X_t=\int_0^tW_s\mathrm ds$$ hence  the process $(r(t))$ is Markov if and only if the process $(X_t)$ is. But the process $(X_t)$ is not Markov (only the bivariate process $(W_t,X_t)$ is).}. The Cheyette separability assumption is well known to ensure a low dimensional Markovian dynamic in the classical framework. 

%\begin{itemize}
%\item We assume spot volatilities are separable: $\sigma_{t,T} = g_th_T$. % = \sigma_t e^{-\int_t^T \beta_v \, dv} $. 
%\item We express the dynamic in terms of correlated one dimensional Ornstein-Ulhenbeck processes driving each risk factor. 
%\end{itemize}

We show that this is still valid in the collateralised multi-curve setting. The resulting models, which we call the Quasi-Gaussian Collateral Models, are the natural generalisation of the Cheyette model to stochastic collateral/currency basis.


\subsection{Model specification}


\subsubsection{Separable volatility condition}%


\begin{defi}
An instantaneous forward rate volatility is \emph{separable} if there exist 
\begin{itemize} 
\item a stochastic process $g : \Omega \times \RR_+ \to M_d(\RR)$ 
\item a deterministic function $h : \RR_+ \to \RR^d$ 
\end{itemize}
such that 
\begin{eqnarray}
  \sigma_{t,T}(\omega)  &=& g(\omega,t)h(T).
  \label{eq:separability_condition}
\end{eqnarray}
\end{defi}

\begin{hypo}[Separable volatility]
The volatility of the following 
\begin{itemize} 
\item The instantaneous forward discount rates under native collateral
\item The instantaneous forward collateral basis spreads.
\end{itemize}
are separable.
\end{hypo}


\subsubsection{HJM dynamic}
Let us  recall the dynamic of the generic HJM model under the  domestic risk neutral measure:
\begin{eqnarray}
  dr^{d|d}_{t,T} &=& \left(\sigma^{d|d,\top}_{t,T} \Sigma^{d|d}_{t,T} \right) dt +  \sigma^{d|d,\top}_{t,T} dW^{d,\RN}_t   \\
  dr^{f|f}_{t,T} &=& \left(\sigma^{f|f,\top}_{t,T} \Sigma^{f|f}_{t,T} \right) dt +  \sigma^{f|f,\top}_{t,T} \left(dW^{d,\RN}_t -\sigma^{df}_t \,dt\right)  \\
  dq^{f|f\to d}_{t,T} &=& \left(\sigma^{f\to d,\top}_{t,T} \Sigma^{f\to d}_{t,T} + \sigma^{f\to d,\top}_{t,T} \Sigma^{f|f}_{t,T} + \sigma^{f|f,\top}_{t,T} \Sigma^{f\to d}_{t,T}\right) dt +  \sigma^{f\to d,\top}_{t,T} \left(dW^{d,\RN}_t -\sigma^{df}_t \,dt\right)
\end{eqnarray}
The drift of the collateral basis is the difference $\sigma^{f|d,\top}_{t,T}\Sigma^{f|d}_{t,T} - \sigma^{f|f,\top}_{t,T}\Sigma^{f|f}_{t,T}$ of the HJM drifts of $r^{f|d}_{t,T}$ and $r^{f|f}_{t,T}$ (both under $\QQ^{f,\RN}$), with $\sigma^{f|d} = \sigma^{f|f} + \sigma^{f\to d}$. In particular the term $\sigma^{f\to d,\top}_{t,T}\Sigma^{f\to d}_{t,T}$ appears once\footnote{An earlier version of this formula had $\sigma^{f|d,\top}_{t,T}\Sigma^{f\to d}_{t,T}$ as last term, which counts it twice.}.
So that, 
\begin{eqnarray}
  \frac{dZ^{d|d}_{t,T}}{Z^{d|d}_{t,T}} &=& r^{d|d}_t  dt - \Sigma^{d|d,\top}_{t,T} dW^{d,\RN}_t  \\
  \frac{dZ^{f|f}_{t,T}}{Z^{f|f}_{t,T}} 
  &=& r^{f|f}_t  dt - \Sigma^{f|f,\top}_{t,T} \left( dW^{d,\RN}_t -\sigma^{df}_t dt\right) \\
  \frac{dZ^{f|d}_{t,T}}{Z^{f|f}_{t,T}} 
  &=&  \left(r^{f|f}_t + q^{f\to d}_t \right)  dt 
  - \left(\Sigma^{f|f}_{t,T} + \Sigma^{f\to d}_{t,T}  \right)^\top \left( dW^{d,\RN}_t -\sigma^{df}_t dt\right) \\
     \frac{dx^{df}_{t}}{x^{df}_{t}} &=& \left(r^{d|d}_t - r^{f|f}_t - q^{f\to d}_t \right) dt + \sigma^{df,\top}_{t} dW^{d,\RN}_t  
\end{eqnarray}

\subsubsection{Notations}
The model is stationary iff the volatilities are a function of time to expiry $\tau = T-t$ rather than both $t$ and $T$. This suggests changing the parametrisation\footnote{This is also sensible for numerical reasons since, under standard parametrisation, the coefficients of $g$ and $h$ tend to have exponential growth.}. 

For each symbol $a$, let us write
\begin{eqnarray}
  H^a_{t} &=& \diag(h^a_t)  \\
  \sigma^a_{z,s} &=& g^a_sH^a_s  \\ 
%  \sigma^a_{s} &=& \diag\left( \sigma^{a,\top}_{z,s}  \sigma^a_{z,s}\right)^{1/2}  \\ 
  M^a_{t,T} &=& H^a_T H^{a,-1}_t \vun = H^{a,-1}_t h^a_T  \\
  G^a_{t,T} &=& \int_t^T M^a_{t,u} \, du = H^{a,-1}_t \int_t^T h^a_u \, du %\\
%  \Omega_t &=& \sigma^{a,-1}_{s}  \left( \sigma^{a,\top}_{z,s}  \sigma^a_{z,s}\right)  \sigma^{a,-1}_{s}   \\
%  \Gamma^\top_t\Gamma_t &=& \Omega_t 
\end{eqnarray}
Let us define the mean reversion speeds as the coefficients of the diagonal matrix: 
\begin{eqnarray}
  \beta^a_{t} &=&  - H^{a,-1}_{t} \frac{d}{dt}H^a_{t}  
\end{eqnarray}
One can represent the matrices more explicitly:
\begin{eqnarray}
  H^a_{t} &=& \diag\left( e^{-\int_0^t \beta^{a,1}_u \, du} , \ldots, e^{-\int_0^t \beta^{a,N_a}_u \, du} \right)  \\
%  \sigma^a_{t} &=& \diag\left( \sigma^{a,1}_t  , \ldots, \sigma^{a,N_a}_t \right)  \\
  M^a_{t,T} &=& \left[ e^{-\int_t^T \beta^{a,1}_u \, du} , \ldots, e^{-\int_t^T \beta^{a,N_a}_u \, du} \right]^\top  \\
  G^a_{t,T} &=& \left[ \int_t^T e^{-\int_t^s \beta^{a,1}_u \, du} \, ds , \ldots,  \int_t^T e^{-\int_t^s \beta^{a,N_a}_u \, du} \, ds \right]^\top  
\end{eqnarray}
%\begin{eqnarray}
%  H^a_{t}  &=& 
%  \begin{bmatrix}
%     e^{-\int_0^t \beta^{a,1}_u \, du}  & 0   &   \\
%      0 &  e^{-\int_0^t \beta^{a,2}_u \, du}  & &    \\
%       &  &  \ddots  &     \\
%     &    & &  e^{-\int_0^t \beta^{a,N_a}_u \, du}    \\      
%   \end{bmatrix}  \\
%   \sigma^a_{t}   &=& 
%   \begin{bmatrix}
%      \sigma^{a,1}_t & 0   &   \\
%      0 &   \sigma^{a,2}_t & &    \\
%       &  &  \ddots  &     \\
%     &    & &   \sigma^{a,N_a}_t   \\      
%   \end{bmatrix}  \\
%   M^a_{t,T}  &=& 
%   \begin{bmatrix}
%      e^{-\int_t^T \beta^{a,1}_u \, du} \\
%      e^{-\int_t^T \beta^{a,2}_u \, du} \\
%      \vdots \\
%      e^{-\int_t^T \beta^{a,N_a}_u \, du} \\
%   \end{bmatrix}  \\
%   G^a_{t,T}   &=& 
%   \begin{bmatrix}
%      \int_t^T e^{-\int_t^s \beta^{a,1}_u \, du} \, ds \\
%      \int_t^T e^{-\int_t^s \beta^{a,2}_u \, du} \, ds  \\
%      \vdots \\
%      \int_t^T e^{-\int_t^s \beta^{a,N_a}_u \, du} \, ds \\
%   \end{bmatrix}  \\
%\end{eqnarray}

\begin{rema}
Moreover, one can reformulate in terms of 1D OU processes by setting:
\begin{eqnarray}
  \sigma^a_{s} &=& \diag\left( \sigma^{a,\top}_{z,s}  \sigma^a_{z,s}\right)^{1/2}  \\ 
  \Omega_t &=& \sigma^{-1}_{s}  \left( \sigma^{\top}_{z,s}  \sigma_{z,s}\right)  \sigma^{-1}_{s}   \\
  \Gamma^\top_t\Gamma_t &=& \Omega_t 
\end{eqnarray}
The factor loading matrix $\Gamma_{t} \in M_{d,k}(\RR)$ is a square root of the instantaneous correlation matrix $\Omega_t$ whose bloc decomposition looks like:
\begin{eqnarray}
  \Omega_t ~=~ \Gamma_{t}^\top\Gamma ~=~ 
  \begin{bmatrix}
     \rho^{d,d}_{t} &   \rho^{d,f}_{t}  &  \rho^{d,c}_{t}  &  \rho^{d,x}_{t}  \\
     \rho^{f,d}_{t} &   \rho^{f,f}_{t}  &  \rho^{f,c}_{t}  &  \rho^{f,x}_{t}  \\
     \rho^{c,d}_{t} &   \rho^{c,f}_{t}  &  \rho^{c,c}_{t}  &  \rho^{c,x}_{t}  \\
     \rho^{x,d}_{t} &   \rho^{x,f}_{t}  &  \rho^{x,c}_{t}  &  \rho^{x,x}_{t}  \\
      \end{bmatrix} 
\end{eqnarray}

This leads to 
\begin{eqnarray}
  \sigma^{a,\top}_{t,T} dW_t 
  &=& H^{a}_T H^{a,-1}_t \sigma^a_{t} \Pi_a^\top \underbrace{\Gamma_t^\top dW_t}_{d\widetilde{W}_t} \\
   \Sigma^{a,\top}_{t,T} dW_t 
  &=& G^{a}_{t,T} \sigma^a_{t} \Pi_a^\top \, d\widetilde{W}_t 
\end{eqnarray}

\end{rema}



\subsubsection{Dynamic}

We use the shorthands $d$, $f$, $c$ and $x$ for the blocks $d|d$, $f|f$, $f\to d$ and $df$, and write $W_t = W^{d,\RN}_t$. For each rates block $a\in\{d,f,c\}$, the volatility of instantaneous forwards is $\sigma^{a}_{t,T} = \sigma^a_{z,t} M^a_{t,T}$ with $\sigma^a_{z,t} = g^a_t H^a_t \in M_{k,N_a}(\RR)$, so that $\Sigma^{a}_{t,T} = \sigma^a_{z,t} G^a_{t,T}$. We also write
\begin{eqnarray}
  \mathcal{M}^a_{u,t} &=& \diag\left(M^a_{u,t}\right) ~=~ H^a_t H^{a,-1}_u , \qquad u\le t,
\end{eqnarray}
for the diagonal transport matrix. It satisfies $\mathcal{M}^a_{u,t}\mathcal{M}^a_{t,T} = \mathcal{M}^a_{u,T}$, $\partial_t \mathcal{M}^a_{u,t} = -\beta^a_t \mathcal{M}^a_{u,t}$ and $\partial_t G^a_{u,t} = \mathcal{M}^a_{u,t}\vun$, and the forward volatilities split as $M^a_{u,T} = \mathcal{M}^a_{u,t} M^a_{t,T}$ and $G^a_{u,T} = G^a_{u,t} + \mathcal{M}^a_{u,t} G^a_{t,T}$.

\begin{theo}[Quasi-Gaussian Collateral Model]\label{theo:qgcm}
Let us consider an arbitrage free model with instantaneous forward rates satisfying the separability condition. Let us define the vector-valued state variables
\begin{eqnarray}
  z^{d}_t &=& \int_0^t \mathcal{M}^d_{u,t} \sigma^{d,\top}_{z,u} \left(  \sigma^{d}_{z,u} G^d_{u,t} \, du + dW_u \right) \\
  z^{f}_t &=& \int_0^t \mathcal{M}^f_{u,t} \sigma^{f,\top}_{z,u} \left(  \sigma^{f}_{z,u} G^f_{u,t} \, du - \sigma^{df}_u \, du + dW_u \right) \\
  z^{c}_t &=& \int_0^t \mathcal{M}^c_{u,t} \sigma^{c,\top}_{z,u} \left(  \sigma^{c}_{z,u} G^c_{u,t} \, du + \sigma^{f}_{z,u} G^f_{u,t} \, du - \sigma^{df}_u \, du + dW_u \right) \\
  z^{x}_t &=& \ln \left( x^{df}_t / x^{df|d}_{0,t} \right)
\end{eqnarray}
together with the matrix-valued processes ($(a,b) \in \{(d,d),(f,f),(c,c),(c,f)\}$, $y^{f,c}_t = y^{c,f,\top}_t$)
\begin{eqnarray}
  y^{a,b}_t &=& \int_0^t \mathcal{M}^a_{u,t} \sigma^{a,\top}_{z,u} \sigma^{b}_{z,u} \mathcal{M}^b_{u,t} \, du
\end{eqnarray}
and the vector-valued process
\begin{eqnarray}
  v^{f,c}_t &=& \int_0^t \mathcal{M}^f_{u,t} \sigma^{f,\top}_{z,u} \sigma^{c}_{z,u} G^c_{u,t} \, du .
\end{eqnarray}
Then, the model is Markovian in these state variables with dynamic under $\QQ^{d,\RN}$
\begin{eqnarray}
  dz^{d}_t &=& \left(y^{d,d}_t\vun_d - \beta^{d}_t z^{d}_t \right)dt +  \sigma^{d,\top}_{z,t} dW_t \\
  dz^{f}_t &=& \left(y^{f,f}_t\vun_f - \sigma^{f,\top}_{z,t}\sigma^{df}_t - \beta^{f}_t z^{f}_t \right)dt +  \sigma^{f,\top}_{z,t} dW_t \\
  dz^{c}_t &=& \left(y^{c,c}_t\vun_c + y^{c,f}_t\vun_f - \sigma^{c,\top}_{z,t}\sigma^{df}_t - \beta^{c}_t z^{c}_t \right)dt +  \sigma^{c,\top}_{z,t} dW_t \\
  dz^{x}_t &=&  \left(\vun_d^\top z^{d}_t - \vun_f^\top z^{f}_t - \vun_c^\top z^{c}_t - \vun_f^\top v^{f,c}_t - \frac{1}{2}|\sigma^{df}_t|^2\right)dt + \sigma^{df,\top}_t dW_t  \\
  dy^{a,b}_t &=&   \left(\sigma^{a,\top}_{z,t} \sigma^{b}_{z,t} - \beta^{a}_t y^{a,b}_t - y^{a,b}_t \beta^{b}_t \right) dt  \\
  dv^{f,c}_t &=&   \left( y^{f,c}_{t}\vun_c - \beta^{f}_{t} v^{f,c}_t\right) dt
\end{eqnarray}
and the following reconstruction formulas:
\begin{enumerate}
\item Instantaneous forward rates:
\begin{eqnarray}
  r^{d|d}_{t,T} &=& r^{d|d}_{0,T} + M^{d,\top}_{t,T}\left( z^{d}_t + y^{d,d}_t G^{d}_{t,T} \right) \\
  r^{f|f}_{t,T} &=& r^{f|f}_{0,T} + M^{f,\top}_{t,T}\left( z^{f}_t + y^{f,f}_t G^{f}_{t,T} \right) \\
  q^{f|f\to d}_{t,T} &=& q^{f|f\to d}_{0,T} +  M^{c,\top}_{t,T}\left( z^{c}_t + y^{c,c}_t G^{c}_{t,T} + y^{c,f}_t G^{f}_{t,T} \right) +  M^{f,\top}_{t,T}\left( v^{f,c}_t + y^{f,c}_t G^{c}_{t,T} \right)
\end{eqnarray}
\item Zero coupon bonds and collateral basis curve $\QZ^{f|f\to d}_{t,T} = Z^{f|d}_{t,T}/Z^{f|f}_{t,T}$:
\begin{eqnarray}
  Z^{d|d}_{t,T} &=& Z^{d|d}_{0,t,T} \exp\left( -G^{d,\top}_{t,T} z^{d}_t -\frac{1}{2} G^{d,\top}_{t,T}  y^{d,d}_t G^{d}_{t,T} \right)  \\
  Z^{f|f}_{t,T} &=& Z^{f|f}_{0,t,T} \exp\left( -G^{f,\top}_{t,T} z^{f}_t -\frac{1}{2} G^{f,\top}_{t,T}  y^{f,f}_t G^{f}_{t,T} \right)  \\
  \QZ^{f|f\to d}_{t,T} &=& \QZ^{f|f\to d}_{0,t,T} \exp\left( -G^{c,\top}_{t,T} z^{c}_t - G^{f,\top}_{t,T} v^{f,c}_t - \frac{1}{2} G^{c,\top}_{t,T}  y^{c,c}_t G^{c}_{t,T} - G^{c,\top}_{t,T}  y^{c,f}_t G^{f}_{t,T}  \right)
  \label{eq:qgcm_collateral_basis_reconstruction}
\end{eqnarray}
\item The FX spot: $x^{df}_{t} = x^{df|d}_{0,t} e^{z^{x}_t}$.
\end{enumerate}
\end{theo}

\begin{rema}
The state variables have a very simple interpretation:
\begin{eqnarray*}
  \vun_d^\top z^d_t &=& r^{d|d}_t - r^{d|d}_{0,t} \\
  \vun_f^\top z^f_t &=& r^{f|f}_t - r^{f|f}_{0,t} \\
  \vun_c^\top z^{c}_t  + \vun_f^\top v^{f,c}_t &=& q^{f\to d}_t - q^{f|f\to d}_{0,t} \\
  z^{x}_t &=& \ln\left( x^{df}_t / x^{df|d}_{0,t} \right).
\end{eqnarray*}
So they are essentially the short rates and FX log moneyness forward spot with their initial conditions factored out.

The additional state variables $y_t$ look like\footnote{Assuming the volatilities are $\sigma_{t,T} = \sigma_t e^{-\int_t^T \beta_v\, dv}$.}
\begin{eqnarray}
  y^{(a,i),(b,j)}_{t} &=& \int_0^t \rho^{(a,i),(b,j)}_s \sigma^{a,i}_s  \sigma^{b,j}_s  e^{-\int_s^t \beta^{a,i}_u + \beta^{b,j}_u \, du}  \, ds
\end{eqnarray}
They correspond to the evolution of covariance in the model and it is necessary to keep track of them in order to keep the dynamic Markovian. However, if the model parameters are deterministic, then so are the $y_t$'s and $v^{f,c}_t$ and they can be pre-computed instead of diffused (see next section).
\end{rema}

\begin{rema}[Quanto drifts]
The quanto drift adjustments $-\sigma^{f,\top}_{z,t}\sigma^{df}_t$ and $-\sigma^{c,\top}_{z,t}\sigma^{df}_t$ are instantaneous: they do not require additional state variables. Only the cross-block covariance $y^{c,f}_t$ between collateral basis and foreign rates has to be tracked, together with $v^{f,c}_t$. The covariances $y^{d,f}_t$, $y^{d,c}_t$ and the FX covariances are not needed.
\end{rema}

\begin{rema}[The collateral basis is not driven by its own Cheyette state]\label{rema:qgcm_vfc}
The answer to the question raised in the proof below, namely whether the collateral basis is driven by an OU process in its own factors only, is negative: the term $\sigma^{f,\top}_{t,T}\Sigma^{c}_{t,T}$ in the drift of $q^{f|f\to d}_{t,T}$ has the maturity dependence $M^{f}_{t,T}$ of the \emph{foreign} factors. This is what the extra finite-variation state $v^{f,c}_t$ accounts for.

However, the stacked state
\begin{eqnarray}
  \hat{z}_t = \begin{bmatrix} z^f_t + v^{f,c}_t \\ z^c_t \end{bmatrix}, \quad
  \hat{G}_{t,T} = \begin{bmatrix} G^f_{t,T} \\ G^c_{t,T} \end{bmatrix}, \quad
  \hat{y}_t = \begin{bmatrix} y^{f,f}_t & y^{f,c}_t \\ y^{c,f}_t & y^{c,c}_t \end{bmatrix} ,
\end{eqnarray}
is exactly a Cheyette state for the cross-currency discount rate $r^{f|d}_{t,T} = r^{f|f}_{t,T} + q^{f|f\to d}_{t,T}$: it has dynamic $d\hat{z}_t = (\hat{y}_t\vun - \hat{\sigma}^\top_{z,t}\sigma^{df}_t - \hat{\beta}_t\hat{z}_t)\,dt + \hat{\sigma}^\top_{z,t} dW_t$ with $\hat{\sigma}_{z} = [\sigma^f_z | \sigma^c_z]$ and
\begin{eqnarray}
  Z^{f|d}_{t,T} &=& Z^{f|d}_{0,t,T} \exp\left( -\hat{G}^{\top}_{t,T} \hat{z}_t - \frac{1}{2} \hat{G}^{\top}_{t,T}  \hat{y}_t \hat{G}_{t,T} \right) .
  \label{eq:qgcm_zfd_reconstruction}
\end{eqnarray}
So the two choices of state considered in the proof below ($z^{c}_t$ driving $\QZ^{f|f\to d}_{t,T}$, or $\hat{z}_t$ driving $Z^{f|d}_{t,T}$) differ only by the shift $v^{f,c}_t$. In the Quasi-Gaussian model $v^{f,c}_t$ is stochastic and must be carried as an extra state in either case, since $Z^{f|f}_{t,T}$ needs $z^f_t$ and $Z^{f|d}_{t,T}$ needs $z^f_t + v^{f,c}_t$. In the LGCM it is deterministic and the two choices carry the same information.
\end{rema}


\begin{rema}
Let us comment our parametrisation choices. 

We chose to keep the quanto-drift adjustments inside the dynamic of the state variables. This ensures the definition of $z^{f|f}_t$ does not depend on the choice of the domestic currency $d$. But it implies that its dynamic under the domestic risk neutral measure does. 

On the other hand, we chose to consider $z^{f|f\to d}_t$ as driving the collateral basis $\QZ^{f|f\to d}_{t,T}$ rather than $\QZ^{d|d\to f}_{t,T}$ or some other observable. This choice is motivated by the calibration of the model to the cross-currency market. Indeed, the most liquid instruments are collateralised in USD so it is natural to chose it as the domestic currency $d$ of the model. But this implies that the first non-native discount curve to calibrate is $Z^{f|d}_{t,T}$. Factoring out the native discount curve $Z^{f|f}_{t,T}$ which is known from the native foreign market, we are led to consider $\QZ^{f|f\to d}_{t,T}$ as containing mostly new information from the cross-currency market.
\end{rema}


\subsection{Proof of the theorem}


\subsubsection{Domestic rates}
Let us inject the separability condition in the dynamic of the instantaneous forward rate.
\begin{eqnarray*}
  r^{d|d}_{t,T} - r^{d|d}_{0,T} &=&  \int_0^t \sigma^{d|d,\top}_{u,T} \Sigma^{d|d}_{u,T} \, du +  \int_0^t \sigma^{d|d,\top}_{u,T} dW^{d,\RN}_u   \\
  &=& \int_0^t h^{d|d,\top}_{T} g^{d|d,\top}_{u} \int_u^T g^{d|d}_{u} h_v \, dv  \, du 
  +  \int_0^t h^{d|d,\top}_{T} g^{d|d,\top}_{u} dW^{d,\RN}_u   \\
  &=& r^{d|d}_{0,T} + h^{d|d,\top}_{T} \int_0^t  g^{d|d,\top}_{u} g^{d|d}_{u} \int_u^T  h_v \, dv  \, du 
  + h^{d|d,\top}_{T} \int_0^t  g^{d|d,\top}_{u} dW^{d,\RN}_u   
\end{eqnarray*}
By writing $h^{d|d}_t =  H^{d|d}_t\vun_d$ and $g^{d|d}_t = \sigma^{d|d}_{z,t} H^{d|d,-1}_t$, this can be decomposed into 
\begin{eqnarray*}
  r^{d|d}_{t,T}  - r^{d|d}_{0,T} 
  &=&  \vun_d^\top H^{d|d}_{T}\int_0^t H^{d|d,-1}_{u}  \sigma^{d|d,\top}_{z,u}   \sigma^{d|d}_{z,u}  H^{d|d,-1}_{u} \int_u^T    h^{d|d}_{v}  \, dv   \, du  \\
  & &   ~~ + ~~ \vun_d^{\top} H^{d|d}_{T}   \int_0^t   H^{d|d,-1}_{u} \sigma^{d|d,\top}_{z,u}  dW^{d,\RN}_u   \\
  &=&  \vun_d^\top H^{d|d}_{T} \int_0^t H^{d|d,-1}_{u} \sigma^{d|d,\top}_{z,u} \sigma^{d|d}_{z,u}  H^{d|d,-1}_{u}    \int_u^T   h^{d|d}_{v} \, dv   \, du  \\
  & &   ~~ + ~~ \vun_d^{\top} H^{d|d}_{T}   \int_0^t   H^{d|d,-1}_{u} \sigma^{d|d,\top}_{z,u}  dW^{d,\RN}_u   \end{eqnarray*}
When $T=t$, we find the expression for the short rate:
\begin{eqnarray*}
  r^{d|d}_{t}  - r^{d|d}_{0,t} 
  &=&  \vun_d^\top H^{d|d}_{t}\int_0^t H^{d|d,-1}_{u}  \sigma^{d|d,\top}_{z,u}   \sigma^{d|d}_{z,u}  H^{d|d,-1}_{u} \int_u^t    h^{d|d}_{v}  \, dv   \, du  \\
  & &   ~~ + ~~ \vun_d^{\top} H^{d|d}_{t}   \int_0^t   H^{d|d,-1}_{u} \sigma^{d|d,\top}_{z,u}  dW^{d,\RN}_u   \\
  &=& \vun_d^{\top} z^{d|d}_t 
\end{eqnarray*}
Now coming back to the forward rate, we isolate the dependence in $T$ and introduce $z^{d|d}_t$: 
\begin{eqnarray*}
  & & r^{d|d}_{t,T}  - r^{d|d}_{0,T}  \\
  &=& \vun_d^\top H^{d|d}_{T}H^{d|d,-1}_{t} H^{d|d}_{t} \int_0^t  H^{d|d,-1}_{u}\sigma^{d|d,\top}_{z,u} \sigma^{d|d}_{z,u} H^{d|d,-1}_{u} \left(\int_u^t   h^{d|d}_{v} \, dv + \int_t^T   h^{d|d}_{v} \, dv \right) \, du \\
  & &   ~~ + ~~ \vun_d^{\top} H^{d|d}_{T}   \int_0^t   H^{d|d,-1}_{u} \sigma^{d|d,\top}_{z,u}  dW^{d,\RN}_u   \\
  &=& \vun_d^\top H^{d|d}_{T}H^{d|d,-1}_{t} \left(  H^{d|d}_{t} \int_0^t  H^{d|d,-1}_{u}\sigma^{d|d,\top}_{z,u} \sigma^{d|d}_{z,u} H^{d|d,-1}_{u} \int_u^t   h^{d|d}_{v} \, dv  \, du \right) \\
  & &  ~~ + ~~ \vun_d^\top H^{d|d}_{T}H^{d|d,-1}_{t} \left( H^{d|d}_{t} \int_0^t  H^{d|d,-1}_{u}\sigma^{d|d,\top}_{z,u} \sigma^{d|d}_{z,u} H^{d|d,-1}_{u}   \, du \, H^{d|d}_{t}\right) \left( H^{d|d,-1}_{t}  \int_t^T   h^{d|d}_{v} \, dv  \right) \\
  & &   ~~ + ~~ \vun_d^{\top} H^{d|d}_{T}   \int_0^t   H^{d|d,-1}_{u} \sigma^{d|d,\top}_{z,u}  dW^{d,\RN}_u   \\
  &=& M^{d|d,\top}_{t,T} \left(  H^{d|d}_{t} \int_0^t  H^{d|d,-1}_{u}\sigma^{d|d,\top}_{z,u} \sigma^{d|d}_{z,u} H^{d|d,-1}_{u} \int_u^t   h^{d|d}_{v} \, dv  \, du + H^{d|d}_{t} \int_0^t   H^{d|d,-1}_{u} \sigma^{d|d,\top}_{z,u}  dW^{d,\RN}_u \right) \\
  & &  ~~ + ~~ M^{d|d,\top}_{t,T} \left( H^{d|d}_{t} \int_0^t  H^{d|d,-1}_{u}\sigma^{d|d,\top}_{z,u} \sigma^{d|d}_{z,u} H^{d|d,-1}_{u}   \, du \, H^{d|d}_{t}\right) \left( H^{d|d,-1}_{t}  \int_t^T   h^{d|d}_{v} \, dv  \right) \\
  &=& M^{d|d,\top}_{t,T} \left( z^{d|d}_{t}  +  y^{d|d}_t G^{d|d}_{t,T}  \right)
\end{eqnarray*}


%We can rewrite this as 
%\begin{eqnarray}
%  r^{d|d}_{t,T}  &=& r^{d|d}_{0,T} + \underbrace{\vun_d^\top H^{d|d}_{T}H^{d|d,-1}_{t}}_{M^{d|d,\top}_{t,T}}\left( z^{d|d}_{t}  +  y^{d|d}_{t} G^{d|d}_{t,T}   \right) \\
%  z^{d|d}_{t}  &=& H^{d|d}_{t} \int_0^t  \widetilde{\sigma}^{d|d,\top}_{u}  \underbrace{\Pi^{d,\top}  \Gamma_t^\top \Gamma_t  \Pi^d}_{\Omega^{d,d}_t}  \widetilde{\sigma}^{d|d}_{u} \int_u^t h^{d|d}_v  \,dv \, du  \\
%  & &   + H^{d|d}_{t} \int_0^t   \widetilde{\sigma}^{d|d}_{u} \underbrace{\Pi^{d,\top} \Gamma_t^\top dW^{d,\RN}_u}_{d\widetilde{W}^{d,\RN,d}_u}   \\
%  y^{d|d}_t &=& H^{d|d}_{t} \int_0^t  \widetilde{\sigma}^{d|d}_{u}  \underbrace{\Pi^{d,\top}  \Gamma_t^\top \Gamma_t  \Pi^d}_{\Omega^{d,d}_t} \widetilde{\sigma}^{d|d,\top}_{u}  \, du\, H^{d|d}_t \\
%  G^{d|d}_{t,T} &=& H^{d|d,-1}_{t}  \int_t^T H^{d|d}_{v} \,dv  \vun
%\end{eqnarray}

\subsubsection{Foreign rates}
Under $\QQ^{d,\RN}$, the foreign forward rate has the extra quanto drift $-\sigma^{f,\top}_{u,T}\sigma^{df}_u$. With $\sigma^{f,\top}_{u,T} = M^{f,\top}_{t,T}\mathcal{M}^f_{u,t}\sigma^{f,\top}_{z,u}$, the same decomposition as for domestic rates gives
\begin{eqnarray*}
  r^{f|f}_{t,T}  - r^{f|f}_{0,T}
  &=& \int_0^t \sigma^{f,\top}_{u,T}\left( \Sigma^{f}_{u,T} - \sigma^{df}_u \right) du + \int_0^t \sigma^{f,\top}_{u,T} dW_u \\
  &=& M^{f,\top}_{t,T} \int_0^t \mathcal{M}^f_{u,t} \sigma^{f,\top}_{z,u} \left( \sigma^{f}_{z,u}\left(G^f_{u,t} + \mathcal{M}^f_{u,t}G^f_{t,T}\right) du - \sigma^{df}_u\, du +  dW_u \right) \\
  &=& M^{f,\top}_{t,T} \left( z^{f}_{t}  +  y^{f,f}_t G^{f}_{t,T}  \right) .
\end{eqnarray*}
The quanto term only contributes to $z^f_t$ and does not depend on $T$ beyond the factor $M^{f}_{t,T}$, which is why it does not require an additional state variable.

\subsubsection{Collateral basis}
The drift of $q^{f|f\to d}_{t,T}$ under $\QQ^{d,\RN}$ is $\sigma^{c,\top}_{u,T}(\Sigma^c_{u,T} + \Sigma^f_{u,T} - \sigma^{df}_u) + \sigma^{f,\top}_{u,T}\Sigma^{c}_{u,T}$. The first part has maturity dependence $M^c_{t,T}$ and is handled exactly like the foreign rate:
\begin{eqnarray*}
  & & \int_0^t \sigma^{c,\top}_{u,T}\left(\Sigma^c_{u,T} + \Sigma^f_{u,T} - \sigma^{df}_u\right) du + \int_0^t \sigma^{c,\top}_{u,T} dW_u \\
  &=& M^{c,\top}_{t,T} \int_0^t \mathcal{M}^c_{u,t}\sigma^{c,\top}_{z,u} \left( \sigma^c_{z,u}(G^c_{u,t} + \mathcal{M}^c_{u,t}G^c_{t,T}) du +  \sigma^f_{z,u}(G^f_{u,t} + \mathcal{M}^f_{u,t}G^f_{t,T}) du - \sigma^{df}_u\, du + dW_u \right) \\
  &=& M^{c,\top}_{t,T} \left( z^{c}_{t}  +  y^{c,c}_t G^{c}_{t,T} + y^{c,f}_t G^{f}_{t,T} \right) .
\end{eqnarray*}
The second part has maturity dependence $M^f_{t,T}$:
\begin{eqnarray*}
  \int_0^t \sigma^{f,\top}_{u,T}\Sigma^c_{u,T} \, du
  &=& M^{f,\top}_{t,T} \int_0^t \mathcal{M}^f_{u,t}\sigma^{f,\top}_{z,u} \sigma^c_{z,u}\left(G^c_{u,t} + \mathcal{M}^c_{u,t}G^c_{t,T}\right) du
  ~=~ M^{f,\top}_{t,T} \left( v^{f,c}_t + y^{f,c}_t G^c_{t,T} \right) .
\end{eqnarray*}
Summing the two gives the reconstruction formula for $q^{f|f\to d}_{t,T}$.

\begin{rema}
An alternative is to assume that $r^{f|d}_{t,T}$ rather than $q^{f|f\to d}_{t,T}$ has separable volatility. Since $\sigma^{f|d}_{t,T} = \sigma^{f}_{t,T} + \sigma^{c}_{t,T} = [\sigma^f_{z,t} | \sigma^c_{z,t}]\, [M^{f}_{t,T} ; M^c_{t,T}]$, the two assumptions are the same. The corresponding Cheyette state for $r^{f|d}_{t,T}$ is the stacked state $\hat{z}_t$ of \autoref{rema:qgcm_vfc}.
\end{rema}

\subsubsection{FX spot}
Since $\ln x^{df|d}_{0,t} = \ln x^{df}_0 + \int_0^t \left( r^{d|d}_{0,u} - r^{f|f}_{0,u} - q^{f|f\to d}_{0,u}\right) du$ and the FX spot has drift $r^{d|d}_t - r^{f|f}_t - q^{f\to d}_t$, we get
\begin{eqnarray*}
  dz^x_t &=& \left( (r^{d|d}_t - r^{d|d}_{0,t}) - (r^{f|f}_t - r^{f|f}_{0,t}) - (q^{f\to d}_t - q^{f|f\to d}_{0,t}) - \frac{1}{2}|\sigma^{df}_t|^2 \right)dt + \sigma^{df,\top}_t dW_t
\end{eqnarray*}
and we conclude with the reconstruction formulas for the short rates ($T = t$).

\subsubsection{Zero coupon bonds}
Since $\partial_T G_{t,T} = M_{t,T}$, for a symmetric matrix $y$ we have $\int_t^T M^\top_{t,u} y G_{t,u}\, du = \frac{1}{2}G_{t,T}^\top y G_{t,T}$, and
\begin{eqnarray*}
  \int_t^T \left( M^{c,\top}_{t,u} y^{c,f}_t G^f_{t,u} + M^{f,\top}_{t,u} y^{f,c}_t G^c_{t,u} \right) du &=& G^{c,\top}_{t,T} y^{c,f}_t G^f_{t,T} .
\end{eqnarray*}
Integrating the instantaneous forward rates from $t$ to $T$ gives the ZCB formulas, with $Z_{0,t,T} = Z_{0,T}/Z_{0,t}$. Multiplying the formulas for $Z^{f|f}_{t,T}$ and $\QZ^{f|f\to d}_{t,T}$ gives \eqref{eq:qgcm_zfd_reconstruction}.

\subsubsection{Dynamic of the state variables}
The dynamic follows by differentiating the integral definitions with respect to $t$, using $\partial_t \mathcal{M}_{u,t} = -\beta_t\mathcal{M}_{u,t}$, $\partial_t G_{u,t} = \mathcal{M}_{u,t}\vun$ and $G_{t,t}=0$. For example,
\begin{eqnarray*}
  \frac{d}{dt}\int_0^t \mathcal{M}^c_{u,t}\sigma^{c,\top}_{z,u}\sigma^f_{z,u} G^f_{u,t}\, du
  &=& -\beta^c_t \int_0^t \mathcal{M}^c_{u,t}\sigma^{c,\top}_{z,u}\sigma^f_{z,u} G^f_{u,t}\, du + \int_0^t \mathcal{M}^c_{u,t}\sigma^{c,\top}_{z,u}\sigma^f_{z,u} \mathcal{M}^f_{u,t}\, du \, \vun_f ,
\end{eqnarray*}
which produces the term $y^{c,f}_t\vun_f$ in the drift of $z^c_t$. This completes the proof of \autoref{theo:qgcm}.

\begin{rema}[Numerical verification]
The theorem was checked in the deterministic-volatility case (where the state is Gaussian) with two domestic factors, one foreign factor, one collateral basis factor and the FX factor. The parameters were piecewise constant with zero, $10^{-7}$ and negative mean reversions on some intervals, under a full $5\times 5$ correlation matrix. The mean and covariance of $(z_t, \int_0^t (r^{d|d}_u - r^{d|d}_{0,u})\,du, \int_0^t (q^{f\to d}_u - q^{f|f\to d}_{0,u})\,du)$ were obtained by integrating the moment ODEs of the dynamic above. The three no-arbitrage identities
\begin{eqnarray*}
  \EE^{d,\RN}\left[e^{-\int_0^t r^{d|d}_u du} Z^{d|d}_{t,T}\right] = Z^{d|d}_{0,T} , \quad
  \EE^{d,\RN}\left[e^{-\int_0^t r^{d|d}_u du} x^{df}_t Z^{f|d}_{t,T}\right] = x^{df}_0 Z^{f|d}_{0,T} , \\
  \EE^{d,\RN}\left[e^{-\int_0^t (r^{d|d}_u - q^{f\to d}_u) du} x^{df}_t Z^{f|f}_{t,T}\right] = x^{df}_0 Z^{f|f}_{0,T} ,
\end{eqnarray*}
then hold to machine precision (relative error below $10^{-15}$) for $t\in\{0.8,2.5,7\}$ and $T-t\in\{0.25,3,15\}$. Dropping the quanto drift of $z^f$, the term $y^{c,f}\vun_f$ or the term $\vun_f^\top v^{f,c}$ breaks the second identity by $6.5\cdot 10^{-2}$, $2.0\cdot 10^{-3}$ and $8.3\cdot 10^{-4}$ respectively.
\end{rema}


\newpage%
\subsection{Reconstruction formulas for market observables}%
All market observables of \autoref{sec:market_observables_list} follow from \autoref{theo:qgcm}. The reconstruction formulas for ZCBs give, for any collateral $i\in\{d,f\}$:
\begin{itemize}
\item FX forwards: $x^{df|i}_{t,T} = x^{df}_t Z^{f|i}_{t,T}/Z^{d|i}_{t,T}$. In particular
\begin{eqnarray}
  \ln \frac{x^{df|d}_{t,T}}{x^{df|d}_{0,T}} &=& z^x_t - \hat{G}^{\top}_{t,T} \hat{z}_t - \frac{1}{2} \hat{G}^{\top}_{t,T}  \hat{y}_t \hat{G}_{t,T}  +  G^{d,\top}_{t,T} z^{d}_t + \frac{1}{2} G^{d,\top}_{t,T}  y^{d,d}_t G^{d}_{t,T} .
\end{eqnarray}
\item Forward discount rates: $1+\theta R^{e|i}_{t,T;\theta} = Z^{e|i}_{t,T}/Z^{e|i}_{t,T+\theta}$.
\item Native forward Libors, under the constant funding spread hypothesis: $1 + \theta L^{e|e}_{t,T;\theta} = (1-\theta S^{e|e,\div}_{0,T;\theta})\, Z^{e|e}_{t,T}/Z^{e|e}_{t,T+\theta}$.
\item Resetting forward Libors under domestic collateral: $1+\theta L^{f/d|d}_{t,T;\theta} = (1-\theta S^{d|d,\div}_{0,T;\theta})\, Z^{f|d}_{t,T}/Z^{f|d}_{t,T+\theta}$.
\item Non-native forward Libors $L^{f|d}_{t,T;\theta} = \EE^{f|d,T+\theta}_t[L^f_{T;\theta}]$ are not functions of the ZCBs. In the Quasi-Gaussian model the collateral convexity adjustment depends on the future path of the volatilities and has to be computed numerically. In the LGCM it is deterministic and given by the collateral convexity formula of the deterministic volatility section:
\begin{eqnarray}
  1+\theta L^{f|d}_{t,T;\theta} &=& \left(1+\theta L^{f|f}_{t,T;\theta}\right) \exp\left\{ - \int_t^T \Sigma^{f|f}_{u,T;\theta}\cdot\Sigma^{f\to d}_{u,T+\theta}\, du \right\} .
  \label{eq:libor_collateral_convexity_lgcm}
\end{eqnarray}
\item Swap rates, currency basis spreads and annuities are then ratios of sums of the above.
\end{itemize}


\newpage%
\subsection{Gaussian par spread approximations}%
The Gaussian approximations for swap rates and currency basis spreads are derived in the LGCM in \autoref{ssec:lgcm_gaussian_par_spread}. They rely on freezing the gradient of the par spread with respect to the state variables. The same technique applies in the Quasi-Gaussian model, where the frozen gradient gives the volatility loading of the swap rate on $\hat{\sigma}_{z,t}$. Since $y_t$ is then stochastic, the resulting swap rate dynamic is a stochastic volatility diffusion rather than a Gaussian one, cf. \cite[\S13]{AndersenPiterbarg}.


\newpage%
\section{The Linear Gaussian Collateral Model: Specification}


\setcounter{subsection}{-1}
\subsection{Introduction}
We now combine our two assumptions:  
\begin{enumerate}
\item spot volatilities are separable. 
\item spot volatilities are deterministic. 
\end{enumerate}
The resulting model, which we call the Linear Gaussian Collateral Model, is the natural generalisation of the classic Black-Scholes and Hull-White models to stochastic collateral/currency basis.

We get the best of both worlds:
\begin{enumerate}
\item The model is Markovian. Moreover, compared to the Cheyette model, the variables $y_t$ are deterministic so the dimensionality of the Markovian diffusion is greatly reduced\footnote{From $\frac{M(M+1)}{2} +N_x$ to $M+N_x$ where $M = N_d+N_f+N_c$.}. 
\item The model is Gaussian. It can be simulated exactly with arbitrary time steps and convexity adjustments are known in closed form.
\end{enumerate}
Obviously, the main limitation is that the model only generates flat smiles\footnote{It is actually trivial to generate skew by replacing lognormal dynamics by shifted ones.}.



\subsection{Model specification}%


\subsubsection{Notations}

For each symbol $a$, let us write
\begin{eqnarray}
  H^a_{t} &=& \diag(h^a_t)  \\
  \sigma^a_{z,s} &=& g^a_sH^a_s  \\ 
  \sigma^a_{s} &=& \diag\left( \sigma^{a,\top}_{z,s}  \sigma^a_{z,s}\right)^{1/2}  \\ 
  M^a_{t,T} &=& H^{a,-1}_t H^a_T  \vun = H^{a,-1}_t h^a_T  \\
  G^a_{t,T} &=& \int_t^T M^a_{t,u} \, du = H^{a,-1}_t \int_t^T h^a_u \, du \\
  \Omega_t &=& \sigma^{a,-1}_{t}  \left( \sigma^{a,\top}_{z,t}  \sigma^a_{z,t}\right)  \sigma^{a,-1}_{t}   \\
  \Gamma^\top_t\Gamma_t &=& \Omega_t 
\end{eqnarray}
Let us define the mean reversion speeds as the coefficients of the diagonal matrix: 
\begin{eqnarray}
  \beta^a_{t} &=&  - H^{a,-1}_{t} \frac{d}{dt}H^a_{t}  
\end{eqnarray}
One can represent the matrices more explicitly:
\begin{eqnarray}
  H^a_{t} &=& \diag\left( e^{-\int_0^t \beta^{a,1}_u \, du} , \ldots, e^{-\int_0^t \beta^{a,N_a}_u \, du} \right)  \\
  \sigma^a_{t} &=& \diag\left( \sigma^{a,1}_t  , \ldots, \sigma^{a,N_a}_t \right)  \\
  M^a_{t,T} &=& \left[ e^{-\int_t^T \beta^{a,1}_u \, du} , \ldots, e^{-\int_t^T \beta^{a,N_a}_u \, du} \right]^\top  \\
  G^a_{t,T} &=& \left[ \int_t^T e^{-\int_t^s \beta^{a,1}_u \, du} \, ds , \ldots,  \int_t^T e^{-\int_t^s \beta^{a,N_a}_u \, du} \, ds \right]^\top  
\end{eqnarray}





\subsubsection{Model dynamic}%

We parametrise the model by a set of correlated one dimensional OU state variables plus the FX state variable.

The model is driven by a $k$-dimensional Brownian motion $W^\PP_t$ under the historical measure, and we write $W_t = W^{d,\RN}_t$ for the corresponding Brownian motion under the domestic risk-neutral measure.
Let us write $\sigma_{z,t} = [\sigma^{d|d}_{z,t} | \sigma^{f|f}_{z,t} | \sigma^{f\to d}_{z,t} | \sigma^{df}_{t} ] \in M_{k,N}(\RR)$, $N = N_d + N_f + N_c + 1$, for the matrix-valued volatility of the state variables. The vols $\sigma^{a}_{z,t}$ are deterministic. Let us set
\begin{eqnarray}
  \sigma_{t}  &=& \diag\left( \sigma_{z,t}^\top \sigma_{z,t} \right)^{1/2} ~=~ \diag\left(\sigma^{d,1}_t,\ldots,\sigma^{d,N_d}_t,\sigma^{f,1}_t,\ldots,\sigma^{c,N_c}_t,\sigma^{x}_t\right) \\
  \Omega_t  &=& \sigma_{t}^{-1}\left( \sigma_{z,t}^\top \sigma_{z,t} \right)\sigma_{t}^{-1} .
\end{eqnarray}
Here $\sigma^{x}_t = |\sigma^{df}_t|$ is the FX spot volatility.

We define a correlated $N$-dimensional Brownian motion $\widetilde{W}$:
\begin{eqnarray}
  d\widetilde{W}_t  &=& \Gamma_t^\top dW_t
\end{eqnarray}
where the factor loading matrix $\Gamma_t$ is such that $\Gamma_t^\top\Gamma_t = \Omega_t$. Its components are denoted by $\widetilde{W}^{(a,i)}_t$, $a\in\{d,f,c\}$, $i=1,\ldots,N_a$, and $\widetilde{W}^x_t$, and $d\langle \widetilde{W}^{(a,i)},\widetilde{W}^{(b,j)}\rangle_t = \rho^{(a,i),(b,j)}_t dt$. In this notation, the forward rate volatilities of \autoref{theo:qgcm} read
\begin{eqnarray}
  \sigma^{a,\top}_{t,T} dW_t ~=~ \sum_{i=1}^{N_a} \sigma^{a,i}_t M^{a,i}_{t,T} \, d\widetilde{W}^{(a,i)}_t , \qquad
  \Sigma^{a,\top}_{t,T} dW_t ~=~ \sum_{i=1}^{N_a} \sigma^{a,i}_t G^{a,i}_{t,T} \, d\widetilde{W}^{(a,i)}_t , \qquad
  \sigma^{df,\top}_t dW_t ~=~ \sigma^x_t \, d\widetilde{W}^x_t .
\end{eqnarray}

We also write $\Pi_a: \RR^N \to \RR^{N_a}$ for the projection onto the bloc $a$. And for $M \in M_{N}(\RR)$, we write $M^{a,b} = \Pi_a M \Pi_b^{\top}$ for its $(a,b)$ bloc. For example, the correlation matrix $\Omega_t$ admits the following bloc decomposition\footnote{We use $c$ as short for $f\to d$ and $x$ as short for $df$.}:
\begin{eqnarray}
  \Omega_t ~=~ \Gamma_{t}^\top\Gamma_t ~=~
  \begin{bmatrix}
     \Omega^{d,d}_{t} &   \Omega^{d,f}_{t}  &  \Omega^{d,c}_{t}  &  \Omega^{d,x}_{t}  \\
     \Omega^{f,d}_{t} &   \Omega^{f,f}_{t}  &  \Omega^{f,c}_{t}  & \Omega^{f,x}_{t}  \\
     \Omega^{c,d}_{t} &   \Omega^{c,f}_{t}  &  \Omega^{c,c}_{t}  &  \Omega^{c,x}_{t}  \\
     \Omega^{x,d}_{t} &  \Omega^{x,f}_{t}  &  \Omega^{x,c}_{t}  &  \Omega^{x,x}_{t}  \\
      \end{bmatrix}
\end{eqnarray}
where each bloc looks like:
\begin{eqnarray}
  \Omega^{a,b}_t ~=~
  \begin{bmatrix}
     \rho^{(a,1),(b,1)}_{t} &   \rho^{(a,1),(b,2)}_{t}   &  \cdots &  \rho^{(a,1),(b,N_b)}_{t}  \\
     \rho^{(a,2),(b,1)}_{t}  &   \rho^{(a,2),(b,2)}_{t}   &   & \rho^{(a,2),(b,N_b)}_{t}   \\
     \vdots &     &  \ddots  &  \vdots  \\
     \rho^{(a,N_a),(b,1)}_{t}  &  \cdots  &  \cdots  &  \rho^{(a,N_a),(b,N_b)}_{t}  \\
      \end{bmatrix}
\end{eqnarray}
Note that we don't require the matrix $\Omega_t$ to be definite: we can have $\Gamma_t \in M_{k,N}(\RR)$ with $k< N$.

\begin{defi}\label{defi:lgcm}
The \emph{Linear Gaussian Collateral Model} (LGCM) is defined by
\begin{enumerate}
\item The following deterministic parameters:
\begin{enumerate}
\item IRD: spot volatilities $(\sigma^{d,j}_t)$ and mean-reversions $(\beta^{d,j}_t)$ for $j = 1,\ldots, N_d$.
\item IRF: spot volatilities $(\sigma^{f,j}_t)$ and mean-reversions $(\beta^{f,j}_t)$ for $j = 1,\ldots, N_f$.
\item CB: spot volatilities $(\sigma^{c,j}_t)$ and mean-reversions $(\beta^{c,j}_t)$ for $j = 1,\ldots, N_{c}$.
\item FX: spot volatility $\sigma^{x}_t$.
\item Correlations: instantaneous correlations between Gaussian drivers $(\rho^{(a,j_a),(b,j_b)}_t)$.
\end{enumerate}
\item The deterministic functions, for $a,b\in\{d,f,c\}$,
\begin{eqnarray}
  y^{(a,i),(b,j)}_t &=& \int_0^t \rho^{(a,i),(b,j)}_u \sigma^{a,i}_u \sigma^{b,j}_u M^{a,i}_{u,t} M^{b,j}_{u,t}\, du  \\
  v^{f,i}_t &=& \sum_{j=1}^{N_c}\int_0^t \rho^{(f,i),(c,j)}_u \sigma^{f,i}_u \sigma^{c,j}_u M^{f,i}_{u,t} G^{c,j}_{u,t}\, du
\end{eqnarray}
(the deterministic versions of $y^{a,b}_t$ and $v^{f,c}_t$ in \autoref{theo:qgcm}).
\item The state process $z_t = [z^d_t|z^f_t|z^{c}_t|z^{x}_t]^\top$, $z_0 = 0$, with the following dynamic under $\QQ^{d,\RN}$:
\begin{eqnarray}
  dz^{d,i}_{t}  &=& \left( \sum_{j=1}^{N_d} y^{(d,i),(d,j)}_t - \beta^{d,i}_t z^{d,i}_t \right) dt + \sigma^{d,i}_t \, d\widetilde{W}^{(d,i)}_t \\
  dz^{f,i}_{t}  &=& \left( \sum_{j=1}^{N_f} y^{(f,i),(f,j)}_t - \rho^{(f,i),x}_t\sigma^{f,i}_t\sigma^{x}_t - \beta^{f,i}_t z^{f,i}_t \right) dt + \sigma^{f,i}_t \, d\widetilde{W}^{(f,i)}_t \\
  dz^{c,i}_{t}  &=& \left( \sum_{j=1}^{N_c} y^{(c,i),(c,j)}_t + \sum_{j=1}^{N_f} y^{(c,i),(f,j)}_t - \rho^{(c,i),x}_t\sigma^{c,i}_t\sigma^{x}_t - \beta^{c,i}_t z^{c,i}_t \right) dt + \sigma^{c,i}_t \, d\widetilde{W}^{(c,i)}_t \\
  dz^{x}_{t}  &=& \left( \sum_{i=1}^{N_d} z^{d,i}_t - \sum_{i=1}^{N_f} \left(z^{f,i}_t + v^{f,i}_t\right) - \sum_{i=1}^{N_c} z^{c,i}_t - \frac{1}{2}\left(\sigma^x_t\right)^2 \right) dt + \sigma^{x}_t \, d\widetilde{W}^x_t .
\end{eqnarray}
\item Reconstruction formulas, with $\hat{z}_t = [z^f_t + v^{f,c}_t ; z^c_t]$, $\hat{G}_{t,T} = [G^f_{t,T};G^c_{t,T}]$ and $\hat{y}_t$ the $(f,c)$ bloc of $y_t$:
\begin{eqnarray}
  Z^{d|d}_{t,T}  &=& Z^{d|d}_{0,t,T} \exp\left( -G^{d,\top}_{t,T} z^{d}_t -\frac{1}{2} G^{d,\top}_{t,T}  y^{d,d}_t G^{d}_{t,T} \right) \\
  Z^{f|f}_{t,T}  &=& Z^{f|f}_{0,t,T} \exp\left( -G^{f,\top}_{t,T} z^{f}_t -\frac{1}{2} G^{f,\top}_{t,T}  y^{f,f}_t G^{f}_{t,T} \right) \\
  Z^{f|d}_{t,T}  &=& Z^{f|d}_{0,t,T} \exp\left( -\hat{G}^{\top}_{t,T} \hat{z}_t - \frac{1}{2} \hat{G}^{\top}_{t,T}  \hat{y}_t \hat{G}_{t,T} \right) \\
  x^{df}_{t}  &=& x^{df|d}_{0,t} \, e^{z^x_t} ~=~ x^{df}_0 \frac{Z^{f|d}_{0,t}}{Z^{d|d}_{0,t}} \, e^{z^x_t} .
\end{eqnarray}
\end{enumerate}
\end{defi}

\begin{rema}
By \autoref{theo:qgcm}, the LGCM is arbitrage free and reproduces the initial curves $Z^{d|d}_{0,\cdot}$, $Z^{f|f}_{0,\cdot}$, $Z^{f|d}_{0,\cdot}$ and the FX spot $x^{df}_0$ for any choice of the parameters. The other discount curves follow from the reconstruction formulas: $Z^{d|f}_{t,T}$ from the change of collateral formula, and $\QZ^{f|f\to d}_{t,T} = Z^{f|d}_{t,T}/Z^{f|f}_{t,T}$. Compared to the Quasi-Gaussian model, the only stochastic state variables are the $N-1$ OU factors $z^{a,i}_t$ and the log FX moneyness $z^x_t$. The $y_t$ and $v^{f,c}_t$ are precomputed by the covariance integral algorithm of \autoref{ssec:covariance_integrals}.
\end{rema}


\begin{rema}[Stationarity]
Remember that, while we specify the dynamic under the domestic risk-neutral measure, we assumed that this measure is equivalent to the historical one. In particular, we expect delivered covariances in a well specified and calibrated model to be close to historically estimated ones.

Since statistical estimation relies on the stationarity of the stochastic process, our model should be stationary. This requires coefficients such mean reversion speeds to be constant. However, for numerical reasons, one often have to rely on term structure of parameters to match market prices. For this reason, we will keep deriving the formulas in the general setting.
\end{rema}



\subsubsection{Dynamic of the the collateral basis}%

We assume a Linear Gaussian dynamic for the instantaneous forward collateral basis $q^{f|f\to d}_{t,T}$, i.e. separable and deterministic volatility
\begin{eqnarray}
  \sigma^{f\to d,\top}_{t,T} dW_t &=& \sum_{i=1}^{N_c} \sigma^{c,i}_t M^{c,i}_{t,T}\, d\widetilde{W}^{(c,i)}_t .
\end{eqnarray}
Equivalently, the cross-currency discount rates $r^{f|d}_{t,T}$ follow a Linear Gaussian dynamic with $N_f + N_c$ factors, the foreign native ones and the collateral basis ones (see the remark on the cross-currency discount rate in the proof of \autoref{theo:qgcm}).

\begin{prop}
In the LGCM, the forward collateral basis has dynamic under $\QQ^{d,\RN}$
\begin{eqnarray}
  dq^{f|f\to d}_{t,T} &=& \sum_{i=1}^{N_c} \sigma^{c,i}_t M^{c,i}_{t,T}\left( \sum_{j=1}^{N_c} \rho^{(c,i),(c,j)}_t \sigma^{c,j}_t G^{c,j}_{t,T} + \sum_{j=1}^{N_f} \rho^{(c,i),(f,j)}_t \sigma^{f,j}_t G^{f,j}_{t,T} - \rho^{(c,i),x}_t\sigma^x_t \right) dt \nonumber \\
  & & + \sum_{i=1}^{N_f} \sum_{j=1}^{N_c} \rho^{(f,i),(c,j)}_t \sigma^{f,i}_t \sigma^{c,j}_t M^{f,i}_{t,T} G^{c,j}_{t,T}\, dt
  + \sum_{i=1}^{N_c} \sigma^{c,i}_t M^{c,i}_{t,T}\, d\widetilde{W}^{(c,i)}_t
\end{eqnarray}
and is reconstructed from the state variables by
\begin{eqnarray}
  q^{f|f\to d}_{t,T} &=& q^{f|f\to d}_{0,T} + \sum_{i=1}^{N_c} M^{c,i}_{t,T}\left( z^{c,i}_t + \sum_{j=1}^{N_c} y^{(c,i),(c,j)}_t G^{c,j}_{t,T} + \sum_{j=1}^{N_f} y^{(c,i),(f,j)}_t G^{f,j}_{t,T} \right) \nonumber \\
  & & + \sum_{i=1}^{N_f} M^{f,i}_{t,T}\left( v^{f,i}_t + \sum_{j=1}^{N_c} y^{(f,i),(c,j)}_t G^{c,j}_{t,T} \right) \\
  \QZ^{f|f\to d}_{t,T} &=& \QZ^{f|f\to d}_{0,t,T} \exp\left( -G^{c,\top}_{t,T} z^{c}_t - G^{f,\top}_{t,T} v^{f,c}_t - \frac{1}{2} G^{c,\top}_{t,T}  y^{c,c}_t G^{c}_{t,T} - G^{c,\top}_{t,T}  y^{c,f}_t G^{f}_{t,T}  \right) .
\end{eqnarray}
In particular, the short collateral basis is $q^{f\to d}_t = q^{f|f\to d}_{0,t} + \sum_i z^{c,i}_t + \sum_i v^{f,i}_t$.
\end{prop}

\begin{proof}
The dynamic is the HJM dynamic of the collateral basis written with the correlated drivers. The reconstruction formulas are those of \autoref{theo:qgcm}.
\end{proof}

\begin{rema}
The collateral basis has three sources of drift: its own convexity ($\Sigma^c$), the convexity with foreign rates ($\Sigma^f$, through both $\sigma^c\cdot\Sigma^f$ and $\sigma^f\cdot\Sigma^c$) and the quanto adjustment ($\sigma^x$). The second one is what makes the deterministic shift $v^{f,c}_t$ appear. It vanishes when foreign rates and collateral basis are uncorrelated, and then $q^{f|f\to d}$ is driven by the $z^{c}$ factors only.
\end{rema}



\subsubsection{Dynamic of the FX spot}%

\begin{prop}
In the LGCM, the FX spot has dynamic under $\QQ^{d,\RN}$
\begin{eqnarray}
  \frac{dx^{df}_t}{x^{df}_t} &=& \left( r^{d|d}_t - r^{f|f}_t - q^{f\to d}_t \right) dt + \sigma^x_t \, d\widetilde{W}^x_t ,
\end{eqnarray}
where the short rates are $r^{d|d}_t = r^{d|d}_{0,t} + \sum_i z^{d,i}_t$, $r^{f|f}_t = r^{f|f}_{0,t} + \sum_i z^{f,i}_t$ and $q^{f\to d}_t$ is given above. The FX forward under collateral $i\in\{d,f\}$ is reconstructed by $x^{df|i}_{t,T} = x^{df}_t Z^{f|i}_{t,T}/Z^{d|i}_{t,T}$, and for $i = d$
\begin{eqnarray}
  \ln \frac{x^{df|d}_{t,T}}{x^{df|d}_{0,T}} &=& z^x_t - \hat{G}^{\top}_{t,T} \hat{z}_t - \frac{1}{2} \hat{G}^{\top}_{t,T}  \hat{y}_t \hat{G}_{t,T}  +  G^{d,\top}_{t,T} z^{d}_t + \frac{1}{2} G^{d,\top}_{t,T}  y^{d,d}_t G^{d}_{t,T} .
\end{eqnarray}
It is a lognormal $\QQ^{d|d,T}$-martingale with volatility
\begin{eqnarray}
  \Sigma^{df,\top}_{t,T} dW_t &=& \sum_{i=1}^{N_d}\sigma^{d,i}_tG^{d,i}_{t,T}d\widetilde{W}^{(d,i)}_t - \sum_{i=1}^{N_f}\sigma^{f,i}_tG^{f,i}_{t,T}d\widetilde{W}^{(f,i)}_t - \sum_{i=1}^{N_c}\sigma^{c,i}_tG^{c,i}_{t,T}d\widetilde{W}^{(c,i)}_t + \sigma^x_t\,d\widetilde{W}^x_t ,
\end{eqnarray}
so that the total implied variance of an FX option of expiry $T$ is $w^{df}_0(T) = \int_0^T|\Sigma^{df}_{u,T}|^2du$, a sum of covariance integrals.
\end{prop}

\begin{proof}
The FX dynamic is that of the HJM section. The FX forward follows from the reconstruction formulas and $x^{df}_t = x^{df|d}_{0,t}e^{z^x_t}$, using $x^{df|d}_{0,T} = x^{df|d}_{0,t}Z^{f|d}_{0,t,T}/Z^{d|d}_{0,t,T}$. Its volatility is given by the proposition on FX forwards in the HJM section.
\end{proof}

\begin{rema}
The FX spot volatility $\sigma^x_t$ is the only FX parameter. The FX forward volatility also contains the rates and collateral basis volatilities through $G^{a,i}_{t,T}$, which grow with the time to expiry. For long expiries, they can account for a large part of the FX implied volatility, which is why bootstrapping $\sigma^x_t$ to FX implied volatilities can fail (variance squeeze, see the remark on the FX options market in the calibration section).
\end{rema}



\subsubsection{Correlation parametrisation}%

The correlation matrix $\Omega_t$ has $N(N-1)/2$ entries, which is too many to be calibrated or estimated freely. We organise them by the blocs of $\Omega_t$ and by the quantities they drive in the model:
\begin{center}
\begin{tabular}{|c|l|}
\hline
Bloc & Main impact \\
\hline
$\Omega^{d,d}$, $\Omega^{f,f}$ & native swaption smiles: decorrelation of swap rates, calibrated in the single currency LGMs \\
$\Omega^{c,c}$ & decorrelation of currency basis spreads of different tenors \\
$\Omega^{f,c}$ & Libor collateral convexity $K_j$, drift of $q^{f|f\to d}$, $v^{f,c}_t$ \\
$\Omega^{d,x}$, $\Omega^{f,x}$, $\Omega^{c,x}$ & quanto drifts, FX forward collateral convexity, FX option variance \\
$\Omega^{d,f}$, $\Omega^{d,c}$ & joint distribution of domestic rates with foreign rates and basis (e.g. CVA) \\
\hline
\end{tabular}
\end{center}
The table follows from the formulas above. For example, $\Omega^{f,c}$ enters the Libor collateral convexity $\exp\{-\int \Sigma^{f|f}_{u,T;\theta}\cdot\Sigma^{f\to d}_{u,T+\theta}du\}$, and $\Omega^{c,x}$ enters the FX forward collateral convexity $\exp\{-\int\Sigma^{df}_{u,T}\cdot\Sigma^{i\to j}_{u,T}du\}$.

A natural parametrisation is the following:
\begin{enumerate}
\item The diagonal blocs $\Omega^{d,d}$ and $\Omega^{f,f}$ are inherited from the single currency LGMs, which are calibrated independently to their native markets.
\item The diagonal bloc $\Omega^{c,c}$ and the off diagonal blocs are marked, or estimated from historical time series of the corresponding market quotes (see the section on historical estimation). They are typically parametrised by a few numbers, e.g. one correlation per pair of blocs applied to the first factors.
\item The resulting matrix is not necessarily positive semi-definite. It has to be projected onto the set of correlation matrices, for example by the nearest correlation matrix in Frobenius norm, or by flooring the negative eigenvalues of $\Omega_t$ at $0$ and renormalising the diagonal. Alternatively, one parametrises the factor loading matrix $\Gamma_t$ directly (e.g. by angles on the unit sphere for each column), which guarantees $\Omega_t = \Gamma_t^\top\Gamma_t$ is a correlation matrix and allows a rank $k<N$.
\end{enumerate}
The projection modifies the native blocs in general. If the native calibrations must be preserved exactly, the projection should keep $\Omega^{d,d}$ and $\Omega^{f,f}$ fixed and only adjust the other blocs.



\subsubsection{Joint dynamic}%

The LGCM is a linear Gaussian model. We write $X_t = [z^d_t|z^f_t|z^c_t|z^x_t]^\top$.

\begin{prop}[Joint law of the state]\label{prop:lgcm_joint_law}
Under $\QQ^{d,\RN}$, $X_t$ solves the linear SDE
\begin{eqnarray}
  dX_t &=& \left( A_t X_t + b_t \right) dt + \sigma_t \, d\widetilde{W}_t
\end{eqnarray}
where $\sigma_t$ is the diagonal matrix of spot volatilities,
\begin{eqnarray*}
  A_t ~=~ \begin{bmatrix} -\beta^d_t & 0 & 0 & 0 \\ 0 & -\beta^f_t & 0 & 0 \\ 0 & 0 & -\beta^c_t & 0 \\ \vun^\top_d & -\vun^\top_f & -\vun^\top_c & 0 \end{bmatrix} , \qquad
  b_t ~=~ \begin{bmatrix} y^{d,d}_t\vun_d \\ y^{f,f}_t\vun_f - \sigma^f_t\Omega^{f,x}_t\sigma^x_t \\ y^{c,c}_t\vun_c + y^{c,f}_t\vun_f - \sigma^c_t\Omega^{c,x}_t\sigma^x_t \\ -\vun^\top_f v^{f,c}_t - \frac{1}{2}(\sigma^x_t)^2 \end{bmatrix} .
\end{eqnarray*}
$X_t$ is Gaussian with mean $m_t$ and covariance matrix $C_t$ given by covariance integrals (with the notations of \autoref{ssec:covariance_integrals}):
\begin{eqnarray*}
  m^{a,i}_t &=& \sum_{j} J_{0,t}[M^{a,i}_{\cdot,t}, G^{a,j}_{\cdot,t}] + \mathbf{1}_{a=c}\sum_{j} J_{0,t}[M^{c,i}_{\cdot,t}, G^{f,j}_{\cdot,t}] - \mathbf{1}_{a\in\{f,c\}} J_{0,t}[M^{a,i}_{\cdot,t}, 1^x] \\
  m^x_t &=& - \frac{1}{2}J_{0,t}[\Sigma^{df}_{\cdot,t},\Sigma^{df}_{\cdot,t}] + J_{0,t}[\Sigma^{df}_{\cdot,t},\Sigma^{d}_{\cdot,t}] \\
  C^{(a,i),(b,j)}_t &=& y^{(a,i),(b,j)}_t , \qquad
  C^{(a,i),x}_t ~=~ J_{0,t}[M^{a,i}_{\cdot,t},\Sigma^{df}_{\cdot,t}] , \qquad
  C^{x,x}_t ~=~ J_{0,t}[\Sigma^{df}_{\cdot,t},\Sigma^{df}_{\cdot,t}] = w^{df}_0(t) .
\end{eqnarray*}
\end{prop}

\begin{proof}
The SDE is \autoref{defi:lgcm} in matrix form. The law is the case $s=0$ of \autoref{prop:lgcm_transition_rn}, where $z_0 = 0$ and $y_0 = 0$.
\end{proof}

\begin{rema}
\begin{enumerate}
\item The covariance of the rates factors is $y_t$ itself: $y_t$ is the covariance matrix of the OU factors, which is why it is called the covariance process in the Quasi-Gaussian model.
\item Equivalently, $z^{a,i}_t = m^{a,i}_t + \xi^{a,i}_t$ where $\xi^{a,i}_t = \int_0^t\sigma^{a,i}_uM^{a,i}_{u,t}d\widetilde{W}^{(a,i)}_u$ are centred correlated one dimensional OU processes, $d\xi^{a,i}_t = -\beta^{a,i}_t\xi^{a,i}_tdt + \sigma^{a,i}_td\widetilde{W}^{(a,i)}_t$. All the drifts are in the deterministic functions $m_t$.
\item The formulas were checked against the moment ODEs of the linear SDE, with the model of the verification of \autoref{theo:qgcm}, at $t=0.8$ and $t=7$: the means and covariances agree to $3\cdot 10^{-16}$ (absolute), for means of order $10^{-2}$ and an FX variance of order $10^{-1}$.
\item The transitions between two dates and the drifts under the forward and jumping num\'eraire measures are given in \autoref{ssec:lgcm_simulation}.
\end{enumerate}
\end{rema}


\newpage%
\subsection{Reconstruction of market observables}
We now derive the dynamic of market observables in terms of correlated Brownian motions. 

\subsubsection{Forward discount rates}

\begin{prop}
The dynamic of forward discount rates is\footnote{Here $\Sigma^{a}_{t,T;\theta} = \Sigma^{a}_{t,T+\theta} - \Sigma^{a}_{t,T}$ has components $\sigma^{a,k}_t G^{a,k}_{t,T;\theta}$ with 
\begin{eqnarray*}
  G^{a,k}_{t,T;\theta} ~=~ G^{a,k}_{t,T+\theta} - G^{a,k}_{t,T} ~=~ \int_T^{T+\theta} e^{-\int_t^u \beta^{a,k}_v\,dv}\,du ~=~ M^{a,k}_{t,T}\, G^{a,k}_{T,T+\theta} .
\end{eqnarray*}
So $G^{a,k}_{t,T;\theta}$ is \emph{not} equal to $G^{a,k}_{T,T+\theta}$: the missing amortising factor is $M^{a,k}_{t,T} = e^{-\int_t^T\beta^{a,k}_v dv}$. An earlier computation found a different expression because it used $\int_t^T e^{-\int_u^T \beta_v dv}du$, which is $\int_t^T M_{u,T}\,du$ rather than $G_{t,T} = \int_t^T M_{t,u}\,du$; the two agree only for constant mean reversion.
}:  
\begin{eqnarray}
  \frac{d (1+\theta R^{e|i}_{t,T;\theta})}{1+\theta R^{e|i}_{t,T;\theta}}  
  &=&  \sum_{k=1}^{N_e}  \Sigma^{e|i}_{t,T;\theta} dW^{e|i,T+\theta,k}_t  \nonumber \\
  &=&  \sum_{k=1}^{N_e}  G^{e|e,k}_{t,T;\theta} \sigma^{e|e,k}_{t}   d\widetilde{W}^{e|i,T+\theta,(e,k)}_t  %\\
  %& & 
  + \sum_{k=1}^{N_{e\to i}} G^{e\to i,k}_{t,T;\theta} \sigma^{e\to i,k}_{t} d\widetilde{W}^{e|i,T+\theta,(e\to i,k)}_t 
\end{eqnarray} 
\end{prop}

\subsubsection{Forward Libors}

[FIXME: Make the parametrisation explicit.]

\begin{prop}
The dynamic of forward Libors is 
\begin{eqnarray}
  \frac{d (1+\theta L^{e|i}_{t,T;\theta})}{1+\theta L^{e|i}_{t,T;\theta}}  
  &=& \Sigma^{e|e}_{t,T;\theta} dW^{e|i,T+\theta}_t  \nonumber \\
  &=& \sum_{k=1}^{N_e}  G^{e|e,k}_{t,T;\theta} \sigma^{e|e,k}_{t}   d\widetilde{W}^{e|i,T+\theta,(e,k)}_t 
\end{eqnarray} 
\end{prop}


\subsubsection{FX forwards}

\begin{prop}
The FX forward has dynamic 
\begin{eqnarray}
  \frac{d x^{df|i}_{t,T}}{x^{df|i}_{t,T}}  
  &=& \left( \Sigma^{d|d}_{t,T} - \Sigma^{f|f}_{t,T}  - \Sigma^{f\to d}_{t,T}  + \sigma^{df}_t \right)   \cdot dW^{d|i,T}_t   \\
  &=& \sum_{k_d=1}^{N_d}  G^{d|d,k_d}_{t,T} \sigma^{d|d,k_d}_{t} \, d\widetilde{W}^{d|i,T,(d,k_d)}_t 
  - \sum_{k_f=1}^{N_f}  G^{f|f,k_f}_{t,T} \sigma^{f|f,k_f}_{t}  \, d\widetilde{W}^{d|i,T,(f,k_f)}_t  \nonumber \\
  & & - \sum_{k_{c}=1}^{N_c}  G^{f\to d,k_c}_{t,T} \sigma^{f\to d,k_c}_{t} \, d\widetilde{W}^{d|i,T,(f\to ,k_c)}_t 
  + \sum_{k_x=1}^{N_x}  \sigma^{df,k_x}_{t}  \, d\widetilde{W}^{d|i,T,(df,k_x)}_t  . \nonumber
\end{eqnarray} 
\end{prop}


\subsubsection{Funding spreads}


\begin{prop}
The funding spreads have dynamic 
\begin{eqnarray}
  \frac{d (1-\theta S^{e|i,\div}_{t,T;\theta})}{1-\theta S^{e|i}_{t,T;\theta}}  
  &=& - \Sigma^{e\to i}_{t,T;\theta}  \cdot dW^{e|i,T}_t  \\
  &=& -  \sum_{k=1}^{N_e}  G^{e\to i,k}_{t,T;\theta} \sigma^{e\to i,k}_{t}   d\widetilde{W}^{e|i,T,(e\to i,k)}_t . \nonumber
\end{eqnarray} 
\end{prop}



\subsubsection{Resetting Forward Libors}



\begin{prop}
The resetting forward Libor has dynamic 
\begin{eqnarray}
  \frac{d (1+\theta L^{f/d|i,r}_{t,T;\theta})}{1+\theta L^{f/d|i,r}_{t,T;\theta}}  
  &=&  \Sigma^{f|d}_{t,T;\theta}  \cdot dW^{f|i,T+\theta}_t
\end{eqnarray} 
\end{prop}


\subsubsection{Currency Basis Spreads}

\begin{prop}
The dynamic of the currency basis spread can be approximated by 
\begin{eqnarray}
  B^{f/d|c,r}_{t,T;\theta} 
  %&=& 
  %B^{f/d|c,r}_{t_0,T;\theta} 
  %+ \int_{t_0}^t \Sigma^{f/d|r}_{u,T;\theta} \, dW^{f|c,T+\theta}_{u} \\
  &\approx& B^{f/d|c,r}_{t_0,T;\theta} 
  +  \frac{1+\theta L^{f/d|c,r}_{t_0,T;\theta}}{\theta} \int_{t_0}^t \Sigma^{f\to d}_{u,T;\theta} \cdot dW^{f|c,T+\theta}_{u} 
  + B^{f/d|c,r}_{t_0,T;\theta} \int_{t_0}^t \Sigma^{f|f}_{u,T;\theta}  \cdot dW^{f|c,T+\theta}_{u}  
\end{eqnarray} 

\end{prop}


\begin{rema}
Note that 
\begin{eqnarray}
  L^{f/d|c,r}_{t_0,T;\theta} 
  &=& L^{f|d}_{t_0,T;\theta} + B^{f/d|c,r}_{t_0,T;\theta}  
\end{eqnarray} 
so that 
\begin{eqnarray}
  L^{f/d|c,r}_{t_0,T;\theta} 
  %&=& L^{f|d}_{t_0,T;\theta} + B^{f/d|c,r}_{t_0,T;\theta}  \\
  &=& B^{f/d|c,r}_{t_0,T;\theta} + \frac{1}{\theta}\left( (1+\theta L^{f|f}_{t_0,T;\theta}) e^{- \int_{t_0}^T \Sigma^{f|f}_{u,T;\theta} \cdot \Sigma^{f\to d}_{u,T+\theta} \, du} - 1 \right)
\end{eqnarray} 
so the volatility of the currency basis spread can be explicitly computed from the foreign Libor, the currency basis spread itself and model parameters.
\end{rema}


\begin{rema}
The first volatility term $ \frac{1+\theta L^{f/d|c,r}_{t_0,T;\theta}}{\theta}  \Sigma^{f\to d}_{u,T;\theta}$ corresponds to the collateral basis volatility. We expect it to be the the most significant one since the coefficient is of the order $1/\theta = 4$. 

However, in periods of stress such as the Argentinian financial crisis, the currency basis spread will increase and the volatility of foreign rates will rise. As a result, the second term will make a material contribution to the realised volatility of the currency basis spread\footnote{Note however that if foreign rates volatility rises too high, the normal approximation for Libors may no longer be accurate.}. 
\end{rema}



\newpage%
\subsection{Gaussian par spread approximations}\label{ssec:lgcm_gaussian_par_spread}%



\subsubsection{Dynamic of swap rates}

%\paragraph{Annuity leg}


\begin{prop}
The annuity has dynamic:
\begin{eqnarray*}
  dA^{x|c}_{t,T;M}  &=&  \sum_{j=1}^{\ell M} \theta_i Z^{x|c}_{t,T_j} \left( r^{x|c}_t \, dt - \Sigma^{x|c,\top}_{t,T_j} dW_t^{x,\RN} \right)
\end{eqnarray*}
\end{prop}


%\paragraph{Floating leg}


\begin{prop}
The floating leg has dynamic:
\begin{eqnarray*}
  d\Fl^{x|c}_{t,T;M}  &=&   \sum_{j=1}^{\ell M} \theta_j Z^{x|c}_{t,T_j} L^{x|c}_{t,T_{j-1};\theta}  \left( r^{x|c}_t \, dt - \Sigma^{x|c,\top}_{t,T_j} dW_t^{x,\RN} \right) \\
  & & ~~+~~ \sum_{j=1}^{\ell M} Z^{x|c}_{t,T_j} (1+\theta_j L^{x|c}_{t,T_{j-1};\theta}) \Sigma^{x|x,\top}_{t,T_{j-1};\theta} dW_t^{x|c,T_j} + O(dt)
\end{eqnarray*}
\end{prop}

\begin{proof}
\begin{eqnarray*}
\Fl^{x|c}_{t,T;M}  &=&   \sum_{j=1}^{\ell M} \theta_j Z^{x|c}_{t,T_j} L^{x|c}_{t,T_{j-1};\theta}  \\
  d\Fl^{x|c}_{t,T;M}  
  &=&   \sum_{j=1}^{\ell M} \theta_j Z^{x|c}_{t,T_j} L^{x|c}_{t,T_{j-1};\theta} \frac{dZ^{x|c}_{t,T_j} }{Z^{x|c}_{t,T_j} }  + \sum_{j=1}^{\ell M} \theta_j Z^{x|c}_{t,T_j} dL^{x|c}_{t,T_{j-1};\theta} + O(dt) \\
   &=&   \sum_{j=1}^{\ell M} \theta_j Z^{x|c}_{t,T_j} L^{x|c}_{t,T_{j-1};\theta}  \left( r^{x|c}_t \, dt - \Sigma^{x|c,\top}_{t,T_j} dW_t^{x,\RN} \right) \\
  & & ~~+~~ \sum_{j=1}^{\ell M} Z^{x|c}_{t,T_j} (1+\theta_j L^{x|c}_{t,T_{j-1};\theta}) \Sigma^{x|x,\top}_{t,T_{j-1};\theta} dW_t^{x|c,T_j} + O(dt)
\end{eqnarray*}
\end{proof}

%\paragraph{Forward swap rates}

\begin{prop}
The forward swap rate has dynamic:
\begin{eqnarray*}
  dY^{x|c}_{t,T;M}  &=&   \frac{1}{A^{x|c}_{t,T;M} }  \sum_{j=1}^{\ell M} Z^{x|c}_{t,T_j}  \left\{ (1+\theta_j L^{x|c}_{t,T_{j-1};\theta}) \Sigma^{x|x}_{t,T_{j-1};\theta}  + \theta_j ( Y^{x|c}_{t,T;M} - L^{x|c}_{t,T_{j-1};\theta} ) \Sigma^{x|c}_{t,T_j}  \right\}^\top dW_t^{x|c,T;M}
\end{eqnarray*}
\end{prop}

\begin{proof}
\begin{eqnarray*}
  dY^{x|c}_{t,T;M}  
  &=&  \frac{d\Fl^{x|c}_{t,T;M}}{A^{x|c}_{t,T;M} } -  Y^{x|c}_{t,T;M} \frac{dA^{x|c}_{t,T;M}}{A^{x|c}_{t,T;M} } + O(dt) \\
  &=& - \frac{1}{A^{x|c}_{t,T;M} } \sum_{j=1}^{\ell M} \theta_j Z^{x|c}_{t,T_j} L^{x|c}_{t,T_{j-1};\theta} \Sigma^{x|c,\top}_{t,T_j} dW_t^{x,\RN}  \\
  & & ~~+~~ \frac{1}{A^{x|c}_{t,T;M} } \sum_{j=1}^{\ell M} Z^{x|c}_{t,T_j} (1+\theta_j L^{x|c}_{t,T_{j-1};\theta}) \Sigma^{x|x,\top}_{t,T_{j-1};\theta} dW_t^{x|c,T_j}  \\ 
  & & ~~+~~  Y^{x|c}_{t,T;M}  \sum_{j=1}^{\ell M} \theta_j Z^{x|c}_{t,T_j} \Sigma^{x|c,\top}_{t,T_j} dW_t^{x,\RN}  + O(dt) \\
  &=& - \frac{1}{A^{x|c}_{t,T;M} } \sum_{j=1}^{\ell M} \theta_j Z^{x|c}_{t,T_j} L^{x|c}_{t,T_{j-1};\theta} \Sigma^{x|c,\top}_{t,T_j} dW_t^{x|c,T;M}  \\
  & & ~~+~~ \frac{1}{A^{x|c}_{t,T;M} } \sum_{j=1}^{\ell M} Z^{x|c}_{t,T_j} (1+\theta_j L^{x|c}_{t,T_{j-1};\theta}) \Sigma^{x|x,\top}_{t,T_{j-1};\theta} dW_t^{x|c,T;M}   \\ 
  & & ~~+~~  Y^{x|c}_{t,T;M}  \sum_{j=1}^{\ell M} \theta_j Z^{x|c}_{t,T_j} \Sigma^{x|c,\top}_{t,T_j} dW_t^{x|c,T;M} \\
  &=& \frac{1}{A^{x|c}_{t,T;M} }  \sum_{j=1}^{\ell M} Z^{x|c}_{t,T_j}  \left\{ \theta_j Y^{x|c}_{t,T;M} \Sigma^{x|c}_{t,T_j}  + (1+\theta_j L^{x|c}_{t,T_{j-1};\theta}) \Sigma^{x|x}_{t,T_{j-1};\theta}  - L^{x|c}_{t,T_{j-1};\theta} \Sigma^{x|c}_{t,T_j}  \right\}^\top dW_t^{x|c,T;M}
\end{eqnarray*}
\end{proof}

\begin{rema}
The component with collateral basis volatility is quite small since it depends on $Y^{x|c}_{t,T;M} - L^{x|c}_{t,T_{j-1};\theta}$ which is zero if the index curve is flat. 
\end{rema}

\newpage%
\subsubsection{Gaussian approximation of forward swap rates and currency basis spreads}\label{sssec:lgcm_gaussian_approx}

We now specialise the weight freezing technique to the LGCM. We consider a schedule $T_0 < T_1 < \ldots < T_n$ with accrual fractions $\theta_j$, in the foreign currency $f$ under domestic collateral $d$. This is the case where collateral basis volatility matters. The native case is recovered by switching off the collateral basis block (see \autoref{rema:gauss_native}). We write
\begin{eqnarray*}
  P_{t,j} = Z^{f|d}_{t,T_j}, \quad
  A_{t} = \sum_{j=1}^n \theta_j P_{t,j}, \quad
  L_{t,j} = L^{f|d}_{t,T_{j-1};\theta_j}, \quad
  L^{r}_{t,j} = L^{f/d|d,r}_{t,T_{j-1};\theta_j}, \quad
  B_{t,j} = L^{r}_{t,j} - L_{t,j} ,
\end{eqnarray*}
for the discount factors, the annuity, the foreign and resetting forward Libors and the swaplet currency basis spreads. The forward swap rate and the forward resetting currency basis spread are
\begin{eqnarray}
  Y_t ~=~ Y^{f|d}_{t,T_0;M} ~=~ \frac{1}{A_t}\sum_{j=1}^n \theta_j P_{t,j} L_{t,j} , \qquad
  B_t ~=~ B^{f/d|d,r}_{t,T_0;M} ~=~ \frac{1}{A_t}\sum_{j=1}^n \theta_j P_{t,j} B_{t,j} .
\end{eqnarray}
Both are martingales under the annuity measure $\QQ^{f|d,A}$. We use the stacked state $\hat{z}_t$, $\hat{G}_{t,T}$, $\hat{y}_t$ of \autoref{rema:qgcm_vfc}. Its volatility is $\hat{\sigma}_{z,t} = [\sigma^f_{z,t}|\sigma^c_{z,t}]$ and its transport matrix is $\hat{\mathcal{M}}_{u,t} = \diag(\mathcal{M}^f_{u,t},\mathcal{M}^c_{u,t})$.

\begin{lemm}\label{lemm:zhat_forward_law}
In the LGCM, under the $T$-forward measure $\QQ^{f|d,T}$, $\hat{z}_T$ is a centred Gaussian vector with covariance $\hat{y}_T$.
\end{lemm}

\begin{lemm}\label{lemm:par_spread_state_functions}
In the LGCM, $Y_t$ and $B_t$ are explicit functions of $\hat{z}_t$:
\begin{eqnarray*}
  P_{t,j} &=& Z^{f|d}_{0,t,T_j} \exp\left( -\hat{G}^{\top}_{t,T_j} \hat{z}_t - \frac{1}{2} \hat{G}^{\top}_{t,T_j}  \hat{y}_t \hat{G}_{t,T_j} \right) \\
  1+\theta_j L_{t,j} &=& (1-\theta_j S^{f|f,\div}_{0,T_{j-1};\theta_j}) \, \frac{Z^{f|f}_{t,T_{j-1}}}{Z^{f|f}_{t,T_j}} \, e^{- \int_t^{T_{j-1}} \Sigma^{f|f}_{u,T_{j-1};\theta_j}\cdot\Sigma^{f\to d}_{u,T_j}\, du}  \\
  1+\theta_j L^{r}_{t,j} &=& (1-\theta_j S^{d|d,\div}_{0,T_{j-1};\theta_j}) \, \frac{P_{t,j-1}}{P_{t,j}}
\end{eqnarray*}
where $Z^{f|f}_{t,T}$ is given by \autoref{theo:qgcm} with $z^f_t = \hat{z}^f_t - v^{f,c}_t$ and $v^{f,c}_t$ deterministic.
\end{lemm}

\begin{prop}[Gaussian approximation of forward swap rates]\label{prop:gauss_swap_rate}
Let us define the deterministic loadings, computed from forward values seen from $t_0 = 0$ ($P_j = Z^{f|d}_{0,T_0,T_j}$, $L_j = L^{f|d}_{0,T_{j-1};\theta_j}$, $Y = Y^{f|d}_{0,T_0;M}$, $A = \sum_j\theta_jP_j$):
\begin{eqnarray}
  q^{f}_Y &=& \frac{1}{A} \sum_{j=1}^n P_j\left( (1+\theta_j Y) G^f_{T_0,T_j} - (1+\theta_j L_j) G^f_{T_0,T_{j-1}} \right) \\
  q^{c}_Y &=& \frac{1}{A} \sum_{j=1}^n \theta_j P_j\, (Y - L_j)\, G^c_{T_0,T_j}
\end{eqnarray}
and $\hat{q}_Y = [q^f_Y ; q^c_Y]$. Then, freezing the gradient of $Y_t$ with respect to the state at its forward value, the swap rate has the Gaussian dynamic under $\QQ^{f|d,A}$
\begin{eqnarray}
  dY_t &\approx& \hat{q}_Y^\top \hat{\mathcal{M}}_{t,T_0} \hat{\sigma}^\top_{z,t} \, dW^{f|d,A}_t
\end{eqnarray}
so that $Y_{T_0}$ is approximately normal with mean $Y_0$ and variance
\begin{eqnarray}
  \VV^{f|d,A}\left[Y_{T_0}\right] &\approx& \hat{q}_Y^\top \hat{y}_{T_0}\, \hat{q}_Y
  ~=~ \sum_{(a,i),(b,j)} q^{(a,i)}_Y q^{(b,j)}_Y \int_0^{T_0} \rho^{(a,i),(b,j)}_u \sigma^{a,i}_u\sigma^{b,j}_u M^{a,i}_{u,T_0}M^{b,j}_{u,T_0}\, du .
  \label{eq:gauss_swap_rate_variance}
\end{eqnarray}
\end{prop}

\begin{prop}[Gaussian approximation of currency basis spreads]\label{prop:gauss_cb_spread}
With the same notations and $B_j = B^{f/d|d,r}_{0,T_{j-1};\theta_j}$, $L^r_j = L_j + B_j$, $B = B^{f/d|d,r}_{0,T_0;M}$, let
\begin{eqnarray}
  q^{f}_B &=& \frac{1}{A} \sum_{j=1}^n \theta_j P_j\left( B\, G^f_{T_0,T_j} - B_j\, G^f_{T_0,T_{j-1}} \right) \\
  q^{c}_B &=& \frac{1}{A} \sum_{j=1}^n P_j\left( (1+\theta_j L_j + \theta_j B) G^c_{T_0,T_j} - (1+\theta_j L^r_j) G^c_{T_0,T_{j-1}} \right) .
\end{eqnarray}
Then $B_{T_0}$ is approximately normal under $\QQ^{f|d,A}$ with mean $B_0$ and variance $\hat{q}_B^\top \hat{y}_{T_0}\, \hat{q}_B$.
\end{prop}

The proofs are in \autoref{sec:gaussian_approx_proofs}.

\begin{rema}[Consistency checks]
\begin{enumerate}
\item For a single swaplet ($n=1$), $G^{a}_{T_0,T_0}=0$ and the propositions give the loadings $q^c_B = \frac{1+\theta L^{r}_{1}}{\theta} G^c_{T_0,T_1}$ and $q^f_B = B\, G^f_{T_0,T_1}$. This is exactly the approximate dynamic of the swaplet currency basis spread in the section on reconstruction of market observables, since $\Sigma^{a}_{u,T_0;\theta} = \sigma^a_{z,u}\mathcal{M}^a_{u,T_0}G^a_{T_0,T_1}$.
\item The collateral basis loading of the swap rate is proportional to $Y - L_j$. As already observed, it vanishes when the index curve is flat. In the numerical example below it contributes $3\cdot 10^{-4}$ of the variance of $Y$.
\item Symmetrically, the foreign rates loading of the currency basis spread is proportional to the level of the basis. The collateral basis loading is of order $(1+\theta L)/\theta$ times the collateral basis volatility of the swaplets.
\end{enumerate}
\end{rema}

\begin{rema}[Native swaps]\label{rema:gauss_native}
For a swap in currency $e$ under native collateral $e$, the collateral basis block disappears, $P_j = Z^{e|e}_{0,T_0,T_j}$, $1+\theta_jL_j = (1-\theta_jS^{e|e,\div})P_{j-1}/P_j$, and \eqref{eq:gauss_swap_rate_variance} reduces to the classic LGM formula $q^{e,\top}_Y y^{e,e}_{T_0} q^e_Y$ with
\begin{eqnarray*}
  q^e_Y &=& \frac{1}{A}\sum_{j=1}^n P_j\left( (1+\theta_j Y) G^e_{T_0,T_j} - (1+\theta_jL_j)G^e_{T_0,T_{j-1}}\right) ,
\end{eqnarray*}
cf. \cite[\S13]{AndersenPiterbarg}. The variance is a quadratic form in the deterministic covariance matrix $y_{T_0}$ of the state, so it is computed by the covariance integral algorithm of \autoref{ssec:covariance_integrals} at no extra cost.
\end{rema}

\begin{rema}[Accuracy]
The approximation is asymptotic in the volatility: it neglects the variation of the gradient $\nabla_{\hat{z}}Y$ along the path. We checked it against the exact distribution of $Y_{T_0}$ and $B_{T_0}$, obtained from \autoref{lemm:zhat_forward_law} and \autoref{lemm:par_spread_state_functions} by 2D Gauss-Hermite quadrature ($160^2$ nodes). The setup was a $5$Y$\times10$Y quarterly swap with one foreign factor ($\sigma^f \approx 1.2\%$, $\beta^f\approx 3\%$) and one collateral basis factor ($\sigma^c\approx 0.4\%$, $\beta^c \approx 15\%$), $\rho^{f,c} = -0.4$, piecewise constant parameters and all volatilities scaled by $\lambda$. The exact pricer reproduces $\EE^{f|d,A}[Y_{T_0}] = Y_0$ and $\EE^{f|d,A}[B_{T_0}] = B_0$ to $10^{-16}$. The relative errors of the approximate normal volatility $\sqrt{\VV/T_0}$ are:
\begin{center}
\begin{tabular}{|c|c|c|c|c|}
\hline
$\lambda$ & normal vol of $Y$ (bp) & rel. error & normal vol of $B$ (bp) & rel. error \\
\hline
$0.25$ & $32.9$  & $2.0\cdot 10^{-4}$ & $4.08$ & $-3.8\cdot 10^{-5}$ \\
$0.5$  & $65.7$  & $8.1\cdot 10^{-4}$ & $8.16$ & $-1.5\cdot 10^{-4}$ \\
$1$    & $131.1$ & $3.3\cdot 10^{-3}$ & $16.3$ & $-6.2\cdot 10^{-4}$ \\
$2$    & $259.0$ & $1.3\cdot 10^{-2}$ & $32.5$ & $-2.6\cdot 10^{-3}$ \\
$4$    & $494.5$ & $5.1\cdot 10^{-2}$ & $64.7$ & $-1.2\cdot 10^{-2}$ \\
\hline
\end{tabular}
\end{center}
The error is multiplied by $4$ each time the volatilities are doubled, i.e. the relative error is $O(\lambda^2)$ (numerical observation). This is consistent with the first-order variation of the gradient averaging out in the variance. In this example the currency basis approximation is about five times more accurate than the swap rate one.
\end{rema}

\begin{rema}[Calibration]
\autoref{prop:gauss_cb_spread} provides the model delivered volatility of currency basis spreads needed by the bootstrapping algorithm of the calibration section. It is explicit in the model parameters and only requires the covariance integrals of \autoref{ssec:covariance_integrals}, so it is cheap to evaluate inside the bootstrapping.
\end{rema}


\newpage%
\section{The Linear Gaussian Collateral Model: Implementation}


%
%Reference is Andersen-Piterbarg, 10.1.6 and 12.1.8.
%
%\subsection{Simulation}
%
%The rates processes are driven by OU processes.
%
%The FX processes require 
%
%\subsection{Num\'eraire}
%
%\subsubsection{Drift under the risk neutral measure}
%
%
%\subsubsection{Drift under the jumping num\'eraire measure}


\subsection{Covariance integrals}\label{ssec:covariance_integrals}

\subsubsection{Introduction}
The expectation and covariance matrix necessary for simulation, as well as all the convexity adjustments in the model, reduce to sums of integrals of the type:
\begin{eqnarray}
  I^{a,b}_{t,S;T_a,T_b} &=& \int_t^S \Sigma^a_{u,T_a} \cdot \Sigma^b_{u,T_b} \, du
  ~=~ \sum_{j_a=1}^{N_a} \sum_{j_b=1}^{N_b} \int_t^S \rho^{(a,j_a),(b,j_b)}_u \sigma^{a,j_a}_u  \sigma^{b,j_b}_u G^{a,j_a}_{u,T_a} G^{b,j_b}_{u,T_b} \, du
\end{eqnarray}
where we set $G^{x}_{u,T} = 1$ for the FX factor. The simulation also requires the same integrals with $G^{a,j}_{u,T}$ replaced by $M^{a,j}_{u,T}$. So we consider the elementary integrals
\begin{eqnarray}
  J_{t,S}\left[\ell^{a,i},\ell^{b,j}\right] &=& \int_t^S \rho^{(a,i),(b,j)}_u \sigma^{a,i}_u  \sigma^{b,j}_u \, \ell^{a,i}_u \, \ell^{b,j}_u \, du
\end{eqnarray}
where each \emph{loading} $\ell^{a,i}_u$ is one of $M^{a,i}_{u,T}$, $G^{a,i}_{u,T}$ or $1$, for some maturity $T\ge S$. Sums of these give all the integrals we need. For example, $\Sigma^a_{u,T;\theta}$ has loadings $G^{a,i}_{u,T+\theta} - G^{a,i}_{u,T}$, and $y^{a,b}_t$ has entries $J_{0,t}[M^{a,i}_{\cdot,t},M^{b,j}_{\cdot,t}]$. For vector loadings $\Lambda_u = \sum_{(a,i)} c_{a,i}\,\sigma^{a,i}_u\ell^{a,i}_u\,\widetilde{e}_{(a,i)}$, we write $J_{t,S}[\Lambda,\Lambda']$ for the corresponding bilinear combination.

\subsubsection{Piecewise constant parameters}
We assume the spot vols, correlations and mean reversion parameters are piecewise constant on a grid $(t_k)$. On a segment $[t_k,t_{k+1}]$ of length $\Delta_k$, we write $\beta = \beta^{a,i}_k$ and $s = t_{k+1}-u \in [0,\Delta_k]$, and introduce
\begin{eqnarray}
  \varphi_\beta(s) ~=~ \int_0^s e^{-\beta v}\,dv ~=~ \frac{1-e^{-\beta s}}{\beta} \qquad (\varphi_0(s) = s).
\end{eqnarray}

\begin{lemm}\label{lemm:affine_loadings}
For $u\in[t_k,t_{k+1}]$ and $T\ge t_{k+1}$, every loading is affine in $\varphi_\beta(t_{k+1}-u)$:
\begin{eqnarray}
  \ell_u &=& p + q\,\varphi_\beta(t_{k+1}-u)
\end{eqnarray}
with
\begin{eqnarray*}
  (p,q) &=& \left( M_{t_{k+1},T}\, ,\, -\beta M_{t_{k+1},T} \right) \quad \textrm{for $\ell = M_{\cdot,T}$,} \\
  (p,q) &=& \left( G_{t_{k+1},T}\, ,\, 1-\beta G_{t_{k+1},T} \right) \quad \textrm{for $\ell = G_{\cdot,T}$,} \\
  (p,q) &=& (1,0) \quad \textrm{for $\ell = 1$.}
\end{eqnarray*}
\end{lemm}

\begin{proof}
We have $M_{u,T} = e^{-\beta s} M_{t_{k+1},T}$ and $G_{u,T} = \varphi_\beta(s) + e^{-\beta s}G_{t_{k+1},T}$, and $e^{-\beta s} = 1-\beta\varphi_\beta(s)$.
\end{proof}

\begin{prop}[Covariance integrals]\label{prop:covariance_integrals}
With the notations of \autoref{lemm:affine_loadings}, for loadings $\ell^{a,i}$ and $\ell^{b,j}$ with coefficients $(p_a,q_a)$, $(p_b,q_b)$ and mean reversions $\alpha = \beta^{a,i}_k$, $\beta = \beta^{b,j}_k$ on the segment,
\begin{eqnarray}
  J_{t_k,t_{k+1}}\left[\ell^{a,i},\ell^{b,j}\right] &=& \rho_k \sigma^{a,i}_k\sigma^{b,j}_k \left( p_ap_b\Delta_k + p_aq_b\Phi_1(\beta) + q_ap_b\Phi_1(\alpha) + q_aq_b\Phi_2(\alpha,\beta) \right)
\end{eqnarray}
where
\begin{eqnarray}
  \Phi_1(\beta) &=& \int_0^{\Delta}\varphi_\beta(s)\,ds ~=~ \frac{\Delta - \varphi_\beta(\Delta)}{\beta} , \\
  \Phi_2(\alpha,\beta) &=& \int_0^{\Delta}\varphi_\alpha(s)\varphi_\beta(s)\,ds ~=~ \frac{\Delta - \varphi_\alpha(\Delta) - \varphi_\beta(\Delta) + \varphi_{\alpha+\beta}(\Delta)}{\alpha\beta} .
\end{eqnarray}
\end{prop}

\begin{proof}
Expand the product of the two affine loadings and use $\varphi_\alpha\varphi_\beta = (1-e^{-\alpha s}-e^{-\beta s}+e^{-(\alpha+\beta)s})/(\alpha\beta)$.
\end{proof}

\begin{rema}[Stable evaluation]
The singularities of $\Phi_1$ and $\Phi_2$ at zero mean reversion are removable but the closed forms lose precision near them. Zero mean reversion is the Ho-Lee case and does occur in practice, for example for the FX factor or a level factor. Both functions are entire:
\begin{eqnarray*}
  \Phi_1(\beta) ~=~ \sum_{n\ge 0} \frac{(-\beta)^n\Delta^{n+2}}{(n+2)!} , \qquad
  \Phi_2(\alpha,\beta) ~=~ \sum_{m,n\ge 0} \frac{(-\alpha)^m(-\beta)^n\Delta^{m+n+3}}{(m+1)!\,(n+1)!\,(m+n+3)} .
\end{eqnarray*}
We use these series when $|\beta|\Delta < 0.1$ (resp. $\max(|\alpha|,|\beta|)\Delta < 0.5$) and the closed forms, computed with \texttt{expm1}, when $\min(|\alpha|,|\beta|,|\alpha+\beta|)\Delta > 0.05$. In the remaining mixed case, where one rate is small and the other is not, we integrate $\varphi_\alpha\varphi_\beta$ by 64-point Gauss-Legendre quadrature. The integrand is entire, so this is exact to machine precision for the values of $|\beta|\Delta$ met in practice.
\end{rema}

\begin{algo}[Covariance integrals]\label{algo:covariance_integrals}
Given the parameter grid $(t_k)$:
\begin{enumerate}
\item Precompute, for each factor, $\int_0^{t_k}\beta_v dv$ on the grid. Then $M_{t,T} = e^{-\int_t^T\beta_v dv}$ for any $t$, $T$.
\item Compute $G_{t_{k},T}$ for all grid points $t_k\le T$ by the backward recursion
\begin{eqnarray*}
  G_{t_k,T} &=& \varphi_{\beta_k}(\Delta_k) + e^{-\beta_k \Delta_k}\, G_{t_{k+1},T} ,
\end{eqnarray*}
starting from the segment containing $T$. It avoids the representation $G_{t,T} = (\Lambda_T - \Lambda_t)/h_t$ with $\Lambda_T = G_{0,T}$, whose difference of two large numbers can lose precision when $t$ is large or the mean reversion is negative.
\item Split $[t,S]$ at the grid points. On each segment compute $(p,q)$ for both loadings by \autoref{lemm:affine_loadings} and the contribution by \autoref{prop:covariance_integrals}. Sum the contributions.
\end{enumerate}
The cost is $O(K)$ per pair of factors, where $K$ is the number of segments in $[t,S]$, i.e. $O(K N_a N_b)$ for $I^{a,b}_{t,S;T_a,T_b}$. All quantities are closed form.
\end{algo}

\begin{rema}[Numerical verification]
We compared \autoref{algo:covariance_integrals} to adaptive quadrature in which $G$ itself is computed by quadrature of $M$. The test used $84$ combinations of loadings ($M$, $G$, $1$), sub-intervals not aligned on the grid and maturities up to $12$Y, with one factor having zero and $10^{-7}$ mean reversion and one factor negative mean reversion on some segments. The maximum relative error is $7\cdot 10^{-16}$. The kernels $\Phi_1$, $\Phi_2$ alone agree with quadrature to $4\cdot 10^{-16}$, including the cases $\alpha=0$, $\alpha = 10^{-9}$, $\alpha+\beta = 0$ and $\beta\Delta = 10$.
\end{rema}


\subsection{Convexity adjustments}
All collateral convexity and quanto adjustments of the deterministic volatility section are exponentials of covariance integrals, computed by \autoref{algo:covariance_integrals}. For example:
\begin{eqnarray*}
  \ln \frac{x^{df|j}_{t,T}}{x^{df|i}_{t,T}} &=& -\int_t^T \Sigma^{df}_{u,T}\cdot\Sigma^{i\to j}_{u,T}\,du , \\
  \ln \frac{1+\theta L^{e|j}_{t,T;\theta}}{1+\theta L^{e|i}_{t,T;\theta}} &=& -\int_t^T \Sigma^{e|e}_{u,T;\theta}\cdot\Sigma^{i\to j}_{u,T+\theta}\,du ,
\end{eqnarray*}
where $\Sigma^{df}_{u,T}$ has loadings $G^{d}_{u,T}$, $-G^f_{u,T}$, $-G^c_{u,T}$ and $1$, and $\Sigma^{e|e}_{u,T;\theta}$ has loadings $G^e_{u,T+\theta} - G^e_{u,T}$.


\subsection{CurveEngine implementation}

In practice, all these computations can be abstracted in a helper class.

\subsubsection{Reconstruction formulas}



\newpage%
\subsection{Monte Carlo simulation}\label{ssec:lgcm_simulation}

In the LGCM, $y^{a,b}_t$ and $v^{f,c}_t$ are deterministic. They are precomputed as $y^{(a,i),(b,j)}_t = J_{0,t}[M^{a,i}_{\cdot,t},M^{b,j}_{\cdot,t}]$ and $v^{(f,i)}_t = \sum_{j} J_{0,t}[M^{f,i}_{\cdot,t},G^{c,j}_{\cdot,t}]$. The state $z_t = [z^d_t | z^f_t | z^c_t | z^x_t]$ of \autoref{theo:qgcm} is Gaussian, and its transition between two dates $s<t$ can be sampled exactly, whatever the step size. For a vector loading $\Lambda$, we write $J[\Lambda,\Lambda'] = J_{s,t}[\Lambda,\Lambda']$ and use the loadings
\begin{eqnarray*}
  \Sigma^a_{\cdot,t}: \ \left(G^{a,j}_{\cdot,t}\right)_j , \qquad
  \Sigma^{df}_{\cdot,t} = \Sigma^d_{\cdot,t} - \Sigma^f_{\cdot,t} - \Sigma^c_{\cdot,t} + \sigma^{df} , \qquad
  \hat{\Sigma}_{\cdot,t} = \Sigma^f_{\cdot,t} + \Sigma^c_{\cdot,t} = \Sigma^{f|d}_{\cdot,t} .
\end{eqnarray*}

\subsubsection{Transition and covariance matrix}

\begin{prop}[Exact transition under the risk-neutral measure]\label{prop:lgcm_transition_rn}
Let $s<t$, $G^a = G^a_{s,t}$ and $\hat{G} = [G^f;G^c]$. Under $\QQ^{d,\RN}$, conditionally on $\mathcal{F}_s$,
\begin{eqnarray*}
  z^a_t &=& \mathcal{M}^a_{s,t}z^a_s + \mu^a_{s,t} + \varepsilon^a , \qquad a\in\{d,f,c\} \\
  z^x_t &=& z^x_s + G^{d,\top}z^d_s - \hat{G}^\top\hat{z}_s + \mu^x_{s,t} + \varepsilon^x \\
  \int_s^t r^{d|d}_u\,du &=& -\ln Z^{d|d}_{0,s,t} + G^{d,\top}z^d_s + \mu^r_{s,t} + \varepsilon^r \\
  \int_s^t q^{f\to d}_u\,du &=& -\ln \QZ^{f|f\to d}_{0,s,t} + G^{c,\top}z^c_s + G^{f,\top} v^{f,c}_s + \mu^q_{s,t} + \varepsilon^q
\end{eqnarray*}
with $\hat{z}_s = [z^f_s + v^{f,c}_s ; z^c_s]$, deterministic means
\begin{eqnarray*}
  \mu^{d,i}_{s,t} &=& \textstyle\sum_{j} J[M^{d,i}_{\cdot,t},G^{d,j}_{\cdot,t}] + M^{d,i}_{s,t}\left(y^{d,d}_sG^d\right)_i \\
  \mu^{f,i}_{s,t} &=& \textstyle\sum_{j} J[M^{f,i}_{\cdot,t},G^{f,j}_{\cdot,t}] + M^{f,i}_{s,t}\left(y^{f,f}_sG^f\right)_i - J[M^{f,i}_{\cdot,t},1^x] \\
  \mu^{c,i}_{s,t} &=& \textstyle\sum_{j} J[M^{c,i}_{\cdot,t},G^{c,j}_{\cdot,t}] + \sum_{j} J[M^{c,i}_{\cdot,t},G^{f,j}_{\cdot,t}] + M^{c,i}_{s,t}\left(y^{c,c}_sG^c + y^{c,f}_sG^f\right)_i - J[M^{c,i}_{\cdot,t},1^x] \\
  \mu^x_{s,t} &=& \frac{1}{2}G^{d,\top}y^{d,d}_sG^d - \frac{1}{2}\hat{G}^\top\hat{y}_s\hat{G} - \frac{1}{2}J[\Sigma^{df}_{\cdot,t},\Sigma^{df}_{\cdot,t}] + J[\Sigma^{df}_{\cdot,t},\Sigma^{d}_{\cdot,t}] \\
  \mu^r_{s,t} &=& \frac{1}{2}G^{d,\top}y^{d,d}_sG^d + \frac{1}{2}J[\Sigma^{d}_{\cdot,t},\Sigma^{d}_{\cdot,t}] \\
  \mu^q_{s,t} &=& \frac{1}{2}G^{c,\top}y^{c,c}_sG^c + G^{c,\top}y^{c,f}_sG^f + \frac{1}{2}J[\Sigma^{c}_{\cdot,t},\Sigma^{c}_{\cdot,t}] + J[\Sigma^{f}_{\cdot,t},\Sigma^{c}_{\cdot,t}] - J[\Sigma^{c}_{\cdot,t},\sigma^{df}]
\end{eqnarray*}
(where $1^x$ denotes the unit loading of the FX factor) and a centred Gaussian noise
\begin{eqnarray*}
  \varepsilon^{a,i} = \int_s^t \sigma^{a,i}_u M^{a,i}_{u,t} \, d\widetilde{W}^{(a,i)}_u , \quad
  \varepsilon^x = \int_s^t \Sigma^{df}_{u,t}\cdot dW_u , \quad
  \varepsilon^r = \int_s^t \Sigma^{d}_{u,t}\cdot dW_u , \quad
  \varepsilon^q = \int_s^t \Sigma^{c}_{u,t}\cdot dW_u ,
\end{eqnarray*}
independent of $\mathcal{F}_s$. Its covariance matrix is $C_{s,t} = \left(J[\Lambda_k,\Lambda_l]\right)_{k,l}$, where $\Lambda_k$ runs over the loadings of the noises: $M^{a,i}_{\cdot,t}$, $\Sigma^{df}_{\cdot,t}$, $\Sigma^d_{\cdot,t}$ and $\Sigma^c_{\cdot,t}$.
\end{prop}

\begin{proof}
\emph{Rates factors.} Integrating the dynamic of \autoref{theo:qgcm}, $z^a_t = \mathcal{M}^a_{s,t}z^a_s + \int_s^t\mathcal{M}^a_{u,t}(\textrm{drift}^a_u\,du + \sigma^{a,\top}_{z,u}dW_u)$. For the term $y^{a,b}_u\vun$ of the drift, Fubini gives
\begin{eqnarray*}
  \int_s^t \mathcal{M}^a_{u,t}y^{a,b}_u\vun\,du ~=~ \int_s^t \mathcal{M}^a_{w,t}\sigma^{a,\top}_{z,w}\sigma^b_{z,w}G^b_{w,t}\,dw + \mathcal{M}^a_{s,t}y^{a,b}_s G^b_{s,t} ,
\end{eqnarray*}
and $\int_s^t\mathcal{M}^a_{u,t}\sigma^{a,\top}_{z,u}\sigma^{df}_u\,du$ gives the terms $J[M^{a,i}_{\cdot,t},1^x]$.

\emph{FX.} The FX forward $x^{df|d}_{u,t}$ is a lognormal $\QQ^{d|d,t}$-martingale with volatility $\Sigma^{df}_{u,t}$ and $x^{df|d}_{t,t} = x^{df}_t$. So $\ln x^{df}_t = \ln x^{df|d}_{s,t} - \frac{1}{2}J[\Sigma^{df},\Sigma^{df}] + \int_s^t \Sigma^{df}_{u,t}\cdot dW^{d|d,t}_u$. We express $x^{df|d}_{s,t} = x^{df}_sZ^{f|d}_{s,t}/Z^{d|d}_{s,t}$ with the reconstruction formulas and use $dW^{d|d,t}_u = dW_u + \Sigma^d_{u,t}du$.

\emph{Numeraire.} Integrating $d\ln Z^{d|d}_{u,t} = (r^{d|d}_u - \frac{1}{2}|\Sigma^d_{u,t}|^2)du - \Sigma^d_{u,t}\cdot dW_u$ from $s$ to $t$ gives $\int_s^tr^{d|d}_u du = -\ln Z^{d|d}_{s,t} + \frac{1}{2}J[\Sigma^d,\Sigma^d] + \varepsilon^r$, and we conclude with the reconstruction formula. The same argument applied to $Z^{f|d}_{u,t}$ and $Z^{f|f}_{u,t}$ under $\QQ^{f,\RN}$, with $dW^{f,\RN} = dW - \sigma^{df}du$, gives $\int_s^t q^{f\to d}_u\,du$.
\end{proof}

\begin{rema}
The covariance matrix $C_{s,t}$ does not depend on the measure since all changes of measure we consider have deterministic Girsanov kernels. It is computed once per time step and its Cholesky (or eigenvalue) decomposition is stored. The state is then simulated as $X_t = F_{s,t}X_s + \mu_{s,t} + C_{s,t}^{1/2}\xi$ with $\xi\sim N(0,I)$. Only the components needed by the chosen num\'eraire are simulated: $\int r^{d|d}$ under $\QQ^{d,\RN}$, nothing under the forward and jumping num\'eraire measures.
\end{rema}

\subsubsection{Drifts under the risk-neutral, $T$-forward and jumping num\'eraire measures}

\begin{prop}[Change of num\'eraire]\label{prop:lgcm_drifts}
Let $\varepsilon^k = \int_s^t\Lambda_{k,u}\cdot dW_u$ be any of the noises of \autoref{prop:lgcm_transition_rn}.
\begin{enumerate}
\item Under the domestic $T$-forward measure $\QQ^{d|d,T}$, $T\ge t$, the transition is the same with means
\begin{eqnarray*}
  \mu^{k,T}_{s,t} &=& \mu^{k}_{s,t} - J[\Lambda_k, \Sigma^d_{\cdot,T}]
\end{eqnarray*}
and the same covariance matrix. For example, for $T=t$ the domestic rates means reduce to $\mu^{d,t}_{s,t} = \mathcal{M}^d_{s,t}y^{d,d}_sG^d_{s,t}$, and the FX mean to $\frac{1}{2}G^{d,\top}y^{d,d}_sG^d - \frac{1}{2}\hat{G}^\top\hat{y}_s\hat{G} - \frac{1}{2}J[\Sigma^{df},\Sigma^{df}]$.
\item The jumping num\'eraire associated with a tenor structure $T_0 = 0 < T_1 < \ldots$ is
\begin{eqnarray}
  N_t &=& Z^{d|d}_{t,T_{q(t)+1}} \prod_{k=0}^{q(t)} \frac{1}{Z^{d|d}_{T_k,T_{k+1}}} , \qquad q(t) = \max\{k : T_k\le t\} .
\end{eqnarray}
It corresponds to rolling an investment in the zero coupon bond maturing at the next tenor date. On $[T_k,T_{k+1}]$ the associated measure $\QQ^{d,J}$ coincides with the $T_{k+1}$-forward measure. So, when the simulation grid is contained in the tenor structure, the means of a step $[t_k,t_{k+1}]$ under $\QQ^{d,J}$ are given by 1.\ with $T = t_{k+1}$.
\end{enumerate}
\end{prop}

\begin{proof}
1. The Girsanov kernel is deterministic: $dW^{d|d,T}_u = dW_u + \Sigma^d_{u,T}du$, so $\varepsilon^k = \int_s^t\Lambda_{k,u}\cdot dW^{d|d,T}_u - J[\Lambda_k,\Sigma^d_{\cdot,T}]$. For the domestic factors with $T=t$, the term $\sum_j J[M^{d,i}_{\cdot,t},G^{d,j}_{\cdot,t}]$ cancels with the shift. For the FX, the shift cancels the last term of $\mu^x$.

2. $N$ is continuous at the tenor dates. Between $T_k$ and $T_{k+1}$ it is proportional to $Z^{d|d}_{t,T_{k+1}}$, with an $\mathcal{F}_{T_k}$-measurable factor, so the density of the measure change is that of the $T_{k+1}$-forward measure on that interval.
\end{proof}

\begin{rema}
The jumping num\'eraire has several advantages:
\begin{itemize}
\item The num\'eraire update $N_{T_{k+1}} = N_{T_k}/Z^{d|d}_{T_k,T_{k+1}}$ only requires the reconstruction formula at $T_k$, whereas $\QQ^{d,\RN}$ requires the simulation of $\int r^{d|d}_u\,du$.
\item On each step, the means are the simplest ones: the domestic factors only carry $\mathcal{M}^d_{s,t}y^{d,d}_sG^d_{s,t}$, and the log FX spot is the log FX forward $\ln x^{df|d}_{s,t}$, lognormally corrected.
\end{itemize}
An earlier version of the formula for $N_t$ used $Z_{t,T_{q(t)}}$, which is not defined since $T_{q(t)}\le t$.
\end{rema}

\begin{rema}[Numerical verification]
We propagated the exact Gaussian law of the state through a quarterly grid up to $3$Y with the transitions above and the model of the verification of \autoref{theo:qgcm}. Under $\QQ^{d,\RN}$, the mean and covariance agree with those obtained by integrating the moment ODEs of \autoref{theo:qgcm} to $4\cdot 10^{-17}$ and $8\cdot 10^{-17}$ (absolute). The martingale identities for $Z^{d|d}_{\cdot,T}$, $x^{df}Z^{f|d}_{\cdot,T}$ and $x^{df}Z^{f|f}_{\cdot,T}$ (under $\QQ^{d,\RN}$), for $Z^{d|d}_{\cdot,T}/Z^{d|d}_{\cdot,T^*}$ and the FX forward (under $\QQ^{d|d,T^*}$, $T^* = 20$Y), and for $Z^{d|d}_{\cdot,T}/N$ and $x^{df}Z^{f|d}_{\cdot,T}/N$ (under $\QQ^{d,J}$) all hold to machine precision. The measure shifts are not small: the means of the simulated variables at $3$Y differ by up to $1.4\cdot 10^{-4}$ between $\QQ^{d,\RN}$ and $\QQ^{d,J}$, and $1.4\cdot 10^{-2}$ between $\QQ^{d,\RN}$ and $\QQ^{d|d,T^*}$.
\end{rema}

\begin{rema}[Initial curves]\label{rema:simulation_initial_curves}
The transitions only involve the initial curves through $Z_{0,s,t}$ at simulation dates and $Z_{0,t,T}$ at payment dates. The instantaneous forward curve $r_{0,\cdot}$ never needs to be evaluated pointwise. In particular, the discontinuities of $r_{0,\cdot}$ created by piecewise flat interpolation have no effect on the simulation (see \autoref{ssec:interpolation_consistency}).
\end{rema}



\newpage%
\section{The Linear Gaussian Collateral Model: Calibration}

\setcounter{subsection}{-1}
\subsection{Introduction}

\subsubsection{Single currency market}% 
The model incorporates a Linear Gaussian Model for each currency. This sub-model is meant to price instruments with native collateral. 

In practice, this means that for each currency, we have an embedded single currency LGM which can be calibrated independently to its native IR market:
\begin{enumerate}
\item Native discount $Z^{x|x}_{t,T}$  and index curves  $L^{x|x}_{t,T;\theta}$ can be built by any standard method (independently of volatilities\footnote{Spot volatilities are included in Future convexity adjustments but they do not require a very precise calibration of the volatility so these convexity adjustments are typically marked externally if not hardcoded. Conversely, the values of the futures in the curve building have almost no impact on the volatilities calibrated to swaption vols.}). 
\item Volatilities $\sigma^{x|x}_{t,T}$  can be bootstrapped on the fly.
\item Mean reversions can be marked/calibrated externally as usual to match other ATM vols or, in the multi-dimensional case, some delivered rates correlations. 
\end{enumerate}

%The LGCM model can take as an input a fully calibrated LGM for each currency. Equivalently, the LGCM calibrator can, as a first step, call a calibrator for each of these sub-models. 
%
%This provides initial native discount and index curves as well as the spot volatilities $\sigma^{x|x}_{t,T}$ calibrated to the corresponding market implied volatilities.



\subsubsection{Cross-currency market}% 
The next step is to calibrate to the cross-currency market. We focus on the G10 currencies where liquid instruments are 
\begin{enumerate}
\item FX forwards (for short maturities)
\item Resetting currency basis swaps 
\item FX options  
\end{enumerate}
We will assume that they are collateralised in currency $d$ (USD in practice).

This leads to the following set of calibration targets for each benchmark expiry:
\begin{enumerate}
\item a resetting currency basis spread (or an FX forward)
\item an FX implied volatility  
\item a resetting cross-currency basis spread (historical) volatility 
\end{enumerate}



\subsection{Calibration to single currency markets}


\subsubsection{Calibration of spot-volatility to the cap market}% 


\subsubsection{Calibration of spot-volatility to the swaption market}% 



\subsubsection{Calibration of mean-reversion to ATM swaptions}% 



\subsubsection{Calibration of mean-reversion to swap rates correlations}% 



\newpage%
\subsection{Calibration to the cross-currency market}


\subsubsection{Necessity for joint calibration}%
Let us consider things from the curve building point of view. On the cross-currency market, the most liquid long-dated linear instruments are resetting currency basis swaps. Say we calibrated the model up to tenor $M$, i.e. up to date $T_{\f M}$ where $f = 4$ is the frequency of the swaps.  We now consider the calibration to the tenor $N$. We have 
\begin{enumerate}
\item one new target: the currency basis spread $B^{x/u|u,r}_{t_0,T;N}$
\item multiple new unknowns: for each $j=\f M+1,\ldots,\f N$
\begin{enumerate}
\item the discount factors $Z^{x|u}_{t_0,T_j}$  
\item the foreign forward rates $L^{x|u}_{t_0,T_{j-1};\theta}$ 
\item the resetting forward rates $L^{x/u|u}_{t_0,T_{j-1};\theta}$ 
\end{enumerate}
\end{enumerate}
The problem is clearly under-determined. 

However, the model provides convexity adjustments that relate the different forward rates to the ones from native single currency markets. But these convexity adjustments depend on the spot volatilities\footnote{as well as the cross-asset correlations. For now, we assume that they are marked but the same approach could be extended to calibrate instantaneous correlations by adding some terminal correlations/covariances as targets.}. Obviously these should be calibrated to match implied volatilities of option prices (FX options and currency basis swaptions) or historical estimations of volatility. 

If we assume that correlations are marked by traders, and MR are marked or calibrated externally, we are left with 3 types of unknown parameters:  
\begin{enumerate}
\item the cross-currency discount curve $Z^{f|d}_{t_0,T}$ 
\item the collateral basis spot volatilities  $\sigma^{f\to d}_{t,T}$ 
\item the FX spot volatilities $\sigma^{df}_{t}$
\end{enumerate}
And since they interact through convexity adjustments, they have to be jointly calibrated. 

Moreover, we want the calibration to be local: for risk management purposes we do not want a 2Y FX option to be exposed to 10Y FX implied vols. This implies that we need to bootstrap the parameters. 


\subsubsection{Bootstrapping}%

\paragraph{Introduction}
Convexity adjustments induce a lot of interdependence between market quotes. For risk management, we still need to be able to localise the risk. This means that the value of a 1Y FX option for example, should not depend on any market quote with expiry/maturity beyond 1Y. 

That is not the case in a global calibration. 

The natural way to achieve that is to "bootstrap" all the curves and term structure of parameters. 

This can be seen as breaking down the global calibration into steps by sorting calibration targets by expiry/maturity\footnote{This can be an interesting feature to add to the CalibrationTarget interface.} and calibrating model parameters iteratively.


%[TODO: Describe the algo!]



\paragraph{Intuition}
Let us assume we observe spot starting currency basis spreads for a given set of maturities $M_i$. 

Let us assume we calibrated the model up to expiry $M_i$. 

Then, we can compute the forward starting spread:
\begin{eqnarray}
  A^{f|d}_{0,T_0,T_{\f N}} B^{f/d|d,r}_{0,T_0,T_{\f N}} 
  &=& A^{f|d}_{0,T_0,T_{\f M}} B^{f/d|d,r}_{0,T_0,T_{\f M}}  + A^{f|d}_{0,T_{\f M},T_{\f N}} B^{f/d|d,r}_{0,T_{\f M},T_{\f N}} \\
  B^{f/d|d,r}_{0,T_0,T_{\f N}}  \sum_{j=1}^{\f N} \theta Z^{f|d}_{0,T_{j}} 
  &=& A^{f|d}_{0,T_0,T_{\f M}} B^{f/d|d,r}_{0,T_0,T_{\f M}}  + \left( \sum_{j=\f M+1}^{\f N} \theta Z^{f|d}_{0,T_{j}} \right) B^{f/d|d,r}_{0,T_{\f M},T_{\f N}} \\
  A^{f|d}_{0,T_0,T_{\f M}} \left(B^{f/d|d,r}_{0,T_0,T_{\f N}} - B^{f/d|d,r}_{0,T_0,T_{\f M}} \right) 
  &=&  \sum_{j=\f M+1}^{\f N} \theta Z^{f|d}_{0,T_{j}}  \left( B^{f/d|d,r}_{0,T_{j-1},T_{j}} -B^{f/d|d,r}_{0,T_0,T_{\f N}} \right)
\end{eqnarray} 
Unknowns are 
\begin{itemize}
\item Currency basis swaplet spreads $B^{f/d|d,r}_{0,T_j;\theta}$ 
\item Discount rates $R^{f|d}_{0,T_{j};\theta}$ 
\end{itemize}

The model formulas relate these to the model parameters of which 
\begin{itemize}
\item Collateral basis spot vol $\sigma^{f\to d}_{t}$, $t\in [T_{\f M},T_{\f N}[$,
\item FX spot vol\footnote{Actually it does not appear under constant funding spreads.} $\sigma^{df}_{t}$, $t\in [T_{\f M},T_{\f N}[$,
\end{itemize}
which we assume to be constant on the segment $[T_{\f M},T_{\f N}[$, are unknown. 

Let us choose a parametrisation for the discount curve $Z^{f|d}$ or at least the collateral spread $\QZ^{f|f\to d}$. For example, we can consider 
\begin{itemize}
\item piecewise constant forwards 
\item quadratic forwards together with continuity and derivability constraint at $T_{\f M}$ 
\end{itemize}
The point is that once the interpolation scheme has been chosen, the value of the discount curve on the segment is entirely specified by the average level of rates on the segment:
\begin{eqnarray}
  \overline{r}^{f|d}_{0,T_{\f M},T_{\f N}} &=& \frac{1}{T_{\f N}-T_{\f M}} \int_{T_{\f M}}^{T_{\f N}} r^{f|d}_{0,u} \, du 
\end{eqnarray} 
So we have reduced the determination of the discount factors $Z^{f|d}_{t_0,T_i}$ to that of a single parameter.

The problem is now well posed since we have 3 unknowns and 3 constraints as listed in \autoref{table:bootstrapping}.
\begin{center}
\begin{table}[h!]
\begin{tabular}{|c|c|}
\hline 
Unknowns & Constraints \\
\hline 
Average discount rate & Currency basis spread \\ 
ColBasis spot vol & Currency basis spread vol    \\
%IRF-ColB spot corr & IRF-CB terminal correlation   \\
\hline
FX spot vol & FX implied vol \\
\hline 
\end{tabular}
\medskip
\caption{Bootstrapping\label{table:bootstrapping}}
\end{table}
\end{center} 

Note that, while mathematically complex, the correspondence between model parameters and calibration targets is very intuitive. 

The only remaining step is a good approximation for model delivered volatilities of currency basis spreads. It is given by \autoref{prop:gauss_cb_spread}.


\subsubsection{Model usage and calibration targets}


\begin{rema}[Model usage and calibration targets]
The bootstrapping algorithm works as long as we have one target per expiry\footnote{Not true! We cannot apply the bootstrapping to a column of currency basis spreads with tenor $>3M$ for example as this would introduce some cyclic dependency between parameters.}. There is still a range of choices for the calibration targets that is up to discussion. The answer should be decided based on the model usage.

The model can be used in two distinct ways:
\begin{enumerate}
\item Valuation and risk management. One should specify a set of hedging instruments and calibrate the model these. When an option market is available and these are included in hedging instruments, the internal parameters should be calibrated to options prices.  When no option market is available as is the case for currency basis volatility risk, the covariance parameters should be calibrated to the historical covariances of the corresponding market quotes, i.e. currency basis spreads in this case. Ideally one would have only one set of parameters so that pricings are consistent across products and break-even levels are well defined. 
\item Benchmarking and reserve usage. One should mark covariance parameters that cannot be implied from market prices conservatively. In particular, these marks should depend on the structure of the portfolio.
\end{enumerate}

Note that since CB delta is going to be hedged with spot-starting resetting cross-currency swaps, one can argue that we should calibrate to their vol. At the same time, it is trivial to synthesise a forward-starting swap. 
\end{rema}




\begin{rema}[FX options market]%
The market ignores collateral convexity for FX options. It seems natural to assume FX forwards and volatilities correspond to USD collateralized products.

Note that in the Linear Gaussian model, the volatility of the FX process is invariant under collateral but the FX forward is not. This means that the calibration targets we should chose may not be the standard ATM options. We may want to calibrate to the volatilities corresponding to the adjusted FX forward. 

Assuming all other parameters are set and we chose target implied volatilities, the FX spot volatility $\sigma^{df}_\cdot$ can easily be bootstrapped to match ATM volatilities. Calibration can only fail if we encounter a variance squeeze.
\end{rema}


%\newpage%
\subsection{Bootstrapping algorithm}%


\subsubsection{Currency basis spread constraint}%

\begin{theo}
Then the model is calibrated to model prices iff the following relations hold:
\begin{enumerate}
\item Currency basis spread:
\begin{eqnarray}
  A^{f|d}_{t,t;M}  \left( B^{f/d|d,r}_{t,t;N} - B^{f/d|d,r}_{t,t;M} \right)
  &=& \sum_{j=\f M+1}^{\f N} \theta Z^{f|d}_{t,T_{\f M}} e^{- \overline{r}^{f|d}_{t,[M,N]} (T_j - T_{\f M})} \left( B^{f/d|d,r}_{t,T_{j-1};\theta} - B^{f/d|d,r}_{t,t;N} \right)
%   &=& \sum_{j=\f M+1}^{\f N} \theta Z^{f|d}_{t,T_{\f M}} Z^{f|d}_{t,T_{\f M},T_j}(\overline{r}^{f|d}_{t,[M,N]}) \left( B^{f/d|d,r}_{t,T_{j-1};\theta} - B^{f/d|d,r}_{t,t;N} \right)
\end{eqnarray}
\item Currency basis spread vol:
\item FX implied volatility:
\begin{eqnarray}
  \sigma^{df|c}_{\BS,t}(T)^2  (T-t)
  &=& \int_t^T \left| \Sigma^{d|d}_{u,T} - \Sigma^{f|f}_{u,T} +\Sigma^{d\to f}_{u,T}  + \sigma^{df}_u  \right|^2 du
\end{eqnarray}
\end{enumerate}
\end{theo}

%
%\begin{prop}
%%
%\begin{eqnarray}
%  A^{f|d}_{t,t;M}  \left( B^{f/d|d,r}_{t,t;N} - B^{f/d|d,r}_{t,t;M} \right)
%  &=& \sum_{j=\f M+1}^{\f N} \theta Z^{f|d}_{t,T_j}  \left( B^{f/d|d,r}_{t,T_{j-1};\theta} - B^{f/d|d,r}_{t,t;N} \right)
%\end{eqnarray}
%\begin{eqnarray}
%  A^{f|d}_{t,t;M}  \left( B^{f/d|d,r}_{t,t;N} - B^{f/d|d,r}_{t,t;M} \right)
%  &=& \sum_{j=\f M+1}^{\f N} \theta Z^{f|d}_{t,T_{\f M}}  e^{- \overline{r}^{f|d}_{t,[M,N]} (T_j - T_{\f M})} \left( B^{f/d|d,r}_{t,T_{j-1};\theta} - B^{f/d|d,r}_{t,t;N} \right)
%\end{eqnarray}
%
%\end{prop}
%

\subsubsection{Currency basis spread volatility constraint}%


\subsubsection{FX implied volatility constraint}%



\subsubsection{Bootstrapping algorithm}%

\begin{theo}
Let $N$ be the next tenor, solve for the average rate and collateral basis volatility satisfying: 
\begin{enumerate}
\item Currency basis spread:
\begin{eqnarray}
  A^{f|d}_{t,t;M}  \left( B^{f/d|d,r}_{t,t;N} - B^{f/d|d,r}_{t,t;M} \right)
  &=& \sum_{j=\f M+1}^{\f N} \theta Z^{f|d}_{t,T_{\f M}} Z^{f|d}_{t,T_{\f M},T_j}(\overline{r}^{f|d}) \left( B^{f/d|d,r}_{t,T_{j-1};\theta} - B^{f/d|d,r}_{t,t;N} \right)
\end{eqnarray}
where each currency basis swaplet is expended as:
\begin{eqnarray}
   B^{f/d|d,r}_{t,T_{j-1};\theta} &=&
\end{eqnarray}
\item Currency basis spread vol:
%\item FX implied volatility:
%\begin{eqnarray}
%  \sigma^{df|c}_{\BS,t}(T)^2  (T-t)
%  &=& \int_t^T \left| \Sigma^{d|d}_{u,T} - \Sigma^{f|f}_{u,T} +\Sigma^{d\to f}_{u,T}  + \sigma^{df}_u  \right|^2 du
%\end{eqnarray}
\end{enumerate}
\end{theo}






\newpage%
\subsection{Off-market collateral}%

%\begin{rema}[Off-market collateral]%
Once the model is calibrated, off-market curves can be computed in closed form using the FX forward convexity adjustment.
%\end{rema}




\newpage%
\subsection{Interpolation and consistency with the model dynamic}\label{ssec:interpolation_consistency}

The bootstrapping requires an interpolation scheme for the initial curves on each segment between calibration pillars. We discuss how this choice interacts with the model dynamic. Interpolation enters at two different levels: the initial curves at $t=0$, and the curves at future dates $t>0$.

\subsubsection{Initial curves}

\begin{rema}[Any interpolation is arbitrage free at $t=0$]
In the HJM framework, the initial curves $Z^{d|d}_{0,\cdot}$, $Z^{f|f}_{0,\cdot}$ and $\QZ^{f|f\to d}_{0,\cdot}$ are free initial conditions: the drifts are determined by the volatilities alone. So for any absolutely continuous interpolated initial curves (in particular piecewise flat forwards), \autoref{theo:qgcm} defines an arbitrage free model that reproduces them exactly. Contrary to Libor market models (see the remark on LMMs with stochastic currency basis), there is no stub interpolation to make consistent with the dynamic: discount factors at any date are given by the reconstruction formulas. Moreover, the exact simulation only uses the initial curves through discount factors at simulation and payment dates (\autoref{rema:simulation_initial_curves}).
\end{rema}

\begin{rema}[Interpolation and convexity adjustments decouple in the LGCM]
In the LGCM, the convexity adjustments (e.g.\ $\exp\{-\int\Sigma^{f|f}_{u,T;\theta}\cdot\Sigma^{f\to d}_{u,T+\theta}du\}$) depend on the model parameters only, not on the initial curves. So, on a bootstrapping segment, the collateral basis spot volatility determines the convexity adjustments of all swaplets of the segment, and the interpolation scheme only determines how the average level of the collateral basis is distributed across these swaplets. The model delivered volatility of \autoref{prop:gauss_cb_spread} depends on the initial curves only mildly, through the weights $P_j$ and the levels $L_j$, $B_j$.
\end{rema}

\begin{rema}[Which curve to interpolate]
On a segment between two currency basis pillars, one can either interpolate the cross-currency curve $Z^{f|d}_{0,\cdot}$ or the collateral basis curve $\QZ^{f|f\to d}_{0,\cdot} = Z^{f|d}_{0,\cdot}/Z^{f|f}_{0,\cdot}$, the native curve $Z^{f|f}_{0,\cdot}$ being known from the native market. Under the constant funding spread hypothesis and in the LGCM, the swaplet basis spreads are
\begin{eqnarray*}
  \theta B_{j} &=& \kappa^d (1+\theta R^{f|d}_{0,T_{j-1};\theta}) - \kappa^f K_j (1+\theta R^{f|f}_{0,T_{j-1};\theta}) 
  ~=~ (1+\theta R^{f|f}_{0,T_{j-1};\theta}) \left( \kappa^d\, \frac{\QZ^{f|f\to d}_{0,T_{j-1}}}{\QZ^{f|f\to d}_{0,T_j}} - \kappa^f K_j \right) 
\end{eqnarray*}
with $\kappa^e = 1-\theta S^{e|e,\div}$ and $K_j$ the (deterministic) Libor collateral convexity.
\begin{itemize}
\item With flat forwards for $r^{f|d}_{0,\cdot}$ on the segment, $R^{f|d}_{0,T_{j-1};\theta}$ is constant and $B_j$ jumps one for one with the native forwards $R^{f|f}_{0,T_{j-1};\theta}$ at the native pillars. The implied collateral basis forward curve $q^{f|f\to d}_{0,\cdot} = r^{f|d}_{0,\cdot} - r^{f|f}_{0,\cdot}$ then carries all the term structure of the native curve with the opposite sign.
\item With flat forwards for $q^{f|f\to d}_{0,\cdot}$ on the segment, $B_j$ only varies through the factor $1+\theta R^{f|f}_{0,T_{j-1};\theta}$ and the convexity $K_j$.
\end{itemize}
Interpolating the collateral basis curve is therefore consistent with the choice of $z^c$ as the state driving $\QZ^{f|f\to d}$ (see the remark on our parametrisation choices in the Quasi-Gaussian section). It also avoids generating a term structure of the short collateral basis $q^{f\to d}_t = q^{f|f\to d}_{0,t} + \vun^\top z^c_t + \vun^\top v^{f,c}_t$, and hence of the FX drift, that only reflects the native curve pillars.
\end{rema}

\subsubsection{Curves at future dates}

\begin{rema}[Calendar-pinned features]
By \autoref{theo:qgcm}, for $t>0$,
\begin{eqnarray*}
  r^{d|d}_{t,T} - r^{d|d}_{0,T} &=& M^{d,\top}_{t,T}\left(z^d_t + y^{d,d}_tG^d_{t,T}\right)
\end{eqnarray*}
and similarly for $r^{f|f}$ and $q^{f|f\to d}$. The right hand side is continuous in $T$, so every discontinuity or kink of the interpolated initial forward curve at a pillar date $T_k$ is still there at the same \emph{calendar} date $T_k$ on every path. On the other hand, a curve re-bootstrapped at $t$ with the same interpolation scheme has its pillars at $t+\tau_k$ for the benchmark tenors $\tau_k$. In general, the two curves differ between the pillars, even if the re-bootstrapping uses the model values of the benchmark quotes.

The practical consequence is that future curves along a simulated path (for exposure, XVA or P\&L explain) should be produced with the reconstruction formulas, not by re-stripping simulated quotes with the curve building code. The latter would price with curves that are not those of the model, which reintroduces exactly the inconsistency the model was designed to remove. For the same reason, the model's carry (roll-down) keeps the initial pillars at fixed calendar dates.
\end{rema}

Whether a parametric interpolation family can be preserved by the dynamic is the consistency problem of Bj\"ork and Christensen \cite{BjorkChristensen}. A family $\mathcal{G}$ of forward curves $\tau\mapsto g(\tau)$, parametrised by time to maturity $\tau = T-t$, is \emph{consistent} with the model if, whenever $\tau\mapsto r_{0,\tau}$ belongs to $\mathcal{G}$, $\tau\mapsto r_{t,t+\tau}$ belongs to $\mathcal{G}$ for all $t$ and all paths.

\begin{prop}[Consistent families for the LGCM]\label{prop:consistent_families}
Assume the parameters of the LGCM are constant and all mean reversions are non zero.
\begin{enumerate}
\item For a native rates block $a$ with mean reversions $\beta^{a,1},\ldots,\beta^{a,N_a}$, the family
\begin{eqnarray*}
  \mathcal{G}^a ~=~ \left\{ \tau\mapsto \zeta_0 + \sum_{i}\zeta_i e^{-\beta^{a,i}\tau} + \sum_{i\le j}\zeta_{ij}e^{-(\beta^{a,i}+\beta^{a,j})\tau} \right\}
\end{eqnarray*}
is consistent with the dynamic of $r^{a}$. If some $\beta^{a,j}$ vanish, the functions $\tau e^{-\beta^{a,i}\tau}$ have to be added.
\item If $\rho^{f,c}\neq 0$, a family of this exponential type is consistent with the dynamic of the collateral basis $q^{f|f\to d}$ only if it contains, in addition to the exponentials of the collateral basis block, the foreign rates exponentials $e^{-\beta^{f,i}\tau}$ and the mixed exponentials $e^{-(\beta^{c,i}+\beta^{f,j})\tau}$ (for generic parameters, i.e. unless their coefficients happen to cancel).
\end{enumerate}
\end{prop}

\begin{proof}
With constant parameters, $M^{a,i}_{t,t+\tau} = e^{-\beta^{a,i}\tau}$ and $G^{a,j}_{t,t+\tau} = (1-e^{-\beta^{a,j}\tau})/\beta^{a,j}$ (or $\tau$ if $\beta^{a,j}=0$). By \autoref{theo:qgcm}, $r^a_{t,t+\tau} = r^a_{0,t+\tau} + \sum_i e^{-\beta^{a,i}\tau}(z^{a,i}_t + \sum_j y^{(a,i),(a,j)}_tG^{a,j}_{t,t+\tau})$. The second term lies in the span of $e^{-\beta^{a,i}\tau}$ and $e^{-(\beta^{a,i}+\beta^{a,j})\tau}$. The first term is in $\mathcal{G}^a$ if $r^a_{0,\cdot}$ is, since $e^{-\beta(t+\tau)} = e^{-\beta t}e^{-\beta\tau}$ and constants are shift invariant. For the collateral basis, the reconstruction formula of \autoref{theo:qgcm} contains $M^{f,\top}_{t,T}(v^{f,c}_t + y^{f,c}_tG^c_{t,T})$ and $M^{c,\top}_{t,T}y^{c,f}_tG^f_{t,T}$. When $\rho^{f,c}\neq 0$, $y^{c,f}_t$ and $v^{f,c}_t$ are non zero for $t>0$, and these terms produce $e^{-\beta^{f,i}\tau}$ with coefficient $v^{f,c}_t + y^{f,c}_t(\beta^c)^{-1}\vun$ and $e^{-(\beta^{c,i}+\beta^{f,j})\tau}$ with coefficient $-y^{(c,i),(f,j)}_t\left(1/\beta^{f,j} + 1/\beta^{c,i}\right)$. The exponential decompositions were checked symbolically (sympy).
\end{proof}

\begin{rema}
\begin{enumerate}
\item A consistent family has at most $1+N_a+N_a(N_a+1)/2$ parameters, so it cannot reproduce an arbitrary set of market quotes. There is a genuine trade-off between exact fit of the initial curves (non-parametric interpolation, inconsistent families) and consistency (parametric curves). In the LGCM the exact fit is preferred, and the inconsistency is harmless as long as future curves are produced by the reconstruction formulas.
\item The proposition is in line with \cite{BjorkChristensen}, where the Nelson-Siegel family is shown to be inconsistent with the Hull-White model but can be augmented to achieve consistency. A one factor Hull-White model generates the exponential $e^{-2\beta\tau}$, which is not in the Nelson-Siegel family. See also \cite{Filipovic:NelsonSiegel}: no non-trivial HJM model is consistent with the Nelson-Siegel family.
\item The second part of the proposition says that, if one uses a parametric family for the collateral basis curve, it must include the foreign rates exponentials. The collateral basis curve inherits part of the shape of the foreign curve through the convexity term $\sigma^{f}_{t,T}\cdot\Sigma^{f\to d}_{t,T}$ in its drift.
\end{enumerate}
\end{rema}


\subsection{Alternative calibrations}
For this purpose, we developed the Non Linear Curve Engine: a generic framework to jointly 
\begin{enumerate}
\item calibrate a stochastic spread model 
\item calibrate discount and index curves taking into account the convexity adjustment induced by the model. 
\end{enumerate}

Calibration targets for this joint calibration are 
\begin{enumerate}
\item FX forwards
\item Cross-currency swaps (resetting and non-resetting)
\item FX options
\item Currency basis swaptions (since these are not liquid we cannot imply volatilities so we target historical estimates instead). 
\end{enumerate}

There are several ways to do the calibration.

\subsubsection{Global calibration}%
We create a function with inputs 
\begin{enumerate}
\item curves parameters
\item model parameters
\end{enumerate}
and outputs 
\begin{enumerate}
\item model induced value of market quotes.  
\end{enumerate}
We run an optimiser to minimise the error between model induced quotes and market quotes.



\subsubsection{Nested calibration}%
We decouple the calibration of model parameters and that of the curves. This leads to the following algorithm.

Until convergence is achieved
\begin{enumerate}
\item Initialise the stochastic model:
\begin{enumerate}
\item Start with an initial guess for $\sigma^{f\to d}_\cdot$ and correlations.
\item Bootstrap the FX spot volatility $\sigma^{df}_\cdot$ to match FX options prices.
\end{enumerate}
\item Build the discount curve $Z^{f|d}_{t_0,\cdot}$ and index curve $L^{f|d}_{t_0,T;\theta}$.
\begin{enumerate}
\item Compute all the convexity adjustments.  
\item Solve for the curve parameters to match currency basis spreads.
\end{enumerate}
\item Use the Gaussian approximation to compute the quadratic covariation of currency basis spreads 
\begin{eqnarray}
  %\Cov_{t_0}\left( B^{f/d|d,r}_{t_0,T_i;M_k}, B^{f/d|d,r}_{t_0,T_j;M_\ell} \right)
  \int_{t_0}^{T_i\wedge T_j} d \left\langle B^{f/d|d,r}_{t_0,T_i;M_k}, B^{f/d|d,r}_{t_0,T_j;M_\ell} \right\rangle
  &=& \int_t^{T_i\wedge T_j}  \ldots \, dt 
\end{eqnarray}
as well as their covariance with other factors.
\item Update $\sigma^{f\to d}_\cdot$ to match historical variances, update correlations to match historical covariances.
\end{enumerate}

Note that the FX forward will move at each iteration. For simplicity, we could assume the FX vol quote do not.

The update (gradient descent) can be implemented in several ways: finite difference, adjoint differentiation etc\ldots




\newpage%
\subsection{Multi-currency model}%


[TODO: This seems really hard to do in closed form! Brute force is having a function that calibrates the $ed$ market and the $df$ market and computes the delivered vols on the cross-FX $ef$. Then run this function, seen as a function of correlations, though an optimiser to reduce the error between model delivered quotes and market observed vols.]

\subsubsection{Odd correlations}


So called odd-correlations such as EUR vs USDJPY correlations cannot be implied from market quotes because there are no liquid market instruments directly exposed to them. The lack of any natural economical interpretation increases the uncertainty associated with historical estimation. This makes this type of correlation unfit as model inputs. On the other hand, options on the FX cross EURUSD are liquid. One can then use a model to imply the odd-correlation EUR vs USDJPY from the natural ones:  EUR vs EURUSD and EUR vs EURJPY.

In the BS-LGM model, this is very easy to do because the ATM vol of the cross is equivalent to a constraint on the FX-FX correlation. 

In the LGCM model, we have many more unknown correlations. Can we still apply the same reasoning?

It seems like we need to bootstrap the FX-FX correlations and FX-CB correlations at the same time. 

\subsubsection{Triangle relationship under stochastic collateral}



%\subsection{Variance of the cross FX-forward}


Let us write $\Sigma^{df}_{t,T}$ for the volatility of the FX-forward process:
\begin{eqnarray}
  \frac{dx^{df|c}_{t,T}}{x^{df|c}_{t,T}} &=& \Sigma^{df}_{t,T} \cdot dW^{d|c,T}_{t} 
\end{eqnarray}
We saw that 
\begin{eqnarray}
  \Sigma^{df}_{t,T} &=& \Sigma^{d|d}_{t,T} - \Sigma^{f|f}_{t,T} - \Sigma^{f\to d}_{t,T}  + \sigma^{df}_t  
\end{eqnarray}
does not depend on the collateral. 

%\begin{coro}
It follows that the total variance of an FX option with maturity $T$ is:
\begin{eqnarray}
  \sigma^{df|c}_{\BS,t}(T)^2  (T-t)
  &=& \int_t^T \left| \Sigma^{d|d}_{u,T} - \Sigma^{f|f}_{u,T} +\Sigma^{d\to f}_{u,T}  + \sigma^{df}_u  \right|^2 du
\end{eqnarray}
regardless of the collateral\footnote{But its moneyness does depend on the collateral.}.
%\end{coro}



Then, from $x^{ef}_{t,T} = x^{df}_{t,T}/x^{de}_{t,T}$, we find 
\begin{eqnarray}
  \Sigma^{ef}_{t,T} &=& \Sigma^{df}_{t,T}  - \Sigma^{de}_{t,T}  .
\end{eqnarray}
It follows that implied vols of the cross-FX are given by:
\begin{eqnarray}
  \sigma^{ef}_{\BS,t}(T)^2  (T-t)
  &=& \int_t^T \left| \Sigma^{df}_{u,T}  \right|^2 \, du
+ \int_t^T \left| \Sigma^{de}_{u,T}  \right|^2 \, du
- 2 \int_t^T \Sigma^{df}_{u,T} \cdot \Sigma^{de}_{u,T} \,du 
\end{eqnarray}
This can be viewed as a constraint on the correlation between the two FX rates. 

When all factors are one-dimensional, assuming all other correlations are known, we can solve for the FX-FX spot correlation to match the delivered ATMF vols on the cross.


In the deterministic spread setting, one only has to mark IR-IR as well as FX-IRD and FX-IRF correlations. In this new context, one has to mark all collateral basis correlations as well. 

This requires to define SDP matrices. 
Note that synchronicity of time series is crucial for correlation estimations.

%\begin{prop}
%The dynamic of the FX-Forward is:
%\begin{eqnarray}
%  \frac{dx^{df|c}_{t,T}}{x^{df|c}_{t,T}} 
%  %&=& \sigma^{df|c}_{t,T}  \cdot dW^{d|c}_{t,T} \\
%  &=& \Sigma^{ef}_{u,T}  \cdot dW^{d|c}_{u,T}  - \Sigma^{ed}_{u,T}  \cdot dW^{d|c}_{u,T} 
%%  &=& \left( \Sigma^{d|d}_{u,T} - \Sigma^{e|e}_{u,T} + \Sigma^{d\to e}_{u,T}  + \sigma^{de}_t   \right) \cdot dW^{d|c}_{u,T}   
%%  - \left( \Sigma^{f|f}_{u,T} - \Sigma^{e|f}_{u,T} +\Sigma^{f\to e}_{u,T}  + \sigma^{fe}_t  \right) \cdot dW^{d|c}_{u,T} 
%\end{eqnarray}
%\end{prop}
%
%\begin{coro}
%It follows that instantaneous correlation between the two FX spot rates is a solution to the following equation:
%\begin{eqnarray}
%  \sigma^{ef|c}_{\BS,t}(T)^2  (T-t)
%  &=& \int_t^T \left| \Sigma^{de|c}_{t,T}  \right|^2 \, du
%+ \int_t^T \left| \Sigma^{fe|c}_{t,T}  \right|^2 \, du
%- 2 \int_t^T \Sigma^{de|c}_{t,T} \cdot \Sigma^{fe|c}_{t,T} \,du 
%\end{eqnarray}
%\end{coro}

\subsubsection{Implied correlations}
We can generalise the reasoning to imply so-called odd-correlations by considering the covariance between the cross-FX and other assets.



\subsubsection{Application to off-market discount curves}%

What does the $Z^{\mathrm{EUR}|\mathrm{JPY}}_{t,T}$ curve look like? 

There is no market for cross-currency swaps collateralized in JPY. So we have to build it from the model calibrated to the EUR/JPY/USD market.




%\newpage%
%\section{The Linear Gaussian Collateral Model: Hedging}%




\newpage%
\section{Remarks}





\subsection{Model usage}

\subsubsection{Break even levels of covariances}%
The model specifies break even levels of covariances. 

This allows the trader to enter trades if the price offered.

Assess whether he is long or short correlations. etc...


\subsubsection{Benchmarking usage}
Because the model is consistent with hedging instruments, it can be used to assess the impact of stochastic currency basis on exotic products. 




\subsection{Stochastic collateral and PnL leaks}%

Q: How does the stochasticity of collateral impact the PnL of a trader hedging with a deterministic collateral model? Does he experience PnL leaks?

We have two impacts: 
\begin{enumerate}
\item Intrinsic value: collateral convexity impacts the forwards
\item Time value: volatility of the collateral changes the volatility of market observables.
\end{enumerate}



\subsection{Materiality}%
Which trading desk would see the largest impacts? 

A guess would be:
\begin{enumerate}
\item DCU/Linear rates
\item Counterparty credit risk
\item Long dated FX
\end{enumerate}



\subsection{Model assessment}


\subsubsection{Pros and cons}
The model has the following:
\begin{enumerate}
\item Pros:
\begin{enumerate}
\item Calibrates directly and consistently to liquid linear market instruments
\item The valuation of market instruments with native collateral is independent of other markets.
\item The valuation of liquid linear market instruments is independent of the diffusion parameters.
\item The model provides relevant intuition for the practitioner.
\item While designed primarily to match linear instruments, the model provides a fully arbitrage-free dynamic of that can be used to price exotics products. 
\end{enumerate}
\item Cons:
\begin{enumerate}
\item Funding spreads are deterministic.
\item Lack of smile.
\item Calibration to historical data is somewhat uncertain.
\end{enumerate}
\end{enumerate}

\subsubsection{Generalisations}%

\paragraph{Quasi-Gaussian model}%
A natural extension of the LGCM model is to allow stochastic volatility. 

Andreasen suggests the following parametrisation. Setting $g_i(\tau) = e^{-\kappa_i \tau}$ and 
\begin{eqnarray*}
  \eta_{t} &=& \begin{bmatrix} g_1(\theta_1) & \ldots & g_n(\theta_1) \\  \\ g_1(\theta_n) & \ldots & g_n(\theta_n) \end{bmatrix}^{-1} \begin{bmatrix} \sigma_1(t,\theta_1) & \ldots & \sigma_n(t,\theta_1) \\  \\ \sigma_1(t,\theta_n) & \ldots & \sigma_n(t,\theta_n) \end{bmatrix}
\end{eqnarray*}
we have,
\begin{eqnarray*}
  dr_{t,t+\theta} &=& \sum_i \sum_j e^{-\kappa_i \tau} \eta^{ij}_t dW^{T;j}_t 
\end{eqnarray*}

\begin{eqnarray*}
  \sigma_j(t,\tau) &=& \sqrt{z_t} \lambda(t,\tau)\left( b(t,\tau) r_t + (1-b(t,\tau)) r_{0,t} - \zeta \right) \rho_j(\tau) \\
  dz_t &=& \theta(1-z_t)dt + \varepsilon(t) \sqrt{z_t} dZ_t 
\end{eqnarray*}
Then, 
\begin{itemize}
\item $\lambda(t,\tau_k)$ controls ATM vol for $r_{t,t+\tau_k}$
\item $b(t,\tau_k)$ controls the skew
\item $\varepsilon(t)$ controls the convexity of the smile
\item $\rho(i)\rho^j(t)$ controls the correlation of forward rates 
\end{itemize}

The dynamic of the swap rate 
\begin{eqnarray*}
  dY_t &=& \sum_i \partial_{x_i} Y \sum_j \eta_{ij}dW^{A,j}_t 
\end{eqnarray*}
is approximated by 
\begin{eqnarray*}
  dY_t &\approx& \sqrt{z_t} \lambda(t)\left( b(t) Y_t + (1-b(t)) Y_{0} - \zeta \right) dW^{A}_t \\
  dz_t &=& \theta(1-z_t)dt + \varepsilon(t) \sqrt{z_t} dZ_t 
\end{eqnarray*}
and one uses the parameter averaging method to get constant parameters followed by a Fourier transform. 



\paragraph{LMM model with stochastic currency basis}
Another alternative is to consider Libor Market Models with stochastic collateral. 

The constant funding spread assumption and many of our computations can carry unchanged. 

One interesting thing is that, as we saw in the calibration to currency basis spreads, we were lead to chose an interpolation scheme for the initial discount curve $Z^{f|d}_{t,T}$.

Similar problems arise in LMM models where one has to specify stub interpolation to create discount factors from the forward rates on a fixed grid. To avoid arbitrage, the interpolation of the dynamic should be consistent with the interpolation scheme for the curves. In the HJM-type models of these notes this issue does not arise, see \autoref{ssec:interpolation_consistency}.




\subsection{Conclusion}
In these notes, 
\begin{enumerate}
\item We characterised absence of arbitrage for instantaneous forward rates models under perfect cash collateral. 
\item We proved that under the Cheyette separability condition, these models are Markovian. 
\item We derived explicit formulas for convexity adjustments. 
\item Finally, we introduced the Linear Gaussian Collateral Model and described efficient algorithms for its calibration and simulation. 
\end{enumerate}







\newpage%


\appendix 
\part*{Appendix}

\section{Risk neutral pricing under perfect collateral}%

\subsection{Risk neutral pricing under perfect collateral}%

We assume perfect collateralization\footnote{Hypothesis can be formulated explicitly.}.

\begin{theo}
Under the previous assumptions, for each currency $d$, there exist a unique measure $\QQ^{d,\RN}$ equivalent to the historical measure $\PP$ such that every domestic self-financed portfolio with domestic collateral rate $r^c$ satisfies the martingale property
\begin{eqnarray}
	\EE^{d,\RN}_t\left[ e^{-\int_t^T r^c_u du} V^{d|c}_T \right] &=& V^{d|c}_t.
\end{eqnarray} 
\end{theo}

\begin{proof}
We follow Piterbarg's reasoning. 

We consider a market with two collateralized assets $S^1$, $S^2$ with collateral rates $r^1,r^2$. An investment strategy is a portfolio: 
\begin{eqnarray*}
	\Pi_t &=& \Delta^1_t S^1_t + \Delta^2_t S^2_t    
\end{eqnarray*} 
At $t$, $A$ enters the trade and at $t+dt$, he closes the position:
\begin{enumerate}
\item at $t$, 
\begin{enumerate}
\item $A$ buys the assets from $B$ paying him the amount $\Pi_t = \Delta^1_t S^1_t + \Delta^2_t S^2_t $\item $B$ immediately posts the amount $\Pi_t$ back to the collateral account (held by $A$ or third party $C$).
\end{enumerate}
\item at $t+dt$, 
\begin{enumerate}
\item $A$ pays to $B$ the interest rate on the collateral posted:
\begin{itemize}
\item $B$ receives: $\Delta^1_t r^1_t S^1_{t} dt + \Delta^2_t r^2_t S^2_{t}dt$. 
\end{itemize}
\item $A$ chooses to close the position:
\begin{itemize}
\item The collateral is returned to $B$ who receives: $\Delta^1_t S^1_t + \Delta^2_t S^2_t$. 
\item $A$ sells the assets which are now valued at $\Delta^1_t S^1_{t+dt} + \Delta^2_t S^2_{t+dt}$. 
\end{itemize}
\end{enumerate}
\end{enumerate}
So $A$'s P\&L is 
\begin{eqnarray*}
	\PnL_{t+dt}
	&=&  \left(\Delta^1_t S^1_{t+dt} + \Delta^2_t S^2_{t+dt}\right) 
	-  \left(\Delta^1_t r^1_t S^1_{t} dt\right)  - \left(\Delta^2_t r^2_t S^2_{t}dt \right)  
	- \left(\Delta^1_t S^1_{t} + \Delta^2_t S^2_{t}\right)    \\
	&=&  \Delta^1_t \left(dS^1_{t} - r^1_tS^1_{t}dt\right)
	+ \Delta^2_t \left(dS^2_{t} - r^2_tS^2_{t}dt\right)    .
\end{eqnarray*} 

Suppose our two assets have dynamic:
\begin{eqnarray*}
	dS^i_t &=& \mu^i_t dt + \sigma^i_t  dW^\PP_t   
\end{eqnarray*} 
driven by the same Brownian motion. Our P\&L is
\begin{eqnarray*}
	\PnL_{t} 
	&=& \left( \Delta^1_t (\mu^1_t-r^1_t) + \Delta^2_t (\mu^1_t-r^2_t) \right)dt 
	+  \left( \Delta^1_t \sigma^1_t + \Delta^2_t \sigma^2_t \right) \cdot dW^\PP_t
\end{eqnarray*} 

We can now choose non-zero weights that cancel the Brownian part: 
\begin{eqnarray*}
	\Delta^1_t \sigma^1_t + \Delta^2_t \sigma^2_t &=& 0
\end{eqnarray*} 
Namely\footnote{This can be generalised to the case of a multi-dimensional Brownian motion and an arbitrary number of assets $S^i$. One simply needs to make the following completeness assumption: the full volatility matrix $\Sigma$ with rows $\sigma^i$ is not invertible.}
\begin{eqnarray*}
	\Delta^1_t &=&  \sigma^2_t \\
	\Delta^2_t &=& -\sigma^1_t 
\end{eqnarray*} 
The P\&L reduces to
\begin{eqnarray*}
	\PnL &=& \left( \Delta^1_t (\mu^1_t-r^1_t) + \Delta^2_t (\mu^2_t-r^2_t) \right)dt .
\end{eqnarray*} 
It is known at $t$ so it has to be\footnote{Why? We should include our funding rate here!} $0$, otherwise we would be able to make a risk-free profit without initial investment by buying or selling the strategy. This implies 
\begin{eqnarray*}
	%\sigma^2_t (\mu^1_t-r^1_t) - \sigma^1_t (\mu^2_t-r^2_t) &=& 0 \\
	\frac{r^1_t - \mu^1_t}{\sigma^1_t} &=& \frac{r^2_t - \mu^2_t}{\sigma^2_t}
\end{eqnarray*} 
Let us write $\lambda_t$ for this quantity. We can define $\QQ^{d,\RN}$ by setting:
\begin{eqnarray*}
	\left. \frac{d\QQ^{d,\RN}}{d\PP} \right|_{\mathcal{F}_t} 
	&=& \mathcal{E}\left(\int_0^\cdot \lambda_u\,dW^\PP_u\right) 
\end{eqnarray*} 
Then, under this measure, 
\begin{eqnarray*}
	dW^{d,\RN}_u &=& dW^\PP_u - \lambda_t dt 
\end{eqnarray*} 
is a Brownian motion and 
\begin{eqnarray*}
	dS^i_t &=& r^i_t dt + \sigma^i_t \cdot dW^{d,\RN}_t   
\end{eqnarray*} 
so the discounted assets are $\QQ^{d,\RN}$-martingales.
\end{proof}


\begin{prop}
For any collateral rate $c$ in currency $f$,
\begin{eqnarray*}
	r^{d|c} &=& \mu^{df} + r^{f|c}.
\end{eqnarray*} 
\end{prop}

\begin{proof}
Enter a collateralized trade at $t$. Close the position at $t+dt$. The value at $t$ in currency $d$ of this strategy is 
\begin{eqnarray*}
	& & \EE_t^{d,\RN}\left[ e^{-r^{d|c}_tdt} \left(V^{d|c}_{t+dt} - x^{df}_{t+dt}(1+r^{f|c}_t)x^{fd}_{t}V^{d|c}_t dt\right) \right] \\
	&=&  \EE_t^{d,\RN}\left[ e^{-r^{d|c}_tdt} \left(V^{d|c}_{t+dt} - (1+r^{f|c}_t)\left(\mu^{df}_t dt + \sigma^{df}_t dW^{d,x}_t \right)V^{d|c}_t dt\right) \right] \\ 
	&=& e^{-r^{d|c}_tdt} \left(\EE_t^{d,\RN}[dV^{d|c}_{t}] - (\mu^{df}_t + r^{f|c}_t)V^{d|c}_tdt\right) \\
	&=& e^{-r^{d|c}_tdt} \left(r^{d|c}_t - \mu^{df}_t - r^{f|c}_t\right)V^{d|c}_t dt
\end{eqnarray*} 
Since we can enter a collateralized trade at no cost, no arbitrage implies this is 0 so $r^{d|c}_t = \mu^{df}_t + r^{f|c}_t$.
\end{proof}

\begin{rema}
Repo rates are equivalent to collateral.
\end{rema}


\begin{prop}
\begin{eqnarray}
	\EE^{d,\RN}_t\left[\frac{d\QQ^{f,\RN}}{d\QQ^{d,\RN}}\right] 
	~=~ e^{-\int_0^t \mu^{df} du} \frac{x^{df}_t}{x^{df}_0} 
	~=~ \mathcal{E}\left(\int_0^\cdot \sigma^x_v dW^x_v \right)_t 
\end{eqnarray} 
\end{prop}

\begin{proof}
No arbitrage between foreign and domestic markets (carry trade). 
\end{proof}

So the change from domestic to foreign measure is the usual one and the quanto adjustments are governed by the correlation with the FX driver.


\begin{prop}[Triangle relation for FX drift]
\begin{eqnarray*}
	\mu^{df} + \mu^{fd} &=& 0 \\
	\mu^{df} + \mu^{fg} + \mu^{gd}  &=& 0 .
\end{eqnarray*} 
\end{prop}

\begin{proof}
Chain the change of measures.
\end{proof}


We can now extend our previous results to arbitrary collateral currency:
\begin{prop}
For any collateral rate $c$ in any currency,
\begin{eqnarray*}
	r^{d|c} &=& \mu^{df} + r^{f|c}.
\end{eqnarray*} 
\end{prop}


\begin{proof}
Let $x$ be the currency of the collateral rate $c$, then
\begin{eqnarray*}
	r^{d|c}  = \mu^{dx} + r^{x|c} = \mu^{dx} + (r^{f|c} - \mu^{fx}) = \mu^{df} + r^{f|c}.
\end{eqnarray*} 
\end{proof}

\begin{prop}
For any currencies $d,f$ and collaterals $c_1,c_2$,
\begin{eqnarray*}
	r^{d|c_1} - r^{d|c_2} &=& r^{f|c_1} - r^{f|c_2}.
\end{eqnarray*} 
\end{prop}

\begin{proof}
\begin{eqnarray*}
	r^{d|c_1} - r^{d|c_2} = (\mu^{df} + r^{f|c_1} ) - (\mu^{df} + r^{f|c_2} ) =  r^{f|c_1} - r^{f|c_2}.
\end{eqnarray*} 
\end{proof}


\subsection{HJM framework: proofs}%

\subsubsection{Dynamic of funding spreads}

\begin{lemm}
Under our funding spread assumption, the multiplicative funding spread
\begin{eqnarray}
	1 - \theta S^{d|c,\div}_{t,T;\theta}
	&=&  \frac{1+\theta L^{d|c,\div}_{t,T;\theta}}{1+\theta R^{d|c,\div}_{t,T;\theta}} 
\end{eqnarray} 
is an exponential martingale under $\QQ^{d|c,T}$:
\begin{eqnarray}
  \frac{d(1 - \theta S^{d|c,\div}_{t,T;\theta})}{1 - \theta S^{d|c,\div}_{t,T;\theta}} 	
  &=& -\Sigma^{d\to c}_{t,T;\theta} \cdot dW^{d|c,T}_t 
\end{eqnarray} 
\end{lemm}


\begin{proof}
From Ito's lemma, we have:
\begin{eqnarray*}
	\frac{d(X/Y)}{X/Y} 
	&=&  \frac{dX}{X} -  \frac{dY}{Y} 
	- \left\langle  \frac{dY}{Y}    ,  \frac{dX}{X} -  \frac{dY}{Y}  \right\rangle 
\end{eqnarray*} 
So, 
\begin{eqnarray*}
	& & \frac{d(1 - \theta S^{d|c,\div}_{t,T;\theta})}{1 - \theta S^{d|c,\div}_{t,T;\theta}} \\
	&=&  \frac{d(1+\theta L^{d|c}_{t,T;\theta})}{1+\theta L^{d|c}_{t,T;\theta}} 
	-  \frac{d(1+\theta R^{d|c}_{t,T;\theta})}{1+\theta R^{d|c}_{t,T;\theta}} 
	- \left\langle   \frac{d(1+\theta R^{d|c}_{t,T;\theta})}{1+\theta R^{d|c}_{t,T;\theta}}   ,  \frac{d(1+\theta L^{d|c}_{t,T;\theta})}{1+\theta L^{d|c}_{t,T;\theta}}  -  \frac{d(1+\theta R^{d|c,\div}_{t,T;\theta})}{1+\theta R^{d|c}_{t,T;\theta}} \right\rangle \\
	&=& \Sigma^{d|d}_{t,T;\theta}dW^{d|c,T+\theta}_t - \Sigma^{d|c}_{t,T;\theta}dW^{d|c,T+\theta}_t - \Sigma^{d|c}_{t,T;\theta}\cdot (\Sigma^{d|d}_{t,T;\theta} - \Sigma^{d|c}_{t,T;\theta})dt \\
	&=& -\Sigma^{d\to c}_{t,T;\theta}dW^{d|c,T+\theta}_t + \Sigma^{d|c}_{t,T;\theta}\cdot \Sigma^{d\to c}_{t,T;\theta} \, dt \\
	&=& -\Sigma^{d\to c}_{t,T;\theta}\left( dW^{d|c,T+\theta}_t - \Sigma^{d|c}_{t,T;\theta} dt\right) \\
	%&=& -\Sigma^{d\to c}_{t,T;\theta}\left( dW^{d|c,T+\theta}_t - \Sigma^{d|c}_{t,T+\theta} + \Sigma^{d|c}_{t,T;\theta} dt\right) \\
	%&=& -\Sigma^{d\to c}_{t,T;\theta}\left( dW^{f,\RN}_t + \Sigma^{d|c}_{t,T;\theta} dt\right) \\
	&=& -\Sigma^{d\to c}_{t,T;\theta} \cdot dW^{d|c,T}_t 
\end{eqnarray*} 
\end{proof}


\subsubsection{Resetting forward Libors}



We first compute the adjustment due to the resetting notional.
\begin{lemm}
Under the constant funding spread assumption,  
\begin{eqnarray}
	 \EE^{d|d,T+\theta}_t\left[ x^{df}_T(1+\theta L^{d}_{T;\theta} )\right] 
	&=& x^{df|d}_{t,T}   \left( 1+\theta L^{d|d}_{t,T;\theta} \right)
\end{eqnarray} 
\end{lemm}


\begin{proof}\label{proof:resetting_notional_libor_lemma}
Using the fact that the funding spreads are deterministic:
\begin{eqnarray*}
	 \EE^{d|d,T}_t\left[ x^{df}_S(1+\theta L^{d|d}_{T;\theta} )\right] 
	&=&  \EE^{d|d,T}_t\left[ x^{df}_S(1+\theta R^{d|d}_{S,T} ) \left(1-\theta S^{d|d,\div}_{T;\theta} \right) \right]  \\ 
	&=& \EE^{d|d,T}_t\left[ x^{df}_S(1+\theta R^{d|d}_{T;\theta} )\right] \left(1-\theta S^{d|d,\div}_{t,T;\theta} \right) \\
	&=& \EE^{d|d,T}_t\left[ \frac{Z^{d|d}_{S,S}}{Z^{d|d}_{S,T}}x^{df}_S \right] \left(1-\theta S^{d|d,\div}_{t,T;\theta} \right) \\
	&=& \frac{Z^{d|d}_{t,S}}{Z^{d|d}_{t,T}}  \EE^{d|d,S}_t\left[ x^{df}_S \right] \left(1-\theta S^{d|d,\div}_{t,T;\theta} \right) \\
	&=&   \left(1+\theta R^{d|d}_{t,T;\theta} \right) x^{df|d}_{t,S}  \left(1-\theta S^{d|d,\div}_{t,T;\theta} \right) \\
	&=&  x^{df|d}_{t,S}  \left( 1+\theta L^{d|d}_{t,T;\theta} \right)
\end{eqnarray*} 
\end{proof}


\begin{prop}
\begin{eqnarray}
  1+\theta L^{f/d|d}_{t,T;\theta} 
  &=& \left( 1+\theta R^{f|d}_{t,T;\theta} \right) \left( 1 - \theta S^{d|d,\div}_{t,T;\theta} \right)
\end{eqnarray} 
\end{prop}


\subsubsection{Resetting cross-currency swaplets}

We now proceed to the proof of \autoref{theo:currency_basis_spread_constant_funding_basis}.
\begin{theo*}
Under deterministic multiplicative funding spread assumptions, the currency basis spread of a resetting cross-currency swaplet is:
\begin{eqnarray*}
	B^{f/d|d,r}_{t,T;\theta} 
	&=& \left(1+\theta R^{f|d}_{t,T;\theta} \right)  \left(S^{f|d,\div}_{t,T;\theta} - S^{d|d,\div}_{t,T;\theta} \right)  .
\end{eqnarray*} 
\end{theo*}


\begin{proof}
By definition, the currency basis satisfies
\begin{eqnarray*}
  B^{f/d|d,r}_{t,T;\theta} 
  &=& L^{f/d|d}_{t,T;\theta}  - L^{f|d}_{t,T;\theta}  \\
  &=& (1+\theta L^{f/d|d}_{t,T;\theta}) - (1+\theta L^{f|d}_{t,T;\theta})  \\
  &=& \left( 1+\theta R^{f|d}_{t,T;\theta} \right) \left( 1 - \theta S^{d|d,\div}_{t,T;\theta} \right) - \left( 1+\theta R^{f|d}_{t,T;\theta} \right) \left( 1 - \theta S^{f|d,\div}_{t,T;\theta} \right) \\
  &=& \left(1+\theta R^{f|d}_{t,T;\theta} \right)  \left(S^{f|d,\div}_{t,T;\theta} - S^{d|d,\div}_{t,T;\theta} \right)  
\end{eqnarray*} 
\end{proof}

%\begin{proof}[Proof of the theorem]
%By definition, the currency basis satisfies
%\begin{eqnarray*}
%	Z^{f|d}_{t,T+\theta}\theta B^{f/d|d,r}_{t,T;\theta} 
%	&=& 
%	x^{fd}_tZ^{d|d}_{t,T+\theta}\EE^{d|d,T+\theta}_t\left[ x^{df}_S(1+\theta L^d_{T;\theta})\right]  
%	- Z^{f|d}_{t,T}\EE^{f|d,T+\theta}_t\left[1+\theta L^f_{T;\theta} \right] .
%\end{eqnarray*} 
%Using the lemma, 
%\begin{eqnarray*}
%	\theta B^{f/d|d,r}_{t,T;\theta} 
%	&=& \frac{x^{fd|d}_{t,T}}{x^{fd|d}_{t,S}} \left(1+\theta L^{d|d}_{t,T;\theta}\right) - \left(1+\theta L^{f|d}_{t,T;\theta} \right)  \\
%	&=& \frac{Z^{d|d}_{t,T;\theta}}{Z^{f|d}_{t,T;\theta}} \left(1+\theta L^{d|d}_{t,T;\theta}\right) - \left(1+\theta L^{f|d}_{t,T;\theta} \right)  \\
%	%&=& Z^{d|d}_{t,T} \frac{Z^{f|d}_{t,S}}{Z^{d|d}_{t,S}}(1+\theta L^{d|d}_{t,T;\theta}) - Z^{f|d}_{t,T}\left(1+\theta L^{f|d}_{t,T;\theta} \right)  \\
%	&=& \left(1+\theta R^{f|d}_{t,T;\theta}\right) \frac{1+\theta L^{d|d}_{t,T;\theta}}{1+\theta R^{d|d}_{t,T;\theta}} - \left(1+\theta R^{f|d}_{t,T;\theta} \right)  \frac{1+\theta L^{f|d}_{t,T;\theta}}{1+\theta R^{f|d}_{t,T;\theta}} \\
%	&=& \left(1+\theta R^{f|d}_{t,T;\theta} \right)  \left(1-\theta S^{d|d,\div}_{t,T;\theta}\right) - \left(1+\theta R^{f|d}_{t,T;\theta} \right) \left(1-\theta S^{f|d,\div}_{t,T;\theta}\right)  
%\end{eqnarray*} 
%giving
%\begin{eqnarray*}
%	B^{f/d|d,r}_{t,T;\theta} 
%	&=& \left(1+\theta R^{f|d}_{t,T;\theta} \right)  \left(S^{f|d,\div}_{t,T;\theta} - S^{d|d,\div}_{t,T;\theta} \right)  
%\end{eqnarray*} 
%\end{proof}




\subsection{Deterministic collateral basis models: proofs}%


\subsubsection{Triangle relation for currency basis}%


\subsection{Deterministic volatility models: proofs}%

\subsubsection{Dynamic of market observables}

\paragraph{Funding basis spreads}

\begin{prop*}[Funding basis spreads]
Under our funding spread assumption, the multiplicative funding spread
\begin{eqnarray}
	1 - \theta S^{d|c,\div}_{t,T;\theta}
	&=&  \frac{1+\theta L^{d|c,\div}_{t,T;\theta}}{1+\theta R^{d|c,\div}_{t,T;\theta}} 
\end{eqnarray} 
is a martingale\footnote{Note that, if $f=d$, the vol is zero and we recover that the spread is constant.}
under $\QQ^{d|c,T}$:
\begin{eqnarray}
  \frac{d(1 - \theta S^{d|c,\div}_{t,T;\theta})}{1 - \theta S^{d|c,\div}_{t,T;\theta}} 	
  &=& -\Sigma^{d\to c}_{t,T;\theta} \cdot dW^{d|c,T}_t 
\end{eqnarray} 
\end{prop*}


\begin{proof}
From Ito's lemma, we have:
\begin{eqnarray*}
	\frac{d(X/Y)}{X/Y} 
	&=&  \frac{dX}{X} -  \frac{dY}{Y} 
	- \left\langle  \frac{dY}{Y}    ,  \frac{dX}{X} -  \frac{dY}{Y}  \right\rangle 
\end{eqnarray*} 
So, 
\begin{eqnarray*}
	& & \frac{d(1 - \theta S^{d|c,\div}_{t,T;\theta})}{1 - \theta S^{d|c,\div}_{t,T;\theta}} \\
	&=&  \frac{d(1+\theta L^{d|c}_{t,T;\theta})}{1+\theta L^{d|c}_{t,T;\theta}} 
	-  \frac{d(1+\theta R^{d|c}_{t,T;\theta})}{1+\theta R^{d|c}_{t,T;\theta}} 
	- \left\langle   \frac{d(1+\theta R^{d|c}_{t,T;\theta})}{1+\theta R^{d|c}_{t,T;\theta}}   ,  \frac{d(1+\theta L^{d|c}_{t,T;\theta})}{1+\theta L^{d|c}_{t,T;\theta}}  -  \frac{d(1+\theta R^{d|c,\div}_{t,T;\theta})}{1+\theta R^{d|c}_{t,T;\theta}} \right\rangle \\
	&=& \Sigma^{d|d}_{t,T;\theta} \cdot dW^{d|c,T+\theta}_t -
	 \Sigma^{d|c}_{t,T;\theta} \cdot dW^{d|c,T+\theta}_t - \Sigma^{d|c}_{t,T;\theta}\cdot (\Sigma^{d|d}_{t,T;\theta} - \Sigma^{d|c}_{t,T;\theta})dt \\
	&=& -\Sigma^{d\to c}_{t,T;\theta} \cdot dW^{d|c,T+\theta}_t + \Sigma^{d|c}_{t,T;\theta}\cdot \Sigma^{d\to c}_{t,T;\theta} dt \\
	&=& -\Sigma^{d\to c}_{t,T;\theta} \cdot \left( dW^{d|c,T+\theta}_t - \Sigma^{d|c}_{t,T;\theta} dt\right) \\
	&=& -\Sigma^{d\to c}_{t,T;\theta} \cdot \left( dW^{d|c,T+\theta}_t - \Sigma^{d|c}_{t,T+\theta}dt + \Sigma^{d|c}_{t,T} dt\right) \\
	&=& -\Sigma^{d\to c}_{t,T;\theta} \cdot \left( dW^{d,\RN}_t + \Sigma^{d|c}_{t,T} dt\right) \\
	&=& -\Sigma^{d\to c}_{t,T;\theta} \cdot dW^{d|c,T}_t 
\end{eqnarray*} 
\end{proof}




\paragraph{Resetting forward Libors}

\begin{prop*}[Resetting forward Libors]
The resetting forward Libors 
\begin{eqnarray}
   1 + \theta L^{f/d|i}_{t,T;\theta} 
   &=& \EE^{f|i,T+\theta}_t\left[ \frac{x^{fd}_{T+\theta}}{x^{fd}_{T}} (1+\theta L^{d}_{T;\theta}) \right] 
\end{eqnarray} 
is an exponential martingale under $\QQ^{f|i,T+\theta}$-martingale:
\begin{eqnarray}
   \frac{d(1 + \theta L^{f/d|i}_{t,T;\theta})}{1 + \theta L^{f/d|i}_{t,T;\theta}}
   &=& \Sigma^{f|d}_{t,T;\theta} \cdot dW^{f|i,T+\theta}_t 
\end{eqnarray} 
So that
\begin{eqnarray}
  1 + \theta L^{f/d|i}_{t,T;\theta} 
   &=& (1+\theta R^{f|i}_{t,T;\theta}) (1-\theta S^{d|i,\div}_{t,T;\theta}) e^{ - \int_t^T \Sigma^{d\to i}_{u,T;\theta} \cdot \Sigma^{df|i}_{u,T} \,du }.
\end{eqnarray} 
\end{prop*}


\begin{proof}
We can rewrite this in terms of in terms of domestic $T$-martingales: 
\begin{eqnarray*}
   1 + \theta L^{f/d|i}_{t,T;\theta} 
   &=& \EE^{f|i,T+\theta}_t\left[ \frac{x^{fd}_{T+\theta}}{x^{fd}_{T}} (1+\theta L^{d}_{T;\theta}) \right] \\
   &=& x^{fd|i}_{t,T+\theta} \EE^{d|i,T+\theta}_t\left[ x^{df}_{T} (1+\theta L^{d}_{T;\theta}) \right]  \\
   &=& x^{fd|i}_{t,T+\theta} \EE^{d|i,T}_t\left[\frac{Z^{d|i}_{T,T+\theta}}{Z^{d|i}_{t,T,T+\theta}} x^{df}_{T} (1+\theta L^{d}_{T;\theta}) \right]  \\
   &=&  \frac{x^{fd|i}_{t,T+\theta}}{Z^{d|i}_{t,T,T+\theta}} \EE^{d|i,T}_t\left[ x^{df}_{T} (1-\theta S^{d|i,\div}_{T;\theta}) \right]  \\
   &=&   \frac{x^{fd|i}_{t,T+\theta}}{Z^{d|i}_{t,T,T+\theta}}  \EE^{d|i,T}_t\left[ x^{df}_{T} (1-\theta S^{d|i,\div}_{T;\theta}) \right]  \\
   &=&   \frac{x^{fd|i}_{t,T+\theta}}{Z^{d|i}_{t,T,T+\theta}}   (1-\theta S^{d|i,\div}_{t,T;\theta}) \frac{ \EE^{d|i,T}_t\left[ x^{df}_{T} (1-\theta S^{d|i,\div}_{T;\theta}) \right]  }  { x^{df|i}_{t,T} (1-\theta S^{d|i,\div}_{t,T;\theta})  }  \\
   &=&  (1+\theta R^{f|i}_{t,T;\theta})  \EE^{f|i,T}_t\left[ 1-\theta S^{d|i,\div}_{T;\theta} \right]  
\end{eqnarray*} 
Now using what we know about the funding spread:
\begin{eqnarray*}
   1 + \theta L^{f/d|i}_{t,T;\theta} 
   &=&  (1+\theta R^{f|i}_{t,T;\theta})  \EE^{f|i,T}_t\left[ 1-\theta S^{d|i,\div}_{T;\theta} \right]   \\
   &=& (1+\theta R^{f|i}_{t,T;\theta})  \EE^{f|i,T}_t\left[ (1-\theta S^{d|i,\div}_{t,T;\theta}) \mathcal{E}\left(- \int_t^T \Sigma^{d\to i}_{u,T;\theta} \cdot dW^{f|i,T}_u \right) \right]   \\
  &=& (1+\theta R^{f|i}_{t,T;\theta}) (1-\theta S^{d|i,\div}_{t,T;\theta}) e^{ - \int_t^T \Sigma^{d\to i}_{u,T;\theta} \cdot \Sigma^{df|i}_{u,T} \,du }  
\end{eqnarray*} 

We can now compute the volatility:
\begin{eqnarray*}
    \frac{d(1 + \theta L^{f/d|i}_{t,T;\theta})}{1 + \theta L^{f/d|i}_{t,T;\theta} } 
   &=& \frac{d(1 + \theta R^{f|i}_{t,T;\theta})}{1 + \theta R^{f|i}_{t,T;\theta}} 
   +  \frac{d(1-\theta S^{d|i,\div}_{t,T;\theta})}{1-\theta S^{d|i,\div}_{t,T;\theta}} + O(dt) \\
   &=& \Sigma^{f|i}_{t,T;\theta} \cdot dW^{f|i,T+\theta}_t  -  \Sigma^{d\to i}_{t,T;\theta} \cdot dW^{d|i,T}_t  + O(dt) \\
   &=& \left( \Sigma^{f|f}_{t,T;\theta} + \Sigma^{f\to i}_{t,T;\theta} -   \Sigma^{d\to i}_{t,T;\theta}\right) \cdot dW^{f|i,T+\theta}_t  \\
   &=&  \Sigma^{f|d}_{t,T;\theta} \cdot dW^{f|i,T+\theta}_t  
\end{eqnarray*} 

\end{proof}



\paragraph{Currency basis spreads}

We consider the extension of the previous theorem to a non-native collateral. 
\begin{theo*}
Let us assume that 
\begin{enumerate}
\item native multiplicative funding spreads are constant
\item volatilities are deterministic.
\end{enumerate}
Then, the currency basis spread of a resetting cross-currency swaplet is:
\begin{eqnarray}
   B^{f/d|c,r}_{t,T;\theta}  
   &=& L^{f/d|c}_{t,T;\theta}  - L^{f|c}_{t,T;\theta}   \\
   \theta B^{f/d|c,r}_{t,T;\theta}  &=& (1+\theta R^{f|c}_{t,T;\theta}) \left\{ (1-\theta S^{d|c,\div}_{t,T;\theta}) e^{ - \int_t^T \Sigma^{df}_{u,T} \cdot \Sigma^{d\to c}_{u,T;\theta} \,du } - (1-\theta S^{f|c,\div}_{t,T;\theta}) \right\}
\end{eqnarray} 
\end{theo*}

%
%\begin{proof}[Proof of the theorem]
%By definition, the currency basis satisfies
%\begin{eqnarray*}
%	Z^{f|c}_{t,T+\theta}\theta B^{f/d|c,r}_{t,T;\theta} 
%	&=& 
%	x^{fd}_tZ^{d|c}_{t,T+\theta}\EE^{d|c,T+\theta}_t\left[ x^{df}_T(1+\theta L^d_{T;\theta})\right]  
%	- Z^{f|c}_{t,T+\theta}\EE^{f|c,T+\theta}_t\left[1+\theta L^f_{T;\theta} \right] .
%\end{eqnarray*} 
%Using the lemma, \footnote{TODO: Finish adjusting this!}
%\begin{eqnarray*}
%	\theta B^{f/d|c,r}_{t,T;\theta} 
%	&=& \frac{x^{fd|c}_{t,T+\theta}}{x^{fd|c}_{t,T}} \left(1+\theta L^{d|c}_{t,T;\theta}\right) - \left(1+\theta L^{f|c}_{t,T;\theta} \right)  \\
%	&=& \frac{Z^{d|c}_{t,T+\theta}/Z^{f|c}_{t,T+\theta}}{Z^{d|c}_{t,T}/Z^{f|c}_{t,T}} \left(1+\theta L^{d|c}_{t,T;\theta}\right) - \left(1+\theta L^{f|c}_{t,T;\theta} \right)  \\
%	&=& \frac{1+\theta R^{f|c}_{t,T;\theta}}{1+\theta R^{d|c}_{t,T;\theta}} \left(1+\theta L^{d|c}_{t,T;\theta}\right) - \left(1+\theta L^{f|c}_{t,T;\theta} \right)  \\
%	&=& \left(1+\theta R^{f|c}_{t,T;\theta}\right) \left(  \frac{1+\theta L^{d|c}_{t,T;\theta}}{1+\theta R^{d|c}_{t,T;\theta}} -  \frac{1+\theta L^{f|c}_{t,T;\theta}}{1+\theta R^{f|c}_{t,T;\theta}} \right) \\
%	&=& \left(1+\theta R^{f|c}_{t,T;\theta} \right) \left( \left(1-\theta S^{d|c,\div}_{t,T;\theta}\right) -  \left(1-\theta S^{f|c,\div}_{t,T;\theta}\right)  \right)  
%\end{eqnarray*} 
%giving
%\begin{eqnarray*}
%	B^{f/d|c,r}_{t,T;\theta} 
%	&=& \left(1+\theta R^{f|c}_{t,T;\theta} \right)  \left(S^{f|c,\div}_{t,T;\theta} - S^{d|c,\div}_{t,T;\theta} \right)  
%\end{eqnarray*} 
%\end{proof}
%
%
%
%
%\begin{rema}
%By construction the currency basis spread is a martingale under the corresponding annuity measure. In the case of swaplet this means $B^{f/d|d,r}_{t,T;\theta}$  is a $\QQ^{f|d,T+\theta}$-martingale. We can check that our formula is correct by computing the dynamic of  
%
%We can focus on the resetting Libor:  
%\begin{eqnarray}
%   1 +\theta L^{d|c,f/d}_{t,T;\theta}  &:=& \frac{ \EE^{d|c,T+\theta}_t\left[ x^{df}_{T} \left( 1 +\theta L^{d}_{T;\theta} \right) \right] } {  \EE^{d|c,T+\theta}_t\left[ x^{df}_{T+\theta} \right]  } 
%\end{eqnarray} 
%Then,
%\begin{eqnarray*}
%   & & x^{df}_{t,T+\theta}  \left( 1 +\theta L^{d|c,f/d}_{t,T;\theta} \right)  \\
%   &=&  \EE^{d|c,T+\theta}_t\left[ x^{df}_{T} \left( 1 +\theta L^{d}_{T;\theta} \right) \right]  \\
%   &=&  \frac{Z^{d|c}_{t,T}}{Z^{d|c}_{t,T+\theta}}\EE^{d|c,T}_t\left[ Z^{d|c}_{T,T+\theta}  x^{df}_T(1+\theta L^{d|c}_{T,T;\theta} )\right]  \\
%  &=&  \frac{Z^{d|c}_{t,T}}{Z^{d|c}_{t,T+\theta}}\EE^{d|c,T}_t\left[ x^{df}_T(1-\theta S^{d|c,\div}_{T,T;\theta} )\right]  \\
%  &=& (1+\theta R^{d|c}_{t,T;\theta}) \EE^{d|c,T}_t\left[ x^{df}_T (1-\theta S^{d|c,\div}_{t,T;\theta} )\mathcal{E}\left( -\int_t^T \Sigma^{d\to c}_{u,T;\theta} \cdot dW^{d|c,T}_u \right)  \right]  \\
%  &=&  (1+\theta L^{d|c}_{t,T;\theta}) \EE^{d|c,T}_t\left[ x^{df}_T \mathcal{E}\left( -\int_t^T \Sigma^{d\to c}_{u,T;\theta} \cdot dW^{d|c,T}_u\right)  \right] \\
%  %&=&  (1+\theta R^{d|c}_{t,T;\theta}) \EE^{d|d,T}_t\left[ x^{df}_T  \right] \\	
%  &=&  (1+\theta L^{d|c}_{t,T;\theta}) x^{df|c}_{t,T} \exp\left\{ \int_t^T \Sigma^{c\to d}_{u,T;\theta}  \cdot \left(  \Sigma^{d|d}_{u,T} - \Sigma^{f|f}_{u,T} + \Sigma^{d\to f}_{u,T} + \sigma^{df}_u \right)  du   \right\} \\
%   &=&  (1+\theta L^{d|c}_{t,T;\theta}) x^{df|c}_{t,T} \exp\left\{ \int_t^T \left( \Sigma^{c\to d}_{u,T+\theta} - \Sigma^{c\to d}_{u,T} \right) \cdot \left(  \Sigma^{d|d}_{u,T} - \Sigma^{f|f}_{u,T} + \Sigma^{d\to f}_{u,T} + \sigma^{df}_u \right)  du   \right\} 
%\end{eqnarray*} 
%
%As a result, we find that $B^{f/d|d,r}_{t,T;\theta}$ is a $\QQ^{f|d,T+\theta}$ martingale as expected:
%\begin{eqnarray*}
%  \EE^{f|d,T+\theta}_t\left[ B^{f/d|d,r}_{T,T;\theta} \right] 
%  &=& \EE^{f|d,T+\theta}_t\left[  \left(1+\theta R^{f|d}_{t,T;\theta} \right)  \left(S^{f|d,\div}_{t,T;\theta} - S^{d|d,\div}_{t_0,T;\theta} \right)  \right]  \\
%  &=& \EE^{f|d,T+\theta}_t\left[  \frac{Z^{f|d}_{T,T}}{Z^{f|d}_{T,T+\theta}}  \left(S^{f|d,\div}_{t,T;\theta} - S^{d|d,\div}_{t_0,T;\theta} \right)  \right]  \\
%  &=&  \frac{Z^{f|d}_{t,T}}{Z^{f|d}_{t,T+\theta}} \EE^{f|d,T}_t\left[ \left(S^{f|d,\div}_{t,T;\theta} - S^{d|d,\div}_{t_0,T;\theta} \right)  \right]  \\
%  &=& \left(1+\theta R^{f|d}_{t,T;\theta} \right)  \left(S^{f|d,\div}_{t,T;\theta} - S^{d|d,\div}_{t_0,T;\theta} \right)  
%\end{eqnarray*} 
%\end{rema}



%\subsubsection{Collateral convexity of currency basis spreads}

%We consider the extension of the previous theorem to a non-native collateral. 
%\begin{prop*}
%Let us assume that 
%\begin{enumerate}
%\item native multiplicative funding spreads are constant
%\item volatilities are deterministic.
%\end{enumerate}
%Then, the currency basis spread of a resetting cross-currency swaplet is:
%\begin{eqnarray*}
%	B^{f/d|c,r}_{t,T;\theta} 
%	&=& \left(1+\theta R^{f|c}_{t,T;\theta} \right)  \left(S^{f|c,\div}_{t,T;\theta} - S^{d|c,\div}_{t,T;\theta} \right)  .
%\end{eqnarray*} 
%\end{prop*}


%\begin{lemm}
%Under the previous assumptions, we have 
%\begin{eqnarray}
%	 & & \EE^{d|c,T+\theta}_t\left[ x^{df}_T(1+\theta L^{d}_{T;\theta} )\right] \\
%	&=& x^{df|c}_{t,T} (1+\theta L^{d|c,\div}_{t,T;\theta})  \EE^{f|c,T}_t\left[  \mathcal{E}\left( -\int_t^\cdot \Sigma^{f\to d}_{t,T;\theta}dW^{d|c,T}_t\right)  \right] \\
%	&=& x^{df|c}_{t,T} (1+\theta L^{d|c,\div}_{t,T;\theta})   \exp\left\{ \int_t^T \Sigma^{d\to f}_{t,T;\theta}  \cdot \left(  \Sigma^{d|d}_{t,T} - \Sigma^{f|f}_{t,T} + \Sigma^{d\to f}_{t,T} + \sigma^{df}_u \right)  du   \right\}
%\end{eqnarray} 
%\end{lemm}
%
%
%\begin{proof}
%Using the fact that the funding spreads are constant:
%\begin{eqnarray*}
%	 \EE^{d|c,T+\theta}_t\left[ x^{df}_T(1+\theta L^{d}_{T;\theta} )\right] 
%	&=&  \frac{Z^{d|c}_{t,T}}{Z^{d|c}_{t,T+\theta}}\EE^{d|c,T}_t\left[ Z^{d|c}_{T,T+\theta}  x^{df}_T(1+\theta L^{d|c}_{T,T;\theta} )\right]  \\
%	&=&  \frac{Z^{d|c}_{t,T}}{Z^{d|c}_{t,T+\theta}}\EE^{d|c,T}_t\left[ x^{df}_T(1-\theta S^{d|c,\div}_{T,T;\theta} )\right]  \\
%	&=& (1+\theta R^{d|c,\div}_{t,T;\theta}) \EE^{d|c,T}_t\left[ x^{df}_T (1-\theta S^{d|c,\div}_{t,T;\theta} )\mathcal{E}\left( -\int_t^\cdot \Sigma^{f\to d}_{t,T;\theta}dW^{d|c,T}_t \right)  \right]  \\
%	&=&  (1+\theta L^{d|c}_{t,T;\theta}) \EE^{d|c,T}_t\left[ x^{df}_T \mathcal{E}\left( -\int_t^T \Sigma^{f\to d}_{t,T;\theta}dW^{d|c,T}_t\right)  \right] \\
%	&=&  (1+\theta L^{d|c}_{t,T;\theta}) x^{df|c}_{t,T} \exp\left\{ \int_t^T \Sigma^{d\to f}_{t,T;\theta}  \cdot \left(  \Sigma^{d|d}_{t,T} - \Sigma^{f|f}_{t,T} + \Sigma^{d\to f}_{t,T} + \sigma^{df}_u \right)  du   \right\}
%\end{eqnarray*} 
%where the last equality is the formula for the expectation of the product of two lognormal variables\footnote{Can't we break this down further as a collateral convexity adjustment on the Libor and one on the FX forward?}.
%\end{proof}



%\subsection{Resetting cross-currency swaplets}
%
%
%\begin{prop}
%The currency basis spread of a swaplet is an exponential martingale\footnote{Move this to Dynamic of market observables!}
%\begin{eqnarray}
%	\frac{d B^{f/d|c,r}_{t,T;\theta}}{B^{f/d|c,r}_{t,T;\theta}}  &=&  \left( \Sigma^{f}_{t,T;\theta} - \Sigma^{d \to f}_{t,T;\theta}\right)  \, dW^{f|d,T}_t 
%\end{eqnarray} 
%\end{prop}
%
%\begin{proof}
%We have
%\begin{eqnarray*}
%  & & \frac{d B^{f/d|c,r}_{t,T;\theta}}{B^{f/d|c,r}_{t,T;\theta}}   \\
%  &=&  \frac{d (1+\theta R^{f|c}_{t,T;\theta})}{1+\theta R^{f|c,r}_{t,T;\theta}} + \frac{d \left( (1+\theta S^{f|c,\div}_{t,T;\theta}) - (1+\theta S^{d|c,\div}_{t,T;\theta})\right)}{\left( (1+\theta S^{f|c,\div}_{t,T;\theta}) - (1+\theta S^{d|c,\div}_{t,T;\theta})\right) }  \\
%  & & + \left\langle \frac{d (1+\theta R^{f|c}_{t,T;\theta})}{1+\theta R^{f|c,r}_{t,T;\theta}} ,  \frac{d \left( (1+\theta S^{f|c,\div}_{t,T;\theta}) - (1+\theta S^{d|c,\div}_{t,T;\theta})\right)}{\left( (1+\theta S^{f|c,\div}_{t,T;\theta}) - (1+\theta S^{d|c,\div}_{t,T;\theta})\right) } \right\rangle \\
%  &=&  \Sigma^{f|c}_{t,T;\theta} \cdot dW^{f|c,T+\theta}_t  
%  + \frac{d \left( (1+\theta S^{f|c,\div}_{t,T;\theta}) - (1+\theta S^{d|c,\div}_{t,T;\theta})\right)}{\left( (1+\theta S^{f|c,\div}_{t,T;\theta}) - (1+\theta S^{d|c,\div}_{t,T;\theta})\right) }  \\
%  & & + \left\langle \Sigma^{f|c}_{t,T;\theta} \cdot dW^{f|c,T+\theta}_t  
% ,  \frac{d \left( (1+\theta S^{f|c,\div}_{t,T;\theta}) - (1+\theta S^{d|c,\div}_{t,T;\theta})\right)}{\left( (1+\theta S^{f|c,\div}_{t,T;\theta}) - (1+\theta S^{d|c,\div}_{t,T;\theta})\right) } \right\rangle \\
%  &=&  \ldots 
%\end{eqnarray*} 
%\end{proof}


\subsubsection{Comparison to the Determinisctic Collateral Basis Model}

\begin{rema}[Comparison to the Deterministic collateral basis model]
According to \autoref{eq:differential_discounting}, in the Deterministic collateral basis models, discount rates are given by:
\begin{eqnarray*}
	R^{f|d,\DCB}_{t,T;\theta} 
	&:=&  L^{f|f}_{t,T;\theta} + B^{f/d|d,r}_{t,T;\theta}   + S^{d|d,-}_{t,T;\theta}
\end{eqnarray*} 
How does our new formula compare to this?

According to our derivation, 
\begin{eqnarray*}
	 \theta B^{f/d|d,r}_{t,T;\theta} 
	 &=& (1+\theta R^{f|d}_{t,T;\theta})\frac{1+\theta L^{d|d}_{t,T;\theta}}{1+\theta R^{d|d}_{t,T;\theta}} - \left(1+\theta L^{f|d}_{t,T;\theta} \right)  \\
	 &=& (1+\theta R^{f|d}_{t,T;\theta})(1-\theta S^{d|d,\div}_{t,T;\theta}) - \left(1+\theta L^{f|d}_{t,T;\theta} \right)  \\
	 B^{f/d|d,r}_{t,T;\theta}  &=& R^{f|d}_{t,T;\theta} -  L^{f|d}_{t,T;\theta}  - S^{d|d,\div}_{t,T;\theta} - \theta R^{f|d}_{t,T;\theta} S^{d|d,\div}_{t,T;\theta}   \\
  \theta B^{f/d|d,r}_{t,T;\theta} 
  &=& R^{f|d}_{t,T;\theta} -  L^{f|f}_{t,T;\theta} - S^{d|d,-}_{t,T;\theta}  \\
  &  & - \left(L^{f|d}_{t,T;\theta} - L^{f|f}_{t,T;\theta}\right) + \left(S^{d|d,-}_{t,T;\theta} - S^{d|d,\div}_{t,T;\theta}\right) -  \theta R^{f|d}_{t,T;\theta}S^{d|d,\div}_{t,T;\theta}
\end{eqnarray*} 
So we find that the correction to the deterministic approximation consists of 
\begin{itemize}
\item a convexity adjustment for the Libor $L^{f|d}_{t,T;\theta} - L^{f|f}_{t,T;\theta}$
\item a second (third?) order term depending on the level of the funding spread $S^{d|d,-}_{t,T;\theta} - S^{d|d,\div}_{t,T;\theta}$
\item a second order term $\theta R^{f|d}_{t,T;\theta}S^{d|d,\div}_{t,T;\theta}$.
\end{itemize}
Neglecting, the higher order terms, we find:
\begin{eqnarray}
  R^{f|d}_{t,T;\theta} &\approx& 
  L^{f|d}_{t,T;\theta} + B^{f/d|d,r}_{t,T;\theta}   + S^{d|d,-}_{t,T;\theta} 
  \label{eq:comparison_to_dd}
\end{eqnarray} 
Note that the convexity adjustment on Libors can be computed from the model parameters.
\end{rema}



The \autoref{theo:currency_basis_spread_constant_funding_basis} and the previous analysis considers the currency basis spread as an output of the model when it is in fact a market observable. It would be more telling to compare discount rates $R^{f|d}_{t,T;\theta}$ when models are calibrated to the same currency basis spreads. 

\begin{prop}[Linear expansion]
Following \autoref{eq:differential_discounting}, let us define  
\begin{eqnarray*}
	R^{f|d,\DCB}_{t,T;\theta} &=&  L^{f|f}_{t,T;\theta} + B^{f/d|d,r}_{t,T;\theta} + S^{d|d,-}_{t,T;\theta}
\end{eqnarray*} 
Then 
\begin{eqnarray*}
	R^{f|d}_{t,T;\theta}  &=& R^{f|d,\DCB}_{t,T;\theta} + \ldots  
\end{eqnarray*} 
[TODO: Finish]
\end{prop}

\begin{proof}
\begin{eqnarray*}
   &  & R^{f|d}_{t,T;\theta} - R^{f|d,\DCB}_{t,T;\theta} \\
   &=& \frac{1}{\theta}B^{f/d|d,r}_{t,T;\theta}\frac{1}{ S^{f|d,\div}_{t,T;\theta} - S^{d|d,\div}_{t_0,T;\theta} }  -  L^{f|f}_{t,T;\theta} - B^{f/d|d,r}_{t,T;\theta}  - S^{d|d,-}_{t,T;\theta} \\
   &=& B^{f/d|d,r}_{t,T;\theta}\left( \frac{1}{ (1+\theta S^{f|d,\div}_{t_0,T;\theta}) - (1+\theta S^{d|d,\div}_{t_0,T;\theta}) }  - 1 \right)  - L^{f|f}_{t,T;\theta} - S^{d|d,-}_{t,T;\theta} \\
   &=& B^{f/d|d,r}_{t,T;\theta} \left(\frac{(1+\theta R^{f|f}_{t_0,T;\theta})(1+\theta R^{d|d}_{t_0,T;\theta}) }{ (1+\theta L^{f|f}_{t_0,T;\theta})(1+\theta R^{d|d}_{t_0,T;\theta}) - (1+\theta R^{f|f}_{t_0,T;\theta})(1+\theta L^{d|d}_{t_0,T;\theta}) } - 1 \right)   -  L^{f|f}_{t,T;\theta} - S^{d|d,-}_{t,T;\theta} \\
   &=& B^{f/d|d,r}_{t,T;\theta} \frac{1 + \theta ( R^{f|f}_{t_0,T;\theta}+ R^{d|d}_{t_0,T;\theta}) + \theta^2  R^{f|f}_{t_0,T;\theta}  R^{d|d}_{t_0,T;\theta} - \theta L^{f|f}_{t_0,T;\theta} + \theta L^{d|d}_{t_0,T;\theta}}{  \theta (L^{f|f}_{t_0,T;\theta} - L^{d|d}_{t_0,T;\theta})  }   -  L^{f|f}_{t,T;\theta} - S^{d|d,-}_{t,T;\theta} \\
   &=& B^{f/d|d,r}_{t,T;\theta} \frac{1 + \theta ( S^{f|f,-}_{t_0,T;\theta}+ S^{d|d,-}_{t_0,T;\theta}) + \theta^2  R^{f|f}_{t_0,T;\theta}  R^{d|d}_{t_0,T;\theta} }{  \theta (L^{f|f}_{t_0,T;\theta} - L^{d|d}_{t_0,T;\theta})  }  -  L^{f|f}_{t,T;\theta} - S^{d|d,-}_{t,T;\theta} \\
\end{eqnarray*} 
\end{proof}








\begin{rema}[FIXME: This is false as it assumes $S^{f|d,\div}_{t_0,T;\theta} = S^{f|f,\div}_{t_0,T;\theta}$, ignoring the collateral volatility!]
\autoref{eq:cb_r_swaplet} is surprisingly simple
\begin{enumerate}
\item The currency basis spread $B^{f/d|d,r}_{t,T;\theta}$ is a $\QQ^{f|d,T+\theta}$-martingale as expected.
\item \autoref{eq:cb_r_swaplet2} shows that $Z^{f|d}_{t,T;\theta}B^{f/d|d,r}_{t,T;\theta} $ is exactly equal to the difference in funding spreads which corresponds to the financial intuition for the existence of currency basis.  
\item Even though we took into account convexity adjustments, the relation between the currency basis $B^{f/d|d,r}_{t,T;\theta} $ and the discount rate $R^{f|d}_{t,T;\theta}$ does not depend on the covariance matrix of the diffusion. This is a small miracle in itself.  It implies that the discount curve $\{Z^{f|d}_{t,T}\}$ can still be bootstrapped from resetting cross-currency swaps without calibrating the model. Only the off-market discount curves would depend on the covariance structure. 
\item The currency basis spread does not depend on the FX spot or its volatility.
\item The volatility of the currency basis spread does depend on the volatility of the foreign rate. 
\item \autoref{eq:cb_r_swaplet} compares nicely to DCB \autoref{eq:differential_discounting} (see below). \end{enumerate}
\end{rema}

\begin{rema}
Let us assume that both the DCB model and the stochastic collateral model are calibrated to the same instruments so that currency basis spreads match. What is the impact of stochastic collateral on discount factors? 

For projection rate pseudo-discount factors:
\begin{eqnarray*}
  \frac{Z^{f|d,\mathrm{proj},\DCB}_{t,S,T}}{Z^{f|d,\mathrm{proj}}_{t,S,T}} 
  &=& \frac{1+\theta L^{f|d}_{t,S,T}}{1+\theta L^{f|d,\DCB}_{t,S,T}} \\
  &=& \frac{1+\theta L^{f|d}_{t,S,T}}{1+\theta L^{f|f}_{t,S,T}} \\
  &=& \exp\left( -\int_t^S \Sigma^{f|f}_{u,S,T} \cdot \Sigma^{f\to d}_{u,T}\, du \right) 
\end{eqnarray*} 
Compounding the forward pseudo-discount factors, we find
\begin{eqnarray*}
  \frac{Z^{f|d,\mathrm{proj},\DCB}_{t,T}}{Z^{f|d,\mathrm{proj}}_{t,T}} 
  &=& \prod_i \frac{Z^{f|d,\mathrm{proj},\DCB}_{t,T_{i-1},T_i}}{Z^{f|d,\mathrm{proj}}_{t,T_{i-1},T_i}}  \\
  &=& \prod_i  \exp\left( -\int_t^{T_{i-1}} \Sigma^{f|f}_{u,T_{i-1},T_i} \cdot \Sigma^{f\to d}_{u,T_i}\, du \right)   \\
  &=&  \exp\left( - \sum_i \int_t^{T_{i-1}} \Sigma^{f|f}_{u,T_{i-1},T_i} \cdot \Sigma^{f\to d}_{u,T_i}\, du \right)   \\
  &=&  \exp\left( - \sum_i \sum_{k=1}^{i-1} \int_{T_{k-1}}^{T_{k}} \Sigma^{f|f}_{u,T_{i-1},T_i} \cdot \Sigma^{f\to d}_{u,T_i}\, du \right)   \\
  &=&  \exp\left( - \sum_{k=1}^{q(T)} \sum_{i=k+1}^{q(T)}  \int_{T_{k-1}}^{T_{k}} \Sigma^{f|f}_{u,T_{i-1},T_i} \cdot \Sigma^{f\to d}_{u,T_i}\, du \right)   \\
  &=&  \exp\left( - \sum_{k=1}^{q(T)} \int_{T_{k-1}}^{T_{k}} \sum_{i=k+1}^{q(T)}  \Sigma^{f|f}_{u,T_{i-1},T_i} \cdot \Sigma^{f\to d}_{u,T_i}\, du \right)   \\
  &=&  \exp\left( - \sum_{k=1}^{q(T)} \int_{T_{k-1}}^{T_{k}} \sum_{i=k+1}^{q(T)}  \Sigma^{f|f}_{u,T_i} \cdot \Sigma^{f\to d}_{u,T_i} - \sum_{i=k+1}^{q(T)}  \Sigma^{f|f}_{u,T_{i-1}} \cdot \Sigma^{f\to d}_{u,T_i} \, du \right)   \\
  &=&  \exp\left( - \sum_{k=1}^{q(T)} \int_{T_{k-1}}^{T_{k}} \sum_{i=k+1}^{q(T)}  \Sigma^{f|f}_{u,T_i} \cdot \Sigma^{f\to d}_{u,T_i} - \sum_{i=k}^{q(T)-1}  \Sigma^{f|f}_{u,T_{i}} \cdot \Sigma^{f\to d}_{u,T_{i+1}} \, du \right)   \\
  &=&  \exp\left( - \sum_{k=1}^{q(T)} \int_{T_{k-1}}^{T_{k}} \Sigma^{f|f}_{u,T} \cdot \Sigma^{f\to d}_{u,T} - \Sigma^{f|f}_{u,T_k} \cdot \Sigma^{f\to d}_{u,T_k} - \sum_{i=k+1}^{q(T)-1}  \Sigma^{f|f}_{u,T_i} \cdot \Sigma^{f\to d}_{u,T_{i-1},T_i} \, du \right)   
\end{eqnarray*} 
\begin{eqnarray*}
  \frac{Z^{f|d,\mathrm{proj},\DCB}_{t,T}}{Z^{f|d,\mathrm{proj}}_{t,T}} 
  &=& \prod_i \frac{Z^{f|d,\mathrm{proj},\DCB}_{t,T_{i-1},T_i}}{Z^{f|d,\mathrm{proj}}_{t,T_{i-1},T_i}}  \\
  &=& \prod_i  \exp\left( -\int_t^{T_{i-1}} \Sigma^{f|f}_{u,T_{i-1},T_i} \cdot \Sigma^{f\to d}_{u,T_i}\, du \right)   \\
  &=&  \exp\left( - \sum_i \int_t^{T_{i-1}} \Sigma^{f|f}_{u,T_{i-1},T_i} \cdot \Sigma^{f\to d}_{u,T_i}\, du \right)   \\
  &=& \exp\left( - \sum_i \int_t^{T_{i-1}} \Sigma^{f|f}_{u,T_i} \cdot \Sigma^{f\to d}_{u,T_i}\, du - \int_t^{T_{i-1}} \Sigma^{f|f}_{u,T_{i-1}} \cdot \Sigma^{f\to d}_{u,T_i}\, du \right)   \\
 \end{eqnarray*} 

\end{rema}




\newpage%
\subsubsection{Resetting cross-currency swaps}


\subsubsection{Dynamic}
From the previous proposition, it is standard to derive the dynamic of the currency basis spread of a swap. It can be written
\begin{eqnarray}
	\frac{dB^{f/d|c,r}_{t,T;M}}{B^{f/d|c,r}_{t,T;M}}  & = &  \left( \sum_{i=1}^{\f M} w^i_t \Sigma^{f|d}_{t,T_{i-1},T_i} \right)  dW^{f|d,T_0,T_M}_t
\end{eqnarray} 
where the stochastic weights are 
\begin{eqnarray}
	w^j_t &=& \frac{B^{f/d|c,r}_{t,T_{j-1},T_j}}{B^{f/d|c,r}_{t,T;M}} \frac{\partial B}{\partial B_j} 
\end{eqnarray} 
We can then freeze those to their initial values.

Since the forward spreads $B^{f/d|c,r}_{t,T_{j-1},T_j}$ are not directly observables, it would be better to express the vol directly as a function of the currency basis spread of swaps: $B^{f/d|c,r}_{t,T_{0},T_k}$. 
This is essentially equivalent to stripping the discount curve $Z^{f|d}_{t,T}$.


\subsubsection{Relation to swaplet spreads}
Let us define  
\begin{eqnarray}
	\omega_j  &=& S^{f|f,\div}_{t_0,T_{j-1},T_j}  -  S^{d|d,\div}_{t_0,T_{j-1},T_j} 
\end{eqnarray}
This is a sort of ``funding basis basis''. We also introduce 
\begin{eqnarray}
	\Omega_n &=& \prod_{j=1}^n \frac{1}{\omega_j}
\end{eqnarray}
so that
\begin{eqnarray*}
	Z^{f|d}_{t,T_0,T_n} 
	= \prod_{j=1}^n Z^{f|d}_{t,T_{j-1},T_j}  
	= \prod_{j=1}^n \frac{1}{\omega_j} B^{f/d|d,r}_{t,T_{j-1},T_j}  
	= \Omega_n \prod_{k=1}^j \frac{1}{B^{f/d|d,r}_{t,T_{j-1},T_j}}  
\end{eqnarray*}


Then the currency-basis spread of a swap reads  
\begin{eqnarray*}
	B^{f/d|d,r}_{t,T;M}  
	&=& \frac{\sum_{j=1}^{M}   Z^{f|d}_{t,T_0,T_j} B^{f/d|d,r}_{t,T_{j-1},T_j}  
	}{\sum_{j=1}^{M}  Z^{f|d}_{t,T_0,T_j}} \\
	&=& \frac{ \sum_{j=1}^{M}   Z^{f|d}_{t,T_0,T_j} \omega_j \left(1+\theta R^{f|d}_{t,T_{j-1},T_j}\right) }{\sum_{j=1}^{M}  Z^{f|d}_{t,T_0,T_j}} \\
	&=& 
	 \frac{ \sum_{j=1}^{M} \omega_j \prod_{k=1}^{j-1} Z^{f|d}_{t,T_{k-1},T_{k}} } {  \sum_{j=1}^{M}  \prod_{k=1}^{j} Z^{f|d}_{t,T_{k-1},T_{k}}  } \\
	&=& 
	\frac{ \sum_{j=1}^{M} \Omega_j \prod_{k=1}^{j-1} \frac{1}{B^{f/d|d,r}_{t,T_{k-1},T_k}}} {  \sum_{j=1}^{M}  \Omega_j \prod_{k=1}^{j} \frac{1}{B^{f/d|d,r}_{t,T_{k-1},T_k}}  } .
	%\label{eq:cb_resetting_swap}
\end{eqnarray*}

\subsubsection{Relations between currency basis spreads}

\begin{eqnarray*}
	\left( \sum_{j=1}^{M+1}  Z^{f|d}_{t,T_0,T_j} \right) B^{f/d|d,r}_{t,T,M+1}  
	&=& \left( \sum_{j=1}^{M}  Z^{f|d}_{t,T_0,T_j} \right) B^{f/d|d,r}_{t,T_0,T_{M}}  + Z^{f|d}_{t,T,M+1} B^{f/d|d,r}_{t,T_{M},T_{M+1}}  
\end{eqnarray*}
So, 
\begin{eqnarray*}
	B^{f/d|d,r}_{t,T_{b},T_{c}} &=&  \frac{A^{f|d}_{t,T_a,T_{c}} B^{f/d|d,r}_{t,T_a,T_{c}}  - A^{f|d}_{t,T_a,T_{b}} B^{f/d|d,r}_{t,T_a,T_b} }
	{  A^{f|d}_{t,T_a,T_c} - A^{f|d}_{t,T_a,T_b} } 
\end{eqnarray*}


As a result,
\begin{eqnarray*}
	B^{f/d|d,r}_{t,T;M}  
	&=& \frac{ \sum_{j=1}^{M}   Z^{f|d}_{t,T_0,T_j} B^{f/d|d,r}_{t,T_{j-1},T_j} }{\sum_{j=1}^{M}  Z^{f|d}_{t,T_0,T_j}} \\
	&=& \frac{ \sum_{j=1}^{M}   Z^{f|d}_{t,T_0,T_j} \omega_j \left(1+\theta Z^{f|d}_{t,T_{j-1},T_j}\right) }{\sum_{j=1}^{M}  Z^{f|d}_{t,T_0,T_j}} \\
	&=& \frac{ \sum_{j=1}^{M} \omega_j \prod_{k=1}^{j-1} Z^{f|d}_{t,T_{k-1},T_{k}} } {  \sum_{j=1}^{M}  \prod_{k=1}^{j} Z^{f|d}_{t,T_{k-1},T_{k}}  }
\end{eqnarray*}




\subsubsection{Dynamic of currency basis spreads}


\begin{prop}
The currency basis spread of a swap has dynamic
\begin{eqnarray}
	dB^{f/d|c,r}_{t,T;M}  & = &  \left( \sum_{i=1}^{\f M} w^i_t \Sigma^{f|d}_{t,T_{i-1},T_i} \right)  dW^{f|d,T_0,T_M}_t
\end{eqnarray} 
where the stochastic weights $w^j_t$ are defined by
\begin{eqnarray}
	w^j_t &=& B^{f/d|c,r}_{t,T_{j-1},T_j} \frac{\partial B}{\partial B_j} \\
	&=&  \ldots 
\end{eqnarray} 
\end{prop}

\begin{proof}
Let us define the ``funding basis basis'' as 
\begin{eqnarray}
	\omega_j  &=& S^{f|f,\div}_{t_0,T_{j-1},T_j}  -  S^{d|d,\div}_{t_0,T_{j-1},T_j} \\
	\Omega_n &=& \prod_{j=1}^n \frac{1}{\omega_j}
\end{eqnarray}
Then, 
\begin{eqnarray*}
	Z^{f|d}_{t,T_0,T_n} 
	&=& \prod_{j=1}^n Z^{f|d}_{t,T_{j-1},T_j}  \\
	&=& \prod_{j=1}^n \frac{1}{\omega_j} B^{f/d|d,r}_{t,T_{j-1},T_j}  \\
	&=& \Omega_n \prod_{k=1}^j B^{f/d|d,r}_{t,T_{j-1},T_j}  
\end{eqnarray*}
Then 
%\begin{eqnarray*}
%	B^{f/d|d,r}_{t,T;M}  
%	&=& \frac{ \sum_{j=1}^{M}   Z^{f|d}_{t,T_0,T_j} B^{f/d|d,r}_{t,T_{j-1},T_j} }{\sum_{j=1}^{M}  Z^{f|d}_{t,T_0,T_j}} \\
%	&=& \frac{ \sum_{j=1}^{M}   Z^{f|d}_{t,T_0,T_j} \omega_j \left(1+\theta Z^{f|d}_{t,T_{j-1},T_j}\right) }{\sum_{j=1}^{M}  Z^{f|d}_{t,T_0,T_j}} \\
%	&=& \frac{ \sum_{j=1}^{M} \omega_j Z^{f|d}_{t,T_0,T_{j-1}} } {  \sum_{j=1}^{M}  Z^{f|d}_{t,T_0,T_{j}} }
%\end{eqnarray*}
%We can rewrite that in terms of currency basis spreads:
\begin{eqnarray*}
	B^{f/d|d,r}_{t,T;M}  
	&=& \frac{ \sum_{j=1}^{M}   Z^{f|d}_{t,T_0,T_j} B^{f/d|d,r}_{t,T_{j-1},T_j} }{\sum_{j=1}^{M}  Z^{f|d}_{t,T_0,T_j}} \\
	&=& \frac{ \sum_{j=1}^{M}   Z^{f|d}_{t,T_0,T_j} \omega_j \left(1+\theta Z^{f|d}_{t,T_{j-1},T_j}\right) }{\sum_{j=1}^{M}  Z^{f|d}_{t,T_0,T_j}} \\
	&=& \frac{ \sum_{j=1}^{M} \omega_j \prod_{k=1}^{j-1} Z^{f|d}_{t,T_{k-1},T_{k}} } {  \sum_{j=1}^{M}  \prod_{k=1}^{j} Z^{f|d}_{t,T_{k-1},T_{k}}  }
\end{eqnarray*}
We have a function 
\begin{eqnarray*}
	B(z_1,\ldots,z_M) &=& \frac{ \sum_{j=1}^{M} \omega_j \prod_{k=1}^{j-1} z_k  } {  \sum_{j=1}^{M}  \prod_{k=1}^{j} z_k }
\end{eqnarray*}
Its derivatives are 
\begin{eqnarray}
	z_i\partial_i B &=& \frac{ \left( \sum_{j=i+1}^{M} \omega_j \prod_{k=1}^{j-1} z_k \right)  \left( \sum_{j=1}^{M}  \prod_{k=1}^{j} z_k \right) 
	- \left( \sum_{j=1}^{M} \omega_j \prod_{k=1}^{j-1} z_k \right)  \left( \sum_{j=i+1}^{M}  \prod_{k=1}^{j} z_k \right) }
	{  \left( \sum_{j=1}^{M}  \prod_{k=1}^{j} z_k \right)^2 }   
\end{eqnarray} 
\end{proof}



\begin{prop}[Gaussian approximations]
The currency basis spread of a resetting swaplet are approximately exponential martingales:
\begin{eqnarray}
	\frac{dB^{f/d|c,r}_{t,T;M})}{B^{f/d|c,r}_{t,T;M}}  &\approx& 
	\widetilde{\Sigma}^{f/d|c,r}_{t,T;M} dW^{f|d,T_0,T_M} .
\end{eqnarray} 
where 
\begin{eqnarray}
	\widetilde{\Sigma}^{f/d|c,r}_{t,T;M}  & = &  \ldots
\end{eqnarray} 
\end{prop}

\begin{proof}
We know that  $B^{f/d|c,r}_{t,T;M}$ is a martingale under the measure corresponding to the annuity num\'eraire $A^{f|d}_{\cdot,T;M}$ so it is enough to compute its volatility. Moreover, we have the relation 
\begin{eqnarray*}
	B^{f/d|c,r}_{t,T;M}  &=& \frac{\sum_{j=1}^{M} \theta^f_j  Z^{f|c}_{t,T_j} B^{f/d|c,r}_{t,T_{j-1},T_j}}{\sum_{j=1}^{M} \theta^f_j Z^{f|c}_{t,T_j}}
\end{eqnarray*}
so
\begin{eqnarray*}
	dB^{f/d|c,r}_{t,T;M}  &=& \sum_{j} \frac{\partial Y}{\partial B_j} d B^{f/d|c,r}_{t,T_{j-1},T_j} + O(dt) \\
	&=& 
\end{eqnarray*}
\end{proof}



\subsection{[TODO: Remove this] Calibration to currency basis swaps}


\subsubsection{Discount curve with domestic collateral}

According to \autoref{eq:comparison_to_dd}, the discount rate with domestic collateral is 
\begin{eqnarray*}
  R^{f|d}_{t,T;\theta} &\approx& L^{f|d}_{t,T;\theta} + B^{f/d|d,r}_{r,T;\theta} + S^{d|d,-}_{t,T;\theta}  \\
  &=& \frac{1}{\theta}\left( \left( 1 + \theta L^{f|f}_{t,T;\theta} \right) e^{-\int_t^S \Sigma^{f|f}_{u,T;\theta} \cdot \Sigma^{f\to d}_{u,T}\, du } - 1\right) + B^{f/d|d,r}_{r,T;\theta} + S^{d|d,-}_{t,T;\theta}
\end{eqnarray*}
It follows that, given the spot volatilities, the discount factors are:
%\begin{eqnarray*}
%  Z^{f|d}_{t,T} 
%  &=& \prod_i \frac{1}{1+\theta R^{f|d}_{t,T_{i};\theta}}  \\
%  &=& \prod_i \frac{1}{1+\theta \left(L^{f|d}_{t,T_{i};\theta} +  B^{f/d|d,r}_{r,T_{i};\theta} + S^{d|d,-}_{t,T_{i};\theta}\right)}  \\
%  &=& \prod_i \frac{1}{\left(1+\theta L^{f|f}_{t,T_{i};\theta}\right)e^{-\int_t^{T_{i}} \Sigma^{f|f}_{u,T_{i};\theta} \cdot \Sigma^{f\to d}_{u,T_{i+1}}\, du } +  \theta B^{f/d|d,r}_{r,T_{i};\theta} + \theta S^{d|d,-}_{t,T_{i};\theta}}  %\\
%  %&=& \prod_i \frac{1}{\left(1+\theta L^{f|f}_{t,T_{i};\theta}\right)e^{-\int_t^{T_{i}} \Sigma^{f|f}_{u,T_{i};\theta} \cdot \Sigma^{f\to d}_{u,T_{i+1}}\, du } +  \theta B^{f/d|d,r}_{r,T_{i};\theta} + \theta S^{d|d,-}_{t,T_{i};\theta}}  
%\end{eqnarray*}
\begin{eqnarray*}
  Z^{f|d}_{t,T} 
  &=& \prod_i \frac{1}{1+\theta R^{f|d}_{t,T_{i};\theta}}  \\
  &=& \prod_i \frac{1}{1+\theta R^{f|d,\DCB}_{t,T_{i};\theta}} \frac{1+\theta R^{f|d,\DCB}_{t,T_{i};\theta}}{1+\theta R^{f|d}_{t,T_{i};\theta}} \\
  &=& Z^{f|d,\DCB}_{t,T} \prod_i \frac{1}{1+\frac{\theta(R^{f|d}_{t,T_{i};\theta}-R^{f|d,\DCB}_{t,T_{i};\theta})}{1+\theta R^{f|d,\DCB}_{t,T_{i};\theta}}}   \\
   &=& Z^{f|d,\DCB}_{t,T} \prod_i \left(1-\frac{\theta(R^{f|d}_{t,T_{i};\theta}-R^{f|d,\DCB}_{t,T_{i};\theta})}{1+\theta R^{f|d,\DCB}_{t,T_{i};\theta}} \right)  \\
   &=& Z^{f|d,\DCB}_{t,T} \prod_i \left(1-\frac{\theta(L^{f|d}_{t,T_{i};\theta}-L^{f|f}_{t,T_{i};\theta})}{1+\theta R^{f|d,\DCB}_{t,T_{i};\theta}} \right)  
\end{eqnarray*}




\begin{rema}[Bootstrapping of cross-currency curves]
The above data allows us to bootstrap curves and price any fixed or Libor cashflow under any collateral. 
The internal parameters are only required for off-market collateral.

%\begin{coro}
The currency basis spread of a resetting cross-currency swap is:
\begin{eqnarray}
	B^{f/d|d,r}_{t,T_M,T_{M+1}} &=& 
	\frac{A^{f|d}_{t,T;M} }{\theta Z^{f|d}_{t,T_{M+1}}} \left( B^{f/d|d,r}_{t,T,M+1} - B^{f/d|d,r}_{t,T;M} \right) + B^{f/d|d,r}_{t,T,M+1}
\end{eqnarray}


%\begin{proof}
We have:
\begin{eqnarray*}
	A^{f|d}_{t,T,M+1}   B^{f/d|d,r}_{t,T,M+1}  &=& 
	A^{f|d}_{t,T;M}   B^{f/d|d,r}_{t,T;M} + Z^{f|d}_{t,T_{M+1}} B^{f/d|d,r}_{t,T_M,T_{M+1}} \\
	\left(A^{f|d}_{t,T_0,T_{M}}  +\theta Z^{f|d}_{t,T_{M+1}} \right) B^{f/d|d,r}_{t,T,M+1}  &=& 
	A^{f|d}_{t,T;M}   B^{f/d|d,r}_{t,T;M} + \theta Z^{f|d}_{t,T_{M+1}} B^{f/d|d,r}_{t,T_M,T_{M+1}} \\
	A^{f|d}_{t,T_0,T_{M}} \left( B^{f/d|d,r}_{t,T,M+1}  - B^{f/d|d,r}_{t,T_0,T_{M}}  \right)  &=& 
	\theta Z^{f|d}_{t,T_{M+1}} \left( B^{f/d|d,r}_{t,T_M,T_{M+1}} - B^{f/d|d,r}_{t,T,M+1} \right) \\
	\frac{A^{f|d}_{t,T_0,T_{M}}}{\theta Z^{f|d}_{t,T_{M}}} \left( B^{f/d|d,r}_{t,T,M+1}  - B^{f/d|d,r}_{t,T_0,T_{M}}  \right)  &=& 
	Z^{f|d}_{t,T_M,T_{M+1}} \left( B^{f/d|d,r}_{t,T_M,T_{M+1}} - B^{f/d|d,r}_{t,T,M+1} \right) \\
	&=& S^{f|f,\div}_{t,T_M,T_{M+1}} - S^{d|d,\div}_{t,T_M,T_{M+1}}  - Z^{f|d}_{t,T_M,T_{M+1}} B^{f/d|d,r}_{t,T,M+1} 
\end{eqnarray*}
\begin{eqnarray}
	Z^{f|d}_{t,T_M,T_{M+1}} B^{f/d|d,r}_{t,T,M+1} 
       &=& S^{d|d,\div}_{t,T_M,T_{M+1}} - S^{f|f,\div}_{t,T_M,T_{M+1}} +  \frac{A^{f|d}_{t,T_0,T_{M}}}{\theta Z^{f|d}_{t,T_{M}}} \left( B^{f/d|d,r}_{t,T,M+1}  - B^{f/d|d,r}_{t,T_0,T_{M}}  \right)
\end{eqnarray}
%\end{proof}

We can write that as a system of linear equations for the discount factors.

Interpolation methods can be included.
\end{rema}





\newpage%
\section{Review of Linear Gaussian Model}



\subsection{Reminder on LGM parametrisations: Single currency case}%
This is detailed in \cite[\S 10]{AndersenPiterbarg}.

\subsubsection{Separable volatilities}%


\begin{hypo}[Separable volatility]
We assume:
\begin{eqnarray}
  \sigma_{t,T}  &=& g(t)h(T)
\end{eqnarray}
where 
\begin{enumerate}
\item $h : \RR_+ \to \RR^d$ is a deterministic differentiable vector valued function.
\item $g : \RR_+ \to M_{k,d}(\RR)$ is a matrix-valued stochastic process satisfying 
\end{enumerate}
such that 
\begin{enumerate}
\item $\sigma_{t,T}^\top \sigma_{t,T}$ is invertible for all $t,T$.  
\end{enumerate}
\end{hypo}

\subsubsection{Parametrisation by a Gaussian martingale}%


\begin{prop}
The model can be described by the martingale state variable: 
\begin{eqnarray}
  Z_t &=& \int_0^t g(u)^\top dW_u 
\end{eqnarray}
and the following reconstruction formulas:
\begin{enumerate}
\item The forward rate is:
\begin{eqnarray}
  r_{t,T} &=& r_{0,t} + \Omega_{t,T} + h_t^\top Z_t 
\end{eqnarray}
where
\begin{eqnarray}
  \Omega_{t,T} &=& h(T)^\top \int_0^t g(s)^\top g(s) \int_0^s h(u) \, du
\end{eqnarray}
\item The ZCB is:
\begin{eqnarray}
  Z_{t,T} &=& Z_{0,t,T} \exp\left( ??   \right) 
\end{eqnarray}
\end{enumerate}
\end{prop}



\subsubsection{Parametrisation by mean reverting state variables}%


\begin{prop}
The model can be described by the variables 
\begin{eqnarray}
  z_t &=& H_t \int_0^t g_s^\top g_s \int_s^t h_u\,du  \, ds +  H_t \int_0^t H_s^{-1}\sigma_{z,s}^\top dW_s \\
  y_t &=& H_t \int_0^t g_s^\top g_s \, ds  H_t
\end{eqnarray}
with dynamic
\begin{eqnarray}
  dz_t &=& \left(y_t\vun - \beta_tz_t\right)dt + \sigma_{z,t}^\top dW_t   \\
  dy_t &=& H_tg_t^\top g_tH_t - \beta_ty_t - y_t\beta_t  
\end{eqnarray}
where
\begin{eqnarray}
  H_{t} &=& \diag(h_t)  \\
  \sigma_{z,s} &=& g_sH_s  \\ 
  M_{t,T} &=& H_T H_t^{-1} \vun \\
  G_{t,T} &=& \int_t^T M_{t,u} \, du
\end{eqnarray}
and the following reconstruction formulas:
\begin{enumerate}
\item The forward rate is:
\begin{eqnarray}
  r_{t,T} &=& r_{0,t} + M_{t,T}^\top\left( z_t + y_t \int_t^T M_{t,u} \, du \right) 
\end{eqnarray}
\item The ZCB is:
\begin{eqnarray}
  Z_{t,T} &=& Z_{0,t,T} \exp\left( -G_{t,T} ^\top z_t -\frac{1}{2} G_{t,T} ^\top y_t G_{t,T}   \right) 
\end{eqnarray}
\end{enumerate}
\end{prop}


\subsubsection{Parametrisation by correlated Brownian motions}%
It is much more intuitive to have a finite number of state variables corresponding to each asset. 
We now explain how OU processes driven by correlated Brownian motions lead to a multi-dimensional OU process. Conversely, we investigate under which conditions a multi-dimensional OU process can de decomposed into one dimensional ones\footnote{It seems that this breaks symmetry: this property is not invariant by change of reference currency.}.


\begin{prop}
The model can be described by the variables 
%\begin{eqnarray}
%  dx^i_t &=& \left(y^i_t\vun - \beta^i_tx^i_t\right)dt + \sigma_{z,t} \Gamma_t^\top d\widetilde{W}_t^i   
%\end{eqnarray}
with dynamic 
\begin{eqnarray}
  dx^i_t &=& \left(\sum_j y^{ij}_t - \beta^i_tx^i_t\right)dt + \sigma^i_{t}  d\widetilde{W}_t^i \\   
  d\langle \widetilde{W}^i, \widetilde{W}^j \rangle_t  &=& \rho^{ij}_t \,dt 
\end{eqnarray}
where
\begin{eqnarray}
  \sigma_{t} &=& \diag(\sigma_{z,t}^{\top}\sigma_{z,t})^{1/2} \\
  \sigma_t^i  &=& \| \sigma_x^{\cdot i}(t)\|_2 = \sqrt{ \sum_j \sigma_{z}^{ji}(t)^2 } \\
  d\widetilde{W}_t^i   &=& \frac{1}{ \sigma_{t}^i } \sum_j \sigma_{z}^{ji}(t)  dW_t^j   \\
  \rho^{ij}(t) 
  &=& \frac{[\sigma_{z}^\top\sigma_{z}]^{ij}_t}{\sqrt{[\sigma_{z}^\top\sigma_{z}]^{ii}_t[\sigma_{z}^\top\sigma_{z}]^{jj}_t}} 
\end{eqnarray}
\end{prop}

\begin{proof}
It is trivial to check that the $\widetilde{W}_t^i$ defined by 
\begin{eqnarray}
  \widetilde{\sigma}_x^i(t)  &=& \| \sigma_x^{\cdot i}(t)\| = \sqrt{ \sum_j \sigma_{z}^{ji}(t)^2 } \\
  d\widetilde{W}_t^i   &=& \frac{1}{ \widetilde{\sigma}_x^i(t) } \sum_j \sigma_{z}^{ji}(t)  dW_t^j   .
\end{eqnarray}
are Brownian motions. Moreover, the instantaneous covariances are given by :
\begin{eqnarray}
  d\langle x^i_t,x^j_t\rangle &=& d\left\langle \sum_k \sigma_{z}^{ki}(t) dW_t^k, \sum_\ell \sigma_{z}^{\ell j}(t) dW_t^\ell \right\rangle     \\
  &=& \sum_k \sigma_{z}^{ki}(t) \sigma_{z}^{kj}(t) dt \\
  &=& [\sigma_{z}^\top\sigma_{z}]^{ij}(t)dt. 
\end{eqnarray}
Hence the expression of the instantaneous correlations.
\end{proof}

\begin{rema}
In vectorial form, the dynamic is 
\begin{eqnarray}
  dz_t &=& \left(y_t \vun - \beta_{t} z_t\right)dt + \sigma_{z,t} \Gamma_t^\top d\widehat{W}^\RN_t 
\end{eqnarray}
where the factor loading matrix $\Gamma_t$ is a square root of the correlation matrix
\begin{eqnarray}
  \Gamma_t^\top \Gamma_t &=& (\rho^{ij}_t) .
\end{eqnarray}
with the same reconstruction formulas:
\begin{enumerate}
\item The forward rate is:
\begin{eqnarray}
  r_{t,T} &=& r_{0,t} + M_{t,T}^\top\left( z_t + y_t \int_t^T M_{t,u} \, du \right) 
\end{eqnarray}
\item The ZCB is:
\begin{eqnarray}
  Z_{t,T} &=& Z_{0,t,T} \exp\left( -G_{t,T} ^\top z_t -\frac{1}{2} G_{t,T} ^\top y_t G_{t,T}   \right) 
\end{eqnarray}
\end{enumerate}
\end{rema}


%
%
%\newpage%
%\subsubsection{Parametrisation by factor weights and correlations}%
%
%The spot volatility vectors of instantaneous forward rates and collateral basis are 
%\begin{eqnarray}
%  \sigma^{f|f}_{t,T}  &=& \sigma^{f|f}_t  
%  \begin{bmatrix}  
%  \alpha^{f|f}_1  e^{- \int_t^T \beta^{f|f}_1(u) \, du } \\ 
%  \vdots \\ 
%  \alpha^{f|f}_{d_f}  e^{- \int_t^T  \beta^{f|f}_{d_f}(u) \, du } 
%  \end{bmatrix} \\
%  \sigma^{i\to j}_{t,T}  &=& \sigma^{i\to j}_t   
%  \begin{bmatrix}  \alpha^{i\to j}_1  e^{-\int_t^T \beta^{i\to j}_1(u) \, du }  \\ 
%  \vdots \\ 
%  \alpha^{i\to j}_{d_{i\to j}}  e^{- \int_t^T \beta^{i\to j}_{d_{i\to j}}(u) \, du } 
%  \end{bmatrix}
%\end{eqnarray}



\subsubsection{Simulation}%


\begin{prop}
Let $z_t$ be a multi-dimensional OU process:
\begin{eqnarray}
  dz_t &=& \left(y_t\vun - \beta_{t}z_t\right)\,dt + \sigma_{z,t}^\top  d\widehat{W}_t 
\end{eqnarray}
then 
\begin{eqnarray}
  \EE[z_{t+h} | z_t ] &=& e^{-\int_t^{t+h} \beta_u \, du}z_t + \int_t^{t+h}  e^{-\int_s^{t+h} \beta_u \, du}y_s \vun \, ds  \\
  \VV[z_{t+h} | z_t ] &=&  \int_t^{t+h} e^{-\int_s^{t+h} \beta_u \, du} \sigma_{z,s}^\top \sigma_{z,s}  e^{-\int_s^{t+h} \beta_u \, du} \,  ds  
\end{eqnarray}
\end{prop}


\subsubsection{Correlation structure of instantaneous forward rates}%



\subsubsection{Correlation structure of forward Libor rates}%



\subsubsection{Correlation structure of forward swap rates}%



%\newpage
\subsection{Reminder on LGM parametrisations: Multi-currency case}%


\subsubsection{FX dynamic}%


\subsubsection{Quanto-drift adjustment}%



\subsubsection{Full dynamic}%


\subsection{[TODO: Remove this]: Parametrisation}



\newpage%
\subsection{Precomputation of model quantities}

Model prescribes quantities that can be precomputed:
\begin{enumerate}
\item Covariance of increments of the state variable
\item Drift adjustments of increments of the state variable
\item Quanto adjustments 
\item Collateral convexity adjustments of market observables.
\end{enumerate}
All of these are expressed as covariance integrals of the form:
\begin{eqnarray}
  I^{a,b}_{t,S;T_a,T_b} &=& \int_t^S \Sigma^a_{u,T_a} \cdot \Sigma^b_{u,T_b} \, du 
\end{eqnarray}

The objective of this section is to describe the algorithm for their computation.


\subsubsection{}

Remember that quantities of interest are 
\begin{eqnarray}
	\sigma^{d|d}_{t,T} 
	&=& \sigma^{d|d}_r(t) \begin{bmatrix} e^{\int_0^t  \beta^1_v \,dv } \\
	\vdots \\ 
	e^{\int_0^t  \beta^d_v \,dv } 
	\end{bmatrix} 
	\begin{bmatrix} e^{-\int_0^T  \beta^1_v \,dv } \\
	\vdots \\ 
	e^{-\int_0^T  \beta^d_v \,dv } 
	\end{bmatrix}  \\ 
	\Sigma^{d|d}_{t,T}  &=& \int_t^T \sigma^{d|d}_{t,u}\,du \\
	&=&   \sigma^{d|d}_r(t) \begin{bmatrix} e^{\int_0^t  \beta^1_v \,dv } \\
	\vdots \\ 
	e^{\int_0^t  \beta^d_v \,dv } 
	\end{bmatrix} 
	\begin{bmatrix} \int_t^T e^{-\int_0^u  \beta^1_v \,dv } \, du \\
	\vdots \\ 
	\int_t^T  e^{-\int_0^u  \beta^d_v \,dv } \, du
	\end{bmatrix}   \\
	&=& g(t)G(t,T)
\end{eqnarray}


\begin{prop}
Let us introduce:
\begin{eqnarray}
	h^i(T) &=& e^{-\int_0^T  \beta^i_v \,dv }  \\
	\Lambda^i(T) &=& \int_0^T e^{-\int_0^u  \beta^i_v \,dv } \, du .
\end{eqnarray}
\end{prop}

The quantities $h^i(t_k)$ and $\Lambda^i(t_k)$ are precomputed and stored. 
Since we assume piecewise constant mean reversion, we compute these integrals:
\begin{eqnarray}
	-\ln h^i(t_{k+1}) &=& -\ln h^i(t_{k}) + \int_{t_k}^{t_{k+1}} \beta^i_v\, dv  \\
	&=& -\ln h^i(t_{k}) +  \sum_{\tau_j\in [t_k,t_{k+1}[} \beta^i(\tau_j) (\min(\tau_{j+1},t_{k+1}) - \tau_j) .    
\end{eqnarray}
We then exponentiate these to get the values of $h^i(t_{k})$.

It could be easier to refine the term structures so that the mean reversion is piecewise constant on all time steps.

Because, $\beta^i$ is piece-wise constant, its integral is piecewise linear so the exponential can be integrated.
\begin{eqnarray*}
	\Lambda^i(t_{k+1})  
	&=& \int_0^{t_{k+1}} e^{-\int_0^u  \beta^i_v \,dv } \, du  \\
	&=& \int_0^{t_{k}} e^{-\int_0^u  \beta^i_v \,dv } \, du  
	+  \int_{t_{k}}^{t_{k+1}} e^{-\int_0^u  \beta^i_v \,dv } \, du  \\
	&=& \Lambda^i(t_{k+1})    
	+  e^{-\int_0^{t_k}  \beta^i_u \,dv }  \int_{t_{k}}^{t_{k+1}} e^{-\int_{t_k}^u  \beta^i_v \,dv } \, du  \\
	&=& \Lambda^i(t_{k+1})    
	+  h^i(t_k) \int_{t_{k}}^{t_{k+1}} e^{-\beta^i_k (u-t_k) } \, du  \\
	&=& \Lambda^i(t_{k+1})  + \begin{cases}
	h^i(t_k) \frac{1-e^{-\beta^i_k(t_{k+t}-t_k)}}{\beta^i_k}   & \textrm{if $\beta^i_k \neq 0$} \\ 
	h^i(t_k) (t_{k+t}-t_k)  & \textrm{else} 
	\end{cases}
\end{eqnarray*}
From these we can compute any requested $G^i(t_j,t_k)$ as 
\begin{eqnarray*}
	G^i(t,T) &=& \int_{t}^T  e^{-\int_0^u  \beta^i_v \,dv } \, du \\ 
	&=& \left(\Lambda^i(T) - \Lambda^i(t)\right) / h^i(t) .
\end{eqnarray*}

Then we need to compute the $y(t_k)$'s.
\begin{prop}
We have
\begin{eqnarray}
	y(t) &=& \left(\Lambda^i(T) - \Lambda^i(t)\right) / h^i(t) .
\end{eqnarray}
\end{prop}



\begin{prop}
We have
\begin{eqnarray}
	G^i(t,T) &=& \left(\Lambda^i(T) - \Lambda^i(t)\right) / h^i(t) .
\end{eqnarray}
\end{prop}



\begin{prop}[]
The short rate is 
\begin{eqnarray}
	r_t &=& r_{0,t} + 1^\top z_t .
\end{eqnarray}
\end{prop}


\begin{prop}[Zero coupon reconstruction]
We have
\begin{eqnarray}
	Z^{d|d}_{t,T} &=& Z^{d|d}_{0,t,T} \exp\left( - G(t,T)^\top z_t - \frac{1}{2} G(t,T)^\top y(t) G(t,T) \right)
\end{eqnarray}
\end{prop}

\begin{prop}[Caplet volatility]
Let us introduce:
\begin{eqnarray}
  w^{d|d}_{t,T;\theta} &=& \int_t^T \left| \Sigma^{d|d}_{u,T;\theta}  \right|^2\, du
\end{eqnarray}
Then 
\end{prop}

\begin{proof}
\begin{eqnarray}
	\frac{d(1+\delta L_{t,T;\delta})}{1+\delta L_{t,T;\delta}} 
	&=& \left( \Sigma_{t,T+\delta} - \Sigma_{t,T} \right)^\top dW^{d|d,T+\delta} 
\end{eqnarray}
So:
\begin{eqnarray}
  (1+\delta L_{S,T;\delta} 
  &=& (1+\delta L_{t,T;\delta})\exp\left( \int_t^S   \Sigma_{u,T;\delta} \cdot  dW^{d|d,T+\delta}_u  - \frac{1}{2} \int_t^S \left| \Sigma_{u,T;\delta}  \right|^2 du \right) 
\end{eqnarray}
and 
\begin{eqnarray}
  \VV_t\left[\ln(1+\delta L_{T,T;\delta}\right] &=&  \int_t^T \left| \Sigma_{u,T+\delta} - \Sigma_{u,T} \right|^2\, du 
\end{eqnarray}
\end{proof}

\begin{prop}[Gaussian swap rate approximation]
Let us introduce:
\begin{eqnarray}
	\VV^{d|d,T}_t[Y^{d|d}_{T,T;M}] &=& 
\end{eqnarray}
\end{prop}


\subsection{Calibration} 


\begin{prop}[Spot volatility bootstrapping]

\end{prop}


\begin{prop}[Mean reversion calibration]

\end{prop}





\subsection{Monte-Carlo simulation} 

\subsubsection{Monte-Carlo simulation of Ornstein-Uhlenbeck processes} 


\begin{eqnarray*}
  dX_{t}  &=& - \beta_t X_t + \sigma_t  \, dW_t \\
  d\left( e^{\int_0^t\beta_u\,du }X_{t} \right) &=&  \sigma_t e^{\int_0^t\beta_u\,du } \, dW_t \\
  e^{\int_0^{t_{k+1}}\beta_u\,du }X_{t_{k+1}}  &=&  e^{\int_0^{t_k}\beta_u\,du }X_{t_{k}} + \int_{t_{k}}^{t_{k+1}} \sigma_u e^{\int_0^u\beta_v\,dv } \, dW_v \\
  X_{t_{k+1}}  &=&  e^{-\int_{t_{k}}^{t_{k+1}}\beta_v\,dv }X_{t} + e^{-\int_0^{t_{k+1}}\beta_v\,dv } \int_{t_{k}}^{t_{k+1}} \sigma_u e^{\int_0^u\beta_v\,dv } \, dW_v \\
   &=&  e^{-\int_{t_{k}}^{t_{k+1}}\beta_v\,dv }X_{t} + e^{-\int_{t_{k}}^{t_{k+1}}\beta_v\,dv } \int_{t_{k}}^{t_{k+1}} \sigma_u e^{\int_{t_{k}}^u\beta_v\,dv } \, dW_v 
\end{eqnarray*}

\begin{prop}
Conditionally on $X_{t_k}$, $X_{t_{k+1}}$ is normal with moments:
\begin{eqnarray}
  \EE[X_{t_{k+1}} | X_{t_k}] &=& e^{-\int_{t_k}^{t_{k+1}}\beta_v\,dv }  X_{t_k}  \\
  \VV[X_{t_{k+1}} | X_{t_k}] &=& \int_{t_{k}}^{t_{k+1}} \left| \sigma_u e^{\int_0^u\beta_v\,dv } \right|^2 dv
\end{eqnarray} 
\end{prop}

This allows for the exact simulation of the process using the scheme: 
\begin{eqnarray*}
  X_{t_{k+1}}  &=&   \alpha_{k+1} X_{t_k} + \alpha_{k+1} \Sigma_{k+1} \xi_{k+1}
\end{eqnarray*}
where 
\begin{eqnarray*}
  \alpha_{k+1} &=& e^{-\int_{t_k}^{t_{k+1}}\beta_v\,dv }  \\
  \Sigma_{k+1} &=& \sqrt{ \int_{t_k}^{t_{k+1}} \sigma_u^2 e^{-2\int_{t_k}^u\beta_v\,dv } \, du} 
\end{eqnarray*}
and $(\xi_k)$ is a Gaussian white noise. 

For correlated OU processes,
\begin{eqnarray*}
  X^i_{t_{k+1}}  &=&  e^{-\int_{t_k}^{t_{k+1}}\beta^i_v\,dv } X^i_{t_k} + (L_k Z_{k+1})^i
\end{eqnarray*}
where 
\begin{eqnarray*}
  \Gamma^{ij}_{k+1} 
   &=&  \int_{t_k}^{t_{k+1}} \rho^{ij}_u \sigma^i_u \sigma^j_u e^{\int_{u}^{t_{k+1}}(\beta^i_v+\beta^j_v)\,dv } \, du 
\end{eqnarray*}
and $L$ is a square root of $\Gamma$, meaning that $L$ is a solution of $LL^\top = \Gamma$ as given by the SVD or Cholesky decomposition.




\subsection{Precomputation of covariance matrix}





\subsection{The Single Currency Linear Gaussian Model: Calibration}



\subsubsection{Calibration of native rates markets}%


\subsubsection{Calibration of spot volatility}%



\subsubsection{Calibration of mean reversion}%





\newpage%
\section{Linear Gaussian Collateral Model: proofs}%


\newpage%
\section{Gaussian approximation of par swap spreads}\label{sec:gaussian_approx_proofs}%

In this section, we prove the Gaussian approximations of forward swap rates and forward currency basis spreads stated in \autoref{sssec:lgcm_gaussian_approx}.

\subsubsection{Proof of \autoref{lemm:zhat_forward_law}}
By Girsanov, $dW^{d,\RN}_t = dW^{f|d,T}_t + \left(\sigma^{df}_t - \Sigma^{f|d}_{t,T}\right) dt$ and $\Sigma^{f|d}_{t,T} = \hat{\sigma}_{z,t}\hat{G}_{t,T}$. So, by \autoref{rema:qgcm_vfc},
\begin{eqnarray*}
  d\hat{z}_t &=& \left( \hat{y}_t\vun - \hat{\sigma}^\top_{z,t}\hat{\sigma}_{z,t}\hat{G}_{t,T} - \hat{\beta}_t \hat{z}_t\right) dt + \hat{\sigma}^\top_{z,t} dW^{f|d,T}_t .
\end{eqnarray*}
The coefficients are deterministic in the LGCM, so $\hat{z}_T$ is Gaussian with covariance $\int_0^T \hat{\mathcal{M}}_{u,T}\hat{\sigma}^\top_{z,u}\hat{\sigma}_{z,u}\hat{\mathcal{M}}_{u,T}\,du = \hat{y}_T$ and mean
\begin{eqnarray*}
  \int_0^T \hat{\mathcal{M}}_{u,T}\left(\hat{y}_u\vun - \hat{\sigma}^\top_{z,u}\hat{\sigma}_{z,u}\hat{G}_{u,T}\right) du ~=~ 0 ,
\end{eqnarray*}
since, by Fubini and $\hat{\mathcal{M}}_{u,T}\hat{\mathcal{M}}_{w,u} = \hat{\mathcal{M}}_{w,T}$,
\begin{eqnarray*}
  \int_0^T \hat{\mathcal{M}}_{u,T}\hat{y}_u\vun\, du ~=~ \int_0^T \hat{\mathcal{M}}_{w,T} \hat{\sigma}^\top_{z,w}\hat{\sigma}_{z,w} \int_w^T \hat{\mathcal{M}}_{w,u}\vun \, du\, dw ~=~ \int_0^T \hat{\mathcal{M}}_{w,T} \hat{\sigma}^\top_{z,w}\hat{\sigma}_{z,w} \hat{G}_{w,T}\, dw .
\end{eqnarray*}

\subsubsection{Proof of \autoref{lemm:par_spread_state_functions}}
The formula for $P_{t,j}$ is \eqref{eq:qgcm_zfd_reconstruction}. Under the constant funding spread hypothesis, $1+\theta L^{f|f}_{t,T;\theta} = (1-\theta S^{f|f,\div}_{0,T;\theta})Z^{f|f}_{t,T}/Z^{f|f}_{t,T+\theta}$, and the foreign Libor under domestic collateral follows from \eqref{eq:libor_collateral_convexity_lgcm}. The resetting forward Libor is given by the proposition on resetting forward Libors in the HJM section.

\subsubsection{Proof of \autoref{prop:gauss_swap_rate}}
Write $\ell_{t,j} = 1+\theta_jL_{t,j}$, so that $Y_t = A_t^{-1}\sum_j P_{t,j}(\ell_{t,j}-1)$. By \autoref{lemm:par_spread_state_functions}, $Y_t = Y(t,\hat{z}_t)$ with
\begin{eqnarray*}
  \nabla_{\hat{z}} \ln P_{t,j} ~=~ -\hat{G}_{t,T_j} , \qquad
  \nabla_{\hat{z}} \ln \ell_{t,j} ~=~ \begin{bmatrix} G^f_{t,T_j} - G^f_{t,T_{j-1}} \\ 0 \end{bmatrix} .
\end{eqnarray*}
Since $Y_t$ is a $\QQ^{f|d,A}$-martingale, its dynamic is $dY_t = \nabla_{\hat{z}}Y(t,\hat{z}_t)^\top\hat{\sigma}^\top_{z,t}dW^{f|d,A}_t$, and a direct computation gives
\begin{eqnarray*}
  \nabla_{\hat{z}} Y &=& \frac{1}{A_t} \sum_{j=1}^n P_{t,j} \begin{bmatrix} (1+\theta_jY_t)G^f_{t,T_j} - \ell_{t,j}G^f_{t,T_{j-1}} \\ \theta_j(Y_t - L_{t,j})G^c_{t,T_j} \end{bmatrix} .
\end{eqnarray*}
For $T\ge T_0$, $G_{t,T} = G_{t,T_0} + \mathcal{M}_{t,T_0}G_{T_0,T}$. The terms in $G_{t,T_0}$ cancel since $\sum_j P_{t,j}\left((1+\theta_jY_t) - \ell_{t,j}\right) = \sum_j\theta_jP_{t,j}(Y_t - L_{t,j}) = 0$, so
\begin{eqnarray*}
  \nabla_{\hat{z}} Y(t,\hat{z}_t) &=& \hat{\mathcal{M}}_{t,T_0}\, \hat{q}_Y(t)
\end{eqnarray*}
where $\hat{q}_Y(t)$ is given by the formulas of the proposition evaluated with the time-$t$ values of the ratios $P_{t,j}/A_t$, $L_{t,j}$ and $Y_t$. The weight freezing approximation consists in replacing $\hat{q}_Y(t)$ by its value $\hat{q}_Y$ computed from forward values seen from $0$. Since $\hat{q}_Y(t)$ only depends on ratios of discount factors, we can use the forward discount factors $Z^{f|d}_{0,T_0,T_j}$. Then $Y_{T_0} - Y_0 = \int_0^{T_0} \hat{q}_Y^\top \hat{\mathcal{M}}_{t,T_0}\hat{\sigma}^\top_{z,t} dW^{f|d,A}_t$ is centred Gaussian with variance $\int_0^{T_0}\hat{q}_Y^\top \hat{\mathcal{M}}_{t,T_0}\hat{\sigma}^\top_{z,t}\hat{\sigma}_{z,t}\hat{\mathcal{M}}_{t,T_0}\hat{q}_Y\,dt = \hat{q}_Y^\top\hat{y}_{T_0}\hat{q}_Y$.

\subsubsection{Proof of \autoref{prop:gauss_cb_spread}}
Write $\ell^r_{t,j} = 1+\theta_jL^r_{t,j}$, so that $B_t = A_t^{-1}\sum_j P_{t,j}(\ell^r_{t,j} - \ell_{t,j})$. By \autoref{lemm:par_spread_state_functions}, $P_{t,j}\ell^r_{t,j} = (1-\theta_jS^{d|d,\div})P_{t,j-1}$ and $P_{t,j}\ell_{t,j}$ is proportional to $P_{t,j}Z^{f|f}_{t,T_{j-1}}/Z^{f|f}_{t,T_j}$, so
\begin{eqnarray*}
  \nabla_{\hat{z}} \left(P_{t,j}\ell^r_{t,j}\right) = - P_{t,j}\ell^r_{t,j} \hat{G}_{t,T_{j-1}}, \qquad
  \nabla_{\hat{z}} \left(P_{t,j}\ell_{t,j}\right) = - P_{t,j}\ell_{t,j} \begin{bmatrix} G^f_{t,T_{j-1}} \\ G^c_{t,T_j} \end{bmatrix} ,
\end{eqnarray*}
and $\nabla_{\hat{z}} A_t^{-1} = A_t^{-2}\sum_j\theta_jP_{t,j}\hat{G}_{t,T_j}$. Using $\ell^r_{t,j} - \ell_{t,j} = \theta_jB_{t,j}$, we get
\begin{eqnarray*}
  \nabla_{\hat{z}} B &=& \frac{1}{A_t} \sum_{j=1}^n P_{t,j} \begin{bmatrix} \theta_j \left( B_t G^f_{t,T_j} - B_{t,j} G^f_{t,T_{j-1}} \right) \\ (\ell_{t,j} + \theta_jB_t) G^c_{t,T_j} - \ell^r_{t,j} G^c_{t,T_{j-1}} \end{bmatrix} .
\end{eqnarray*}
The terms in $G_{t,T_0}$ cancel in both blocks since $\sum_j\theta_jP_{t,j}(B_t - B_{t,j}) = 0$, and we conclude as for swap rates.

\begin{rema}
The weight freezing could be applied to other representations of the same par spread, for example the representation of $B_t$ as a weighted sum of differences of funding spreads in the corollary on resetting cross-currency swaps. All representations have the same exact gradient at $t$, so they lead to the same frozen loadings when the gradient is frozen at the forward values. They differ only if one freezes the weights rather than the gradient.
\end{rema}



%\newpage%
%\section{Historical estimation of parameters}
%
%
%\subsection{Static vs rolling instruments}%
%We observe quotes for rolling instruments but are trying to calibrate the model at time $t_0$ to a static set of instruments. What is the dynamic of rolling currency basis spreads in our model?   
%
%
%
%\subsection{Historical distribution of rolling currency basis spreads}%
%
%
%
%
%\subsection{Liquidity}%
%
%
%
%\subsection{Jumps}%
%It seems that most jumps are associated with political and economical announcements. They should not be included in the estimation of volatility as this would lead to an overestimation of volatility with no practical relation to daily delta hedging. Instead, the traders should adjust for these events.
%
%
%\subsection{Estimation methodology}%
%
%We consider a model for the forward (swap rate, currency basis, etc\ldots)
%\begin{eqnarray}
%  dF_{t,T}  &=&  \sigma_{t,T}  \, dW^{F}_t 
%\end{eqnarray}
%under its natural martingale measure (typically the annuity measure). 
%
%We make the stationarity assumption $\sigma_{t,T} = \nu_{T-t}$. 
%
%From the daily time series of rolling instruments, we build an estimator of spot volatility:
%\begin{eqnarray*}
%  \widehat{\nu}_\tau  &=& A \mathrm{StdDev}(F_{t_k,t_{k-1}+\tau} - F_{t_k,t_{k-1}+\tau})
%\end{eqnarray*}
%where $A$ is an annualisation constant, e.g. $A = \sqrt{252}$.
%
%Then, the total variance of the forward is
%\begin{eqnarray}
%  \VV[F_{T,T}]  ~=~ \int_0^T \sigma_{t,T}^2  \, dt 
%  ~=~  \int_0^T \nu_{T-t}^2  \, dt 
%\end{eqnarray}
%Discretising this integral and replacing the spot vols by their estimators leads to the following estimator
%\begin{eqnarray}
%  \widehat{w} &=& \sum_i  \widehat{\nu}_{\tau_i} \frac{\widehat{\nu}_{\tau_{i}} + \widehat{\nu}_{\tau_{i-1}}}{2} (\tau_{i} - \tau_{i-1}).
%\end{eqnarray}
%of total variance. 




\newpage%
\section{Normal approximations}

\subsection{Shifted lognormal distributions}
In the Linear Gaussian Model, the simply compounded forward rates (Libors) follow shifted lognormal distributions. But, when calibrating the model to G10 market data, one finds that caplet smile is flat which means that the model distribution of Libors is very close to be Gaussian. 

In this section, we derive Gaussian approximations of shifted lognormal distributions.  

\subsubsection{Moment matching}

\paragraph{Random variables}

Let $Z$ be a standard Gaussian variable and let $X$ be a shifted lognormal variable with distribution:
\begin{eqnarray}
  1 + \theta X &=& (1+\theta \mu_{\mathrm{LN}} )e^{\sigma_{\mathrm{LN}} Z - \frac{1}{2}\sigma_{\mathrm{LN}}^2} 
\end{eqnarray}
where $Z$ is standard Gaussian variable. 

We can compute the moments of $X$: 
\begin{eqnarray}
  \EE[X] &=& \mu_{\mathrm{LN}} \\ 
  \VV[X] &=& \left(\frac{1+\theta \mu_{\mathrm{LN}}}{\theta}\right)^2 \left( e^{\sigma_{\mathrm{LN}}^2} - 1 \right) 
\end{eqnarray}
It follows that $X$ can be approximated by 
\begin{eqnarray}
  Y &=& \mu_{\mathrm{N}} +  \sigma_{\mathrm{N}}Z \\
  \textrm{where} \qquad \mu_{\mathrm{N}} &=& \mu_{\mathrm{LN}} \\ 
  \sigma_{\mathrm{N}} &=& \frac{1+\theta \mu_{\mathrm{LN}}}{\theta} \sqrt{e^{\sigma_{\mathrm{LN}}^2} - 1}
\end{eqnarray}



\paragraph{Martingales}
The previous result easily extends to Gaussian martingales.

\begin{prop}
Let $X_t$ be a shifted exponential martingale:
\begin{eqnarray}
  \frac{d(1 + \theta X_t)}{1+\theta X_t}  
  &=& \sigma_{\mathrm{LN},t} dW_t  
\end{eqnarray}
Then, when the lognormal volatility of marginal distribution is low, $X_t$ is well approximated by a Brownian martingale:
\begin{eqnarray}
  dY_t   &=& \sigma_{\mathrm{N},t} dW_t  
\end{eqnarray}
where 
\begin{eqnarray}
   Y_0 &=& X_0 \\
   \sigma_{\mathrm{N},t} &=& \frac{1+\theta X_0}{\theta}  \sigma_{\mathrm{LN},t}
\end{eqnarray}
\end{prop}

\begin{proof}
We clearly match expectations so let us match total variance, that is that variance of each terminal distribution $X_T$:
\begin{eqnarray}
   \int_0^T \sigma_{\mathrm{N},t}^2 \, dt &=&  \left(\frac{1+\theta X_0}{\theta} \right)^2 \left( e^{\int_0^T \sigma_{\mathrm{LN},t}^2\, dt} - 1 \right)
\end{eqnarray}
Taking the derivative wrt $T$, we find:
\begin{eqnarray}
   \sigma_{\mathrm{N},T}^2  &=&  \left(\frac{1+\theta X_0}{\theta} \right)^2 \sigma_{\mathrm{LN},T}^2 \left(  e^{\int_0^T \sigma_{\mathrm{LN},t}^2\, dt} - 1 \right). 
\end{eqnarray}
And, at first order, 
\begin{eqnarray}
   \sigma_{\mathrm{N},t}  & \approx &  \frac{1+\theta X_0}{\theta}  \sigma_{\mathrm{LN},t}.
\end{eqnarray}
\end{proof}




\subsubsection{Cumulative distribution}


Let $Z$ be a standard Gaussian variable and $1 + \theta X = (1+\theta \mu)e^{\sigma Z - \frac{1}{2}\sigma^2}$. The CDF of $X$ is 
\begin{eqnarray*}
  F_X(x)  &=& \PP\left( 1 + \theta X \le 1+\theta x \right) \\
  &=&  \PP\left( (1+\theta \mu)e^{\sigma Z - \frac{1}{2}\sigma^2} \le 1+\theta x \right) \\
  &=&  \PP\left( e^{\sigma Z - \frac{1}{2}\sigma^2} \le \frac{1+\theta x}{1+\theta \mu} \right) \\
  &=&  \PP\left( \sigma Z - \frac{1}{2}\sigma^2 \le \ln \left(\frac{1+\theta x}{1+\theta \mu} \right)  \right) \\
  &=&  \PP\left( Z \le \ln \left(\frac{1+\theta x}{1+\theta \mu} \right) + \frac{1}{2}\sigma^2   \right) \\
  &=&  \Phi\left( \ln \left(\frac{1+\theta x}{1+\theta \mu} \right) + \frac{1}{2}\sigma^2  \right) 
\end{eqnarray*}
We can then expand\footnote{
The lognormal random variable with parameters $\mu, \sigma$ is 
$e^{\mu + \sigma Z}$, where $Z$ is standard normal, and so its CDF is
$F(x) = \Phi\left(\frac{\ln(x) - \mu}{\sigma}\right)$ for $x > 0$, where $\Phi$ is
the standard normal CDF.  When $x = e^{\mu} (1+ t)$ is near $e^{\mu}$, 
$$\ln(x) = \mu + \ln(1 + t) = \mu + t  + O(t^2)$$
so that
$$F(x) \approx \Phi\left(\frac{t }{\sigma}\right)$$
In fact
    $$F(x) - \Phi\left(\frac{t }{\sigma}\right) = \Phi\left(\frac{\ln(1+t )}{\sigma}\right) - \Phi\left(\frac{t}{\sigma}\right) $$
    
This will be a very good approximation for $t$ near $0$. For fixed $t$, it
goes to $0$ as $\sigma \to 0$ because both $\Phi$'s go to  $0$ if $t < 0$
and both $\Phi$'s go to $1$ if $t > 0$.  When $\sigma$ is small there appear to be two local
minima of $F(x) - \Phi\left({t }/{\sigma}\right) $, at approximately
$t = \pm \sqrt{2} \sigma$, where
$F(x) - \Phi\left(t/\sigma\right) \approx  -\sigma/(\sqrt{2\pi} e) \pm
\sigma^2/(6 \sqrt{\pi} e)$.  Thus the maximum error in $F$ is
approximately $$ \frac{\sigma}{\sqrt{2\pi} e} + \frac{\sigma^2}{6 \sqrt{\pi} e}$$

For example, for $\sigma = 1/2$, the actual maximum error is approximately
$0.0817231495769266$, and the  above approximation to the maximum error is $0.08202940444$.
As $\sigma$ gets smaller, the approximation should get better. 
}
around $x = \mu$.

\subsection{Lognormal mixtures}

\subsubsection{Additive mixtures of lognormal variables}

Additive mixtures appear in options on baskets in the BS model. One can moment match a shifted lognormal distribution to the first 3 moments of the basket. This can give very good results at very low cost. 


\subsubsection{Difference of lognormal variables}
We consider 
\begin{eqnarray}
  X &=& e^{\mu^1 + \sigma^1 Z^1} - e^{\mu^2 + \sigma^2 Z^2} 
\end{eqnarray}
where $(Z_1,Z_2)$ is Gaussian vector with unit variance but non-zero correlation. 

For processes, the analog is:
\begin{eqnarray}
  X_t &=& e^{\mu^1_t +  \int_0^t \sigma^1_u dW^1_u} - e^{\mu^2_t + \int_0^t \sigma^2_u dW^2_u} 
\end{eqnarray}
where $W^1,W^2$ are correlated Brownian motions. 

This is the case of the currency basis which is the difference of two lognormal processes. Note that this implies that the model assigns a non-zero probability to the basis changing sign. Moreover the marginal distributions have fatter tails than a Gaussian. 

Nevertheless, we can approximate this process by a normal or shifted lognormal process by moment matching of the marginal distributions.

\paragraph{Normal approximation}

\begin{eqnarray}
  \EE[X] &=& e^{\mu^1 +  \frac{1}{2} (\sigma^1)^2} - e^{\mu^2 +  \frac{1}{2} (\sigma^2)^2}  \\
  \VV[X]  &=& e^{\mu^1} e^{\mu^2} e^{\rho \sigma^1 \sigma^2}  ????
\end{eqnarray}

So $X$ can be approximated by $\mathcal{N}(\mu,\sigma^2)$ where:
\begin{eqnarray}
  \mu &=& e^{\mu^1 +  \frac{1}{2} (\sigma^1)^2} - e^{\mu^2 +  \frac{1}{2} (\sigma^2)^2}  \\
  \sigma  &=& e^{\frac{1}{2}(\mu^1 + \mu^2 + \rho \sigma^1 \sigma^2)}   ????
\end{eqnarray}


\paragraph{Shifted lognormal approximation}

We have 
\begin{eqnarray}
  \EE[X] &=& e^{\mu^1 +  \frac{1}{2} (\sigma^1)^2} - e^{\mu^2 +  \frac{1}{2} (\sigma^2)^2}  \\
  \VV[X]  &=& e^{\mu^1} e^{\mu^2} e^{\rho \sigma^1 \sigma^2}  \\
  \mathbb{SK}[X]  &=& ?? 
\end{eqnarray}

Comparatively, the moments of a shifted lognormal $Y = a + e^{\mu + \sigma Z}$ are: 
\begin{eqnarray}
  \EE[Y] &=& a + e^{\mu + \frac{1}{2}\sigma^2} \\
  \VV[Y]  &=& e^{2\mu +  \frac{1}{2}\sigma^2}  \\
  \mathbb{SK}[Y]  &=& ??
\end{eqnarray}

We can solve this system of equations:



\subsection{Second order mixtures}

In the case of an option on a non-resetting cross-currency swaplet on a non-USD pair, the underlying is a basket with 3 terms:
\begin{enumerate}
\item $B^{d|c,nr}_{T,T;\theta}$: Normal 
\item $S\left(\frac{x^{df}_{T}}{x^{df}_{t}} - 1\right)$: Lognormal.
\item $x^{df}_{t}B^{f|c,nr}_{T,T;\theta}$: product of Normal and Lognormal
\end{enumerate}

Let us moment match this distribution to a shifted lognormal.

[TODO: Write down the computation]



\newpage%
\section{Estimation of historical volatilities and correlations}%




\subsection{Stationarity assumption}%

\subsubsection{Market quotes and stationarity}


\subsubsection{Statistical evidence}

\subsection{Model}%

\subsubsection{Statistical evidence}


\subsection{Estimator}%



\subsubsection{Unexplained}%

\begin{rema}
There is a significant gap between implied and realised vols. Here are a list of possible explanations:
\begin{enumerate}
\item The model assumes Gaussian returns while we know they are fat-tailed. 
\item The implied volatility quoted in the market account for transaction/liquidity costs that the model ignores.  
\end{enumerate}
\end{rema}

\begin{rema}[Rough volatility]
Empirical studies show that volatility is rough, i.e. similar to a fractional Brownian motion with Hurst exponent $H \approx 0.1$. 
\end{rema}


\newpage%
\begin{thebibliography}{99}
%
\bibitem{AndersenPiterbarg}
\textsc{L. Andersen \& V. Piterbarg,}
\textit{Interest rate modelling}.
%
\bibitem{BjorkChristensen}
\textsc{T. Bj\"ork \& B. J. Christensen,}
\textit{Interest rate dynamics and consistent forward rate curves},
Mathematical Finance 9(4), 323--348, 1999.
%
\bibitem{Filipovic:NelsonSiegel}
\textsc{D. Filipovi\'c,}
\textit{A note on the Nelson-Siegel family},
Mathematical Finance 9(4), 349--359, 1999.
%
\bibitem{GeometryOfMultiCurve}
\textsc{C. Fontana, G. Lanaro, A. Murgoci,}
\textit{The geometry of multi-curve interest rate models},
SSRN.
%
\bibitem{McCloud:CollateralConvexity}
\textsc{P. McCloud,}
\textit{Collateral Convexity Complexity},
Risk Magazine, 2011.
%
\bibitem{McCloud:CollateralVolatility}
\textsc{P. McCloud,}
\textit{Collateral Volatility},
SSRN.
%
\bibitem{Piterbarg:Cooking}
\textsc{V. Piterbarg,}
\textit{Cooking with collateral}.
%
\end{thebibliography}








\end{document}






%%%%%%%%%%%%%%%%%%%%%%%%%%%%%%%%%%%%%%%%%%%%%%%%%%%%%%%%%%%%%%%%%%%%%%%%%%%%%%%%%%%%%%%%%%%%var